% source: https://learning.edx.org/course/course-v1:HarvardX+TinyML3+1T2022/home
% Applications of TinyML
% google log in

% Bilder TinyML/Deployment

% Reihenfolge
%1. TinyML/Fundamentals
%2. TinyML/Applications
%3. TinyML/deployment
%4. TinyML/MLOps for Scaling TinyML

\chapter{Deploying TinyML}


\section{Objectives}

The goal of this course is to teach learners how to engineer end-to-end tinyML applications using TensorFlow Micro. We teach learners how to program in TF Micro, and use it to deploy real-world applications.

\section{Prerequisites}

\subsection{Prerequisites}

\begin{itemize}
  \item  Applications of TinyML (Course 2)
  \item Basic programming in C++
  \item TinyML course kit
\end{itemize}

\subsection{Learning Objectives/Focus}

\begin{itemize}
  \item Discover the rich landscape of embedded ML applications
  \item Understand how to design and deploy end-to-end TinyML applications
  \item Learn to deploy TinyML models using TFLite Micro on embedded systems
\end{itemize}

\section{Welcome Message from VJ}

VIJAY JANAPA REDDI: Welcome to deploying TinyML, Course 3 of the TinyML edX
series.
I'm excited to see you back here and I hope
you had a really good time in Course 2.
As always, I'd like to start off by reminding us
why we are all excited about TinyML.
TinyML is one of the fastest growing fields of machine learning, which
brings together algorithms, hardware, and software so
that we can do on device sensor analytics
right where the data resides in real time.
If we can do this at ultra low power consumption using
small little microcontrollers, then we unlock
the ability to run machine learning models at all times.
And if so, we can run them on battery powered devices,
and this would unlock a whole new range of applications which I've been psyched
about in both Course 1 in Course 2.
And of Course 3, we're going to touch on a whole new range
of different applications.
Now that said, let's take a quick recap on where we came from
and where we're going.
In Course 1, we focused on the language of machine learning.
We exposed you to the fundamental building blocks of machine learning.
We taught you how to program in TensorFlow.
We taught you about Colab and how you can use Colab as a programming
environment free of your own resources.
In doing so we expose you to things like stochastic gradient descent,
we expose you to loss, we expose you to cost functions,
and so forth, all the building blocks you
need to have in order to build machine learning models.
Then in Course 2, we got excited about different kinds
of applications of TinyML.
And then we use the knowledge that we got from Course 1 in Course two,
specifically we learned how to take the models that we trained in TensorFlow
and shrink them down using the TF Lite converter so that we
can have quantized models that are more easier
to run on small embedded devices.
In doing that, we also learned the art of training
neural networks faster for instance using things like transfer learning.
We did all of this in the context of three different applications ranging
from keyword spotting, to visual wake words,
to anomaly detection, which was all about unsupervised learning.
So if you take a step back and you kind of
reflect on what you have covered just between Course 1 and Course 2,
it's a wealth of knowledge surrounding machine learning.
Now finally in Course 3, this is where the pedal meets the metal my friend.
This is where we're going to take those small compressed
models that we had trained and learn how to deploy them onto a microcontroller.
One of the key things that you will learn here
is that simply having a small little neural network
is not enough in terms of deploying it.
You actually have to understand the art of packaging up
the entire application as a whole binary and being able to push that down.
And when you're running it, you need to really think
about what it means to invoke a model.
So we'll talk about things like pre processing, post processing,
and of course we'll also learn how to write code in TensorFlow Lite
Micro, which is specifically designed for microcontrollers
so that you can actually take a model, a TF Lite model,
and then be able to run it effectively on a microcontroller.
So one of the core things that we really want
to get through by the end of this course is teach you
how to write an end to end ML application.
As I've said before in Course 2, I do not
care if you know how to write a model.
A model alone is a part of the bigger picture.
It really is about orchestrating the entire flow, which
includes the neural network computation as well as all the non neural network
computation.
That is what a real application is, and that
is really going to be the big takeaway that I'm going to push hard with you
in this particular course.
And that my friend, is exactly why when you finish this course,
you can call yourself a TinyML engineer.
Because you will actually have the hands on knowledge
of working with the physical microcontroller
and learning how to program it with a production grade inference framework,
which is TensorFlow Lite Micro.
It has got new APIs, and it's about the state of the art inference
engine that is out there.
So you will be way ahead of the curve.
So stick with me and let's get through Course 3.

\section{Course 1 \& 2 Recap}

\subsection{TinyML1 and TinyML2 Recap}

\includegraphics[width=\textwidth]{TinyML/3 Development/img1-re}

Over the first two courses, you have been exposed to various facets of the machine learning ecosystem. You have seen how models are developed using the machine learning workflow, from the data collection phase through to the training and deployment phases. You have seen how models are produced using Tensorflow and then ported into a lighter-weight framework such as Tensorflow Lite. You have also seen how these models can interact with various forms of input data, such as time-series data from sensors and microphones, to images from an onboard camera. Throughout both courses, we also made sure to keep in mind the core principles of responsible AI. We talked you through how to design products with all users in mind. We discussed various error types that (can) occur when we are not careful in designing and developing AI models. We explored these concepts using various real-world case studies. We then considered how to ensure that models and datasets are designed with the goals of maximizing fairness and minimizing bias in mind. Finally, we explored Google’s What-If Tool which allowed us to explore and analyze fairness and bias issues in real datasets.

In this course, you will go one step further still, and learn to deploy TinyML models responsibly on your own embedded system. Let’s first revise some of the key concepts in more detail.

\subsection{Course 1: Fundamentals of TinyML}

Before we could understand TinyML, we needed to understand its underpinnings, namely, machine learning. The utility of TinyML only becomes apparent once we understand the machine learning workflow, the different supervised learning algorithms, how they are structured, and the tasks they are able to solve (and what data they act on). To this end, the first course provided a crash course in these topics, showing that models can be trained on existing data, allowing the model to learn important associations between its parameters and the data. These trained models could then be exposed to new information, and make predictions about the data based on these associations. The two main tasks that these models perform are called regression and classification. Regression algorithms try to approximate the mapping function (f) from the input variables (x) to numerical or continuous output variables (y), whereas classification algorithms try to approximate the mapping function from the input variables to different discrete or categorical output variables.

\subsection{Machine Learning}

We have seen that machine learning entails the development of data-driven models to make predictions or decisions for a particular task. Supervised learning takes labeled data, meaning that both the input and output variables are known, and attempts to produce a model that can manipulate the input variables in such a way that the output variable can be effectively predicted.

Depending on the curated dataset, there may be a large number of input variables or very few; there may be very many data points for us to train our model with or only a small number; we may need to preprocess the data in order to extract more relevant feature information, known as feature engineering, or simply leave the input data unperturbed. Thus, the first important steps are dataset collection and dataset preprocessing. The collection stage is often the most time-consuming unless an existing dataset is available. Determining how much data is necessary to perform a robust analysis, the best method to obtain the data (e.g., crowd-sourcing, external data providers, simulated data), as well as the number of features and their relevance, are challenging steps that will often require multiple iterations to get right.

Once our dataset is curated and preprocessed, we are ready to begin developing and training machine learning models. The training process is where a model tunes its parameters based on associations learned from a particular dataset. Different tasks lend themselves better to certain machine learning methods. For example, convolutional neural networks often work best for image data - a standard fully connected neural network can also be used, but will likely perform worse at the same task. Usually, models are selected based purely on their performance, but other metrics may also be important, such as interpretability and model complexity.

Different machine learning algorithms have different hyperparameters, pre-specified values which characterize the configuration of the algorithm. During the training of a machine learning model, often through the use of a stochastic gradient descent algorithm on a chosen loss function, these parameters are often altered to find the values that lead to the highest-performing model. This process is known as hyperparameter optimization. Some platforms, known as AutoML tools, have automated this activity, making it easier for the user to optimize their models during the training process (e.g., Google’s Cloud AutoML). The optimization process performs the model evaluation, which assesses the performance of a model using metrics of interest, such as false-positive rate and F1 score. Cross-validation is often used to improve the robustness of the algorithm by preventing overfitting from occurring.

Recall that once our model has been suitably trained and optimized, we test our model on a hold-out dataset known as the test set. This simulates an unseen dataset, like our model would see in the production environment. If our model works on this set of data, it should work in our production environment! After a positive result here, we are then ready to deploy our model, which, depending on the application, may involve integration into existing infrastructure. Models are often monitored after deployment to ensure that they are functioning as expected. 

\subsection{Embedded Systems}

The main embedded system used in this course is the Arduino Nano 33 BLE Sense, which is a relatively small microcontroller equipped with a 64 MHz processor, 1 MB of flash memory, and 256 KB of SRAM. While this system is highly resource-constrained, it is low cost and can be purchased for about \$30. Clearly, this system is very different from the device you are using to view this reading. Laptops and smartphones have processors in the range of GHz, which is hundreds, if not thousands of times faster than the Arduino Nano. Furthermore, the Nano does not have any input or output peripherals that we are mostly accustomed to, such as a trackpad, keyboard, and screen. Finally, the small amounts of program memory and RAM available mean that the system has very little capacity to perform complex tasks. In fact, our particular embedded system is designed to only run one task at a time, as compared to your laptop which probably currently has multiple tabs of Chrome open!

However, we have seen that there are multiple benefits despite the severe resource constraints. Since there is no (or a very limited and lightweight) operating system, we have minimal computational overhead, which improves inference speed (i.e., how many times our model can run per second) and reduces system latency (i.e., communication delay between sending and receiving data). 

Privacy is another important benefit. Embedded systems that can perform machine learning do not need to communicate data, which reduces network loads and also prevents the possibility of man-in-the-middle attacks of transmitted data. Perhaps more important, only localized data is stored on the devices, meaning that there is no risk of data being stolen from a central repository.

Perhaps the greatest benefit of embedded systems is their small power consumption. Modern devices such as laptops and smartphones consume a great deal of power, and running large machine learning algorithms often has a high computational cost. In comparison, embedded systems use very little power, typically on the order of mW. This is no small feat since such devices can run on a simple coin cell battery for an entire year without needing to be recharged - just imagine not having to charge your laptop for a year!
  
\includegraphics[width=\textwidth]{TinyML/3 Development/img2-re}  

\subsection{Course 2: Applications of TinyML}

After learning the ropes of machine learning and motivating the utility of performing machine learning on embedded systems, we started to look at some archetypal examples of TinyML applications and some of the specifics of how models are ported from frameworks such as TensorFlow to embedded systems. We studied all of this in the context of the machine learning flow moving from collecting and preprocessing data to designing, training, evaluating, and converting models. We’ll cover deployment in this course!

\includegraphics[width=\textwidth]{TinyML/3 Development/img3-re}  

The first example we looked at was keyword spotting, which involved extracting the presence of specific keywords from a short voice recording. This example already exists in smartphones such as with Apple’s ``Hey Siri'' and Google’s ``OK Google''. We performed feature extraction of the voice recording by using spectrograms, which were then used to train our model.



\includegraphics[width=\textwidth]{TinyML/3 Development/img4-re}  

Following feature extraction, we performed post-training quantization of model weights and inference calculations to allow our model to run using the 8-bit arithmetic available on most embedded systems¹. In the penultimate stage, we took the quantized model and converted it to a more suitable file format, which, in this case, was a TFLite model. After converting this model to a binarized format, the model was then ready to be deployed in an embedded system.

We then went on to look at further examples, such as visual wake words. In the visual wake words application, we trained a model to determine the presence of a person in an image and saw how transfer learning could be used to retrain models for similarly related applications, such as detecting the presence of masks on a person’s face.

In the final parts of Course 2, we looked at more advanced industry-oriented applications that began to move towards unsupervised learning algorithms, such as anomaly detection.

We have already done a considerable amount in these first two courses. In Course 3, we will go one step further still, and provide you with all the tools necessary to develop, build, and troubleshoot your own TinyML systems. Let’s jump in!

%We must also note that we explored how Quantization Aware Training can reduce the errors introduced by quantization by allowing the model to adjust the weights to the quantization scheme during training.


\section{TinyML Application Deployment Preview}

VIJAY JANAPA REDDI: Hi.
In this video, I want to give you a little bit of a sneak preview into the kind of applications that we will actually be deploying in course 3. In order to do that, let me start off by reminding you about the ML workflow. If you recall from course 2, the ML workflow is all about how you collect your data, how you pre-process that data so that you can then prepare that data to be ingested by a neural network that you're designing, which you're going to train using tools like TensorFlow and then evaluate for both performance and accuracy.
And when you're happy with it, you use a converter to actually convert that model into something that can be deployed in a wide range of devices, be it a microcontroller or mobile device.
In this particular course, we'll focus a lot on the deployment aspects of things, which we are yet to cover, because, in course 1, we focused very much on helping you understand how to actually design a model and how to train a model.
We kind of gave you the pre-existing data sets that you could bootstrap from.
But, in course 2, we took it all the way from collecting your own data to actually deploying the model or coming close to deploying the model by actually doing [? tier ?] flight conversions.
In course 3, we're going to spend the entire gamut going all the way from collecting data to actually deploying the converted model onto an embedded microcontroller and actually making inferences.
And the way we're going to do that is by physically deploying the models onto an embedded microcontroller that comes coupled with a couple of different sensors that allow you to have that full experience.
Now this is a kit that I will talk about in more detail in the upcoming videos.
And one of the other things that's interesting is that this kit, while that's the hardware portion, coupled very elegantly with TensorFlow Lite for microcontrollers, which is a state-of-the-art inference framework from Google.

Now TensorFlow Lite Micro runs in a wide variety of different embedded microcontrollers.
This just happens to be one microcontroller or one embedded system that we happen to choose.
I'll give you the rationale for this particular setup later on.
But using this setup, we're going to help you cover that ML workflow.
Specifically, we want to help you understand how to build applications like keyword spotting, which have a lot of core fundamental attributes in terms of engineering a solid end-to-end pipeline.
Now you might be thinking, hey, yes, I know about keyword spotting.
I've heard about it in course 1.
And hey, in course 2, I did train a small, little with spotting model.
So hey, what more do I need to know?
I'm going to tell you that there is much more than just getting the model right.
I hinted at this in course 2.
But when it comes to actually deploying the model on device, you really have to think about where is the audio stream coming from.
It's not a file on your computer.
It literally streams in from the input sensor.
And how do you couple that input sensor with your neural network and make sure all of these things run in a timely fashion?
And there's also the aspect of post processing.
If some of these things seem new or a little complicated, yes, they should be because they're new concepts really.
So we'll talk about that.
While dealing with one sensor of audio sensor is interesting, with Visual WakeWords, we dealt with another sort of sensor which is the visual instance or the camera.
In this particular course, we will actually use a camera module in order to implement Visual WakeWords.
But hey, that's not it.
We're actually going to focus on taking the audio sensor and the visual sensor and doing sensor fusion where we can actually run the two models together so that we can actually do audio visual recognition.
In other words, how can I get the machine to respond only when it recognizes that I'm physically in its space, which is going to make the TinyML device much more intelligent?
And in showing its application, there is a set of fundamental new things or concepts that we're going to teach you about, like model multitenancy, which is very important and rarely talked about in TinyML.
But it's one of the upcoming areas that is something of quite a lot of value.

Finally, we'll expose you to this super awesome and exciting application called the Magic Wand, where you move the device in thin air.
And you can see in real time the device recognizing what you're actually doing.

For example.
Here, the machine will actually recognize that the letter is Z. This sort of gesture recognition capability has many different applications in many different areas, especially in biomedical applications.
So more on this later on.
I'm very excited to be talking to you about these different applications.
One of the things that we will do is go all the way from understanding the characteristics of the sensors and learning how to couple them onto the embedded device, which is then running the neural network and learning how to assess whether the neural network that is running on devices is actually performing well or not.
And that, my friend, is going to give you the end-to-end picture, which, as I've always maintained, is extremely important to be a very good TinyML engineer.
With that said, I'll see you in the next video where I actually introduce the TinyML kit.

\section{The TinyML Kit}

\subsection{The TinyML Kit}
If you haven’t already ordered your tinyML Kit, it is available from Arduino for \$49.99 at this link. The kit is the easiest and most reliable way to obtain all of the parts necessary for this course. If you already have your kit or are planning on purchasing one then you can move onto the next reading.

That said, we understand that for some learners it may be beneficial, or even necessary, to obtain some or all of these parts via other means. As such, the remainder of this document details a bill of materials for the components you need to complete all of the planned exercises alongside potential vendors and functional alternatives. Note that availability may depend on your location and other shipping dynamics limitations.

\begin{figure}
  \includegraphics[width=\textwidth]{TinyML/3 Development/kit1-re.png}  
  
  \caption{Image of the tinyML Kit including the box, Arduino Nano 33 BLE sense, Arduino Tiny Machine Learning Shield, Arducam OV7675, and USB microB cable}
\end{figure}

\subsection{tinyML Kit BOM}

In the picture above you can see all of the components included in the purpose-designed kit that we have developed with Arduino. Fortunately, there aren’t too many parts, to begin with, and most of them are generally accessible on a global scale, with the exception of Arduino’s Tiny Machine Learning Shield, or PCB cradle (in the middle), for the Nano 33 BLE Sense AI-enabled microcontroller board (on the left) that serves to breakout the MCU’s IO to connectors that permit easy, reliable connections with the camera module (on the right), as well as other sensing modules via connectors included preemptively for expansion to new areas of application, calling upon the Grove system for nearly plug-and-play compatibility with a range of transducers, available from SeeedStudio¹. Fortunately, while the shield can certainly make your life easier and lends a hand in creating reliable connections with off-board sensors/actuators, we can readily replace its function with other off-the-shelf components (e.g., a breadboard and some wires as listed below).

The table below is a list of parts that are comparable with suggested vendors alongside alternatives. To be clear, the edX course staff will only support the official kit we have developed in partnership with Arduino.

\begin{tabular}{l|l|l|l}
  Description & Part   & Vendors	& Alternatives \\
              & Number & &  \\ \hline
Nano 33 BLE Sense with headers 
See note (1)&	ABX00035	&Arduino, Amazon, DigiKey, Newark, RS Components, Arrow, Mouser, Verical, Farnell, element14, Arrow (cn)	& NA \\
USB microB cable, often to type A or C &	Varied &	DigiKey, Amazon, and many others	& See note (2) \\
Breadboard &	Varied	 & DigiKey, Amazon, and many others & 	See note (3) \\
Jumper wires‘Male to female' & Varied & DigiKey, Amazon, and many others & See note (4) \\
Camera sensor with breakout PCB	& OV7675  RobotShop, Botland, Uctronics, and others	& OV7670 See note (5) \\
\end{tabular}

\subsection{Notes}

\begin{enumerate}
  \item Make sure that you select the variant of the Nano 33 BLE Sense with headers, which means that the board will arrive with 0.1” pitch pins pre-soldered, enabling you to quickly seat the microcontroller board into a breadboard without additional equipment.
  \item The USB cable specification need only be data carrying and with connector type micro-B (plug) on one end — the other end will depend on the port available on your computer. Type A is the most common connector type, whereas modern laptops are increasingly featuring Type C connectors. Note that not all USB cables support data exchange and some are intended for power delivery only. Be sure to select a cable with a description that explicitly mentions data transfer or equivalent.
    \begin{itemize}
      \item Type C cables tend to be more expensive than Type A, so it may be comparable or even cheaper to source a USB A to C adaptor if your computer only features a Type C port. These adaptors are generally available.
      \item For both cables and adaptors, as applicable (see 1a), think carefully about the necessary designation of plug vs. receptacle on each end. For cables, you’ll likely want plugs on both ends. For a typical USB A to C adaptor, you’ll find a Type A receptacle that gives way to a Type C plug, to be inserted into your PC.
    \end{itemize}
  \item Be sure that the breadboard you purchase is both solderless and equal to or greater in spatial dimensions than the standard ``half'' size (5.5 $\times$ 8.5 cm), which we generally recommend for use in this context.
  \item Here, ``jumper wires'' refer to bundles of adjoined or individual lengths of wire that terminate in one of two possible forms: a ``male'' pin (M) or a ``female'' receptacle (F). To connect the camera module to our microcontroller board via the breadboard, you'll want to pick up about 25 or more (individual wires) of the M-F variety. If you intend to go on to connect other off-board sensor modules, it may be helpful to have M-M jumpers on hand. Jumper wires can be purchased in various lengths. We suggest between 3 to 6 inches, here.
  \item In searching for the OV7675 camera module, you may come across vendors selling the sensor itself, with only a ribbon cable breakout. Be sure to pick up a PCB breakout module that terminates in a 2x20 array of pins, like in the suggested link. If for some reason, you can’t obtain the OV7675 module, we have previously had success in using the related OV7670 module, but note, as above, that the edX course staff will only officially support the official kit we have developed in partnership with Arduino, including the requisite code for the OV7675 modules. Any adjustments to said code to call on the OV7670 module technically fall outside of the officially supported scope. With that said, we have noted in the course the required changes (fortunately there are very few) as we recognize that many students are finding it hard to source the OV7675.
\end{enumerate}

%You’ll want to check module voltage requirements and IO diagrams to verify Grove compatibility.

\section{TinyML Course Kit Overview}


This is a kit that we developed in conjunction with Arduino and is widely available to everyone across the world. One of the cool things about this kit is the fact that we designed this kit to be extensible.
In other words, while this kit serves as the basis for you to experiment with everything that we teach you in this course, it's also got capabilities that allow you to expand outside of this course, using the learnings that you get from this course.
Let me introduce you to the kit components so that you get familiar with what's in it.
First and foremost is the Arduino Nano 33 BLE Sense.
Now, this is the brains of your TinyML Kit.
Now, you know the Nano 33 by now because I've shown it to you a thousand times on my slides.
But hopefully, you're finally excited to hold it in your hand.
This is a small little microcontroller development board which has got all the elements of a modern embedded computing system on it.
For example, it's got a rich number of sensors.
It's got a microphone, it's got an IMU, it's got a light sensor, it's got an intensity sensor, and it's also got actuators in the form of RGB LED indicators that allow you to understand if the machine is trying to communicate something to you, according to what you've written in the code.
It's also got a communication hardware module, specifically in the form of Bluetooth, which allows you to transmit data without having a cable tethered to it, which makes it quite exciting.
And it's something we will be using for the magic wand example.
We can actually interface that Nano 33 with a camera module, as shown here.
The OV7675 camera module is a low voltage color CMOS image sensor, which has got a couple of breakout pins on it.
Now, this camera module outputs a resolution of about 0.3 megapixels, which is sufficient for the kind of things that we will be doing for our TinyML use cases.
Now, because it's got so many pins, one of the problems that we often tend to have is that students will incorrectly wire the different pins.
And that leads to a lot of problems in debugging.
In order to get rid of that problem and to make things easier, we designed the TinyML Shield.
The TinyML Shield is basically a daughter PCB board that's designed to have many IO pins that allow you to easily interface the Nano 33 and the camera module together.
But not only that, it's also got a lot of flexibility built into it.
Specifically, if you look at the bottom of the board, that's where your Nano 33 is going to easily plug in.
Now, if you look at the top of the board, that's where the pins of the OV7675 camera module get plugged in.
That's it.
You don't have to worry about any wires in order to connect those bazillion pins that are staring at you on the camera module.
No mistakes.
We'll be able to get started immediately.
Another nice thing about this particular TinyML Shield is the fact that it allows you to interface the Nano 33 BLE Sense with many other kinds of sensors.
You're not restricted to just the sensors that are on the board.
So on the right-hand side of the board, you'll find Standard Grove connectors that allow you to have serial communication with other types of sensors.
On the left-hand side of the board are also another set of Grove connectors that allow you to break out with analog and digital GPIO pins.
In other words, you will be able to use these GPIO pins to interface with LEDs, buzzers, buttons, and so forth.
Here are some examples of the kind of sensors that you might be able to interface with.
For example, you might be able to interface with a thermocouple sensor.
Or you might be able to interface with an optical rotary encoder.
Now, these are sensors that you would buy separately.
And they come with a cable that can easily plug into the sides of that board.
So therefore, the TinyML Shield allows you to explore much more than what we do in this particular class.
Finally, right on the bottom, there's a programmable button that's on the left-hand side.
And it's got a screw terminal for external power on the right.
The screw terminal allows you to easily and safely connect a battery to make the system fully independent.
In other words, you can completely disconnect all the wires and be able to interact with the device while it's still on using battery power.
In other words, you might be able to hold the device up and take a selfie just like that.
So with all that said, I hope you're excited about the TinyML Kit.
We spent a lot of time thinking about what's the most easiest and convenient way of packaging up the materials for you so that you can focus on what really matters, and not get distracted by the physical components and the mechanics that are needed to get the whole thing running.

\section{How the Course is Structured}

VIJAY JANAPA REDDI: Hi.
In this video, I want to talk to you a little bit about how
Course 3 is structured because Course 3 is different from Course 1
and Course 2 in that we actually have a hands on component.
Now, some parts are going to stay the same.
For example, I'm going to continue giving you
video lectures in order to help you realize the important concepts that you
have to keep in mind.
Then you're going to have readings.
Now you have certainly had many readings through both Course 1 and Course 2.
One of the interesting things that we're doing in Course 3
is that we're going to embed a couple of GIFs and videos
in order to spruce up the content in order
to help you navigate the programming experience and the development
side of the experience a bit more smoothly.
And then finally you're going to have labs.
Now these labs will require you to actually do the deployment cycle.
Now, in doing deployment you might run into a couple of issues.
For example, in the magic [? wand ?] ask you to move the device.
Well, it sounds simple, but might not be that easy
for you to understand exactly what I'm asking you to do.
Therefore, I've asked Brian, my TA, to actually join us in helping us out.
He's going to be doing hands-on demonstrations,
walking me through the collabs, the compilation process, as well
as the deployment process, all the way from getting started,
making sure your kit works, to actually seeing TinyML in action.
In addition to that, you're going to have quizzes, your tests,
your assignments, any discussion [? forums. ?]
My hope is that all of this material put together
will help you have a great experience with the TinyML kit
as well as building the actual applications that go onto this device.
With that said, I'm going to leave you in the good hands of Brian, who's
going to introduce himself next and help you get started with the TinyML kit.

% 1.2: Getting Started

\section{Introduction to the Lab Sections}

BRIAN PLANCER: Hi, I'm Brian Plancher, and I'm
a PhD candidate at Harvard University.
And I'll be here to help you through the lab components of this course.
We've prepared detailed readings with lots of screenshots and short video
walk-throughs to help make sure you are all set to deploy your TinyML
applications onto your Arduinos.
In this first lab section, we're going to make sure all of your software
gets installed and all of your hardware gets assembled and set up correctly.
Then we'll run a series of test programs to make sure that you are all
set for the rest of this course.
Because TensorFlow Lite for microcontrollers,
and Arduino code in general, is written in C++,
we'll start with a quick introduction to C++ for Python programmers.
Let's dive in.

\section{C++ for Python Users}

If you are comfortable with C/C++, please feel free to move on to the next reading. If you are new to C++, we hope this introductory material will be helpful for you. If you are comfortable with C/C++, please feel free to move on to the next reading. If you are new to C++, we hope this introductory material will be helpful for you.

Python, the language you have been using in all of your Colabs, is a dynamically-typed, “high-level” language that is interpreted at runtime. C++ (also written as Cpp), on the other hand, is a statically-typed, “low-level” language that is pre-compiled before running, allowing for very compact code. 

The good news is that since we are using the Arduino platform, we won’t have to deal with much of the complications of C++ as that is taken care of for us by the many libraries and board files we will be able to leverage (more on that soon). We’ll then just have to pay attention to the changes in syntax which are more cosmetic than functional. That is, all of the main loops and conditional statements (e.g., for, if-else) remain the same functionally!

In order to make sure you feel comfortable we’ve put together a set of short appendix items and have collected some other good resources to walk you through the main changes between Python and C++ for Arudino which you can find below. As you go through the course, if you have good ideas for other introductory material, changes to the current material, or additional resources, please let us know and we’d be very excited to add them to the list for future learners!

For more information, take a look at our short \href{https://github.com/tinyMLx/appendix/blob/main/CppForPythonProgrammers.md}{appendix} document that covers:

\begin{itemize}
  \item Data Types
  \item Scope, Parentheses, and Semicolons
  \item Functions
  \item Libraries, Header Files, \PYTHON{\#include}
  \item Other General Syntax Points
\end{itemize}


\bigskip

Additional resources:

\begin{itemize}
  \item \href{https://www.cplusplus.com/doc/tutorial/}{The C++ Language Tutor}
  \item \href{https://google.github.io/styleguide/cppguide.html}{The Google C++ Style Guide}
  \item \href{https://www.arduino.cc/reference/en/}{The Arduino Standard Library Language Reference}
\end{itemize}


\section{Setting up your Hardware}

This reading will walk you through setting up your hardware whether you purchased the TinyML Kit or your own off-the-shelf parts.


This document will help you set up your hardware, whether you purchase the Tiny Machine Learning Kit or your own off the shelf parts.
We've provided detailed instructions below with lots of pictures, and pinout diagrams, and everything else you're going to need.
I thought I'd also provide a visual example of setting up the kit, as well as some tips and tricks for wiring up off the shelf parts.
So let's begin with the kit.
To set up your Tiny Machine Learning Kit what you need to do is, first, take out all the parts.
So we have Arduino Nano 33 BLE.
We have our Tiny Machine Learning Shield.
We have the OV7675 Camera Module, which is a little tricky to get out of the protective plastic.
And we have a USB cable.
And since we only have four parts, the kit's actually quite simple to assemble.
So what we're going to do is, first, we're going to take the shield, and we're going to place it on a flat surface with this marking that says USB pointing forward.
We're then going to take the Arduino and make sure it's USB port is also pointing forward in the same direction of that USB marking on the shield.
And then we're going to place it -- these male headers here -- into these ports here on the shield.
And it should sit in, but we're going to have to push it the rest of the way in.
So the way to do that is to grab it by the outside and sort of gently rock it back and forth.
And you'll feel it press all the way down.
To attach the camera module, we're going to want to use a similar procedure.
We're going to want to hold it so that the camera itself, and the words OV7675 are pointing away from the Arduino on the board.
And then we can just line up the headers again into these ports here, which might take a little finagling.
And then we can use that same rocking motion to push it all the way in.
And that's it.
Your kit is fully assembled.
All we need to do then is undo the twist tie, and we can attach to the USB cable to our computer.
And so the USB cable itself will attach to the front of your Arduino.
And then the other end will just go straight into your computer.
And that's it.
Your kit is all set up and ready to go.
Now, if you're using off the shelf parts, the process is at a high level very similar.
But you're going to have to do a bit more wiring.
So setting this aside for a second.
What we suggest using is a solderless breadboard.
And the reason we suggest that is because it'll make it a little bit easier to do all of the wiring itself.
And so what you're going to want to do is you're going to want to take your Arduino.
And hopefully, you purchased it with the headers attached.
If the headers came separate, like this one here, where you can see the two headers are sort of on the side, you're going to have to find a friend with a soldering iron and get those soldered on, because the headers need to be rigidly attached for the camera to work correctly.
So assuming you have one with the headers on, you can do a similar procedure as what we did to seed the Arduino onto the shield, and seed it into the middle of this breadboard, making sure one set of the headers is on one side of this main channel and the other on the other.
And once you do that, you can now use the holes in the breadboard for all your wires.
And I would look in our document for the pinout tables and wiring diagrams to make sure you know which wire goes from which hole here on the breadboard, which connects to the pins on the Arduino, to which pin on your camera module.
And it's going to take a bit of wiring.
There's 18 wires you're going to have to do.
But after a couple of minutes of wiring, you should be all set and ready to go.
And so I'll see you all in the software install.

\subsection{Setting up the TinyML Kit}

Certainly, the easiest and most reliable way to obtain all of the parts necessary for this course is to procure the purpose-designed kit that we have developed with Arduino, available at this link for \$49.99. If you ordered the official kit, then getting setup will be quick and easy. We outline the required steps below:
    
\begin{enumerate}    
  \item Slot the Nano 33 BLE Sense board into the Tiny Machine Learning Shield.
    
  \includegraphics[width=\textwidth]{TinyML/3 Development/ss1.jpg}

  \includegraphics[width=\textwidth]{TinyML/3 Development/ss2.jpg}
    
    
  A mapping of the OV7675 pins to the Arudino Nano 33 BLE board. If we number the Arduino pins from top left (1) to bottom left (16) to bottom right (17) to top right (32). The mappings are as follows (Ardunio Pin, OV7675 Pin). 2,1; 4,7; 5,6; 6,19; 7,20; 8,4; 9,3; 14,2; 18,15; 19,14; 22,13; 23,12; 24,9; 25,11; 26,10; 28,5; 29,8; 30,16.
    
    The pinout (designation) for the Nano 33 BLE Sense development board. For more information and documentation of the board and its pinouts, please visit the official Arduino site.
    
  You’ll want to target the pair of spatially separated 1x15 female headers. Carefully align the pins of the microcontroller board with the headers below and then gently push down until the board is seated flush against the top of each header. The downward facing pins should no longer be visible. As best you can, avoid touching the components atop the board to prevent inadvertently damaging the surface mount devices. Pay attention to the orientation of the board so that the indication of the USB port on the PCB silkscreen matches the physical port on the board itself.
    
  \item Slot the OV7675 camera module into the shield using the same technique.

  \includegraphics[width=\textwidth]{TinyML/3 Development/ss3.jpg}
    
  The pinout (designation) for the Nano 33 BLE Sense development board. For more information and documentation of the board and its pinouts, please visit the official Arduino site.
    
  You’ll want to target the 2x10 female header. Carefully align the pins of the camera module with the headers and then gently push down until the board is seated flush against the top of each header. The downward facing pins should no longer be visible. As best you can, avoid touching the camera module atop the board to prevent inadvertent damage. Pay attention to the orientation of the camera module so that the camera sensor is to the right of the header array (as shown), further from the microcontroller board than the header array.
    
    Image of the shield with the Arduino Nano 33 BLE sense in place as well as the Arducam OV7675 slotted into the remaining 2x10 female header.

  \includegraphics[width=\textwidth]{TinyML/3 Development/ss4.png}
    
  \item Finally, use the provided USB cable (type-A to microB) to connect the Nano 33 BLE Sense development board to your machine. If your PC only features type-C USB ports, you will need to obtain an adaptor.
    
   \includegraphics[width=\textwidth]{TinyML/3 Development/ss5.jpg}
    
    Image of the tinyML shield with the Arduino Nano 33 BLE sense and Arducam OV7675 slotted in place and the USB cable connected to the Arduino.
    
    Note that if you are only calling upon hardware found on the Nano 33 BLE Sense development board (say the MCU and IMU), you could forgo connecting it to the Tiny Machine Learning Shield. If you need to remove either the Nano board or the camera module from the shield, grip each side of whichever board and pull back with a gentle rocking motion (back and forth) to work the pins out from the headers below.
\end{enumerate}

   \textbf{And that’s it! If you purchased the official kit you can move onto the next reading!}
   
    
\subsection{Wiring Up Off-the-Shelf Parts}

While we generally encourage all to consider Arduino’s Tiny Machine Learning kit as the primary and recommended hardware option, we understand that in certain circumstances, it may be beneficial, or even necessary, to obtain some or all of these parts via other means. If you have bought off-the-shelf components using our guide, we outline the steps you are likely to take in getting setup, below.

\begin{enumerate}
    \item Slot the Nano 33 BLE Sense board into a solderless breadboard.
    
  \includegraphics[width=\textwidth]{TinyML/3 Development/ss6.jpg}
    
    Image of the Arduino Nano 33 BLE sense next to a solderless breadboard.
    
  \includegraphics[width=\textwidth]{TinyML/3 Development/ss7.jpg}

    Image of the Arduino Nano 33 BLE sense slotted into a solderless breadboard. Note that one set of headers is to the left of the main channel that runs down the center of the board and the second set of headers is to the right of that main channel.
    
    If you aren’t familiar with how breadboards can be used as a substrate for connecting various electrical components, you can take a moment to review the mechanism, here. As shown in the pictures above, we’ve targeted the left-hand side of the breadboard and seated the microcontroller board into the solderless breadboard along its length, with each of the 1x15 arrays of male pins on opposite sides of the ‘trough’ that runs through the middle of the breadboard. Given the board’s width, one side will have three rows of receptacles beside it along the board while the other will have two. Which side is of no consequence. Note also that we’ve positioned the USB port at the end of the breadboard, to facilitate connections to your PC. So, altogether, you’ll want to carefully align the pins of the microcontroller board with the receptacles below and then gently push down until the board is seated flush against the breadboard. The downward facing pins should not be visible. Avoid touching the components atop the board to prevent inadvertently damaging the surface mount devices.
    
  \item     With female-to-male jumper wire, use the following Fritzing (wiring) diagram, pinout diagrams, and connection table to link the OV7675 camera module to the microcontroller board via the solderless breadboard.
    
    You can start by isolating (tearing away, as applicable) 20 adjoined female-to-male jumper wires as shown in the first photo above. For clarity, keep the desired 20 wires connected together, but tear away any number above this. The sequence of colors is of no consequence. Next connect the female side of this wire assembly to the male pins of the camera module in an alternating pattern. To keep track of what each color represents, it may be easiest to have the left or right-most color in your wire assembly connect to pin 1, where the next color in your sequence connects to 2, and so on. That way, the camera module’s pinout is now encoded by the sequence of colors.
  
    \includegraphics[width=\textwidth]{TinyML/3 Development/ss8.jpg}
    
    Image of a the set of jumper wires connected to the Arducam OV7675.
    
  \item Our next step will be to connect the camera module pinout (1-20) to specific pins on the microcontroller board via the solderless breadboard.
    
    Below we’ve mapped these OV7675 module pin numbers onto a fritzing (wiring) diagram for the Arduino Nano 33 BLE sense (assuming it is placed into a breadboard):
    
  \includegraphics[width=\textwidth]{TinyML/3 Development/ss9.png}
    
    The pinout (designation) for the Nano 33 BLE Sense development board. Note that 
    
    
    In case also helpful, we have also included the full pinout (designation) for the Nano 33 BLE Sense development board below as well as a table explaining the full pin connections needed from the OV7675 to the Arduino Nano 33 BLE Sense:
    
  \includegraphics[width=\textwidth]{TinyML/3 Development/ss10.png}
    
    The pinout (designation) for the Nano 33 BLE Sense development board. Note that
    
\begin{tabular}{lll}    
    Description	& Camera Module Pin	& Microcontroller Board Pin \\
    VCC / 3.3V  &	1	& 3.3V \\
    GND &	2	& GND \\
    SIOC / SCL&	3&	SCL / A5\\
    SIOD / SDA&	4&	SDA / A4\\
    VSYNC / VS&	5&	D8\\
    HREF / HS&	6&	A1\\
    PCLK&	7&	A0\\
    XCLK&	8&	D9\\
    D7&	9&	D4\\
    D6&	10&	D6\\
    D5&	11&	D5\\
    D4&	12&	D3\\
    D3&	13&	D2\\
    D2&	14&	D0/RX\\
    D1 (may be labeled D0)&	15&	D1/TX\\
    D0 (may be labeled D1)&	16&	D10\\
    NC	&17	&--\\
    NC	&18	&--\\
    PEN / RST&	19&	A2\\
    PWDN / PDN	&20&	A3\\
\end{tabular}
    
    So, for example, the black wire shown in the top right of the image of the set of wires connected to the OV7675 above, is connected to the camera module's first pin (pin \#1). We can directly see on the fritzing diagram that pin \#1 from the OV7675 should be connected to the second to the top left pin on the Arduino. We can also see that according to the table above, we need to connect that to the Arduino’s 3.3V pin, which we can find in the pinout diagram above as being the second pin on the top left.
    
    Another example: the second wire (white) should then connect to either the second pin from the bottom on the left (as we marked on the fritzing diagram), or the fourth pin from the bottom on the right hand side of the MCU board, as either spot² will provide a connection to GND.
    
    Continue down the table (or around the fritzing diagram) until all pins have been connected, outside of pins 17 and 18, which are left unconnected. We choose to include wires for these pins, even though they serve no purpose, because their presence in connecting to the camera module makes the entire cable assembly physically more robust, which is favorable in maintaining a solid connection.
    
    We want to underline that you should take your time in this process and verify each connection as you go, because most (nearly all) of the issues reported for this camera module stem from a wire that is connected to the wrong pin - perhaps only one location off or so. Precision here for the entire table is required for functionality.
    
    If you’re interested, the communication between the controller and the camera module is standardized as a Serial Camera Control Bus (SCCB) as well as a Camera Parallel Interface (CPI) or, equivalently, Digital Video Port (DVP).
    
    When you are done wiring up the camera your setup should like something like the following:

  \includegraphics[width=\textwidth]{TinyML/3 Development/ss11.jpg}
    
    The Arduino Nano 33 BLE sense on a breadboard connected to the OV7675 camera through jumper wires
    
\item  Finally, use the provided USB cable (type-A to microB) to connect the Nano 33 BLE Sense development board to your machine. If your PC    only features type-C USB ports, you will need to obtain an adaptor.
    
  \includegraphics[width=\textwidth]{TinyML/3 Development/ss12.jpg}
    
    The Arduino Nano 33 BLE sense on a breadboard connected to the OV7675 camera through jumper wires showing a USB cable attached to the Arduino.
    
    It’s probably apparent why we are so excited that Arduino has developed the Tiny Machine Learning Shield for our course kits! :)
    
    If you’ve sourced your jumper wires from SparkFun, you may be able to call on the same sequence of colors as we have above.
\end{enumerate}    

\subsection{Changes for the OV7670 Camera Module}


    While the course staff will not officially support the OV7670 camera module, it exists as a viable alternative if the OV7675 is unavailable in your region. Fortunately, there is an Arduino blog post concerning using the OV7670 module for tinyML.
    
    \textbf{IMPORTANT NOTE:} The connection table in the Arduino blog post is out of date, you should use the same connection table as for the OV7675 listed above. The one change is that the OV7670 only has 18 pins. Therefore, you need to omit the two NC pins between D0/D1 and RST/PWDN.
    
    ¹Be careful, the silkscreen label on some OV7675 camera modules mistakenly swaps D0 and D1!
  
    ²Note that GND, or ground, is the only example where a pin designation is repeated.
    
    
    \section{Hardware Setup Forum}
    
    If you are having trouble with the hardware setup please post to this forum. If you successfully set up the hardware please try to help your fellow students get their setup working!
    
\section{Setting up your Software}

In this reading, we will walk through setting up the software you’ll need for this course, the Arduino integrated development environment (IDE). We will be using the Arduino IDE to program your microcontroller.

Note: you do not need your Aduino connected to your computer for this step so you can install the software even if you do not yet have a kit!

BRIAN PLANCHER: Hi, everyone.
In this reading, we're going to get your software all set up for this course.
We're going to be working with the Arduino Integrated Development
Environment or IDE for short.
The Arduino IDE has a ton of built-in functionality,
and we'll be using a lot of it in this course.
For example, will be loading in libraries, writing code,
deploying code, monitoring output over the serial port,
and a host of other things.
By the end of this reading, your software
will be all ready to deploy some TinyML applications onto your Arduino.
Now before we dive in, I do want to note that I
will be using the latest versions of all libraries, board files,
et cetera in this reading.
However, I do imagine that over time those versions may become out of date.
We'll make sure to update the reading with the appropriate versions
that you should be using.
However, I am not sure that will be able to update the videos.
So please, please, please use whatever version is written in the reading,
and use that version when you are installing the various libraries.
Thank you.

\subsection{Introduction}

\subsubsection{What is Arduino?}

    Arduino describes itself as “an open-source electronics platform based on easy-to-use hardware and software.” In large part, this description is fitting, as the company designs and sells a collection of microcontroller development boards that simplify deployment of embedded hardware alongside a software framework that abstracts away all, but the most relevant considerations for your application. Perhaps, what’s under-represented by this description is the role that the surrounding community plays in enabling many plug-and-play experiences given the number and quality of auxiliary hardware modules, support libraries, and tutorials that have and continue to be produced within Arduino’s ecosystem.
    
    One thing we’d like to acknowledge here is that Arduino’s mission of creating easy-to-use hardware and software has the necessary tradeoff of limiting the feature set of its development environment to the essentials.
    
\subsubsection{What is an Integrated Development Environment (IDE)?}
    As is true for all forms of programming, an integrated development environment, or IDE, is an application with a feature set that facilitates software development generally or within a particular niche. The Arduino Desktop IDE is highly specific in the sense that it is intended to facilitate software development for a specific set of microcontroller boards, in C++.
    
    Within the niche of IDEs for embedded software, there is a noticeable dichotomy between light-weight applications, (like Arduino’s IDE) that minimize functionality and abstract away details in the name of simplicity and full-featured IDEs (like Keil µVision and Visual Studio Code) that exist to support industry professionals. 
    
    If you’re interested in finding a happy medium to use in future embedded projects, our staff have had good experience using VS Code with the hardware specific extension PlatformIO, that together enable nice-to-have features like line completion, reference tracking, etcetera, without all of the complexity that more advanced IDEs introduce. Having said this, this particular combination of VS Code + PlatformIO will not be officially supported within this course.
    
\subsubsection{What is the Arduino IDE?}

    As you might expect given the description above, the Arduino IDE is a light-weight development environment with features that permit you to very quickly manipulate microcontroller development boards. While there is a cloud-based offering (the Arduino Create Web Editor) as well as a so-called Arduino Pro IDE that is in development (specifically in the alpha stage at the time of writing), we are going to use the standard Arduino Desktop IDE in this course.
    
\subsubsection{Downloading and Installing the Arduino IDE}

    Note: Our staff will only officially support the the Arduino Desktop IDE v1.8.15 or later (tested up to v1.8.19). If you have an older version of the Desktop IDE, we recommend that you update your software.

\begin{enumerate}
  \item Navigate to arduino.cc and at the top of the page select Software and click Downloads.
  \item Click on the download link appropriate for your machine.
  \item There is no obligation to contribute, you can click ‘Just Download’ to proceed.
\end{enumerate}    

\bigskip

    Operating specific installation instructions vary, so if what follows isn’t self evident, you can find guidance, by OS, at the following links:

\begin{itemize}
  \item Windows
  \item Mac OS X
  \item Linux
\end{itemize}    

\section{A Quick Tour of the IDE}

  \includegraphics[width=\textwidth]{TinyML/3 Development/a1.png}


    A Screenshot of the Arduino IDE showing the menus and toolbar across the top highlighted in orange, the file tabs below that highlighted in yellow, the large code area in the middle of the screen highlighted in blue, and the console at the bottom highlighted in purple.
    
    When you first open the IDE, you will see a screen that looks something like the above screenshot.
    
    Up at the top of the IDE, you will find the menus and the toolbar, which we will explore as we work through the rest of this document and when we run the tests later in this section. 
    
    Below that, you will find the file tabs. For now, there will only be one file open that is most likely named \PYTHON{sketch\_<date>}, however, later in the course when we work through more complicated examples, you will see all of the various files that make up the project across the file tab area. In “Arduino speak,” a sketch is a simple project/application.
    
    In the middle of the screen, you will see the large code area. Each sketch can consist of many files and use many libraries, but at its core, each sketch is made of two main functions, “setup” and “loop”, which you can see pre-populated in the code area. We’ll explore what those functions are used for through the Blink example in a future video. For now, let's continue orienting ourselves with and setting up the IDE.
    
    Finally, at the bottom of the screen, you’ll find the console. This is where you will see debug and error messages that result from compiling your C++ sketch and uploading it to your Arduino.
    
\section{Installing the Board Files for the Nano 33 BLE Sense}


    0:00 / 0:00Klicken Sie auf die Pfeiltaste nach oben, um das Geschwindigkeitsmenu aufzurufen. Regeln Sie dann mit den Hoch und Runter Pfeiltasten die Geschwindigkeit. Klicken Sie Enter, um die jeweilige Geschwindigkeit festzulegen.Geschwindigkeit1.0xKlicke auf diesen Knopf und diess Video stumm oder nicht stumm zu schalten, oder drücke die HOCH oder RUNTER Knöpfe um die Lautstärke zu erhöhen oder abzusenken.Maximal Lautstärke.Videotranskript
    Beginn des Transkriptes, zum Ende springen.
    Brian Plancher: Now that you have the Arduino IDE installed,
    let's work on setting it up for the rest of this course.
    We're going to start by installing the board files we need
    for do our Arduino Nano 33 BLE Sense.
    To do that, we're going to go to the menu.
    Although, if you did plug in your Arduino,
    you might have noticed that there's this little pop-up on the bottom left that
    says, install this package to use your Arduino Nano 33 BLE board.
    By clicking that link, it'll actually get you to the same place exactly,
    where we want to go.
    But let's assume you haven't plugged in the board or you happened to X it out
    and don't know how to get that link back.
    Well, what you're going to want to do is go
    to Tools and then to Board, where you might
    see Arduino Uno already populated.
    And then, we're going to open the Boards Manager.
    Inside the Boards Manager, we're going to go up to this top
    here and search for Nano 33 BLE.
    And when we do that, up at the top, we're
    going to find this Arduino MBed-enabled boards, board file library.
    And that's what we want.
    And we want to make sure to install version--
    well, at the time of this recording, it's version 1.3.1,
    but please check the reading to make sure that hasn't changed.
    And then click Install.
    And once you click Install, you should see
    this blue bar going across the bottom.
    And it should start to move quite quickly.
    And then, as it gets to the end, it might hang for a long time.
    Don't be alarmed.
    Don't worry.
    Just let it run.
    It could take 5 minutes, 10 minutes.
    Sometimes, what will happen is that on Windows machines,
    you might need to approve an administrator pop-up to allow it
    to install some USB drivers or whatnot.
    And so just be patient with it, and let it do its thing.
    Mine seems to be hanging.
    This is a Windows machine.
    And so, yeah, just be patient.
    And eventually, it'll complete.
    And then you'll be ready to move on to the next step.
    So I'm going to hang here for a little while and wait for it to finish.
    All right, and so now, after a little bit of waiting,
    your board files will be installed.
    And you can see this teal INSTALLED in all caps,
    next to the version of your board files.
    And you should be all set.
    The last thing we're going to do before we move on is actually show you
    how to select the board.
    So if you go, again, to Tools and to Board,
    instead of using the Arduino Uno, we can go under the MBed OS boards,
    and we can select the Nano 33 BLE.
    And now, we should be all configured with the board and libraries installed
    and that board selected.
    And you can actually see, down on the bottom right, it should say,
    Arduino Nano 33 BLE.
    And so we should be ready to go for the rest of the course.
    Ende des Transkriptes, zum Anfang springen
    Downloads und Transskripte
    Video
    Videodatei herunterladen
    Transskripte
    Download SubRip (.srt) file
    Download Text (.txt) file

\section{Installing the Board Files for the Nano 33 BLE Sense}


    One of the primary advantages that the Arduino ecosystem affords is the portability of code you write for one or another board within their line-up or even in porting code to affiliate boards. This is made possible by the support files organized in the Boards Manager, which coordinates a download and installation of files that detail the Arduino functions (sometimes ‘core’) that are defined for that particular board (which is how hardware differences between boards are abstracted) as well as compiler or linker details specific to the given board.
    
    To install the board files that you will need for your Arduino Nano 33 BLE Sense, please do the following:

\begin{enumerate}
    \item 
    Open the Boards Manager, which you can find via the Tools drop-down menu. Navigate, as follows: Tools → Board → Boards Manager. Note that the Board may be set to “Arduino UNO” by default.
    
  \includegraphics[width=\textwidth]{TinyML/3 Development/a2-re.png}
    
    A screenshot of the top left corner of the Arudino IDE showing the top menu of “File, Edit, Sketch, Tools, Help” with “Tools” highlighted. This opens up a large submenu and about 2/3rds of the way down we have highlighted “Board: Arduino Uno” which has opened up another sub menu to the right and we have highlighted “Boards Manager..” and selected that.
    
\item     In the Boards Manager dialog box, use the search bar at the top right to search for “mbed nano” which should bring up a few results. We’re interested in the first result (as shown), named “Arduino MBED OS Nano Boards,”. Make sure you select Version 2.3.1 or newer (tested up to Version 3.0.1) and then click “Install.” As the install process progresses, you will see a blue completion bar work its way across the bottom of the Board Manager window. Be patient, you may need to install USB drivers, which requires you to approve an administrator privileges popup which can take a couple of minutes to appear. 

  \includegraphics[width=\textwidth]{TinyML/3 Development/mbedNanoBoards.png}


    A screenshot of the Boards manager showing an input bar across the top where we have typed in “mbed nano” which has next to it a dropdown for “type” that we have left to the default “All”. We see that two results have appeared. The first titled in board “Arduno MBed OS Nano Boards” and the second “[Deprecated - Please install standalone packages] Arduin Mbed OS Boards”. Inside each, there is a description of the package and its creator as well as links for help and more info. By selecting the first package (the non-deprecated one), on the right of that row appears a dropdown for a version (we have selected 2.3.1) as well as an Install button you'll want to click.
    
    After you have successfully installed the board, if you exit and re-open the Board Manager and search again for “mbed nano” you will now see INSTALLED next to the library and the option to “Remove” the library or install a different version.

  \includegraphics[width=\textwidth]{TinyML/3 Development/mbedNanoBoards_installed.png}
    
    A screenshot of the Boards manager after installation. Everything is as described in the previous image except now there is additional text under the title which says version 2.3.1 INSTALLED. Where “INSTALLED” is highlighted in teal.
\end{enumerate}    


\section{Installing the Libraries Needed for this Course}

    0:00 / 0:00Klicken Sie auf die Pfeiltaste nach oben, um das Geschwindigkeitsmenu aufzurufen. Regeln Sie dann mit den Hoch und Runter Pfeiltasten die Geschwindigkeit. Klicken Sie Enter, um die jeweilige Geschwindigkeit festzulegen.Geschwindigkeit1.0xKlicke auf diesen Knopf und diess Video stumm oder nicht stumm zu schalten, oder drücke die HOCH oder RUNTER Knöpfe um die Lautstärke zu erhöhen oder abzusenken.Maximal Lautstärke.Videotranskript
    Beginn des Transkriptes, zum Ende springen.
    All credits for copyrighted materials are listed at the end of this work.
    BRIAN PLANCHER: All right, and so the next thing that will need to do
    is to install all of the libraries that we'll be using in this course,
    We'll use the Arduino IDE built-in library manager to do this.
    You can find that under the tools menu by clicking manage libraries.
    This operates very similar to the boards Manager
    that we just used, and will have a similar pop-up that
    shows up in the middle of the screen with a search bar across the top.
    So, for example, in that search bar, we could type TensorFlow,
    and the first thing that will show up is the Arduino TensorFlow
    library for microcontrollers.
    And so we can then go over here and find the correct version.
    Again, please check the reading, for the most up-to-date version.
    At the time of this recording, it is 2.4.0-alpha.
    However, it may change over time.
    So again, please check the reading below for the correct version.
    Anyway, once you have the correct version, you can click install
    and it will install the library across the bottom
    with another one of those blue bars, and this
    should be pretty quick in this case.
    Yes, it's done.
    And when it's done, you'll see the installed word and field
    next to the library.
    So, please look to the reading below, and make
    sure you install all of the libraries that we'll need for this course,
    and install the correct version of those libraries Thanks.
    Ende des Transkriptes, zum Anfang springen
    Downloads und Transskripte
    Video
    Videodatei herunterladen
    Transskripte
    Download SubRip (.srt) file
    Download Text (.txt) file

\section{Installing the Libraries Needed for this Course}

    Another advantage of the Arduino ecosystem is the availability of a wide array of libraries for performing various tasks, such as interfacing with a sensor module or manipulating data using common algorithms. There are many libraries that can be accessed from within the Library Manager in the Arduino IDE as described below. Check here for a complete list.
    
    For this course, we are going to need four libraries. To install the libraries, please do the following and make sure to install the version specified in the reading below or the tinyML applications will not work:

\begin{enumerate}    
\item     Open the Library Manager, which you can find via the Tools drop-down menu. Navigate, as follows: Tools → Manage Libraries.
    
  \includegraphics[width=\textwidth]{TinyML/3 Development/a5-re.png}
    
    Screenshot of opening the library manager using the menus on the Arduino IDE
    
    
    \item Then, much like for the Boards Manager, in the Library Manager dialog box, use the search bar at the top right to search for the following libraries, one at a time. Note that like with the Board manager, a blue completion bar will appear across the bottom of the Library Manager window.
    
    \begin{enumerate}
        \item   The Tensorflow Lite Micro Library:
          \begin{itemize}
              \item     Search Term: Tensorflow
              \item Library Name: Arduino\_TensorFlowLite
              \item    Version: 2.4.0-ALPHA
              
                \includegraphics[width=\textwidth]{TinyML/3 Development/a6-re.png}
    
    Screenshot of using the library manager to search for and install the  Arduino\_TensorFlowLite library.

          \end{itemize}
    
   \item  The \PYTHON{Harvard\_TinyMLx} Library we put together for this course!
   
     \begin{itemize}
       \item Search Term: TinyMLx
       \item Library Name: Harvard\_TinyMLx
       \item Version: 1.1.0-Alpha
     \end{itemize}
            
\end{enumerate}
\end{enumerate} 

\subsection{Preparing for Deployment}

Use a USB cable to connect the Arduino Nano 33 BLE Sense to your machine. You should see a green LED power indicator come on when the board first receives power.
Open the Blink.ino sketch, which you can find via the File drop-down menu. Navigate, as follows: File → Examples → 01.Basics → Blink. You’ll notice that Arduino has provided a wealth of examples to choose from should you like to explore the board more on your own outside of the course material. There is great documentation about those examples on the Arduino website.

A screenshot of the top left of the Arduino IDE showing the navigation to the blink example which opens up a series of submenus below and then to the right of the choices. The correct path is as follows: File (first option on the top menu) → Examples (about half way down) → 01.Basics (first option) → Blink (third option). You’ll notice that Arduino has provided a wealth of additional examples to choose from should you like to explore the board more on your own outside of the course material. Click the Blink example to open up the correct code in the IDE. Do note that Arduino sometimes changes the ordering of menu items during updates to the IDE so please follow the words and do not be alarmed if they are in slightly different locations.
Use the Tools drop-down menu to select appropriate Port and Board. This is important as it is telling the IDE which board files to use and on which serial connection it should send the code. In some cases, this may happen automatically, but if not, you’ll want to select:
Select the Arduino Nano 33 BLE as the board by going to Tools → Board: <Current Board Name> → Arduino Mbed OS Boards (nRF52840) → Arduino Nano 33 BLE. Note that on different operating systems, the exact name of the board may vary but/and, it should include the word Nano at a minimum. If you do not see that as an option, then please go back to Setting up the Software and make sure you have installed the necessary board files.
Another screenshot of the IDE showing another set of menus and submenus that need to be traversed to select the correct board. The correct path is: Tools (4th option) → Board: <Current Board Name> (about 2/3rds of the way down)→ Arduino Mbed OS Boards (nRF52840) (3rd option)→ Arduino Nano 33 BLE. (3rd option). Do note that Arduino sometimes changes the ordering of menu items during updates to the IDE so please follow the words and do not be alarmed if they are in slightly different locations.

Then, select the USB Port associated with your board. This will appear differently on Windows, macOS, Linux, but will likely indicate ‘Arduino Nano 33 BLE” in parenthesis. You can select this by going to Tools → Port: <Current Port (Board on Port)> → <TBD Based on OS> (Arduino Nano 33 BLE). Where <TBD Based on OS> is most likely to come from the list below where <\#> indicates some integer number.

\begin{itemize}
  \item Windows → COM<\#>
  \item macOS → /dev/cu.usbmodem<\#>
  \item Linux → ttyUSB<\#> or ttyACM<\#>
\end{itemize}

Another screenshot of the IDE showing another set of menus and submenus that need to be traversed to select the correct port. The correct path is Tools (4th option) → Port: <Current Port (Board on Port)> (about 3/4 of the way down)→ <TBD Based on OS> (Arduino Nano 33 BLE) (location depends on your machine). Do note that Arduino sometimes changes the ordering of menu items during updates to the IDE so please follow the words and do not be alarmed if they are in slightly different locations.
Finally, use the checkmark button at the top left of the UI to verify that the code within the example sketch is valid. 

A screenshot of the top left of the Arduino IDE showing the verify checkmark right below the top menu items all the way on the left.
Verification will compile the code, so take note of the status and results indicated in the black console at the bottom of the IDE. The level of detail presented here will depend on whether or not you have enabled ‘verbose output during compilation’ in Preferences. You should most likely at the end see a final output indicating how much memory the sketch will take on the Arduino once it is uploaded. Something like: “Sketch uses 86568 bytes (8\%) of program storage space. Maximum is 983040 bytes. Global variables use 44696 bytes (17\%) of dynamic memory, leaving 217448 bytes for local variables. Maximum is 262144 bytes.” As you can see, this is a very simple example and does not take up much space. While, as you’ll see in a moment, uploading your code will also verify your code automatically, it is often helpful to verify your code first as you can iron out any compilation errors without having any hardware on hand.

\subsection{Deploying (Uploading) the Sketch}
    
Once we know that the code at hand is valid, we can ‘flash’ it to the MCU:

Use the rightward arrow next to the ‘compile’ checkmark to upload / flash the code.
Note that pragmatically, this step will re-compile the sketch before flashing the code, so that in the future if you intend to sequentially compile and flash a program, you need to only press the ‘upload’ arrow.

A screenshot of the top left of the Arduino IDE showing the download right arrow right below the top menu items second from the left and one to the right of the verify checkmark we used in the last step.

As before, take note of the status and results indicated in the black console at the bottom of the IDE. The level of detail presented here will depend on whether or not you have enabled ‘verbose output during compilation’ in Preferences, accessible via the File drop-down menu in the IDE.

You’ll know the upload is complete when you see red text in the console at the bottom of the IDE that shows 100\% upload of the code and a statement that says something like “Done in <\#.\#> seconds.” Again, if this is the first time you are uploading a sketch to an Arduino, the upload may hang for a little while until you get another administrator approval popup and approve it. Don’t worry, this is just a one time thing.

If you receive an error, you will see an orange error bar appear and a red error message in the console (as shown below). Don’t worry -- there are many common reasons this may have occurred. To help you debug, please check out our FAQ appendix with answers to the most common errors!

A screenshot of the Arudino IDE showing the blink sketch code at the top and the terminal output at the bottom. The terminal has an orange highlight that says “An error occuried while uploading the sketch” and then has a series of somewhat cryptic error messages in the text. As noted in the text if you see this orange error bar -- check out FAQ appendix first as we have documented the common errors. https://github.com/tinyMLx/appendix/blob/main/ArduinoFAQ.md
At this point, you’ll want to look to the board itself. The orange LED opposite the green LED power indicator about the USB port should now be blinking!

A GIF of the tinyML kit running the blink sketch. We can see that the orange LED opposite the green LED power indicator about the USB port is blinking.
As a final note, if you’d like to learn more about how the process of flashing code to a microcontroller works, please check out this appendix document.

\subsection{Understanding the Code in the Blink Example}

0:00 / 0:00Klicken Sie auf die Pfeiltaste nach oben, um das Geschwindigkeitsmenu aufzurufen. Regeln Sie dann mit den Hoch und Runter Pfeiltasten die Geschwindigkeit. Klicken Sie Enter, um die jeweilige Geschwindigkeit festzulegen.Geschwindigkeit1.0xKlicke auf diesen Knopf und diess Video stumm oder nicht stumm zu schalten, oder drücke die HOCH oder RUNTER Knöpfe um die Lautstärke zu erhöhen oder abzusenken.Maximal Lautstärke.Videotranskript
Beginn des Transkriptes, zum Ende springen.
BRIAN PLANCHER: Let's now take a look at the code behind the [? blink ?]
example.
If you look, it's actually mostly a comment.
So let's up that text size, and just scroll down to the bottom here.
And what we'll find is that the Arduino blink example is pretty much just two
very short functions.
The setup function, and the loop function.
And it turns out that any Arduino sketch boils down to a setup and a loop,
whether it's this simple blink example, or a full keyword spotting application.
In the setup function, we, well, set things up.
You can imagine in a keyword spotting application, that
could be as complicated as turning on the microphone
and initializing a neural network.
Here, it's as simple as just telling the device that we're
going to use the LED as an output.
And then in the loop function, we're going
to run whatever we need to run over and over and over again.
So in the keyword spotting example, that's
recording audio, doing preprocessing, passing it through the neural network,
getting the output, and then doing whatever
we want to do with that output.
In this case, it's just turning the LED on, waiting a second,
and then turning it off.
And in general, it's just to remember that every Arduino sketch is
going to do some amount of setup, and then it's
going to do something in a loop, as fast as it can, over and over again.
And so you just need to get into a mindset of like,
that's how your code will work.
It's, set something up once, and then infinite loop forever.
So hopefully that explains the code a little bit,
and I'll see you all in the next test.
Ende des Transkriptes, zum Anfang springen
Downloads und Transskripte
Video
Videodatei herunterladen
Transskripte
Download SubRip (.srt) file
Download Text (.txt) file
Understanding the Code in the Blink Example
Now that you have gotten the blink example deployed to your microcontroller, let's explore the code as shown below.

You’ll notice that it consists of two functions: setup, and loop. As we mentioned before, this is the standard setup for an Arduino sketch. This is because when the Arduino turns on it runs the setup() function ONCE to initialize (aka setup) the sketch. Then it runs the loop() function infinitely many times (aka it runs as an infinite loop) to execute the sketch. This works well for most tinyML applications as they are designed to respond to continuous sensor input. You can imagine in the case of Keyword spotting that we need to initialize the neural network and the microphone and then in a loop we want to listen to audio and trigger (or not) depending upon the output of the neural network!

\begin{lstlisting}
// the setup function runs once when you press reset or power the board
void setup(){
    // initialize digital pin LED_BUILTIN as an output.
    pinMode(LED_BUILTIN, OUTPUT);
}

// the loop function runs over and over again forever
void loop(){
    digitalWrite(LED_BUILTIN, HIGH); // turn the LED on (HIGH is the voltage level)
    delay(1000);                     // wait for a second
    digitalWrite(LED_BUILTIN, LOW);  // turn the LED off by making the voltage LOW
    delay(1000);                     // wait for a second
}
\end{lstlisting}

You’ll also notice that both functions have the void return type as they do not ever return anything but instead have side effects.

For example, in the setup function the \PYTHON{LED\_BUILTIN} (which is a shortcut name for the pin that controls the voltage to the LED) is set to be an output for the duration of the loop. In general, you will need to set all of the pins you use as either inputs or outputs during the setup function. If you wired up the camera yourself, you have already explored a lot of the special names reserved for the pins as shown on the pinout diagram.

In the loop function, you’ll notice that we are alternating between writing a HIGH (aka turning on the LED) and writing a LOW (aka turning off the LED). The delay of 1000 milliseconds (1 second) between each step is crucial as otherwise the light would turn on and off so fast that it would be imperceptible. In fact, if you make the delay too short the light will simply seem to be dim. This is a trick called Pulse Width Modulation that is actually used often in industry to e.g., control motors. If you’d like, feel free to experiment with modifying these delays and redeploying the code to see its effect!            
    \PYTHON{TinyMLx} Library selection in the Library Manager
    \begin{itemize}

   \item The library that supports the accelerometer, magnetometer, and gyroscope on the Nano 33 BLE sense:
   
   \begin{itemize}
      \item Search Term: LSM9DS1
      \item Library Name: Arduino\_LSM9DS1
      \item    Version: 1.1.0

       \includegraphics[width=\textwidth]{TinyML/3 Development/a8-re.png}
    
    Screenshot of using the library manager to search for and install the \PYTHON{Arduino\_LSM9DS1} library.
      \end{itemize}
   
   \item     ArduinoBLE:
     \begin{itemize}
       \item Search Term: ArduinoBLE
       \item Library Name: ArduinoBLE
       \item Version: 1.2.1 or newer (tested up to 1.2.2)
       
       \includegraphics[width=\textwidth]{TinyML/3 Development/ArduinoBLE.png}
       
    ArduinoBLE Library selection in the Library Manager
    
     \end{itemize}
 \end{itemize}

\section{Setting your Preferences}

    You can adjust the preferences set for the Arduino Desktop IDE via the File drop-down menu, File → Preferences. There are a few preferences that we recommend enabling to make the Arduino IDE a little easier to use, namely:

\begin{enumerate}
  \item Show verbose output during: compilation and upload.
  \item Enable code folding.
  \item Display line numbers.
\end{enumerate}    
    
       \includegraphics[width=\textwidth]{TinyML/3 Development/a10.png}
    
    On final note, if you don’t like the default theme for the Arduino Desktop IDE, there is a nice tutorial for a dark theme you can find here. Also, if you would like to learn more about the IDE, check out Arduino's documentation.
    
    And that’s it! Your Arduino IDE should be all configured for this course. Now that you have all of the necessary board files and libraries installed, it’s time to explore more of the features of the IDE available under the “Tools” menu and start to test out your Arduino by deploying the Blink example!
    
    
\section{The Arduino Blink Example}

    In this reading, we will deploy the Arduino Blink example to make sure everything is working properly and to give you your first experience deploying code to your Arduino!
    
Brian Plancher: All right, now it's time to begin a series of tests
to make sure that everything is working with your hardware and software setup.
We'll start with the Arduino Blink example.
That example does kind of what it sounds like it would do.
It blinks the lead in the top-left corner of your Arduino.
Mine is currently just off.
So to get to that example, we're going to use the File menu.
Under File, we're going to go to Examples.
Then, under Built-in Examples, we're going to select Basics.
And under Basics, we're going to select Blink.
And this will open up a new Arduino IDE window for the blink
example, which we can maximize.
We'll talk about the code in a minute.
But first, let's just deploy it to your Arduino
to make sure everything's working.
To do that, we're going to need to go to the Tools menu
and make sure that the board says Arduino Nano 33 BLE.
If it doesn't, we can get to that by going over to Mbed OS boards
and then selecting the correct board.
We also need to make sure the port is selected.
Currently, mine is empty.
Now, when you go to Port, you might actually
see a couple of different things here, and exactly what you see
will depend on what operating system you're using.
We've described that in detail below, so please check the reading
depending on your operating system.
I'm using Windows, so usually, the serial ports show up as COM something.
And the important one we're looking for is
we're looking for the port that has Arduino Nano 33
BLE in parentheses next to it.
Again, you might have multiple ports there,
so please select the one that says Arduino.
Once you do that, you can just click the Upload button on the top over here,
and that will start to compile the sketch.
You can see this compiling happening over here.
And you can see the debug output in your console
if you selected your preferences to have the verbose debug.
If you didn't, you might not see as much.
Once it's done compiling, it will start to upload the code to our Arduino.
And at that point, the text will change from white to orange.
And you can see that here, as we're having our percentage uploading going.
And there, we have done in 3.642 seconds.
And if we go over and look at our Arduino,
we can see that that top-left light is blinking on and off.
And that's exactly what we want to see.
You've now deployed your first Arduino sketch to your Arduino.
And you know that, at least at that level, that's working.
And so I'll see you in the next test.

\subsection{Preparing for Deployment}

\begin{enumerate}
  \item Use a USB cable to connect the Arduino Nano 33 BLE Sense to your machine. You should see a green LED power indicator come on when the board first receives power.
  \item Open the Blink.ino sketch, which you can find via the File drop-down menu. Navigate, as follows: File → Examples → 01.Basics → Blink. You’ll notice that Arduino has provided a wealth of examples to choose from should you like to explore the board more on your own outside of the course material. There is great documentation about those examples on the Arduino website.

       \includegraphics[width=\textwidth]{TinyML/3 Development/b1-re.png}

  A screenshot of the top left of the Arduino IDE showing the navigation to the blink example which opens up a series of submenus below and then to the right of the choices. The correct path is as follows: File (first option on the top menu) → Examples (about half way down) → 01.Basics (first option) → Blink (third option). You’ll notice that Arduino has provided a wealth of additional examples to choose from should you like to explore the board more on your own outside of the course material. Click the Blink example to open up the correct code in the IDE. Do note that Arduino sometimes changes the ordering of menu items during updates to the IDE so please follow the words and do not be alarmed if they are in slightly different locations.

  \item Use the Tools drop-down menu to select appropriate Port and Board. This is important as it is telling the IDE which board files to use and on which serial connection it should send the code. In some cases, this may happen automatically, but if not, you’ll want to select:
  
  \begin{itemize}
      \item Select the Arduino Nano 33 BLE as the board by going to Tools → Board: <Current Board Name> → Arduino Mbed OS Boards (nRF52840) → Arduino Nano 33 BLE. Note that on different operating systems, the exact name of the board may vary but/and, it should include the word Nano at a minimum. If you do not see that as an option, then please go back to Setting up the Software and make sure you have installed the necessary board files.

       \includegraphics[width=\textwidth]{TinyML/3 Development/b2-re.png}

       Another screenshot of the IDE showing another set of menus and submenus that need to be traversed to select the correct board. The correct path is: Tools (4th option) → Board: <Current Board Name> (about 2/3rds of the way down)→ Arduino Mbed OS Boards (nRF52840) (3rd option)→ Arduino Nano 33 BLE. (3rd option). Do note that Arduino sometimes changes the ordering of menu items during updates to the IDE so please follow the words and do not be alarmed if they are in slightly different locations.

    \item Then, select the USB Port associated with your board. This will appear differently on Windows, macOS, Linux, but will likely indicate ‘Arduino Nano 33 BLE” in parenthesis. You can select this by going to Tools → Port: <Current Port (Board on Port)> → <TBD Based on OS> (Arduino Nano 33 BLE). Where <TBD Based on OS> is most likely to come from the list below where <\#> indicates some integer number.

      \begin{itemize}
        \item Windows $\rightarrow$ COM<\#>
        \item macOS $\rightarrow$ /dev/cu.usbmodem<\#>
        \item Linux $\rightarrow$ ttyUSB<\#> or ttyACM<\#>

       \includegraphics[width=\textwidth]{TinyML/3 Development/b3-re.png}

       Another screenshot of the IDE showing another set of menus and submenus that need to be traversed to select the correct port. The correct path is Tools (4th option) → Port: <Current Port (Board on Port)> (about 3/4 of the way down)→ <TBD Based on OS> (Arduino Nano 33 BLE) (location depends on your machine). Do note that Arduino sometimes changes the ordering of menu items during updates to the IDE so please follow the words and do not be alarmed if they are in slightly different locations.
    \end{itemize}
  \end{itemize}

  \item Finally, use the checkmark button at the top left of the UI to verify that the code within the example sketch is valid. 

       \includegraphics[width=\textwidth]{TinyML/3 Development/b4-edited.png}

   A screenshot of the top left of the Arduino IDE showing the verify checkmark right below the top menu items all the way on the left.

   Verification will compile the code, so take note of the status and results indicated in the black console at the bottom of the IDE. The level of detail presented here will depend on whether or not you have enabled ‘verbose output during compilation’ in Preferences. You should most likely at the end see a final output indicating how much memory the sketch will take on the Arduino once it is uploaded. Something like: “Sketch uses 86568 bytes (8\%) of program storage space. Maximum is 983040 bytes. Global variables use 44696 bytes (17\%) of dynamic memory, leaving 217448 bytes for local variables. Maximum is 262144 bytes.” As you can see, this is a very simple example and does not take up much space. While, as you’ll see in a moment, uploading your code will also verify your code automatically, it is often helpful to verify your code first as you can iron out any compilation errors without having any hardware on hand.

\end{enumerate}

\subsection{Deploying (Uploading) the Sketch}

Once we know that the code at hand is valid, we can ‘flash’ it to the MCU:

  \begin{enumerate}
    \item Use the rightward arrow next to the ‘compile’ checkmark to upload / flash the code.

           Note that pragmatically, this step will re-compile the sketch before flashing the code, so that in the future if you intend to sequentially compile and flash a program, you need to only press the ‘upload’ arrow.

       \includegraphics[width=\textwidth]{TinyML/3 Development/b5-edited.png}

      A screenshot of the top left of the Arduino IDE showing the download right arrow right below the top menu items second from the left and one to the right of the verify checkmark we used in the last step.

As before, take note of the status and results indicated in the black console at the bottom of the IDE. The level of detail presented here will depend on whether or not you have enabled ‘verbose output during compilation’ in Preferences, accessible via the File drop-down menu in the IDE.

You’ll know the upload is complete when you see red text in the console at the bottom of the IDE that shows 100\% upload of the code and a statement that says something like “Done in <\#.\#> seconds.” Again, if this is the first time you are uploading a sketch to an Arduino, the upload may hang for a little while until you get another administrator approval popup and approve it. Don’t worry, this is just a one time thing.

If you receive an error, you will see an orange error bar appear and a red error message in the console (as shown below). Don’t worry -- there are many common reasons this may have occurred. To help you debug, please check out our FAQ appendix with answers to the most common errors!

       \includegraphics[width=\textwidth]{TinyML/3 Development/b6-re.png}

  A screenshot of the Arudino IDE showing the blink sketch code at the top and the terminal output at the bottom. The terminal has an orange highlight that says “An error occuried while uploading the sketch” and then has a series of somewhat cryptic error messages in the text. As noted in the text if you see this orange error bar -- check out FAQ appendix first as we have documented the common errors. \url{https://github.com/tinyMLx/appendix/blob/main/ArduinoFAQ.md}
  
  \item At this point, you’ll want to look to the board itself. The orange LED opposite the green LED power indicator about the USB port should now be blinking!

       \includegraphics[width=\textwidth]{TinyML/3 Development/GIF1}% .gif}

A GIF of the tinyML kit running the blink sketch. We can see that the orange LED opposite the green LED power indicator about the USB port is blinking.

\end{enumerate}

As a final note, if you’d like to learn more about how the process of flashing code to a microcontroller works, please check out this appendix document.


    
\subsection{Understanding the Code in the Blink Example}


0:00 / 0:00Klicken Sie auf die Pfeiltaste nach oben, um das Geschwindigkeitsmenu aufzurufen. Regeln Sie dann mit den Hoch und Runter Pfeiltasten die Geschwindigkeit. Klicken Sie Enter, um die jeweilige Geschwindigkeit festzulegen.Geschwindigkeit1.0xKlicke auf diesen Knopf und diess Video stumm oder nicht stumm zu schalten, oder drücke die HOCH oder RUNTER Knöpfe um die Lautstärke zu erhöhen oder abzusenken.Maximal Lautstärke.Videotranskript
Beginn des Transkriptes, zum Ende springen.
BRIAN PLANCHER: Let's now take a look at the code behind the [? blink ?]
example.
If you look, it's actually mostly a comment.
So let's up that text size, and just scroll down to the bottom here.
And what we'll find is that the Arduino blink example is pretty much just two
very short functions.
The setup function, and the loop function.
And it turns out that any Arduino sketch boils down to a setup and a loop,
whether it's this simple blink example, or a full keyword spotting application.
In the setup function, we, well, set things up.
You can imagine in a keyword spotting application, that
could be as complicated as turning on the microphone
and initializing a neural network.
Here, it's as simple as just telling the device that we're
going to use the LED as an output.
And then in the loop function, we're going
to run whatever we need to run over and over and over again.
So in the keyword spotting example, that's
recording audio, doing preprocessing, passing it through the neural network,
getting the output, and then doing whatever
we want to do with that output.
In this case, it's just turning the LED on, waiting a second,
and then turning it off.
And in general, it's just to remember that every Arduino sketch is
going to do some amount of setup, and then it's
going to do something in a loop, as fast as it can, over and over again.
And so you just need to get into a mindset of like,
that's how your code will work.
It's, set something up once, and then infinite loop forever.
So hopefully that explains the code a little bit,
and I'll see you all in the next test.

Ende des Transkriptes, zum Anfang springen
Downloads und Transskripte
Video
Videodatei herunterladen
Transskripte
Download SubRip (.srt) file
Download Text (.txt) file

\subsection{Understanding the Code in the Blink Example}

Now that you have gotten the blink example deployed to your microcontroller, let's explore the code as shown below.

You’ll notice that it consists of two functions: setup, and loop. As we mentioned before, this is the standard setup for an Arduino sketch. This is because when the Arduino turns on it runs the setup() function ONCE to initialize (aka setup) the sketch. Then it runs the loop() function infinitely many times (aka it runs as an infinite loop) to execute the sketch. This works well for most tinyML applications as they are designed to respond to continuous sensor input. You can imagine in the case of Keyword spotting that we need to initialize the neural network and the microphone and then in a loop we want to listen to audio and trigger (or not) depending upon the output of the neural network!


\begin{lstlisting}
// the setup function runs once when you press reset or power the board
void setup(){
    // initialize digital pin LED_BUILTIN as an output.
    pinMode(LED_BUILTIN, OUTPUT);
}

// the loop function runs over and over again forever
void loop(){
    digitalWrite(LED_BUILTIN, HIGH); // turn the LED on (HIGH is the voltage level)
    delay(1000);                     // wait for a second
    digitalWrite(LED_BUILTIN, LOW);  // turn the LED off by making the voltage LOW
    delay(1000);                     // wait for a second
}
\end{lstlisting}

You’ll also notice that both functions have the void return type as they do not ever return anything but instead have side effects.

For example, in the setup function the \PYTHON{LED\_BUILTIN} (which is a shortcut name for the pin that controls the voltage to the LED) is set to be an output for the duration of the loop. In general, you will need to set all of the pins you use as either inputs or outputs during the setup function. If you wired up the camera yourself, you have already explored a lot of the special names reserved for the pins as shown on the pinout diagram.

In the loop function, you’ll notice that we are alternating between writing a HIGH (aka turning on the LED) and writing a LOW (aka turning off the LED). The delay of 1000 milliseconds (1 second) between each step is crucial as otherwise the light would turn on and off so fast that it would be imperceptible. In fact, if you make the delay too short the light will simply seem to be dim. This is a trick called Pulse Width Modulation that is actually used often in industry to e.g., control motors. If you’d like, feel free to experiment with modifying these delays and redeploying the code to see its effect!    


\section{Testing the Tensorflow Install}

\subsection{Testing the TFLM Installation}

In this reading, we are going to walk through deploying the ‘Hello World’ example provided in the TensorFlow Lite for Microcontrollers (TFLM) library. The Hello World application runs a neural network model that can take a value, x, and predict its sine, y. Visually, the sine wave is a pleasant curve that runs smoothly from –1 to 1 and back. This makes it perfect for controlling a visually pleasing light show! We’ll be using the output of our model to power a smoother blink of your Arduino LED.

       \includegraphics[width=\textwidth]{TinyML/3 Development/Circle_cos_sin-re.jpg}

       \includegraphics[width=\textwidth]{TinyML/3 Development/IMG_0084__2_}% .gif}

On the left a sine wave and on the right the LED fading in and out in a sine wave pattern.

While the result won’t differ greatly from what we’ve accomplished with standard Arduino Blink sketch we just deployed (in either case, an LED will be turning on and off), in this instance, we are going to be deploying a model. Previously, we have learned about how neural networks can learn to model patterns they see in training data to go on to make predictions. Here, the model we deploy will call on maps an input (x) to an output y = sin(x). While this is admittedly contrived, it will allow us to quickly evaluate if the TFLM stack is functional on the hardware in front of us by using this output y to control the intensity of light emitted by the on-board LED. As promised, below, you’ll find instructions and screenshots and short videos walking you through the process!

\subsection{TFLite Micro Hello World}

0:00 / 0:00Klicken Sie auf die Pfeiltaste nach oben, um das Geschwindigkeitsmenu aufzurufen. Regeln Sie dann mit den Hoch und Runter Pfeiltasten die Geschwindigkeit. Klicken Sie Enter, um die jeweilige Geschwindigkeit festzulegen.Geschwindigkeit1.0xKlicke auf diesen Knopf und diess Video stumm oder nicht stumm zu schalten, oder drücke die HOCH oder RUNTER Knöpfe um die Lautstärke zu erhöhen oder abzusenken.Maximal Lautstärke.Videotranskript
Beginn des Transkriptes, zum Ende springen.
BRIAN PLANCHER: Next, we'll be using the Hello World example from the TensorFlow
Lite for Microcontrollers library to test out your TensorFlow
install and make sure that all worked correctly.
To do that, we're going to open the Hello World
example by going to File, Examples, Arduino TensorFlow Lite, Hello World.
And again, we're going to have a new window [? COM ?] that we can open up,
which is the Hello World example.
Now, this one has multiple different files.
And I'll mention some of them briefly in the next short video.
But for now, let's get this deployed to make sure that it will deploy correctly
on your device.
To do that, again, we're going to go up to Tools
and make sure that the board is set to the Arduino Nano 33 BLE,
and that the port is set to the port that the Arduino is on.
Again, look to the reading for what the port should
look like depending on if you're a Windows, Mac, or Linux operating system
user.
For me, I'm on Windows, so it's COM, some number.
And in this case, it's 7.
And then again, just like before, we can click the Upload
button to start compiling the code.
And again, if you have set your preferences
to show the full debug output, you're going to see a lot of things go here.
And this might take a couple of minutes to compile.
But once it's done, we'll see the output on the board.
Now, just as a comparison, so you know what's coming, right now,
I have the Arduino Blink example going on my board.
And you can see the light turning on and off, blinking.
When we get the Hello World example from the TensorFlow
Lite for Microcontrollers library running,
what we're going to see there is that, instead of blinking on and off,
the light is going to follow a sine wave, which is predicted
by the neural network on the board.
Therefore, it will dim on and dim off according to that sine wave, which
again, is learned by a neural network.
So we're doing the super fancy TensorFlow version of Blink
on this board.
I'll give my computer a minute to finish compiling and uploading the code.
And then we'll look at what the output looks like.
OK, now the Hello World example has finished compiling
and has been uploaded to my board.
We can look over to the board.
And we can see that the LED is now dimming
in and out at a much faster rate than it was before with the Blink example.
And that's exactly what you should expect from the Hello World example.
We'll explore the code in a little bit more detail in the next video.
See you there.
Ende des Transkriptes, zum Anfang springen
Downloads und Transskripte
Video
Videodatei herunterladen
Transskripte
Download SubRip (.srt) file
Download Text (.txt) file

\subsection{Preparing for Deployment:}

\begin{enumerate}
  \item Use a USB cable to connect the Arduino Nano 33 BLE Sense to your machine.
  \item Open the  \FILE{hello\_world.ino} sketch, which you can find via the File drop-down menu. Navigate, as follows: File → Examples → \FILE{Arduino\_TensorFlowLite} → \FILE{hello\_world}.

       \includegraphics[width=\textwidth]{TinyML/3 Development/c1-re.png}

A screenshot of the top left of the Arduino IDE showing the navigation to the \FILE{Arduino\_TensorFlowLite} \FILE{hello\_ world} example. The correct path is as follows: File (first option on the top menu) → Examples (about half way down) → \FILE{Arduino\_TensorFlowLite} (about half way down) → \FILE{hello\_world} (first option). Do note that Arduino sometimes changes the ordering of menu items during updates to the IDE so please follow the words and do not be alarmed if they are in slightly different locations.

  \item Just like with the Arduino Blink example (and as we will do for every deployment), use the Tools drop-down menu to select appropriate Port and Board.

  \begin{itemize}
      \item Select the Arduino Nano 33 BLE as the board by going to Tools → Board: <Current Board Name> → Arduino Mbed OS Boards (nRF52840) → Arduino Nano 33 BLE. Note that on different operating systems the exact name of the board may vary but/and it should include the word Nano at a minimum. If you do not see that as an option, then please go back to Setting up the Software and make sure you have installed the necessary board files.

       \includegraphics[width=\textwidth]{TinyML/3 Development/c2-re.png}

      Another screenshot of the IDE showing another set of menus and submenus that need to be traversed to select the correct board. The correct path is: Tools (4th option) → Board: <Current Board Name> (about 2/3rds of the way down)→ Arduino Mbed OS Boards (nRF52840) (3rd option)→ Arduino Nano 33 BLE. (3rd option). Do note that Arduino sometimes changes the ordering of menu items during updates to the IDE so please follow the words and do not be alarmed if they are in slightly different locations.

     \item Then select the USB Port associated with your board. This will appear differently on Windows, macOS, Linux, but will likely indicate ‘Arduino Nano 33 BLE” in parenthesis. You can select this by going to Tools → Port: <Current Port (Board on Port)> → <TBD Based on OS> (Arduino Nano 33 BLE). Where <TBD Based on OS> is most likely to come from the list below where <\#> indicates some integer number.
     
     \begin{itemize}
       \item Windows $\rightarrow$ COM<\#>
       \item macOS $\rightarrow$ /dev/cu.usbmodem<\#>
       \item Linux $\rightarrow$ ttyUSB<\#> or ttyACM<\#>
       
              \includegraphics[width=\textwidth]{TinyML/3 Development/c3-re.png}
       
        Another screenshot of the IDE showing another set of menus and submenus that need to be traversed to select the correct port. The correct path is Tools (4th option) → Port: <Current Port (Board on Port)> (about 3/4 of the way down)→ <TBD Based on OS> (Arduino Nano 33 BLE) (location depends on your machine). Do note that Arduino sometimes changes the ordering of menu items during updates to the IDE so please follow the words and do not be alarmed if they are in slightly different locations.
      \end{itemize}
  \end{itemize}

  \item As always, it is best practice to then verify that the code is valid. Note that this may take a while the first time you run this as all of TFLM is being compiled (it will compile much faster in the future). Also, you may, or may not, see a series of compiling warnings (but not errors) as the code compiles, this is entirely normal either way and is the result of TFLM being bleeding edge software.
  
       \includegraphics[width=\textwidth]{TinyML/3 Development/c4-re.png}
  
A screenshot of the top left of the Arduino IDE showing the verify checkmark right below the top menu items all the way on the left.
\end{enumerate}


\subsection{Deploying (Uploading) the Sketch}

Once we know that the code at hand is valid, we can flash it to the MCU:

\begin{enumerate}
  \item Use the rightward arrow next to the ‘compile’ checkmark to upload / flash the code. 

       \includegraphics[width=\textwidth]{TinyML/3 Development/c5-re.png}

   A screenshot of the top left of the Arduino IDE showing the download right arrow right below the top menu items second from the left and one to the right of the verify checkmark we used in the last step.

  You’ll know the upload is complete when you see red text in the console at the bottom of the IDE that shows 100\% upload of the code and a statement that says something like “Done in <\#.\#> seconds.”

  If you skipped the verification step, do note that this may take a while the first time you run this as all of TFLM is being compiled (it will compile much faster in the future). Also, you may, or may not, see a series of compiling warnings (but not errors) as the code compiles, this is entirely normal either way and is the result of TFLM being bleeding edge software.

If you receive an error, you will see an orange error bar appear and a red error message in the console (as shown below). Don’t worry -- there are many common reasons this may have occurred. To help you debug, please check out our FAQ appendix with answers to the most common errors!

  \item At this point, you'll want to look to the board itself. The orange LED opposite the green LED power indicator about the USB port should now be fading at the rate set by the sinusoidal model. Note that the default rate is quite fast so it should almost appear to be blinking very quickly. However, if you look closely you’ll notice that it is fading in and out vs. blinking.

       \includegraphics[width=\textwidth]{TinyML/3 Development/GIF2}% .gif}

  A GIF of the tinyML kit running the \FILE{Arduino\_TensorFlowLite} \FILE{hello\_world}  sketch. We can see that the orange LED opposite the green LED power indicator about the USB port is fading in and out at the rate set by the sinusoidal model. Note that the default rate is quite fast so it should almost appear to be blinking very quickly. However, if you look closely you’ll notice that it is fading in and out vs. blinking.

\end{enumerate}

Congratulations! At this point, if all has gone well, you can be sure that your embedded microcontroller is configured properly to run TensorFlow Micro.

\subsection{Understanding the TFLM Hello World Example}

0:00 / 0:00Klicken Sie auf die Pfeiltaste nach oben, um das Geschwindigkeitsmenu aufzurufen. Regeln Sie dann mit den Hoch und Runter Pfeiltasten die Geschwindigkeit. Klicken Sie Enter, um die jeweilige Geschwindigkeit festzulegen.Geschwindigkeit1.0xKlicke auf diesen Knopf und diess Video stumm oder nicht stumm zu schalten, oder drücke die HOCH oder RUNTER Knöpfe um die Lautstärke zu erhöhen oder abzusenken.Maximal Lautstärke.Videotranskript
Beginn des Transkriptes, zum Ende springen.
BRIAN PLANCHER: Now that you got the Hello World example deployed,
I thought I would talk a little bit about how the code is structured.
At a high level, it's going to operate very similarly to the blink example,
but with a bit of complication.
So if we go back to the code, you'll notice that on the blink example,
we just had one file, this blink file, with that setup
and loop function we talked about earlier.
Now, on our Hello World example, we have a whole series
of files across the top and a lot more code in here
before we get to that set up.
And a lot more code in the setup function
before we get to the loop and then a lot more code in that loop function.
But at a high level we're doing the same thing.
Now all these files across the top are simply just helper files, many of which
you'll see included here.
Like the model.h file, we have the \FILE{output\_handler.h} file,
these are just some of these files that are up here.
And as you saw in the intro to C++, generally you just include the h files
or in practice you just include the h files.
And then during compilation, the c files,
or the C++ files are also compiled.
So just as an example, the model.h file, is describing our model
and there's very little there.
It just says you know, there is going to be something
called \PYTHON{g\_model} and \PYTHON{g\_model len} and that is actually--
let me Zoom in a little bit for you, --that
is actually just the model itself and the size of the model.
And if we go to the C++ file, and what we find is that there is a model file,
which is this giant array of bytes, which is our actual TensorFlow Lite
for microcontrollers model in the length of the bottom, and that's it.
And many of these files are relatively simple,
but it's a nice way to sort of compartmentalize
the code into various files that describe exactly what is in that file.
And I'll let Pete and Vijay get into more detail on that.
But just going back to the overall flow, we
have these includes up front, include the TensorFlow library--
it's actually Zoom in a little bit, --incclude the TensorFlow library,
include all those helper files and then define some global variables.
This is important because we need to define something out
of the context of the setup or loop function,
if we want them to be used by both the setup and loop function.
So in this case things like the Tensor arena,
which we'll talk about a lot more, but it's where we'll store the model.
We need to define that once so we can initialize
the model into there in the setup function
and then use it in the loop function.
And then, the Setup and loop will do somewhat what you would think.
So Setup function will--
I mean set up a bunch of things, but also really
at a high level, the important thing it does here, is that initializes
our neural network, which it does in a couple of steps.
But essentially we're going to build an interpreter, allocates some memory,
get some pointers and essentially by doing all of that,
we will have defined our neural network.
And again Vijay and Pete will describe more what that exactly means.
And then our loop function will loop and it will take in some x value.
Which is this x value we're getting here and then
it will run it through the model for a little quantize it.
It'll run through the model by running inference.
And then at the back end, we'll get some y value out,
y equals sine of x that is sort of what we're computing.
Well that is exactly what we're computing,
although it's an approximation, because we're using
a neural network for the sine function.
And then we'll take that y and we'll use some output Handler to then
dim or brighten the LED accordingly.
And this is the sort of flow that you'll see in kind of all of the examples.
And again I'll let Vijay and Pete get into the details.
But I want to kind of show you that flow,
because even in keyword spotting, in person detection and all
these sorts of things that we'll do in the rest of this course,
there will be some set of includes, some set of global variables.
In the setup, we will initialize the neural network among other things.
And in the loop, we will do some work to pre-process and get some raw data
and then pass it to our neural network.
And then once we get something out, we will do some post-processing
and handle the output and then have some effect.
And so that's sort of the flow that you'll see over and over again.
Now, I know you're itching to get into the details of TensorFlow
Lite from microcontrollers and embedded systems in general,
but there's one more set of tests I need to run with all of you
before I can pass you on to Vijay.
And that is to make sure that all of your sensors
are actually working correctly.
So I'll see you in the sensor test reading
and then after that, you should be all set for the rest of this course.
Ende des Transkriptes, zum Anfang springen
Downloads und Transskripte
Video
Videodatei herunterladen
Transskripte
Download SubRip (.srt) file
Download Text (.txt) file


\subsection{Understanding the TFLM Hello World Example}

The first thing you will notice when looking at this example is that there is a lot more code in this TFLM version of blink! However, at a high level the code is structured the same. There is still a setup() function that runs once to initialize things and a loop() function that runs continuously to blink the LED.

As far as the high level differences go, in the main \FILE{hello\_world} file, before either of those functions, there are a set of \#include lines and global variable definitions. The \#include lines are used to (in Python speak) import the various libraries we need for this example. In this case it is the TFLM library and some other helper functions. The global variables are defined outside of the setup()and loop() functions, so that there are references to variables that can be used in both the \PYTHON{setup()} function as well as in every loop of the \PYTHON{loop()} function. For example, the global \PYTHON{tensor\_arena} memory is used to initialize and store our neural network once and then is used to execute the neural network every loop iteration. This is much more efficient than re-initializing the neural network in each loop iteration.

One other high level difference is that you’ll notice that there are a series of tabs across the top of the IDE window. These tabs each represent different files that make up this \FILE{hello\_world} project. Each holds some helper functions or definitions that our main \FILE{hello\_world} file uses.

       \includegraphics[width=\textwidth]{TinyML/3 Development/c6.png}


  A screenshot of the top of the Arduino IDE showing that now below the veryify and upload icons there is a ersies of tabs across the IDE window that represetn the different files used by this sketch. The first is the \FILE{hello\_world.ino} main file and then there are a series of others.

For example, the model.cpp file contains a large array which contains the binary data for the neural network model (in fact, it is simply the binary version of the TensorFlow Lite file) as well as an integer, which tells TFLM how long that array is. We’ll explore how that binary file is generated later in the course!

       \includegraphics[width=\textwidth]{TinyML/3 Development/c7.png}

A screenshot of the top of the Arduino IDE showing that now below the veryify and upload icons there is a ersies of tabs across the IDE window that represetn the different files used by this sketch. The first is the \FILE{hello\_world.ino} main file and then there are a series of others. In this image we have selected the \FILE{model.cpp} file far to the right. We see that this file contains a large amount of hexidecimal values (e.g., 0x00). We'll explore these later in the course but as a preview, they described the actual neural network model values.

Getting back to the main \FILE{hello\_world} file, as mentioned above, the setup() function has a decent amount of code in it, but at a high level, it does exactly what you would expect, it initializes our neural network.

The loop() function also has a decent amount of code in it, but again at a high level, all it does is keep track of a counter over time for the x-value and uses the neural network to compute the approximate y-value of y = sin(x) and then set the brightness of the LED according to that approximate y value.

Now, there are a lot of details that go into making all of those helper functions work, and excitingly, Pete Warden, the Technical lead of TensorFlow Mobile and Embedded team at Google will dig deeper into many of these topics in more detail later on. But before that, you’ll need to run one more set of sketches to test your sensors!


% 1.2: Getting Started 13/15

\section{Testing the Sensors}

With tests of the microprocessor itself and TFLM stack completed, we can turn our attention to verifying that the sensors required for this course are functioning properly. To do so, we’ll be calling on a predefined script that the course staff have put together. By the end of this reading, you will have seen live data streams for the on-board microphone, the STMicroelectronics, MP34DT05, the on-board internal measurement unit (IMU), the STMicroelectronics, LSM9DS1, as well as at least a static image return from the off-board OV7675 camera module, the OmniVision OV7675. We will also detail the optional extensions you’ll need to do to obtain a live video feed from the camera as well.

\subsection{The Serial Monitor \& Plotter}

The Serial Monitor \& Plotter are two core tools that you can use to get information from your Arduino when it is connected to your computer with a data USB cable. In fact, as you develop your own Arduino applications, you’ll probably find yourself opening up the Serial Monitor and using the \PYTHON{Serial.println()} command in your Arudino sketches much like how you have used the \PYTHON{print()} function in Python or \PYTHON{printf()} function in C++. Since the Arduino sketch runs in an infinite loop, you may find that it may be easier to plot the graph of the data you are sending back from the Arduino instead of continuously printing the raw values, and that is where the Serial Plotter will come in.

While the primary function of the Serial Monitor (and the Serial Plotter) is to view data incoming from the MCU, the Serial Monitor also features a text entry field that we will use to issue pre-determined commands conveyed in the console.

Below we have included a screenshot of the Serial Monitor elements to differentiate where we expect to see incoming data and where we can enter commands to send to the microcontroller:

       \includegraphics[width=\textwidth]{TinyML/3 Development/d1.png}

A screenshot of the serial monitor showing that data is displayed in the middle and that you can enter text/commands at the top. Next to the text/ command input area is a send bytton which will send that command. At the bottom of the window are a series of options. The first two on the left are for autoscroll (checked by default) and to show timestamps (unchecked by default). On the right there is a dropdown to select the kind of line endings to send a command with (by default no line ending), a dropdown to select the baud rate of the serial connection (we will use the defauly 9600 for this course), and a button to clear the output in the middle of the screen.

You can open up the Serial Monitor (and Serial Plotter) by navigating through the menu to Tools → Serial Monitor or Tools → Serial Plotter.

       \includegraphics[width=\textwidth]{TinyML/3 Development/d2-re.png}

A screenshot of the Arduino IDE showing how to access the Serial Monitor using the menus by navigating to Tools (4th option) → Serial Monitor (4th option) or Serial Plotter (5th option).

\subsection{Testing the Microphone}

0:00 / 0:00Klicken Sie auf die Pfeiltaste nach oben, um das Geschwindigkeitsmenu aufzurufen. Regeln Sie dann mit den Hoch und Runter Pfeiltasten die Geschwindigkeit. Klicken Sie Enter, um die jeweilige Geschwindigkeit festzulegen.Geschwindigkeit1.0xKlicke auf diesen Knopf und diess Video stumm oder nicht stumm zu schalten, oder drücke die HOCH oder RUNTER Knöpfe um die Lautstärke zu erhöhen oder abzusenken.Maximal Lautstärke.Videotranskript
Beginn des Transkriptes, zum Ende springen.
SPEAKER: The first sensor that we're going to test, is the microphone.
To get to that testing sample sketch, we're going to go to File, examples
and we're going to scroll all the way down and keep
going past this arrow down at the bottom, until we get to Harvard TinyMLx
and then we will open the test microphone sketch.
As always, we need to make sure that the board is set to the Arduino Nano 33
and that the port is set to whatever port your Arduino is on.
And then we need to click Upload, to start compiling and uploading the code.
And once it's uploaded, we'll talk about how we run this test.
All right.
So now that you've deployed the code, let's use it to test your microphone.
To do that, we're going to use the serial monitor and serial
plotter, which allow us to communicate with the Arduino
over the universal serial bus or USB protocol, with which we've used
to connect our Arduino to our computer.
To get to the serial monitor, we're going to go to Tools.
And one quick tip, make sure the port is still selected.
Sometimes after the board is uploaded with new code, it will reset the port
and it will sort of deselect.
And so you may have to reselect it there.
So if you run into an error just go back to error and just check.
So let's open up that serial monitor.
And what we'll see, is it'll be a message from our Arduino.
It says welcome to the microphone test.
Well thank you very much Mr Arduino.
Use the on-shield button or send the command click,
to start and stop an audio recording.
So what we can do is up at the top here, there's this text field
and we can either send the command click there,
or we can click the button on our shield.
So let's just start by sending the command here.
Now, one thing to be careful of, is you have
to make sure you select NL and CR, both NL and CR for this
to work with our code.
So if we send that command, suddenly numbers
are going to be dropping all over the screen here.
And all of these numbers, this is actually just the waveform
coming off of the microphone.
And I can stop it by sending that click Command again.
Or I can start it and stop it, by clicking the button.
Now, this means that there is some communication coming
from the microphone there, there is data coming off.
But these numbers are a little hard to interpret.
So let's open up the serial plotter, which
should help us better visualize this.
Now the serial plotter, assumes the numbers coming in
are some sort of time series data and plots it
to the screen, which is kind of exactly what we want here.
So let's send that click Command over the serial plotter,
noting that now the text entry field is on the bottom.
And now we'll see a waveform coming across
and now this looks a lot more like somebody talking.
If I get closer to the thing and get louder,
you'll see the scale numbers go up.
And if I get a lot quieter when I move over here,
the numbers go a little bit down.
So it seems like at a high level, that's kind of working.
Now one last test you can do and this is a test
that I am very bad at because I am not a great singer.
Is that if you can actually hum or sing a single pitch,
you should end up with a perfect sine wave.
And however far off you are from the perfect pitch,
it's a less perfect sign wave but what we're
going for is something that's kind of repetitive.
So for an example, let me give it my best shot here
and try to make some tone.
So like, [MUSIC PLAYING] And so that doesn't look like a sine wave,
I'm not a great singer, but it is a nice repetitive senesh--
we can call it thing.
And so I feel pretty confident that my microphone is roughly
picking up my mediocre singing and delivering
some sort of repetitive wave, which which is exactly what we want.
And so, you should try something else yourself at home.
If you're a great singer, if you have perfect pitch,
I challenge you to see how perfect your sine wave can be.
Mine will not be much better than that, I've
actually tried doing this a couple times on a couple of different recordings.
And I'm going to call that good for me for now.
So I'll see you in the next test.
Ende des Transkriptes, zum Anfang springen
Downloads und Transskripte
Video
Videodatei herunterladen
Transskripte
Download SubRip (.srt) file
Download Text (.txt) file

\subsection{Testing the Microphone}

\begin{enumerate}
    \item Use a USB cable to connect the Arduino Nano 33 BLE Sense to your machine. You should see the green LED power indicator come on when the board first receives power.
  \item Open the \FILE{test\_microphone.ino} sketch, which you can find via the File drop-down menu. Navigate, as follows: File → Examples → \FILE{Harvard\_TinyMLx} → \FILE{test\_microphone}.
  \item As always, use the Tools drop-down menu to select appropriate Port and Board.
  
  \begin{itemize}
    \item Select the Arduino Nano 33 BLE as the board by going to Tools → Board: <Current Board Name> → Arduino Mbed OS Boards (nRF52840) → Arduino Nano 33 BLE. Note that on different operating systems, the exact name of the board may vary, but/and, it should include the word Nano at a minimum. If you do not see that as an option, then please go back to Setting up the Software and make sure you have installed the necessary board files.
     \item Then, select the USB Port associated with your board. This will appear differently on Windows, macOS, Linux, but will likely indicate ‘Arduino Nano 33 BLE” in parenthesis. You can select this by going to Tools → Port: <Current Port (Board on Port)> → <TBD Based on OS> (Arduino Nano 33 BLE). Where <TBD Based on OS> is most likely to come from the list below where <\#> indicates some integer number.
       \begin{itemize}
           \item Windows → COM<\#>
           \item macOS → /dev/cu.usbmodem<\#>
           \item Linux → ttyUSB<\#> or ttyACM<\#>
       \end{itemize}
    \end{itemize}
  \item Use the rightward arrow next to the ‘compile’ checkmark to upload / flash the code.
You’ll know the upload is complete when you see red text in the console at the bottom of the IDE that shows 100% upload of the code and a statement that says something like “Done in <#.#> seconds.”

If you receive an error, you will see an orange error bar appear and a red error message in the console (as shown below). Don’t worry -- there are many common reasons this may have occurred. To help you debug, please check out our FAQ appendix with answers to the most common errors!

  \item Open the Serial Monitor. You can accomplish this in three different ways as shown below. If the Serial Monitor fails to open, check to make sure the appropriate Port is selected. Sometimes, the port is reset during the upload process.
  
  \begin{itemize}
      \item Use the menu to select: Tools → Serial Monitor.
    \item Click on the magnifying glass at the top right of the Arduino Desktop IDE.
    \item Use the keyboard shortcut: CTRL + SHIFT + M or CMD + SHIFT + M depending on your operating system.
  \end{itemize}
  \item When the Monitor opens, you should see the following text appear: 
Welcome to the microphone test for the built-in microphone on the Nano 33 BLE Sense.

Use the on-shield button or send the command 'click’ to start and stop an audio recording Open the Serial Plotter to view the corresponding waveform.

  \item Press the on-shield button to start an audio recording. After you do, a stream of data should appear in the Serial Monitor. Press the button again to stop the recording. If the numbers were changing, then your microphone is most likely correctly recording audio! Congratulations! If you do not have the shield, you can also control the starting and stopping of the waveform by sending the command click through the serial monitor. In the drop-down menu that reads “No line ending” by default at the bottom right of the Serial Monitor, you need to select “Both NL \& CR”. This tells the Serial Monitor to send both a NL = new line character as well as a CR = carriage return character every time you send a message. This is important for our application as our Arduino sketch is looking for that to differentiate between each input.

          \includegraphics[width=\textwidth]{TinyML/3 Development/d3-re.png}
   
A screenshot of the Serial Plotter showing how the data comes into the large white area at the top and then on the bottom there is the dropdown to select the baud rate and an input box in which we have input the “click” command. Before we press the “send” button next to the input on the dropdown all the way on the bottom right we select the “Both NL \& CR” option to send the correct line ending.

  \item To help visualize this a little better, we are going to use the Serial Plotter. Note that you cannot have both the Serial Monitor and Serial Plotter open at the same time as the ARduino can only communicate over a single serial port. Importantly, this also means that you cannot upload new code to the Arduino if either the Serial Monitor or Plotter is open! You can open the Serial Plotter in two ways:
  
  \begin{itemize}
      \item 
Use the menu to select: Tools → Serial Plotter.
  \item Use the keyboard shortcut: CTRL + SHIFT + L or CMD + SHIFT + L depending on your operating   \end{itemize}
system.
\item Perform the same experiment, press the on-shield button (or send the serial command click) to start an audio recording and again to stop it. You should see a waveform appear on the Serial Plotter. This is simply a graphical representation of the numbers you saw earlier on the Serial Monitor. Note that the vertical axis of the Serial Plotter scales dynamically, so that the magnitude of the data conveyed is not fixed with respect to the enclosing window, but rather scaled to occupy most of the headroom.Like the previous figure this is a a screenshot of the Serial Plotter showing how the data comes into the large white area at the top and then on the bottom there is the dropdown to select the baud rate and an input box in which we have input the “click” command. Before we press the “send” button next to the input on the dropdown all the way on the bottom right we select the “Both NL \& CR” option to send the correct line ending. In this image we see the data of a wave form in the white area as it is showing an audio recording.

  \includegraphics[width=\textwidth]{TinyML/3 Development/d4-re.png}
  
\item Now, record some audio again. This time, try to whistle or hum a constant tone. If you can keep the tone constant, you should see the waveform start to look like a sinusoid. An example that the course staff made is below. Don’t worry if you can’t get this to work perfectly. You should in general find that a constant tone has a more constant pattern. If when you make different sounds and you find that the waveform changes, you can be confident that your microphone is working.
 
             \includegraphics[width=\textwidth]{TinyML/3 Development/d5.png}
 
Like the previous two figures this is a a screenshot of the Serial Plotter in this one we see a much clearer sinusoidal waveform created by trying to hum a constant tone. The course staff are not trained singers and so a break in the wave is seen in the middle (we tried our best).
\end{enumerate}


\subsection{Testing the Inertial Measurement Unit}
0:00 / 0:00Klicken Sie auf die Pfeiltaste nach oben, um das Geschwindigkeitsmenu aufzurufen. Regeln Sie dann mit den Hoch und Runter Pfeiltasten die Geschwindigkeit. Klicken Sie Enter, um die jeweilige Geschwindigkeit festzulegen.Geschwindigkeit1.0xKlicke auf diesen Knopf und diess Video stumm oder nicht stumm zu schalten, oder drücke die HOCH oder RUNTER Knöpfe um die Lautstärke zu erhöhen oder abzusenken.Maximal Lautstärke.Videotranskript
Beginn des Transkriptes, zum Ende springen.
SPEAKER 1: The next sensor that we're going to test
is your inertial measurement unit.
We're going to open the example again, in the same way, by going to File,
Examples, Harvard TinyMLx, test IMU.
And again, as always, we're going to make sure that the Board and the Port
are correctly selected.
And we're going to go and upload and compile that code.
Once that's all up on your device, we can
talk about how to actually use it to test your sensor.
All right, now that you've deployed your code, let's use it to test the IMU.
To do that, we're going to use the Serial Monitor and Serial Plotter,
just like in the last test.
So let's start by opening up that Serial Monitor.
And again, make sure the Port is still selected.
So when we open up the Serial Monitor, we're
going to get a window with some instructions
about three tests we can use.
a, g, and m, which will display the accelerometer, the gyroscope,
and the magnetometer, which are the three different sensors that
make up an IMU.
Now, just like in the last test, if we do it in the Serial Monitor,
we're just going to see those numbers getting displayed to the screen.
So let's move over to the Serial Plotter,
where we can actually get a graph of that output.
And I'm going to use g, for the gyroscope,
because I find that that's the easiest one to test and the easiest one
to visualize.
What the gyroscope is asking for, or displaying,
is the angular acceleration.
And that is the rate of change of the orientation of the Board.
And since the Board is flat on the table,
we see basically, just a little bit of noise around zero.
But if we pick it up and start to yaw it about one axis,
suddenly, we're going to see a huge spike and trough in the green line.
And if we change the axis about, which we are rotating it,
so let's say we start to pitch it up and down.
Now we can get the orange line excited.
And if we say, roll it back and forth, now we can get the blue line excited.
And so that is a really great test to see if your IMU is working.
If you can just yaw it, pitch it, and roll it, and get the blue, green,
and orange lines to be excited individually,
then your IMU should be working great.
You can also test the accelerometer and the magnetometer,
as well, if you'd like.
And once you're done with that, I will see you in the camera test.
Ende des Transkriptes, zum Anfang springen
Downloads und Transskripte
Video
Videodatei herunterladen
Transskripte
Download SubRip (.srt) file
Download Text (.txt) file

\subsection{Testing the Inertial Measurement Unit}

\begin{enumerate}
  \item Now, let's open the \FILE{test\_IMU.ino} sketch, which you can find via the File drop-down menu. Navigate, as follows: File → Examples → \FILE{Harvard\_TinyMLx} → \FILE{test\_IMU}.
  \item As always, use the Tools drop-down menu to select appropriate Port and Board.
  \item Then use the rightward arrow next to the ‘compile’ checkmark to upload / flash the code. If you get the error “\FILE{Arduino\_LSM9DS1.h}: No such file or directory,” then please go back to Setting up the Software and make sure you install the accelerometer, magnetometry, and gyroscope library! To help you debug, please check out our FAQ appendix with answers to the most common errors!
  
  \item Open the Serial Monitor and you should see the following. As a reminder, you can open the serial monitor in one of three ways as shown below. If the Serial Monitor fails to open, check to make sure the appropriate Port is selected. Sometimes, the port is reset during the upload process.
   
    \begin{enumerate}
      \item Use the menu to select: Tools → Serial Monitor.
      \item Click on the magnifying glass at the top right of the Arduino Desktop IDE.
      \item Use the keyboard shortcut: CTRL + SHIFT + M or CMD + SHIFT + M depending on your operating system.
      
      
        Welcome to the IMU test for the built-in IMU on the Nano 33 BLE Sense.

        Available commands:
        
          \begin{description}
              \item[a] - display accelerometer readings in g's in x, y, and z directions
              \item[g] - display gyroscope readings in deg / s in x, y, and z directions
              \item[m] - display magnetometer readings in uT in x, y, and z directions
          \end{description}

  \item Like with the microphone, in the drop-down menu that reads “No line ending” by default at the bottom right of the Serial Monitor, you need to select “Both NL \& CR”.
  \item Now, you can enter one of the possible arguments (a, g, or m) into the text entry bar across the top of the Serial Monitor and click “Send.” As mentioned in the Serial Monitor, this will start to print the data coming from the [a]ccelerometer (computing acceleration in the x, y, or z direction), the [g]yroscope (computing the angular velocity -- the change in the roll, pith, and yaw), or the [m]agnetometer (computing the magnetic forces present on the IMU). Depending on your selection, you should now see a stream of corresponding data. However, like with the microphone, this data is hard to interpret in raw form. So, instead, lets us open the Serial Plotter.
  \item Close the Serial Monitor and open the Serial Plotter. As a reminder, you can open the Serial Plotter in two ways:
  
    \begin{itemize}
      \item Use the menu to select: Tools → Serial Plotter.
      \item Use the keyboard shortcut: CTRL + SHIFT + L or CMD + SHIFT + L depending on your operating system.
    \end{itemize}

  \item With the Plotter open, you should see the same stream of data presented graphically. The most interesting one to plot is the [g]yroscope. Again, make sure you have selected “Both NL \& CR” and type “g” into the text entry box (now at the bottom of the window) and click “Send.”
  
  \includegraphics[width=\textwidth]{TinyML/3 Development/d6-re.png}
  
  Like the previous figures we see another screenshot of the Serial Plotter. In this case we have again selected “Both NL \& CR” on the bottom right dropdown to send the correct line ending and have sent the message “g” for “gyroscope”. We see threee waveforms in blue, orange, and green representing the different axis of measurements. Since we tried to rotate first along one axis and then on a second axis we see that the orange wave (representing rotation on the thrid axis) is always a small wave, while the blue wave at first shows large rotations and then is small while we see at first small and then large rotations on the green wave.

  \item Move the board around and observe the response. Can you figure out which axis of rotation is the x, the y, and the z (also known as roll, pitch, and yaw)? If by moving the board around in different directions you get different responses from the various lines in the Serial Plotter, you can feel confident that you have a working IMU.

  \end{enumerate}
\end{enumerate}

\subsection{Testing the OV7675 Camera}

0:00 / 0:00Klicken Sie auf die Pfeiltaste nach oben, um das Geschwindigkeitsmenu aufzurufen. Regeln Sie dann mit den Hoch und Runter Pfeiltasten die Geschwindigkeit. Klicken Sie Enter, um die jeweilige Geschwindigkeit festzulegen.Geschwindigkeit1.0xKlicke auf diesen Knopf und diess Video stumm oder nicht stumm zu schalten, oder drücke die HOCH oder RUNTER Knöpfe um die Lautstärke zu erhöhen oder abzusenken.Maximal Lautstärke.Videotranskript
Beginn des Transkriptes, zum Ende springen.
BRIAN PLANCHER: All right, it's time for the last sensor test.
Let's test your camera.
To do that we're going to go to File, Examples, \FILE{Harvard\_TinyMLx}, \FILE{test\_camera}.
And again, like the rest of them, we're going
to make sure to set the board and the port by going to Tools, Board,
and make sure we have the Arduino Nano 33 selected.
And on the port, we'll make sure we have the serial port with our Arduino.
And then we can click Upload to compile and upload the code.
Once your code is uploaded to your board, your camera is now active.
Before we start taking pictures, let's make
sure we know which way to hold the camera,
regardless of whether you have the TinyML shield
or you wired up your own off-the-shelf parts.
And the important thing to keep in mind here
is that when you are taking a picture with the camera,
you want the letters OV and numbers 7675 to be facing up and towards you.
As long as those are facing at the top and you can read them,
then you are taking the correct selfie.
And if you were, say, facing the camera the other way,
you'd want them oriented in the other direction.
So that's how you want to hold the camera, with the OV7675 up.
If you wanted to take a selfie, you'd point it at you, then, like that.
So let's go to the Serial Monitor now to start to test the camera.
So we'll open up the Serial Monitor.
And I forgot to select the port again, so I will select the port.
Again, sometimes the ports disconnect after it uploads.
So let's open that Serial Monitor.
And on the Serial monitor, we will see the welcome message here
telling us that we have entered the OV7675 test.
And we have three commands available-- "single," "live," and "capture."
We'll leave "live" for you if you would like later.
You'll see some notes in the reading.
It doesn't really work on all operating systems,
and so we won't use that as our main test.
The main test we're going to use is the single camera,
which is just to take a single picture.
So let's type in "single" here and make sure we have "Both NL \& CR" there.
And then, when we want to take a picture,
we can either type in "capture" or just click the button on the shield.
So let's send that single, and now the camera is in single mode.
It's ready to take a picture.
And so one note--
the camera field of view is very, very small.
So if you are taking a selfie, you're going
to want to hold that camera really, really far away from you.
In fact, it probably works better if you prop it up on something,
like maybe onto your computer monitor.
But then it's kind of hard to press the button.
So I'm just going to hold it as far as I can, make a ridiculous face,
and hope I can capture myself in the picture.
So after you press the button, you'll see a note
up here that says "image data will be printed out in 3 seconds,"
and then, suddenly, a lot of hexadecimal values will be coming at you.
And that's because this is the raw image coming off of the camera.
Unfortunately, as you see in the reading, there's a colab that we have,
which will take this raw data and display it for you.
So what you'll need to do is you'll need to copy all of that raw data.
So you can just highlight all of it, and if you just pull it down,
it should highlight the whole row.
And I will copy that.
And then I will open up that colab.
I'll make sure to run this first cell to import the correct packages,
and then what I'll do here is I will highlight the existing hexadecimal bits
and I will paste in my new data so we can run that.
And then we can run this cell below here, which will take that raw data
and display it down here for you.
And so let's see if I captured myself.
And I kind of got my face.
I kind of missed a little bit--
a little too high.
Again, it's pretty zoomed in, but I'm pretty sure the camera is working.
It's not the most flattering picture of me, but it is a picture of me.
So you can do the same thing to test your camera.
And that's it.
You're done with all the sensor tests, and I'll turn it over to VJ and Pete
to talk more about embedded systems and [? TensorFlow ?] [? Lite ?]
for [? microcontrollers. ?]
Ende des Transkriptes, zum Anfang springen
Downloads und Transskripte
Video
Videodatei herunterladen
Transskripte
Download SubRip (.srt) file
Download Text (.txt) file


\subsection{Testing the OV7675 Camera}

\begin{enumerate}
  \item Now, let’s open the \FILE{test\_camera.ino} sketch, which you can find via the File drop-down menu. Navigate, as follows: File → Examples → \FILE{Harvard\_TinyMLx} → \FILE{test\_camera}.
    
  \item As always, use the Tools drop-down menu to select appropriate Port and Board.

        Note: if you did not purchase the Tiny Machine Learning Kit and sourced an OV7670, as compared to the OV7675 that comes with the kit, you will need to make one software changes as detailed at the end of this document.

  \item Then, use the rightward arrow next to the ‘compile’ checkmark to upload / flash the code. If you get the error “\FILE{Arduino\_OV767X.h}: No such file or directory,” then please go back to Setting up the Software and make sure you install the accelerometer, magnetometry, and gyroscope library! To help you debug, please check out our FAQ appendix with answers to the most common errors! Note: if you get an error message that contains something like 'OV7675' was not declared in this scope this means you already had a conflicting copy of the \FILE{Arduino\_OV767X} library installed on your system. We have bundled a forked copy of it with our library so please uninstall the conflicting version.

  \item Open up the Serial Monitor and you should see:
  
         Welcome to the OV7675 test

         Available commands:

         \begin{description}
             \item[single] - take a single image and print out the hexadecimal for each pixel (default)
             \item [live] - the raw bytes of images will be streamed live over the serial port
             \item [capture] - when in single mode, initiates an image capture

         \end{description}

  \item Type “single” in the text entry field and press send or hit the enter / return key. The camera is now primed to take a picture. Now, you can either text “capture” or press the on-shield button to take a selfie (or photo). To get the photo oriented correctly, make sure the “OV7675” text on the camera module (or the “Tiny Machine Learning Shield” on the shield) is oriented such that it would be readable by the subject of the photo (e.g., you if you are taking a selfie). The camera does not have a large field of view, so if you are taking a selfie, make sure to hold the camera far away from your face.

  \item To view your image, you will need to copy the sequence of hexadecimal digits that will print to the Serial Monitor and paste them into this Google Colab that we created to display the image (link included below as well). Note: in some versions of the IDE the large amount of digits printed will cause an error in the Serial Monitor and the text will seem to disappear. Don't worry the text is still there (its just white on a white background). Just highlight the whole line below "Image data will be printed out in 3 seconds..." and copy and paste it into the Colab and the digits will appear!

  \url{https://colab.research.google.com/github/tinyMLx/colabs/blob/master/4-2-12-OV7675ImageViewer.ipynb}
\end{enumerate}  

\subsection{(Optional) Real-Time Video with Processing}

If you are interested in seeing a live stream capture from the OV7675, you can do that, but not through the Arduino IDE. You’ll have to use an additional piece of software, Processing. You can download the application for your OS, here (we’ve tested with version 3.5.4). We would like to note that we have found this to be buggy on Windows, which is why we provide this as an optional extension and not the default suggested process.

\begin{enumerate}
    \item Launch the \FILE{test\_camera} example through the Arduino IDE, open the Serial Monitor, and then type “live” in the text entry field and press send or hit the enter / return key. Your camera will now be live streaming the image to the computer at 1 frame per second. Make sure to then close the Serial Monitor to free up the serial port as we will be using it again later.

    \item With Processing installed (which for some operating systems is as simple as unzipping the pre-built application folder), open the Processing Application. You’ll find that the application looks a lot like the Arduino Desktop IDE.
 
   \item Then find your Arduino folder. On most operating systems it is located either in your home or Documents directory.

  \item In Processing, open the file Arduino → Libraries → \FILE{Harvard\_TiynMLx} → extras → CameraVisualizerRawBytes → \FILE{CameraVisualizerRawBytes.pde}. 

  \item Go to lines 34-36 and update them to be the name of the serial port on which your Arduino is currently connected.

  \item Now, click the “run” play button at the top left of the UI. A pop-up should appear streaming the camera image. Again, we have found this to be buggy on Windows, so don’t be surprised if the image looks odd on Windows.

  \includegraphics[width=\textwidth]{TinyML/3 Development/d7-re.png}

  Screenshot of the Processing IDE showing the play button needed to launch the file and start visualizing the camera output on the top left of the image. This button is in the same place as it would be for the Arduino IDE and the menus across the top are also almost all the same.

\end{enumerate}

\subsection{Changes for the OV7670}

To use the OV7670 instead of the OV7675, simply change the fourth argument of the call to \PYTHON{Camera.begin()} on line 25 of the \FILE{test\_camera} sketch, as highlighted below, from OV7675 to OV7670. That’s it! The library will handle the rest!

\begin{lstlisting}
// Initialize the OV7675 camera
void if(!Camera.begin(QCIF, RGB565, 1 OV7675)) {
    Serial.println(``Failed to initialize camera'');
    While(1);   
}
\end{lstlisting}

\subsection{Additional Resources}
If you’d like to learn more about debugging microcontrollers, we’ve put together a short guide in this appendix document.
If you’d like to learn more about powering your Arduino using a battery, we’ve put together another short guide in this appendix document.


\section{Appendix - Powering Your TinyML System}

Now that you have your Arduino development board up and running, let's talk about how you could deploy it independent of your computer! While some embedded systems call on AC-DC converters (or “wall warts,” colloquially) to provide low voltage power to their electronics (the Google home speaker, for example), others are battery powered. Both of these paradigms are applicable to real-world deployment of tinyML and both are achievable using your TinyML kit.
    
\subsection{USB Power Delivery}

To this point, we have provided power over USB to our microcontroller via the microB port on the Nano 33 BLE Sense. The 5V that USB carries is then down regulated on the development board to 3.3V, the logical reference for the MCU. While there’s nothing wrong with this in development, a prototype of your application ought not depend on drawing power from your computer. Instead, you could call on an AC-DC converter with USB output. This has the added benefit of likely raising the current capacity of the 5V power rail, from at most 500 mA via a computer to whatever the specification happens to be for a given wall-bound converter. This could be meaningful for driving certain power hungry actuators, like speakers.
    
\section{Battery}

While the above solution removes the need for the computer, you’re still tethered to a wall. To go fully mobile you’ll want to call on a battery. So what are our options? The idea of using a voltage regulator (specifically, the MPM3610) to cut down an input voltage to a nice, stable 3.3V level applies here as well. If you were to take a closer look at the linked datasheet, you’d find that the MPM3610 accepts input voltages from 4.5V to 21V. The 5V delivered over USB is within this range, and any compatible battery will need to be as well. This unfortunately eliminates the possibility of calling on single cell 3.7V lithium batteries, but makes the selection of a 9V alkaline battery fairly obvious.
    
    You might be wondering how any battery might connect to the boards in front of you, but never fear, we’ve got you covered. At one corner of the Tiny Machine Learning Shield you’ll find a green terminal block with silkscreen labels that read, “VIN” and “GND,” where GND is our reference voltage and as such should be connected to the negative terminal of any compatible battery. This green terminal block is where you’d want to screw in wires carrying 4.5V to 21V, and we’ll add that 9V clip, like this, that terminates in pre-stripped hook-up wire makes this quite easy!

\bigskip
    
    Assembly Steps

\begin{enumerate}
  \item Screw down a wire leading from the negative battery terminal (black) to GND (Most < 3mm flat-head screwdrivers will suffice here).
  \item Repeat this process for the positive battery terminal (red) to VIN. And that's it you're all set to power your Arduino from a battery!
\end{enumerate}


\bigskip

Important Notes
    
\begin{enumerate}
  \item While there is clever circuitry on board to handle such an exception, it is generally good practice to avoid having competing power sources, so we’d recommend that you unplug the Nano 33 BLE Sense from USB power before connecting a battery
    
  \item With about 550 mAh capacity, a 9V battery can source 15 mA for about 37 hours before you will need to get a new one.
\end{enumerate}
    
  \includegraphics[width=\textwidth]{TinyML/3 Development/batteryScrew.png}
    
    
    An image of screwing down the black negative terminal to attach a wire.

  \includegraphics[width=\textwidth]{TinyML/3 Development/batteryDone.png}
    
    An image of the attached battery powering the device.
    
    
%1.3: Embedded Hardware and Software    

\section{Embedded System}
Embedded System
0:00 / 0:00Klicken Sie auf die Pfeiltaste nach oben, um das Geschwindigkeitsmenu aufzurufen. Regeln Sie dann mit den Hoch und Runter Pfeiltasten die Geschwindigkeit. Klicken Sie Enter, um die jeweilige Geschwindigkeit festzulegen.Geschwindigkeit1.0xKlicke auf diesen Knopf und diess Video stumm oder nicht stumm zu schalten, oder drücke die HOCH oder RUNTER Knöpfe um die Lautstärke zu erhöhen oder abzusenken.Maximal Lautstärke.Videotranskript
Beginn des Transkriptes, zum Ende springen.
SPEAKER 1: Hi.
I hope you had a good time with Brian getting your TinyML system all set up
and ready to go.
In this video and the next couple of videos, what I want to do
is I want to talk to you a little bit about embedded systems at large.
And specifically, I want to help you understand how the embedded system
that I'm holding in my hand--
the one that you are going to be working with--
actually has the right building blocks that compare it
with a production grade system.
The reason I want to make this comparison
is because I want you to feel confident that you're actually
learning with a real embedded platform.
And the lessons that you pick up in this particular course
are actually something that are transferable
well outside of the course itself.
And that's important, because I do want you to be a TinyML
engineer when you finish the course.
So let's start off by talking about what is an embedded system.
I talked about this in course one.
Specifically, any given embedded system has got three fundamental building
blocks.
The first is the input stage, where you get input from the sensors.
The second is a digital phase, where you actually
process all of that analog signal and try and make some decisions about it.
And then the third stage is the actual output stage.
Typically in an embedded system, this might be analog output,
or it might be a digital output.
Let's take some real world examples to actually understand what these are.
For example, if you look at a Fitbit.
A Fitbit might be looking at the angle of rotation of the wrist--
based on IMU data that's coming in-- or it might be looking at the temperature,
or so forth.
And that data is then sent over to your processing system, which
then actually runs the routines that were engineered by the Fitbit engineers
in order to make some sort of decision and then act on it.
And that actuation might be in the form of actually generating output
to the screen, or it might be recording the data into the actual subsystem.
Let's take something that we find around houses, a thermostat.
A Nest thermostat, for instance, has a temperature sensor
that detects what the temperature is around it,
or it might also have a rotational sensor which
detects how much you're actually changing the dial so that it can,
in reaction to that, actually update the screen so that it tells you what
new setting you're reaching out for.
Now these are production grade systems.
And what I want to help do is actually help
you make the connection between these production grade systems
and the TinyML system in your hand.
Let's take an Echo Dot, since we are often talking about keyword spotting.
An Amazon Echo Dot reacts to Alexa.
How does that happen?
Well, if we take it apart, we see a set of different components in here.
In this teardown, we're seeing that the microphone array is on the PCB.
There are four different microphones here
that actually are the source of input signals into the computing subsystem.
These Texas Instruments microphones take the data
and hand it off to the processing system.
From a processing standpoint, you actually
see two different computing engines here on the two different boards
that are inside the Echo Dot.
This actually should remind you of the cascade architecture, which is
something I talked about in course two.
Now, if you don't remember cascade architectures,
I strongly encourage you to go back and look it up.
Cascade architectures are extremely critical for deploying keyword
spotting efficiently.
And the third stage is obviously the process, or the output stage.
Once you finish processing the data, you actually
actuate on that input process data.
So how does this compare with your embedded system?
On the left-hand side is the Amazon Echo Dot PCB.
On the right-hand side is your embedded system.
One of the key differences is that on the left-hand side,
the Amazon Echo Dot is specialized for one particular task.
It does one thing.
It just listens to analog input signals, processes it, and then
generates some sort of output on its speakers.
Your microcontroller actually has a bunch of different capabilities.
It's more general purpose.
That said, there are very strong similarities between the two of them,
specifically from the microcontroller standpoint.
All the limitations of memory compute and so forth that we've talked about
are actually very much identical.
So that should give you confidence that when
you are working with this microcontroller,
you are very much working on a production grade system.
So let's go a little bit deeper into the developer embedded system
and contrast it with the real embedded system.
The punch line is really this.
A real embedded system is designed to perform one particular task.
A developer embedded system has a lot more bells and whistles
because it is a developer platform and allows you to explore different things.
While you have the processor and the Bluetooth,
you have a microphone on your embedded system.
You have an IMU on it.
You have a temperature and humidity sensor.
And you also got a physical I/O port there so
that you can actually interface that embedded system with your computer.
You don't have this sort of flexibility and features--
general-purpose flexibility-- on the Echo Dot.
Why?
Because this is a developer system.
So you got a couple of extra bells and whistles, all of which
actually lead to an increase in the price of the board.
Now think about it.
An Echo Dot doesn't need to be general purpose.
It's got one functionality.
So you tend to strip all these things down and make it very focused
so you can keep the cost as low as possible,
because it just needs to perform the task it needs to perform.
Now, while we are talking about this one particular microcontroller board here,
there are many different embedded systems.
Later on I will explain why we have chosen the Arduino microcontroller
board.
For example, here is a Himax AI development board.
Much like the Arduino, it's a development board.
It's got a camera module on it, it's got a microcontroller on it,
and it's got an I/O subsystem that allows
you to interface with the computer.
Here's another example.
This is the SparkFun Edge 2, which has got an AMBIQ microcontroller, which
can also interface with your computer.
And it also has the exact same constraints of an embedded system.
Or, the ESP-EYE.
It's got a camera on it.
It's also got a small microcontroller.
And the form factors are all about the same size
as the Arduino microcontroller.
Why am I talking about all these different devices?
The reason is because these are all different kinds of embedded platforms.
And when you are actually going to work on these platforms,
you need to know a little bit about them so that you're actually
knowledgeable when something does not work and you know how to fix them.
All of these systems come with a wide range
of capabilities, specifically at the board level--
the embedded system board level.
What you find is that there are differences in the kind of sensors
they support.
For example, the Himax has an accelerometer, microphone,
and a camera, whereas the Arduino on the second row
has got a wide range of different sensors,
as I was already hinting earlier.
So these are some of the differences that we end up having.
In addition, there might be differences in the kind of communication
module [? some might ?] provide.
For example, the Himax does not provide a radio option,
whereas Arduino and the other devices have
the Bluetooth capability, which allows us to do the magic wand, for example.
In addition, the embedded system also has the microcontroller,
which has got both the processing engine and the memory all tightly coupled
together.
Now there are differences in there, as well.
As you can see, the clock speeds go all the way from 400 megahertz
to as low as 48 megahertz on the SparkFun Edge 2.
That said, they are roughly in the same ballpark,
in the order that they're all about megahertz, not gigahertz.
As for the memory, they are also very limited.
You have very little bit of flash, you have very little bit of RAM.
So these different embedded systems are roughly the same
in terms of the constraints, though they have different levels of flexibility.
We happened to choose the Arduino microcontroller because it gives you
a lot of different sensors, which allows you to expand well
outside the scope of the course, while also
having the right set of constraints.
It gives you that practical hands-on feel
about dealing with the microcontroller.
With that said, in the next video, I'm going
to go a little deeper into understanding the characteristics
of the microcontroller itself.
And after that, I'll talk a little bit about the software and then
I'll kind of bring it all together so you--
at the end of the day-- can feel confident
that you are, in fact, working with what is a production grade
microcontroller that's actually found in real world devices.
Ende des Transkriptes, zum Anfang springen
Downloads und Transskripte
Video
Videodatei herunterladen
Transskripte
Download SubRip (.srt) file
Download Text (.txt) file

\section{Diversity of Embedded Systems}

\subsection{Embedded Systems}

In this reading, we will explore the diversity of embedded systems.
    
If you take a few minutes to look around the room you’re in, you’ll no doubt recognize the proliferation of embedded hardware within our world. Embedded systems are increasingly ubiquitous because they serve to complete specific computational tasks within their environment, allowing us to not only collect all sorts of data from distributed sensors but also to process the resulting information locally and act accordingly. That we can accomplish this in situ, or within the environment that this hardware is embedded within, at a fraction of the cost and physical scale of most general-purpose computing hardware, opens the door to many advancements in automation and generally to the creation of smart, connected things.
    
    While the specifications and capabilities of a given embedded system ought to be tailored for its application, in general, we can look to the common construct for an embedded system that links sensing to processing and later actuation through the process of transduction, the conversion of one form of energy into another, for throughlines. In sense, we convert energy from a physical phenomenon into an electrical signal we can go on to digitize and compute with. In actuation, the converse. In this course, we’ll focus predominantly on the central element of said construct: the computing hardware, and review specifications and technical considerations that vary between implementations of embedded systems.

  \includegraphics[width=\textwidth]{TinyML/3 Development/pasted_image_IMU}
    
    An image of the Arduino Nano 33 BLE Sense noting the location of the IMU, USB port, Microphone, temperature + humidity sensors, and processor + bluetooth module.
    
\subsection{Development Boards}

 In the context of this course, we have talked about and deployed a microcontroller unit (or MCU) development board from Arduino, the Nano 33 BLE Sense. Here, we distinguish the nRF52840 MCU from the development board it is soldered to. What’s interesting about this board is that despite its relatively small size, it is jam-packed with many of the representative elements of an embedded system, all on one printed circuit board (PCB): a collection of surface-mount sensors, a form of an actuator in the on-board programmable RGB LED, and an integrated MCU-BLE module, where BLE is an acronym for a branch of the Bluetooth standard called Bluetooth Low Energy. Further still, the PCB that the MCU is soldered to serves to ‘break out’ additional IO from the controller to provide interfaces for external, off-board modules and connections to daughter boards. So, while in many ways, the Nano 33 BLE Sense is representative of an embedded system, it is separated from commercial implementations in that it serves no particular purpose, other than the open-ended development of an application that is. Put another way, a development board has the advantage of enabling many potential use cases. Further, you don’t have to go through the process of designing circuitry, capturing schematics, or physical layout for such a board to get started prototyping a concept that could go on to be manufactured with specific intent later. The tradeoffs you make in employing a development board often involve board size and cost.
    
\subsection{Board Size}

    Board size, and shape, have fairly meaningful implications for use in the field, given these simple parameters constrain where and how a system can be deployed. Some environments or contexts within which we’d like to embed a system are forgiving (like the Google Home, say), but other manifestations (smart glasses, for example) depend wholly on the requisite hardware being quite tiny and perhaps of unique form. Somewhat obviously, board size is determined by the number and physical size, or package, of the components that live upon it. At a minimum, an embedded system must include a MCU chip alongside a power source, often a lithium polymer battery, and power circuitry. The nRF52840 MCU on the Nano 33 BLE Sense lives within a MCU-BLE module, the U-Blox NINA-B306, which spans 10 by 15 mm, dimensions that include a trace antenna. The MPM3610 step-down converter used to down-regulate the 5V delivered to the board over USB occupies 3 by 5 mm, alongside small SMD passives. The entire board, meanwhile, spans 18 by 45 mm. While impressive, clearly a purpose-designed PCB could remove some of the sensors and IO breakout unnecessary for a specific application to reduce the overall scale even further, perhaps by about 60 to 70%.
    
\subsection{Cost}

    Unsurprisingly, the cost of a development board scales with its MCU’s compute capability and general feature set. In selecting the appropriate MCU for an application, it is important to remember that MCUs are often deployed to complete specific tasks, where the complexity they will face is predetermined or constrained. For very simple tasks, we highlight the ATTINY85 (an 8-bit processor with 8 kB of flash) that can, depending on the package selected, cost well less than \$1 EA and even less at scale, perfect for simple computing requirements. In the context of TinyML, the AI-capable NINA-B306 module featuring the Nordic Semiconductor nRF52840 chip (an ARM Cortex M4) and all-in-one BLE hardware (including antenna) costs about \$10 EA and is more generally representative. In view of the on-board sensors and design costs, the Arduino Nano 33 BLE Sense costs just over \$30, not even an hour's time of an electrical engineer who might design such a board — a tremendous value. In general, the boards we originally considered for this class span from \$2 (BluePill) to \$54 (Disco-F746NG).
    
    You can find a list of the boards we considered in the table below. Ultimately, our staff selected the Arduino Nano 33 BLE Sense for its versatility, in providing a large selection of on-board sensors, accessible IO via breakout pins, and a BLE module for projects that involve wireless communication, with reasonably representative compute specifications for resource-constrained hardware. We’ll explore the compute specifications in a bit more detail in the next video and reading.
    
\subsection{TinyML Development Board Comparison   - officially TFLM supported   - unofficially compatible boards}

\begin{tabular}{llll llll}
    Board &	MCU &	CPU	& Clock &	Memory&	IO	& Sensor(s)&	Radio\\
    Arduino Nano 33 BLE Sense &	Nordic nRF52840 &	32-bit ARM Cortex-M4F &	64 MHz&	1 MB flash 256 kB RAM	& x8 12-bit ADCs x14 DIO UART, I2C, SPI &	Mic, IMU, temp, humidity, gesture, pressure, proximity, brightness, color	& BLE \\
    Espressif ESP32-DevKitC	& ESP32 D0WDQ6 &	32-bit, 2-core Xtensa LX6&	240 MHz&	4 MB flash 520 kB RAM &	x18 12-bit ADCs x34 DIO** UART, I2C, SPI &	Hall effect, capacitive touch***&	WiFi, BLE\\
    Espressif EYE&	ESP32 D0WD&	32-bit, 2-core Xtensa LX6&	240 MHz&	4 MB flash* 520 kB RAM &	SPI via surface pads&	Mic, camera&	WiFi, BLE\\
    Teensy 4.0&	NXP iMXRT1062&	32-bit ARM Cortex-M7&	600 MHz&	2 MB flash 1 MB RAM&	x14 10-bit ADCs x40 DIO** UART, I2C, SPI&	Internal temperature, capacitive touch&	None\\
    MAX32630FTHR&	Maxim MAX32620	&32-bit ARM Cortex-M4F&	96 MHz&	2 MB flash 512 kB RAM&	x4 10-bit ADCs x16 DIO UART, I2C, SPI&	Accelerometer, gyroscope&	BLE\\
\end{tabular}
    
    *This board also features 4 MB flash and 8 MB of PSRAM external to the MCU, 
    
    **shared programmable functions, 
    
    ***with external touchpads 
    
\begin{tabular}{llll llll}
    Board	&ASIC&	DSP	&Clock&	Memory&	IO&	Sensor(s)	& Radio\\
    Himax WiseEye WE-I Plus EVB&	HX6537-A&	32-bit ARC EM9D DSP&	400 MHz&	2 MB flash 2 MB RAM&	x3 DIO I2C&	Mic, accelerometer, camera	& None\\
\end{tabular}

    We include the Himax WiseEye as an officially supported example of hardware optimization for TensorFlow Lite that calls on an application-specific integrated circuit (ASIC).
    
\section{Embedded Computing Hardware}

Embedded Computing Hardware

VIJAY JANAPA REDDI: Hi, in this video I want to talk to you a little bit about the embedded computing hardware, that you find on an embedded system. Recall, that there are many different kinds of developer embedded systems out there.
And these developer embedded systems, come with very different kinds of capabilities in terms of sensors, the communication modules like the radio and so forth.
Similarly, they also come with different types of microcontrollers on them.
We are specifically going to focus on the Arduino subset.
Specifically, we are going to focus on the capabilities of the nano 33 BLE sense.
We're going to try and understand what its computing capabilities are, with respect to the general class for application processors that are out there.
The reason I want to touch on this, is because I want you to understand what are the pros and cons of what are the trade offs that we're making here, as we are going down to TinyML.
By and large, when you look at the class of processors that are out there, there are many different types of processors.
I'm going to talk specifically about the ARM Cortex profile.
The ARM Cortex profile of processors, are widely used for mobile applications, ranging from smartphones, all the way down to small little microcontrollers that we find everywhere.
So, by and large, the're three tiers I would say.
The application of processors that you find in things like your smartphone for instance, that are capable of running complicated tasks.
Then, you have real-time processors, that you tend to find for example in automotive vehicles.
And then you have microcontrollers, that are pervasive where you tend to find everywhere.
Now, all of these tiers come with many different variants in each tier.
For example, the application processors go from A5, A9 to A65 and so forth.
The real time processors have, again, a different tier or class of processors in there.
And Similarly, in the microcontroller space, you have a different set of tiers.
One of the reasons that you tend to have these tiers of processors, is because really you're trying to shape power off.
For example, as we go from application processors that consume anywhere between like 3 to 4 watts, when you go down to microcontrollers, these consume something in the order of a few milliwatts.
And what's the trade-off?
Well, you simplify the functionality.
By renewing functionality, by limiting their ability, you actually lower the power down.
Why do we care about having low power?
Well, it's because it's TinyML.
We want TinyML machine learning to be running at all times.
And we want to do that on battery based devices, things like a coins or battery.
And therefore we are trying to focus on microcontrollers.
That's why in this particular course, we're going to use small O, M class microcontroller processors.
Now, in addition to just sort of understanding that piece of the lay of the land, I also want you to understand what are the differences within the microcontroller space itself.
Even in the microcontroller space, there are extremely low and simple microcontrollers, like the Cortex M0 that we're looking at here, which are capable of doing General Data processing and general input output control tasks.
Then, you might have the M3 for example, which does a bit more advanced data processing because it has instruction set architecture or instruction support for things like bit manipulations, which allow you to do arithmetic a bit more efficiently.
Then you have something like the M4 that we will be using, which has what we call as the floating point unit, which allows us to do floating point arithmetic on device.
Which is a great thing to have, especially when we're dealing with things like keyword spotting and other types of applications, where we might have to do a little bit of complicated math.
It's also got capabilities, it's what we call the DSP or the Digital Signal Processor, where it can actually execute multiple instructions all as one instruction.
Some of these details are reserved more for embedded systems class, but the core thing that I'm trying to get to here, is that even when you pick a small microcontroller class, there are many different variants in there and you need to be conscious about that.
Because, that choice of microcontroller is going to tremendously influence how you are actually going to be able to build your application.
For example, if we didn't have the FPU, the Floating Point Unit, we would not be able to do floating point arithmetic.
Which might actually limit some of the things that we would be able to do, in terms of application and design.
So, keep some of these details in mind, because it's extremely important to have an appreciation for how
our applications will work.
I will come back and touch on some of these concepts, as I go through the application examples in the code itself later on.
So bottom line what I'm trying to say here, is that this small little microcontroller that we have chosen has got all the fundamental nuggets being limited in terms of its performance and memory, because of the power efficiency that we get from it.
And it's something that you actually find for example in devices, like the Amazon's Echo Dot that we did a Teradyne of earlier.
Now, in addition to this, there are a couple of other things that we need to touch on, such as what is the software stack that actually enables these embedded systems to run code effectively.
And we need to understand that as well and that's where I'll be taking your next.
And I'll see you there.


\section{Diversity of Embedded Microcontrollers}

\subsection{Embedded Microcontrollers}
In this reading, we will explore the diversity of embedded microcontrollers.

Microcontrollers sit at the bottom of the computing hardware hierarchy, at least in terms of computational capability. You might be wondering what differentiates an MCU from the CPU in the machine in front of you. Well, in some ways an MCU and CPU are quite similar, namely because an MCU is the integration of a low-performance CPU and other peripherals within a single chip. In this way, MCUs are also similar to, albeit less sophisticated than, a System-on-a-Chip (SoC). Typical peripherals integrated into an MCU include memory, analog-to-digital converters (ADC), digital-to-analog converters (DAC), timers, counters, general-purpose IO (GPIO), pulse-width modulation (PWM), direct memory access (DMA), interrupt managers, and serial protocol controllers. We can look at this integration and draw two conclusions: MCUs emphasize connections to the environments they live within and require that all on-board hardware be as tightly packaged as possible. Much like a CPU is seated upon a motherboard, MCUs are chips with packages soldered to much smaller PCBs that, at a minimum, must have supporting power circuitry and a programming interface. The latter prerequisite can be avoided if the MCU is programmed prior to assembly, but this is unusual, at least in development. Beyond this, the specifics of a custom embedded system are contingent upon the needs of the application. In leveraging a development board, you’ll want to center your selection on the meeting, if not exceeding, the same requirements. We’ll walk you through a comparison of MCU development boards and specifications in a moment.

To return to the computing hardware hierarchy: MCUs are clearly outstripped by the performance of your average personal computer, as well as computing clusters. They do fall somewhat adjacent to, if not arguably short of, what we’ll call ‘intermediate’ computing in the form of single-board computers (SBCs), like the Raspberry Pi. SBCs and MCU development boards are similar enough conceptually and in terms of computing power but are differentiated in that SBCs aim to be low-power, single board computers with proper operating systems and IO typical of personal computers, but often lack the peripherals we’ve listed above, making them ill-suited in the context of embedded systems design. Looking at the computing hierarchy we’ve put forward (clusters > computers > intermediate computing > microcontrollers), you may wonder why exactly so much emphasis is placed on developing advanced software solutions for MCUs, like deploying deep learning models at the edge. Conventional wisdom would suggest that increasingly powerful computing hardware is required for and enables increasingly complex software applications. In recent years, however, there is an incentive for us to develop solutions that allow microcontrollers to take on such complex tasks to realize visions of distributed sensing and computation, where we need to strike a balance between cost (which can beget ubiquity), performance, and power efficiency. To put a name to such an incentive: we can compare the power required to transmit a stream of data from some remote controller to a server uplink and find that this energy requirement greatly outstrips the power we would need to perform a similar analysis on board. Further, there are complicated ethical and policy issues enveloping the deployment of distributed sensors that capture and then share streams of possibly identifying information, rather than de-identified reports of model returns. As such, we have seen MCUs take on increasingly complex tasks through the clever implementation of software, coupled with ever-advancing architectures that enable more powerful processing in even tinier packages.

\subsection{Compute Capability}

Microcontroller units have varied CPU architectures. Of these, chipmakers can license the 32-bit ARM Cortex M architecture, making modifications from this common starting point to introduce features and differentiate their products. For instance, the Nano 33 BLE Sense calls on a U-Blox NINA-B306, which is a stand-alone MCU-BLE module, that integrates the nRF52840 MCU from Nordic Semiconductor, which licensed the underlying architecture from ARM, as the nRF52840 is an ARM Cortex-M4. While this is generally representative of the modern controller, it’s worth noting that there is a great deal of variation in compute capability, with processing cores that step all the way down to 8-bits, carrying modern relevance in simple applications.

As with other computing hardware, a faster MCU clock enables greater temporal performance, at the cost of diminishing power efficiency. Typical clock speeds sit between 50 and 500 MHz. For example, the Nano 33 BLE Sense is clocked at 64 MHz. We’ve deliberately said ‘MCU clock,’ here, rather than CPU clock, since the core clock of an MCU enables functions beyond just the CPU and trickles down to other peripherals (ADCs, DACs, et cetera) via scalars so that a faster clock can not only enable faster processing but also create headroom (bandwidth) for non-compute functions as well. For instance, while a less than 0.5 GHz CPU clock might not sound impressive, this can translate to theoretical sample intervals for analog-to-digital conversion on the order of tens of nanoseconds, which is more than sufficient for most physical phenomena of interest. Furthermore, lower speeds may actually be beneficial, somewhat counterintuitively, for timing longer intervals accurately. 

Ultimately, the best processing core for your application will depend on the complexity of the task, its performance requirement, and temporal dynamics. One thing that MCUs are known for is their role in facilitating real or near real-time responses to events via interrupts and event systems. The latency of these responses will of course be tied to the performance specifications for the controller. Another thing that MCUs are known for is lying largely dormant in the environment until they are called upon. Dynamic clocking can ensure that an MCU sips as little current as possible until an interrupt wakes the MCU from this low-power mode to respond to whichever event. Consider a TV remote that sits at home while you’re at work that can respond to a button press, wake its controller, and begin emitting the necessary sequence of electromagnetic pulses from its IR transmitter. The ability to adjust the system clock to reduce power consumption is a powerful technical feature — pun intended.

\subsection{Memory \& Storage}

An important consideration to make when approaching MCU memory and storage for the first time is that the scales you may be used to working with on computers will likely not apply. Microcontrollers call on flash as a non-volatile option for program memory so that the very same program and initialization will occur at each reset. Typically, MCUs feature about 1 MB of program memory. Microcontrollers call on random access memory (RAM) to store working data, but this is of course volatile, so data from one power cycle cannot persist to the next. In the context of this course, it is interesting to consider the size of various models. One of the primary constraints for TinyML is ensuring that MCUs, as resource-constrained hardware, can fit a desired model in the first place. Compact speech recognition models, for instance, require 20-30 kB, whereas more complex models like those required for visual feature detection, require hundreds of kilobytes, really at the extreme of what most MCUs provide.

We want to highlight here that some microcontroller boards interface with SD cards as a way to extend the storage capacity of their system as well as to create a non-volatile record of data, for applications that require it. For what it’s worth, EEPROM is another form of on-board or on-chip memory that some embedded applications (smart cards, for example) leverage to store programmable, persistent data.

In the end, you should consider the scale of your application and the role that volatility might play in limiting or enabling the function you require.

\section{Embedded I/O}


VIJAY JANAPA REDDI: Hi.
One of the key things that we often deal with when we are dealing with embedded systems is that we tend to interface them with different types of sensors.
Now, if you remember, when I talked about the Tiny Machine Learning shield, which is a daughter PCB design with breakout for input/output pins that allow you to interface the Nano 33 with many different kinds of external sensors, for example, like the camera sensor or an LED and so forth.
Now, this particular TinyML shield has got a set of pins.
If we recall, on the left-hand side, we have general-purpose input/output pins that allow you to interface with digital and analog components.
And, on the right-hand side, we have pins that interface with I2C protocol-based sensors.
Now, I want to talk a little bit about this because, on the TinyML shield, we are, in fact, using these capabilities.
And so I want you to understand how we're actually interfacing them so that, later on, when you expand yourself outside the scope of whatever we're teaching, you will know how to integrate new sensors.
And, for example, in the I2C, you can actually integrate a Seeed Grove sensor, like the thermocouple that we're showing here, using a simple cable.

Now, there are many other types of sensors you might be able to interface-- thermocouple, a Coulomb counter, a rotary encoder.
Or pick your favorite one from the website shown below.
The key thing to understand is how these sensors are actually interfacing.
I'm not going to be able to go really deep in here because this is not an embedded systems class, but, that said, I want to give you just enough of the pointers in terms of interfacing these embedded components, just as we have done with the camera module in here, so that you actually have a little bit of inclination on what's going on.
And then the readings will supplement some of this material.
And we will give additional pointers so you can build outward from there.
The key thing is that, when you're interfacing these sensors, you have to interface them through some sort of communication module.
For example, here is an old-school parallel interface, as we used to call it, where you have a transmitter on the left-hand side, and you have a receiver.
Think of it as the transmitter being the camera, for instance, which is picking up data, the pixel data, and it's transmitting that over to the receiver, which is the Arduino Nano 33 BLE sensor.
It's got to receive data.
So you can have parallel interfaces like this, or you can have serial interfaces like this.
In parallel interfaces, you tend to have a whole bunch of wires that are going from the transmitter to the receiver. 
Now, the poor thing about this is that it uses up a lot of space and area.
The nice thing about a serial interface, like we're showing here, is that all of that data is communicated effectively over one line in some sort of intelligent sequence that is predetermined.
And there are many different types of communication protocols out there.
The three primary ones you need to know about when you're interfacing with sensors is generally the UART, or the Universal Asynchronous Receive Transmit protocol, or the I2C, the Inter-Integrated Circuit protocol, or SPI, the Serial Peripheral Interface.
These are all serial communication protocols.
In other words, you have a single line through which you're effectively communicating from between the transmitter and the receiver.
Now, the readings will supplement the knowledge that I have here, but I'm going to specifically focus on the I2C protocol because that's one of the most common ways to interface with sensors.
And, in fact, it's one of the things that we are actually using.
So I want you to understand a little bit about what's actually going on underneath the hood.
So an I2C protocol basically has two critical ways of communication channels.
One is called an SDA, which is the serial data channel where data actually moves, and the serial clock channel, which actually controls when the data is actually coming about.
These two wires are already sort of there.
You just need to know that they exist and what their role is because often, when you're integrating sensors, you're dealing with these two pins.
The nice thing about these pins that you can actually connect different sensors on this.
So you can take the Nano 33, which, for example, will be your mother down here, and connect it to the camera, which might be your son here.
Typically, these are called master-slave circuits, but I prefer not to use that terminology to be progressive.
Now, let's be a bit more specific.
What happens here is that the mother, like the Nano 33, is able to transmit data.
And the son is able to receive the data.
On the other hand, the son might be able to transmit the data over to the mother.
Being a bit more specific, in our case, for example, the Nano 33 BLE is the mother, which is sort of the brains.
And the sensors are like the camera module or some other Seeed Grove sensor that you might choose to interface.
Basically, all three components or n number of these different sensors that you might have connect over to the SDA and SCL.
And, by connecting them and then programmatically setting up a mechanism to communicate between them, you can get data in from the sensor, or you can push data to another device.
So just a very, very high-level skim of what's actually going on in terms of how a sensor communicates with the Nano 33 BLE sensor.
There's often a start and a stop.
Think of it as a handshake or saying hi and bye.
So, at the start, you have a special signal where you can see that the SDA line goes low, and the SCL clock signal stays high.
When that happens that basically indicates to the device that you're about to start transmitting data.
And the first thing you will do is you kick off that start signal, and then you start transmitting your bits.
Now, the first thing that happens is that you're actually telling what address you want to send to.
So, for example, let's say you want to pull data in from the IMU.
Now, an IMU might have an accelerometer.
It might have a gyroscope.
And you want to know which one of those do you actually want to read.
So, in that particular case, you will give the address to either one of them.
And then you will say whether you want to read or write.
In our case, for example, we might want to read the data.
So you will say, I want to read the data.
And then, after that, you will get an acknowledgment saying, yes, things are ready to be read.
And then the data starts coming in from the sensor.
And the data comes in a chunk of 1 byte, which is 8 bits of data.
Every time 8 bits of data are passed back, you get an acknowledgment.
And then this process repeats until there's a stop signal.
The readings are going to go into much more detail.
I'm simply trying to pick out the key points that I want you to know.
At this point, I just want you to know that we're trying to understand that there's an I2C protocol that has some way of communicating with the sensor where the sensor might actually have different kinds of pieces of data that you might want to get to, like the accelerometer or the gyroscope data and so forth.
So, therefore, you have an address.
And, once you program it, you are able to read and write from it using I2C.
How does this pertain to what we are doing in our class?
Well, the OV7675 camera module that we are interfacing with, for example, uses
a variant of an I2C.
Remember, I just said a variant.
That's because the OV module has a proprietary I2C interface.
But, by and large, you will find that, if you look at the pins,
the pins are the SDCs and the SCLs that we talked about earlier.
And then you tend to have some additional pins.
And those pins are basically transmitting pixel data back to the Arduino Nano 33.
And, of course, then you have the power pins that actually need to provide power to the OV module.
Just the way the OV module interfaces through an I2C interface onto the Nano 33 BLE, the IMU that is on the board that is already physically integrated onto the board also communicates with the processor using I2C.
Now, you are starting to understand why I'm actually emphasizing the need to understand these protocols.
Later on, when you're integrating things, you will need to know what type of sensor you're actually dealing with.
So this kit that we have put together for you is basically using fundamental embedded system principles, which will be common to no matter what device you pick up and you play with.
Hopefully, at this point, you have an idea about why we chose this particular microcontroller and specifically why we chose this particular embedded system platform and why or how the different sensors that we have provided you interface with the board itself.
There is much more to this.
And this class unfortunately cannot do due diligence in explaining every single thing around embedded systems [INAUDIBLE]..
But what I'm trying to hint to you is to keep in mind the language that I'm getting you exposed to.
And, once you kind of know a couple of those terms, then you'll be able to go on and build outward from that particular point.
With that said, in the next two videos, I'm going to focus very much on the software side of how we get the embedded system to actually do what we want it to do.
And I'll see you there.

\section{Transducer Modules and Wireless Communication}

\subsection{Transducer Modules \& Communication}

    While some of the TensorFlow Lite Micro capable development boards feature onboard sensors, others do not and you may find yourself interested in connecting an external sensor or actuator transducer module. Further, your use case may require you to periodically transmit results to other devices using wireless communication protocols like Bluetooth. In this reading, we briefly discuss each of these topics at a high level, leaving more thorough explanations to external readings you can complete at your discretion.
    
\subsection{Sensor and Actuator Modules}

    When you write code to compute the average of a set of variables, you rarely if ever add the terms and divide by their number. Instead, you call on a pre-built average function. Much in the same way, many of the desired circuits required to achieve sensing and actuation of various kinds have been commoditized to the point where it isn’t a question of designing these circuits again from first principles but interfacing the modules with an MCU using an API mindset if you will. The inputs and outputs of various modules are our concern.
    
    We’ll turn to consider possible interfaces in just a moment, but first, we want to highlight that there is a catalog of largely plug-and-play sensors and actuators that will work with the Arduino Tiny Machine Learning Kit, recalling that the role that each category plays in a canonical embedded system (sense → process → actuate), converting energy from and to a physical phenomenon, respectively, in a process called transduction.
    
    The modules we mention are available from SeeedStudio, given their compatibility with the Grove connectors on the Tiny Machine Learning Shield and are as such the easiest way to extend the functionality of your system, but there’s no reason why you couldn’t leverage a solderless breadboard and some hookup wire to unlock compatibility with an enormous number of available modules. In the US, we like to source such modules from Adafruit and SparkFun, but they are widely available if you do some digging.
    
\subsection{IO \& Serial Protocols}

    A microcontroller has three main types of IO: digital, analog, and serial, listed in order of relative complexity. The foremost is very simple, as digital pins have only two states, high (or equivalently 1 or true) and low (or 0, false), and can be written to as outputs or read from as inputs. For clarity, these binary states are physically conveyed in voltage extremes, if you will: either 0V with respect to circuit “ground” or the full-extent of the logical reference voltage for the given microcontroller, for the nRF52840, this is 3.3V.
    
    Analog inputs are a bit more complex in that they seek to translate a continuous representation of voltage in time and amplitude into a binary representation, that requires discretization and quantization, respectively. For more on analog-to-digital conversion (ADC) and digital-to-analog conversion (DAC), including approximations made by pulse-width modulation (PWM) for slow-response loads or circumstances, you can read on at the links above.
    
    Finally, we turn our attention to serial communication protocols, of which we’ll highlight three: universal asynchronous receive transmit (UART), inter-integrated circuit (I2C), and serial peripheral interface (SPI). We’ve provided a brief appendix document describing them in more detail and note at a high level that the various serial protocols are designed for transmitting longer data messages with different applications in mind (e.g., one to one vs. one to many communication). We also note that variants of the I2C protocol are used by both the IMU on your Arduino Nano 33 BLE Sense and between the Arduino and the OV7675 camera module.
    
\subsection{Efficient Wireless Communication}

    A major incentive for deploying deep learning models locally within an embedded system is to avoid the energetic cost of transmitting a constant stream of data from a sensor via some wireless communication interface, like WiFi or Bluetooth. Indeed, in most cases, it is favorable to compute a result that you might later transmit periodically, cutting down on the overall energy expenditure. This paradigm fits well within the dogma that led to the division of Bluetooth specification in 2011 and the creation of a new fork dedicated to transmitting information as needed, rather than in a continuous stream, Bluetooth Low Energy (BLE).
    
    There is a BLE module, the U-Blox NINA B306, on your Nano 33 BLE Sense development board, which we will use in the Magic Wand example later in this course! If you’d like to learn more about BLE you can check out this short appendix document we created.
    
    
    
\section{Embedded System Software}

Embedded System Software
0:00 / 0:00Klicken Sie auf die Pfeiltaste nach oben, um das Geschwindigkeitsmenu aufzurufen. Regeln Sie dann mit den Hoch und Runter Pfeiltasten die Geschwindigkeit. Klicken Sie Enter, um die jeweilige Geschwindigkeit festzulegen.Geschwindigkeit1.0xKlicke auf diesen Knopf und diess Video stumm oder nicht stumm zu schalten, oder drücke die HOCH oder RUNTER Knöpfe um die Lautstärke zu erhöhen oder abzusenken.Maximal Lautstärke.Videotranskript
Beginn des Transkriptes, zum Ende springen.
VIJAY JANAPA REDDI: Hi.
So far, we've been talking about the embedded system hardware.
What I want to do in this video and the next one
is talk a little bit about the embedded system software.
Specifically I want to talk about the system software,
and then I want to talk a little bit about the ML software.
Now I'm going to focus very much on just the system side of things,
I'm going to just kind of give you a little bit of a preview about what's
kind of wrapping the whole embedded hardware in order
to make it functional.
As with any system if you look at it, there
are two core components in any computing system.
There's a hardware, there's a software.
If you open up the software black box, there's usually an operating system,
there's libraries, and then you have your application code.
And now specifically if you look at the operating systems,
by and large you are used to running things on Linux, or Mac, or Windows,
and you're probably using one of these operating systems to watch this video.
Or perhaps you're on your phone right now watching it on a mobile OS,
like the iOS or Android based device.
And the question is, what happens when you go to an embedded system?
Embedded systems don't have those kinds of rich operating system capabilities.
In fact, they're quite lean and mean.
So things like the Arm Mbed OS, something that we are actually
going to be using without you having to get
into the details too much of course, or things like the FreeRTOS.
These are very lean and mean embedded system operating systems.
Quite often you don't have an operating system,
but sometimes you might have an operating system
because operating systems raise the level of abstraction,
so it's much more easier for users to actually program and interface
with the device.
So when you are dealing with your microcontroller,
by and large, when you're programming it you're
going to be writing code in the Arduino processing IDE.
This is going to be your bread and butter for all the examples
that we're going to do.
However, I do want you to understand what's in that software box.
So for example here, on top of the Arduino hardware,
which is a software ecosystem that's cooked with the Mbed operating
system at the base layer.
The operating system provides a lot of functionalities, as I'll show you.
On top of that, we have the Arduino, which
provides a set of functions that hide a lot of the complexity of the operating
system.
And then, you use the Arduino APIs to interface with the sensors,
and so forth to build your TensorFlow Lite Micro application.
So when you are writing a code, this is what you will be dealing with.
You'll be using the Arduino IDE and writing your TensorFlow Lite Micro
application using TensorFlow Lite Micro specific APIs calls.
We'll go into that much more detail.
But all of that is being orchestrated by some sort of low level code,
for example, the Mbed OS has got digital interface drivers,
it's got the ability to interface with the analog and digital pins,
it's got timers, and it's got a hardware abstraction layer.
So there's a lot of complexity in terms of actually running a microcontroller,
all of which is hidden away from you.
Why am I talking about all of this?
Two reasons.
One, I don't like learners to just blindly be using something.
I want you to understand the kind of ecosystem
you are sitting in so that you know how to navigate that space as challenges
come about, so you can actually build on top of the knowledge
that I'm sharing with you.
That's one.
The second reason I'm talking about all of this
very much is from the resource constraints perspective.
Now remember, memory is extremely constrained on our embedded systems.
For example, if you just look at the raw amount of memory that you have,
it's only about 1 megabyte of flash storage
and about 256 kilobytes of RAM on this particular device.
So therefore we have to be very careful about how memory is being used.
Now of course when you're developing neural network models,
you are thinking about the model size, because of course,
you're in a TinyML course.
In course two we talked a lot about the model sizes.
We did quantization in order to optimize a model from an FP32
to get it down to INT8, because we said that's a 4x reduction in the model
size.
Well that's great, but in this course that's not enough.
You need to know more, and learn how to optimize even more.
For instance, in order to run the model you've
got an interpreter, or TF Lite Micro, that actual TensorFlow Lite Micro.
Well, that's going to need a little bit of memory.
Not only that, then you've got the application code
that I just gave you a screenshot of earlier, through which you
will actually be writing and interfacing with the TensorFlow Lite Micro
interpreter, which then runs the model.
Now all of that is going to need some variables and temporary state, which
is going to take up memory.
In addition to all of this, you've got sensors which
also need to store data someplace.
So all of that might actually eat up more memory.
Keep in mind though that the numbers that I'm
showing you here a hypothetical.
They're just illustrative, but they get the message across to you
that you need to be quite aware and conscious about where memory is really
going.
And you need to keep that in mind because if you're not
conscious about it, given that we have such limited resources,
you can very quickly run out of memory.
So let me give you a real world example.
Mbed OS is an operating system, and over time it's been evolving.
So here's Mbed OS 2, and I'm showing you in the two bars
how much of Flash memory we are consuming,
and how much of static RAM that the operating system
itself is consuming when it's running.
Now this is Mbed OS 2, now on the right hand side, I'm showing you Mbed OS 5,
it's an improved version of the Mbed OS.
And what do you see?
In both the static RAM and the flash, you
see nearly a doubling of the requirements.
For example, in static RAM, the amount of RAM that is being used
has gone from five kilobytes to 13 kilobytes.
That's more than twice the amount of memory that has been needed by the OS.
Now that might not seem very big, but remember, at least on the Arduino Nano
33 BLE Sense, you only have 256 kilobytes.
And this is just the operating system that's taking up 13 kilobytes.
Now if you look at the Flash, which is what you usually
use to store your application on, you can
see that the operating system version 2 went from 38 kilobytes to 57 kilobytes.
It's a 1.5x increase in OS 5.
The point I'm trying to make is that in addition to focusing on the model,
you also need to be thinking about the system.
And that's why I want you to understand how the applications are
built and so forth.
Now all that said, we will be focusing on learning how to optimize
the whole system throughout the course.
How can I reduce this, for instance?
One of the key things you have to keep in mind is that as OS 5 for instance,
is growing, you need to understand why it's growing
and what can you do to simplify things.
So let me give you a real world example of what Arm actually
did to simplify things.
Now the Mbed OS has a lot of different things that it comes with.
It gives you a lot of flexibility and capability
to operate the device in a very easy and efficient way.
But in doing that, you have a lot of different features and capabilities
that add up, that take up a lot of Flash and SRAM.
And you want to detect whether you really need them.
For example, the features and frameworks folder
is actually something that's actually not needed.
There are many unnecessary test tools that are built into the actual Mbed OS.
Do you really need it?
Well it takes up one kilobyte of RAM and it takes about eight kilobytes
of Flash, which is very precious.
So if you were to take this out you can save a fair bit of memory.
So for example here, now we have gone down in terms of the static RAM usage
and in terms of the Flash.
We've actually cut down our flash usage from 57 kilobytes to just 49 kilobytes.
And in the RAM case, we've gone down from 13 kilobytes 12 kilobytes.
That's still important.
Point is really that TinyML applications are more than just the TinyML models.
You've got to think about the system stack, because all of that
is going to influence how well we can actually
run the end to end TinyML application.
Having a TinyML little model will not get you anywhere.
Having a TinyML application mindset, that's going to get you somewhere.
So in the next video, let's push a little bit
further in terms of understanding the software ecosystem
before we start building our applications out.
I'll see you there.
Ende des Transkriptes, zum Anfang springen
Downloads und Transskripte
Video
Videodatei herunterladen
Transskripte
Download SubRip (.srt) file
Download Text (.txt) file 

\section{Arduino Cores, Frameworks, mbedOS, and ‘Bare Metal’}

In working with the Nano 33 BLE Sense, you have interacted with some of the underlying layers between your application and the board’s physical hardware. Here, we quickly define some key terms (Arduino core, embedded frameworks, mbedOS, RTOS, and bare-metal) before discussing the implications of these concepts to your application.

\subsection{What is an Arduino Core?}

An Arduino “core” is central to a particular microcontroller’s compatibility with the Arduino framework. You can think of it as a low-level software API for a specific chip or family, like AVR or SAMD. Since each microcontroller has its own architecture, the concomitant core acts as an abstraction layer between your application and the physical hardware. In this context, libraries can be thought of as extensions to the core that are portable between chips or board variants for a particular chip.

More granularly, each Arduino core contains a pair of generic files: arduino.h that handles declarations for what are often called ‘built-in’ functions like \PYTHON{pinMode()}, \PYTHON{digitalWrite()}, \PYTHON{analogRead()}, \PYTHON{delay()}, et cetera as well as macros for constants, like \PYTHON{HIGH} or \PYTHON{INPUT\_PULLUP}, and operations like \PYTHON{bitToggle()} or \PYTHON{abs()}; \FILE{main.cpp} declares the basic program structure for a sketch, as it relates to \PYTHON{setup()} and \PYTHON{loop()}. We should note that \FILE{arduino.h} also sets the stage for board-specific pin mappings defined in \FILE{pins\_arduino.h} at \PYTHON{<architecture>/variants/<board name>}. 

The remainder of a core is largely architecture-specific. There is a collection of C files with names wiring.c and \FILE{wiring\_<type>.c} that define hardware initialization, timing, and functions related to IO of various types, like \PYTHON{analogRead()}. This nomenclature is a nod to the Wiring API that Arduino utilizes. There are also the header and C++ files for serial classes Stream, Print, and HardwareSerial, alongside a light USB stack. If you’re interested in digging deeper into a specific core, check out this appendix document.

\subsection{What is an Embedded Framework?}

Loosely, to account for variation between manifestations, an embedded framework is a collection of development tools for microcontrollers that may or may not be portable and often serve as some combination of low-level API, software development kit (SDK), and hardware abstraction layer (HAL). For instance, the Arduino framework comprises its integrated development environment (IDE), built-in functions (defined by cores), as well as library extensions. As an open-source project, the Arduino framework is perhaps a bit more difficult to put bounds around, given its reliance on related projects, namely Wiring and its predecessor Processing, born out of the MIT Media Lab. Here are some other examples:

CMSIS - Arm Cortex Microcontroller Software Interface Standard

FreeRTOS - a real-time operating system (RTOS) kernel for embedded devices

Mbed - an embedded RTOS for Internet-of-Things, drivers for I/O devices

STM32Cube - HAL between STM32 devices and other libraries

If you work with more than one framework, we recommend PlatformIO, an extension of Visual Studio Code, that enables you to call on many frameworks, with a unified workflow.

\subsection{What is an RTOS?}

Okay, so two of the examples we provide for embedded frameworks make mention of a real-time operating system or RTOS. In a circular definition, an RTOS is an operating system that serves real-time applications. More helpfully, an embedded RTOS is extremely light-weight system software that organizes the attention of microcontroller processing cores, often singular, to minimize unnecessary task switching and to prioritize time-sensitive computation so that the processing time required for a given task is ideally less than the interval to the next input. We won’t dwell on the details here, but you can learn more about task scheduling, apparent multitasking (threads), and other RTOS fundamentals, here.

\subsection{What are Mbed and mbedOS?}

Mbed is a competing embedded framework to Arduino, with a fair bit of overlap. A glance at the ‘What is Mbed?’ section of their landing page illustrates this. Both are centered around a C++ API and feature IDEs, example code, libraries, and drivers for common components. Perhaps most striking is the difference in accessibility and other professional-grade considerations Mbed makes, like security. Notably, while mbedOS can be and usually is an RTOS, for applications that require thread management, there is also a limited ‘bare metal’ profile that you can utilize to cut down on the memory utilization tied to RTOS overhead.

Why are we talking about mbedOS? Well, the Nano 33 BLE Sense runs on it. Here is an Arduino blog post from Martino Facchin, the head of Arduino’s firmware development team, that explains the Nano 33 BLE Sense unique core, and its reliance on mbedOS.

\subsection{What is Bare Metal?}

You may happen upon variations of this definition, but at a high, inclusive level, ‘bare metal’ programming in an embedded systems context means utilizing microcontroller hardware independent of an operating system or scheduler, at the least. Typically, there is a base super loop, not unlike the Arduino sketch structure, and all processing is defined by your application. In many cases, this also means your code is written independently of pre-existing frameworks. This is where you’ll find some variation in definitions. The colloquial spirit of bare-metal programming is counter to most developer mindsets, in that because microcontroller tasks are often simple but mission-critical, their developers want strict control over processing that precludes reliance on abstraction.

A ‘true’ bare-metal project, then, will have a developer that either writes or, if supplemented by a light-weight framework, fully understands each layer, from the physical hardware on up, using an MCU/SoC datasheet as their guide to registering definitions and hardware architecture.

\subsection{Implications for your Application}

As you can imagine, your choice of the framework (or lack thereof) should stem from the needs of your application. If you have no hard time-constraints, an RTOS could be eating away at precious resources that may be constrained. However, there are obvious challenges to developing a bare-metal solution, which can largely be summarized by ‘recreating known solutions.’ If your task is simple, time-insensitive, mission-critical, or needs to be deployed on very constrained hardware, a bare-metal approach may be appropriate, assuming you have the knowledge and time required for the development. If you have a complex web of tasks in front of you that could be organized by a scheduler, an RTOS is really a no-brainer. In most cases, there is room for innovation in the middle, where you place trust in proven frameworks and libraries, but manage processing attention in a simple base loop, perhaps organized as a finite state machine or FSM.


\section{Embedded ML Software}

    
    VIJAY JANAPA REDDI: Hi, In this video, I'm
    going to touch on Embedded Machine Learning Software
    frameworks a little bit.
    The reason I want to do this, is because there
    are many different machine learning frameworks out there
    and need to understand how to choose the right one for the right problem,
    that you're trying to solve.
    But before we get there, let's talk about the general ecosystem
    of what's involved with the machine learning frameworks.
    And what are some of the challenges as we try and take those frameworks
    and cram them down, onto something that would
    fit in the constraints of like a one megabyte flash
    or just 256 kilobytes of RAM.
    Let's take TensorFlow and open up the hood and see what's inside.
    Here is a general architecture of TensorFlow.
    As you can see, it looks quite complicated.
    There are many different pieces to it.
    For example there's a lot of high level code that you might write in Python,
    or you might use language like C++, Java, Go and so forth,
    in order to express your program or the neural network itself.
    And then you want to actually train it.
    So you have a distributed execution engine,
    which is capable of using different computing systems all concurrently,
    in order to very quickly train the model.
    In addition to having capability of using different systems,
    each of these systems might have different types of hardware.
    They might have a CPU, they might have a GPU
    and you've got to orchestrate all that execution.
    You might also have a TPU, Google's Tensor Processing
    Unit, which is custom designed for just running machine learning tasks.
    Well, you've got to manage all of that capability in execution heterogeneity
    data.
    So, naturally these systems end up being quite complicated.
    Just by that nature or by that fact, you end up
    having to rely on big operating systems.
    So these frameworks like TensorFlow or [? pintor, ?]
    rely on operating system capabilities.
    Operating systems like Windows and Linux for instance,
    are very well capable of running these systems.
    However, these operating systems come with a fair bit of complexity
    themselves.
    For example, they have to support virtual memory,
    they have to support the ability to manage multiple threads that
    are concurrently running on the system.
    Why am I talking about all of this?
    Well, remember we are talking about embedded systems.
    How can we take for instance, our embedded system which
    has got a very limited amount of capability, in terms of compute,
    in terms of memory, in terms of storage, in terms of power.
    Given that we're orders of magnitude smaller,
    we cannot have the flexibility of supporting a massive complex machine
    learning framework on top of a very simple device.
    Therefore, we need to think about how can we
    reengineer the machine learning frameworks,
    so they're more aptly suited to run all the embedded device.
    So typically, when you're thinking about building an inference framework,
    there are three things that you would need to consider;
    what kind of model features do you need to support,
    what kind of software capabilities do you need to have
    and what kind of hardware are you trying to support.
    If we break this down from a model standpoint,
    do you need to support training?
    The answer is no.
    We are doing only on device inference.
    Do you need a support inference?
    Obviously yes.
    You need to support a wide range of ops.
    For example, TensorFlow has over 1,400 ops,
    but in reality we often use only about 50 different ops.
    By the way, when I'm saying ops, I'm talking about things like convolutions
    and Macs pulling, those kinds of operators.
    So just from the model standpoint, you need
    to understand that when we're deploying things on to the embedded device,
    we only need some subset of the big features
    that the engines actually offer.
    Furthermore, when we think about software capability,
    should the toolchain actually support quantization?
    The answer is absolutely yes.
    We don't want to typically be running f.p.
    32 models, we want to run integer eight models.
    Small models that take up less space, where the image or arithmetic can also
    be done efficiently.
    Then, we cannot assume that we're going to have complex operating systems,
    so therefore we cannot rely on things like virtual memory support.
    That is something that we take for granted to the big operating
    systems or the big frameworks.
    In addition, when it comes to the hardware, we need to think about,
    do we support a diverse range of processor or hardware device.
    Well, the answer is no, because typically
    on a simple embedded microcontroller, you're
    looking at a very small subset of processing capability,
    so you need to have.
    And do you need support many different platforms and features?
    The answer is yes, you need to be able to build
    a framework that can easily target many different types of embedded hardware.
    But on a single device itself, the heterogeneity is quite limited.
    You need to keep these things in mind, as we try and build
    an inference framework that actually runs efficiently
    on the embedded microcontroller.
    This has led to the proliferation of many different types of machine
    learning frameworks.
    So for example, TensorFlow Lite Micro which we are going to learn how to use,
    is one amongst the ecosystem that is starting to thrive.
    The reason I'm touching on this ecosystem,
    is because there's a TinyML engineering you
    need to know that there's more than just TensorFlow Lite Micro, which
    you're learning about in this course.
    So there is tvm, there's uTensor, there's GLOW, there's STM Cube AI.
    These are all different types of machine learning frameworks
    that support the different class of hardware that we're seeing.
    Let's take for example, a particular framework
    and see how they go about deploying a model that's on the left hand side,
    all the way down to a custom embedded system on the right hand side.
    Let's take STM for example, which is a company that is well
    known in the microcontroller space.
    What they often do, is like they'll take a neural network model that comes in
    and then they will generate some sort of code in their own proprietary way,
    to run on the embedded system.
    Let's open up that box and understand this a little bit further.
    Typically you have a network that's in petrovitch in TensorFlow.
    You'll take that into an importer, which imports the model in and builds
    an internal representation of what the network looks like.
    For example here, PINNR is some sort of internal representation
    that's used by the software framework.
    Then, a compiler will actually generate optimized code for the target device.
    It does this because it knows the capabilities of the target
    that you're looking for.
    It'll bring in the best capabilities needed, in order
    to achieve the best performance, from optimized kernel
    libraries that are available.
    Remember, I talked about how the M0, the M3 the M4,
    a range of different microcontrollers all have different kinds
    of instructional support.
    The reason I was showing you that is because now you can make the connection
    that.
    Depending on the type of microcontroller,
    you might want to choose the right type of instructional level support
    in order to maximize the performance of that particular device.
    So you can run inference much faster.
    And finally, you generate executable code
    that runs with the microcontroller.
    This is a very typical flow.
    Implementations tend to vary quite a bit.
    For example here, if you take a model like this on the left hand side
    and you run it through STM cube, as well as TensorFlow Lite Micro,
    there are differences.
    TensorFlow Lite Micro for example here, consumes a little bit more
    flash, a little bit more RAM, therefore it also
    consumes a little bit more runtime.
    While this is not always the case, it is very subjective and dependent
    upon the model you're using and the type of optimizations that are supported
    on the platform of your liking.
    Of course there are also differences where TensorFlow Lite is positive,
    for example it's an open source project.
    Does that matter?
    Well that depends on your needs.
    But the general point that I'm getting across,
    is that it's important to understand that there many different frameworks
    that are out there.
    And TensorFlow Lite Micro is one that we will be going deeper into.
    How do you choose the right framework is an important question
    that you should really think about.
    And you have to think about it from the hardware standpoint,
    you think about it from the embedded inference standpoint
    and you also have to understand about things
    like the accessibility of the experience as a developer.
    Now, we are going to try and convey some of those lessons
    through the lens of TensorFlow Lite.
    Specifically Pete Warden, who is the technical lead for TensorFlow Mobile
    and Embedded, is actually going to talk us through the implementation nuances
    and details of TensorFlow Lite micro, all of which
    you need in order to be able to start building the actual applications.
    So with that said, in the next couple of videos,
    I'm going to hand it off to Pete and I'll
    pick up after you're done with that, where we
    start talking about keyword spotting.
    I'll see you there.

%%%%%%%%%%%%%%%%%%%%%%%%%%%%%%%
    
    
\section{What is TensorFlow Lite for Microcontrollers?}

As a future TinyML engineer, it is essential to understand the inner workings of the software you use to know its capabilities and limitations. So, in the upcoming series of videos and readings, you will learn more about the challenges that led to the development of TF Lite Micro, straight from the source. Pete Warden from Google, who leads the team that works on TF Lite Micro, will introduce TF Lite Micro and give us a sneak peek into its internal workings. Here is a preview of what is to come next.

TensorFlow has become the most popular deep learning framework, superseding other popular frameworks such as PyTorch and Keras. TensorFlow, developed by Google, contains a Python frontend with highly optimized C++ code at its core, making it simple to program, fast, and efficient. The library has a large developer community and is now seen as the de facto standard for most machine learning applications.

Despite this, TensorFlow is not suitable for every scenario. The standard TensorFlow library is ~400 MB in size, and even running a relatively small model (e.g., 200 MB) can take up a considerable amount of random access memory (> 1 GB). Such large storage and memory requirements make running simple models on lightweight systems largely intractable.

Recognizing this issue, Google developed a more lightweight framework, TensorFlow Lite, also sometimes referred to as TensorFlow Mobile. The TFLite binary is approximately 1 MB in size, considerably more compact than the original library, making it possible to run deep learning models on mobile devices such as smartphones. This compression was achieved by removing superfluous functionality that is largely unnecessary for mobile deployment. 

While this is an improvement, our problem still remains: even TFLite is not suitable for every scenario. Many important deep learning applications exist at the microcontroller-level, which are significantly more resource-constrained than mobile devices, often equipped with less than 1 MB of storage and 256 KB RAM. Clearly, deploying TFLite models is not feasible for microcontrollers, so an alternative solution was needed.

Enter TF Lite Micro. TF Lite Micro takes the compression of the TensorFlow library to the extreme, removing all but essential functionality. In fact, the core runtime of the library takes up only 16 KB, several orders of magnitude smaller than TFLite. With such a small memory footprint, this lightweight framework makes it possible to deploy deep learning models on the smallest of microcontrollers, such as an Arduino Nano.

However,  this is not without its complications.  Deploying models with TF Lite Micro is fraught with new and unique challenges when building models. For example, since all functionality for plotting and debugging is removed, troubleshooting model issues is difficult. Additionally, since many microcontrollers do not have floating-point units or use 8-bit arithmetic, the model weights and activations must be suitably quantized on the microcontroller system. Since model training requires near-machine precision to perform gradient descent, this largely precludes on-device training. Thus, TF Lite Micro models must first be trained on a device with greater computational resources before being ported to the microcontroller, adding an additional stage to the machine learning workflow. 

Despite this, the benefits provided by TF Lite Micro - the ability to perform machine learning inference on microcontroller devices - far exceed the challenges, heralding a new era of machine learning that is often referred to as tiny machine learning.





\section{TFMicro: The Big Picture}

TFMicro: The Big Picture
0:00 / 0:00Klicken Sie auf die Pfeiltaste nach oben, um das Geschwindigkeitsmenu aufzurufen. Regeln Sie dann mit den Hoch und Runter Pfeiltasten die Geschwindigkeit. Klicken Sie Enter, um die jeweilige Geschwindigkeit festzulegen.Geschwindigkeit1.0xKlicke auf diesen Knopf und diess Video stumm oder nicht stumm zu schalten, oder drücke die HOCH oder RUNTER Knöpfe um die Lautstärke zu erhöhen oder abzusenken.Maximal Lautstärke.Videotranskript
Beginn des Transkriptes, zum Ende springen.
PETE WARDEN: Hi, I'm Pete Warden, and today I'm
going to be giving you an introduction to Google's Open Source Machine
Learning Framework for Embedded Systems--
TensorFlow Lite Micro.
I'll be starting with a broad overview, the big picture of what
the framework is for, and we'll be diving more deeply into each topic
I mention in later videos.
So don't worry if you don't understand everything I'm covering immediately.
So what is TensorFlow Lite Micro?
It's a software library that aims to make
it easy to run machine learning models efficiently on embedded hardware.
To understand what that means in more depth,
I'd like to show the constraints of the embedded platforms it
has to run on so hopefully, the design decisions make a bit more sense.
There are two broad categories of constraints
that drive the library's design.
Embedded systems face extremely tough requirements for power usage, price,
form factor, and compute performance.
These requirements mean that there are massively
diverse sets of solutions available because different products require
radically different trade-offs.
The pressure to meet different trade-offs
shows up in the hardware that's available.
There are a lot of different embedded chip architectures in active use,
far more than in any other computing domain.
The cloud has x86 CPUs, GPUs from a couple of vendors,
and recently ARM chips have been emerging.
Smartphones almost exclusively run on ARM Cortex
A chips, but by contrast, embedded systems have ARM Cortex M
microcontrollers, [INAUDIBLE] DSPs, Qualcomm DSPs, [? Seever ?] DSPs,
Synopsis microcontrollers, FPGAs, and with many other architectures
like Texas Instruments dominating in particular markets.
These are more than just manufacturers.
Most of the chips have entirely unique instruction sets
that require their own compilers and tool chains.
The instruction sets are often incompatible
even between different versions of a chip in the same family
from a single manufacturer.
This variety is because there are so many different trade-offs that
make sense for different situations.
So many different solutions exist for all the various types of problems
that embedded systems try to solve.
A consequence of this is that the framework
has to be radically portable, able to compile and run with no hitches
across a wide range of computer architectures.
At the same time, compute latency is often an important requirement, too,
so there has to be an easy path to optimizing
the code for the particular features offered by a platform.
Another consequence of the tough requirements embedded systems face
is that they often have comparatively small amounts of memory and low clock
rates.
These memory constraints make library design particularly hard
because, in an ideal world, you want a framework that
can be a closed box with an abstraction layer
where you don't have to worry about the internals,
but, because the implementation details of the library
directly affect its binary footprint, how much memory the code takes up,
users of it do have to be aware of how it works.
It also means that the library code has to be
very careful about how much memory it uses for other purposes,
too, with a focus on memory minimization algorithms.
The constraints of embedded requirements also
extend to the system software that's available.
If you use the more relaxed environments,
the lack of a heap and dynamic memory allocation on embedded systems
will come as shock.
Dynamic memory allocation using malloc or [INAUDIBLE]
is a productivity-enhancing feature of most modern computing environments,
but it implicitly assumes that there is a large enough pool of memory
available, possibly through virtual memory on disk, that running out of it
is unlikely in normal operation.
Because embedded systems may only have tens of kilobytes of memory
and may need to run continuously for months or years,
developers have to be a lot more careful and explicit about the allocation
strategy.
TensorFlow Lite Micro avoids using any dynamic memory allocation internally
and just asks for an arena of contiguous memory from the application
to use for its variables, which is not altered after initialization time
so that products can have confidence that the framework won't
crash due to heat fragmentation when running for long periods.
Embedded systems also often lack other facilities
we take for granted, like file system calls or even functions like printf
that we might expect to always be available.
To cope with this, TensorFlow Lite Micro takes pains to avoid using any C or C+
library calls that require operating system support.
This includes features like the STL containers
that rely on dynamic memory allocation as well as printf and other POSIX
functions that need low level support.
Even where these are available, they often
require additional space in the binary so they're avoided for size reasons,
too.
The end result of all this work is a software library
that's able to run across a wide variety of different boards.
We're using the Arduino Nano BLE Sense 33 primarily in this course,
but the same application code you're using to ML models
should work portably and performantly across a wide range of other devices.
Hopefully I've given you a bit of an idea of what TensorFlow Lite Micro is,
but how can you use it yourself as a developer?
The basic flow as shown here, and you can see the same pattern
in all of our code examples.
Don't worry if these particular terms don't make sense yet.
This is just an overview, and I'll be diving
into each stage in much more detail later in this section.
The first stage is specifying what ops you're using in your model
using an OpResolver.
This is important because it ensures only the implementations of those ops
will be included in the binary, keeping the space used as small as possible.
Then you load the model into an interpreter object, which handles
all of the initialization needed.
You copy the input data you want to feed into the model into the model inputs.
You call invoke on the interpreter to run the MODEL running the input
data through all of the operations to calculate the final result.
You then read that final result from the model outputs,
and presumably, you're going to be taking some action based on the model
outputs.
You're using the ML Model to decide something,
so your application should respond to what's being calculated.
And then usually, you'll go right back to reading some new input data
and running the cycle again.
I know this has been a very quick overview of a very complicated topic,
but I hope it's at least giving you a flavor of what TensorFlow Lite
Micro is all about.
To sum up, it's a library that's able to take in models from the TensorFlow
training environment and run them on embedded systems
with very little memory and no operating system facilities.
To work well in these environments, it focuses on memory size optimization,
avoiding dynamic memory allocation, using
a very minimal set of standard library functions,
and keeping its code as portable as possible.
Ende des Transkriptes, zum Anfang springen
Downloads und Transskripte
Video
Videodatei herunterladen
Transskripte
Download SubRip (.srt) file
Download Text (.txt) file

\section{How to Use TFMicro: HelloWorld Tour}

Beginn des Transkriptes, zum Ende springen.
PETE WARDEN: In this video, we will be looking
at the code that's at the heart of the Hello World example,
where we call into Tensorflow Lite Micro to actually run the machine learning
model we've trained.
This code is contained in a loop function that's
called repeatedly by the Arduino framework,
and it fits onto a single screen.
So we'll take a few minutes to step through the APIs calls involved.
As a reminder, the model we've trained takes an input value
and calculates an estimate of the sine function on that value.
We're incrementing the input value over time to help pulse and LED.
So in this code to get started, we figure out
the float version of that incrementing value,
and store it in a variable called x.
The Hello World model has been optimized use 8-bit quantization.
So we need to convert that float value into the quantized equivalent.
During initialization, we stored the structure
that holds information about the input, which is a TFLite tensor in a variable.
So we use those parameters to perform the conversion.
To go from a floating point value to 8 bits,
we need to do a linear transformation controlled by the scale and zero point
settings that are stored in the model.
You'll always want to do something similar
whenever you're feeding in quantized values as inputs, scaling
and subtracting an offset.
To be extra safe, you could also make sure the float
values don't overflow the minimum and maximum of the expected range.
But for simplicity's sake in this example, we don't do that.
Now instead of passing values into a function call,
TFL Micro sets up areas in memory to hold inputs and outputs ahead of time,
and the application developer is expected
to fill in the contents of the inputs before calling the invoke function.
In our case, the input is just a single 8-bit number.
So we use the input data structure we set up during initialization to write
into the correct memory location.
The data member of TFLite tensor holds the address of the array,
and it's a C union, so that we can have type safe access.
Since a model is expecting a single signed 8-bit value,
we use the int8 accessor, and write the first entry in the array.
The actual model calculation is triggered
by calling the invoke function on the interpreter.
This runs through all the operations defined in the model,
reading the inputs from the memory area that we stored
the quantized x value into, and writing to the memory areas reserved
to hold output.
You might be surprised that the invoke call doesn't take any arguments.
But all the parameters are either defined as part of the model data,
or taken from the input memory areas we just filled out.
It also doesn't return any calculation values,
since those are written into the areas of memory that
are reserved for outputs.
It does return an error code, which in normal operation is K TFLite OK.
Any other result indicates that something went wrong
with the calculation process.
Just as there is a reserved area for input values,
the invoke call writes the results of the model calculations
into memory areas that can be accessed once the function is complete.
In this case, we've stored the data structure
holding information about the output, and like we
do to write the input value, we read a single 8-bit value
from the output array.
Since it's a lot easier to work with a floating point number, rather
than the 8-bit encoded value we got out from the model,
we do the reverse process we did for the input,
to convert the output back into a floating point value.
Again, we're reading the conversion factors from the model information,
and it's good practice to follow this pattern whenever you need
to convert quantized output values.
Hopefully this run through has given you some ideas
on how the model calculation process works with TensorFlow Lite Micro.
You'll be seeing more complex models in the future.
But they'll all follow this basic pattern
of writing input values into memory arrays, calling invoke,
and then reading the results.

  \includegraphics[width=\textwidth]{TinyML/3 Development/TFLiteHelloWorld}
  
\section{TFLite Micro: Interpreter}


    PETE WARDEN: Hi, today I'm going to talk about an important part
    of the TensorFlow lite micro libraries design, Its interpreter.
    So what is an interpreter?
    Interpretation means executing a program by following a data structure stored
    in memory, as opposed to a purely compiled approach,
    where the execution is determined by instructions ahead of time.
    In our case, this is the difference between storing a model as data
    and looping through the operations of runtime,
    calling the appropriate functions based on the OP type versus examining
    the model at compile time and omitting code that
    makes direct call to the functions.
    TFlite Micro is based on an interpreted approach.
    But an interpreter slow, the most common place
    you might have encountered interpreters and compilers in the past,
    is on desktop computers.
    Python is an interpreted language and C+ is compiled for example.
    Broadly speaking, desktop interpreted languages are easier to use
    but run more slowly.
    Is this true for embedded ML too?
    Running a procedural language like Python through an interpreter
    is inherently challenging, because each line of code
    probably only takes a few assembly instructions to execute.
    But passing that line, may take thousands of instructions.
    This gives a big performance advantage to compiled languages like C++,
    which do the parsing ahead of time.
    Machine learning is different though.
    In machine learning, each layer like convolutional softmax,
    is likely to take tens of thousands, even millions of cycles to execute.
    This means any passing overhead is comparatively small.
    As a result, there isn't a noticeable speed difference between the approaches
    for real models in our experiments.
    As a practical example, we run an experiment
    to count the number of instructions that were taken up
    with interpreting the model, versus calculating the operations.
    For a large model, it's a vanishingly small proportion of the time taken.
    And even very small models, it's only a few of the cycles.
    And there are some very important advantages to using an interpreter,
    it's possible to change the model without recompiling the code,
    though the lack of fowl system on most embedded systems
    can make this hard in practice.
    The same operator code can be used across multiple different models,
    in the same system.
    The same portable model serialization format
    can be used across a lot of different systems.
    So, how does interpreting a model actually work?
    The model contains a list of operators and arguments.
    For example, CONV 2D, FULLY\_CONNECTED softmax.
    The interpreter loops through this list of runtime
    and calls the appropriate operator implementations or kernels.
    You can think of the interpreter as pseudocode
    that loops through a list of ops and for each one decides what function to call,
    based on the ops type.

  
\section{MCU Memory Hierarchy}

 
\subsection{MCU Memory Hierarchy Micro?}

    Given the limited resources available on microcontroller units (MCUs), it is paramount for us as future tinyML engineers to understand the different types of on-device memory, including their interplay and performance capabilities. This understanding will become vital for optimizing machine learning models on given hardware architecture.
    
\subsection{What is Memory?}

    Computer memory stores information as binary values, known as bits, which can take on a value or one or zero. However, storing a single bit is not particularly informative, and so multiple bits are often grouped together into quantities known as bytes, which consist of 8 bits. A byte can represent 256 distinct values, which can be used to represent anything from letters of the alphabet to the weights of a neural network.
    
    Microcontrollers read from and write to memory using bytes, which are located by a memory address. This memory address is a hexadecimal value that tells the microcontroller where the memory it is looking for is located, much like a postal address. Microcontrollers have a limited amount of memory, which fundamentally limits the complexity of programs and computations they can run. Different forms of memory may also take longer to access or lose data when they are not constantly powered, leading to the use of multiple forms of memory for various purposes.
    
\subsection{RAM and Flash Memory}

    The two most important forms of memory are flash memory and random access memory (RAM). Our Arduino Nano 33 BLE Sense is equipped with 1 MB of flash memory and 256 KB of static RAM.
    
    Flash memory is non-volatile, meaning that data persists even when the system is turned off. This memory is where we save our model weights and program code, and is what we are transferring when we “flash” our model onto the device. The process of saving to flash memory is slow and gradually degrades the memory over time. Consequently, to minimize this degradation, flash memory is effectively used as read-only memory and only overwritten during reprogramming, not when storing intermediate results during program execution.
    
    RAM, however, is volatile, and thus data is lost when the device is powered off. Accordingly, RAM is used for the temporary storage of variables, such as input and output buffers, and intermediate tensors. RAM is much faster to read and write than flash memory, making it ideal as the primary memory used during program execution. RAM comes in two forms: static and dynamic.
    
    Dynamic RAM (DRAM) uses a single transistor and capacitor to store a bit, but the capacitor quickly loses charge and must be periodically refreshed to prevent the stored information from being lost. Static RAM (SRAM) uses 6 transistors to store each bit but is able to maintain the stored information without refreshing. SRAM is more expensive due to the additional transistors needed but is faster and less power-intensive because no time or energy is spent performing a DRAM refresh between reading and write operations. These differences make DRAM a suitable candidate for main memory in modern computers, and SRAM more suitable for caches. In our microcontroller, we have access to SRAM and use it as our main memory during program execution.
    
    The third type of memory is also used in microcontrollers but is often abstracted from the user; this type of memory is known as a register. Typically, special purpose registers exist that store various values needed to perform certain low-level computing functions, such as the program counter and stack pointer. In contrast, general-purpose registers are used to store values and memory addresses. It is unlikely you will need to play around with registers as a tinyML engineer unless you are working at the assembly code level.
    
\subsection{Relevance to TF Lite Micro}

    Why do we care about introducing computer memory to you? Well, the size of your model weights and code cannot exceed the size of the available MCU flash memory (i.e., flash memory minus the 2 KB bootloader), and the size of input buffers and intermediate tensors cannot exceed the available SRAM on the device. These are hard constraints, and as tinyML engineers, we must be wary when working with large models that might exceed these memory constraints. In such cases, an MCU platform with additional resources or a simpler model may be necessary.  
    
\section{TFLite Micro: Model Format / FlatBuffer}

    0:00 / 0:00Klicken Sie auf die Pfeiltaste nach oben, um das Geschwindigkeitsmenu aufzurufen. Regeln Sie dann mit den Hoch und Runter Pfeiltasten die Geschwindigkeit. Klicken Sie Enter, um die jeweilige Geschwindigkeit festzulegen.Geschwindigkeit1.0xKlicke auf diesen Knopf und diess Video stumm oder nicht stumm zu schalten, oder drücke die HOCH oder RUNTER Knöpfe um die Lautstärke zu erhöhen oder abzusenken.Maximal Lautstärke.Videotranskript
    Beginn des Transkriptes, zum Ende springen.
    SPEAKER 1: Hi so we've talked about how an interpreter needs
    a list of operators, weights, and other information to run a model.
    But how is this data stored?
    The model format is the link between training a model
    and running it on device.
    So in this video, I'll run through the format, and how we use it.
    If you look near the start of the Hello World sketch,
    you'll see this section of code.
    This is how we set up the model data structure.
    The most important call is this line.
    The g model variable is an array of bytes
    that we've exported from TensorFlow, held as a CLI in a source file.
    This array of bytes is the raw data for the model,
    and the GetModel call returns a C++ class that lets us access
    the model's data elements.
    As the comment mentions, this call doesn't
    do any data movement or copying, so it's very memory efficient.
    It just creates a thin wrapper to read the information in place.
    Once we have this C++ class, we can read the operations and weights
    to understand how to run the model calculations.
    The underlying bytes of the model are stored
    in a read-only array, which means it can be held in flash memory
    on devices that have it.
    If you are used to programming desktop computers,
    you can think of this array of bytes as the equivalent of a file on disk.
    But since embedded systems usually don't have file systems,
    we store the data in read-only memory.
    Serialization turns a data structure in a language,
    like Python, into a form that can be stored as an array of bytes,
    and then understood by many other languages.
    You might be familiar with other ways of serializing data,
    like JSON or ProtoBuffers.
    For TensorFlow Lite Micro, we're using FlatBuffers,
    and the key reason for that is that they're very memory efficient.
    This matters because there's often very little memory available.
    A typical embedded chip might only have a megabyte of flash memory
    and 256 kilobytes of SRAM.
    This means we need a library that uses very little memory
    and doesn't need operating system support,
    since that's often left out to save space.
    So we have this model stored as an array of bytes,
    the equivalent of a file, that holds all the information we
    need about a machine learning model.
    The FlatBuffers format allows us to keep that information in place,
    without needing any additional memory, and also has
    the benefit of not requiring any operating system support.
    The simulation is defined as a text file called a schema
    and specifies what elements are held within the data structure.
    The FlatBuffer compiler takes a schema definition and creates C++ code
    to access all of the data held in a serialized byte array.
    To give you an idea of what's held in the model data,
    the main parts are metadata about the model, like the version, quantization
    parameters, inputs and outputs, the list of operations to perform,
    and the size and values for all of the weights.
    If you want to get a better feel for what's inside of FlatBuffer model
    structure, we've created a co-lab that lets you explore reading, modifying,
    and writing out a model in Python.
    This will let you change the stored weights, for example.
    So it's a great way to, kind of, get your hands dirty with this.
    Ende des Transkriptes, zum Anfang springen
    Downloads und Transskripte
    Video
    Videodatei herunterladen
    Transskripte
    Download SubRip (.srt) file
    Download Text (.txt) file  
    
    \section{TensorFlow Lite Flatbuffer Manipulation Example}    
    
  In this Colab, we’ll dig a little deeper into the TensorFlow Lite Flatbuffer file format. We’ll load in the Flatbuffer library and the TensorFlow Lite Schema, which will allow us to access and directly modify the weights in the pretrained KWS model manually. You’ll then get to see how your direct modifications impact the final model accuracy using the Speech Commands dataset!
  
  https://colab.research.google.com/github/tinyMLx/colabs/blob/master/4-4-8-Flatbuffers.ipynb
 
 \section{TFLite Micro: Memory Allocation (a.k.a Tensor Arena)}
 
 
     0:00 / 0:00Klicken Sie auf die Pfeiltaste nach oben, um das Geschwindigkeitsmenu aufzurufen. Regeln Sie dann mit den Hoch und Runter Pfeiltasten die Geschwindigkeit. Klicken Sie Enter, um die jeweilige Geschwindigkeit festzulegen.Geschwindigkeit1.0xKlicke auf diesen Knopf und diess Video stumm oder nicht stumm zu schalten, oder drücke die HOCH oder RUNTER Knöpfe um die Lautstärke zu erhöhen oder abzusenken.Maximal Lautstärke.Videotranskript
     Beginn des Transkriptes, zum Ende springen.
     PETE WARDEN: Hi.
     In this video, I'm going to talk about how TensorFlow Lite Micro deals
     with one of the most challenging areas of embedded application
     development, memory management.
     The reason memory is so challenging is embedded systems
     typically only have hundreds or even just tens of kilobytes of RAM.
     This means it's very easy to hit memory limits when building an application.
     This means any framework that integrates with embedded products
     needs to offer a lot of control over the memory it uses,
     because developers simply won't be able to use
     a library that takes up too much space.
     There's also a more subtle problem that shows up
     when an embedded application needs to run continuously for months, or even
     years.
     Even if you're sure that you'll never ask for more memory
     than the system contains, allocating a new area
     requires as a big enough gap between other allocations.
     Over time, it's possible for small allocations
     to litter the available address space, so that there were no gaps big enough.
     This state is known as heap fragmentation.
     It's extremely hard to predict how heap usage will evolve over
     long periods of usage, and the resulting lack of available memory
     is likely to cause the application to crash.
     The standard way to avoid these kinds of memory fragmentation issues
     with long-running programs, is to avoid dynamic memory allocation altogether,
     at least after a startup phase.
     Any libraries used by applications with this requirement
     also need to follow the same rules.
     Because memory is so scarce, a lot of devices
     ship with no operating system at all.
     They just run on bare metal.
     This also makes memory handling a lot more difficult,
     since there were no facilities like a heap or memory virtualization
     to help with allocation.
     The Arduino board we're using does have an operating system, mbed OS,
     but that's not something that that TensorFlow Lite Micro
     library can rely on existing on all the platforms it supports.
     So far I've talked a lot about the challenges of our memory.
     But how does TensorFlow Lite Micro deal with them?
     The most important concept to understand as an application developer using
     the library is the arena that you supply.
     The library needs memory to store variables and intermediate calculation
     values.
     But since it can't rely on standard operating system calls, like Malloc,
     we asked developers to pass in an area of memory
     to an interpreter when it's created.
     In this code snippet from an example, you
     can see how we create a global byte array and pass it into the interpreter
     when it's created.
     To avoid the problem of fragmentation, the interpreter
     only allocates from this arena at startup time.
     It will never do any allocations during invoke,
     so there's no chance of failing due to memory fragmentation.
     This API also ensures that the memory usage of the library is very clear,
     and how it's allocated is up to the application developer.
     The only requirements are that the memory area is available for as long
     as the interpreter is in use, and that it's large enough
     for everything the model requires.
     The memory space itself is divided up into various sections internally.
     You don't need to understand what going on under the hood
     to use this arena interface.
     But the array is used to hold calculated values the operators need,
     the internal state of the interpreter, input and output
     arrays for all of the operators, and various other things.
     Earlier, I mentioned that you need to make sure the arena is large enough.
     But what size should it be?
     Unfortunately, this is actually a very tricky question to answer.
     The total size needed for the arena depends
     on what operations are in the model, and what the operator parameters are.
     It's also possible for optimized implementations of operators
     to acquire scratch memory, which has to be allocated from the arena.
     The size needed for operator inputs and outputs
     doesn't depend on which platform you're running on,
     but different devices can have different operator implementations.
     So the scratch buffer sizes can vary.
     This means it's hard to come up with the exact size needed
     for the arena ahead of time.
     So what can you do?
     Luckily, there is an arena used by its function in the interpreter.
     So I recommend you start off with as large an arena as you can create,
     load your model, and then log the number of bytes actually used.
     Ideally you should do this on the device you
     want to deploy on, since you may have optimized operator implementations that
     use scratch memory.
     Once you have the actual number, set the arena to that size and recompile.
     Hopefully, this has given you a better understanding
     of how TensorFlow Lite Micro deals with the challenges of memory
     on an embedded platform, and will give you a head start building
     your own memory conscious applications.
     Ende des Transkriptes, zum Anfang springen
     Downloads und Transskripte
     Video
     Videodatei herunterladen
     Transskripte
     Download SubRip (.srt) file
     Download Text (.txt) file   
 
 \section{TFLite Micro: NN Operator Support (OpsResolver)}
 
     0:00 / 0:00Klicken Sie auf die Pfeiltaste nach oben, um das Geschwindigkeitsmenu aufzurufen. Regeln Sie dann mit den Hoch und Runter Pfeiltasten die Geschwindigkeit. Klicken Sie Enter, um die jeweilige Geschwindigkeit festzulegen.Geschwindigkeit1.0xKlicke auf diesen Knopf und diess Video stumm oder nicht stumm zu schalten, oder drücke die HOCH oder RUNTER Knöpfe um die Lautstärke zu erhöhen oder abzusenken.Maximal Lautstärke.Videotranskript
     Beginn des Transkriptes, zum Ende springen.
     PETE WARDEN: Hi.
     In this talk, I'm going to talk about one
     of the most mysterious parts of the TensorFlow Lite Micro interface.
     Op resolvers.
     Op resolvers exist to help reduce the amount of space
     that the library code takes up.
     The functions that make up the library compile down
     into machine code instructions, and these all take up space.
     Typically, in Flash.
     But Flash is a limited resource on embedded devices.
     There are often just tens of kilobytes available.
     This means if a library's compiled code is too large,
     it won't be usable by applications.
     To help minimize the binary footprint of the library's executable code,
     TensorFlow Lite Micro is broken up into two main sections.
     The core functionality includes things like model loading and the interpreter,
     and is always required.
     But we aim to keep it's size below 20 kilobytes.
     The second section contains the code that implements the operating
     calculations, also known as kernels.
     These are divided up by layer, and designed
     so that they can be independently included or excluded
     from the compiled library.
     The reason for this design is there are potentially
     a lot of different operations that you might
     need to support, with over 1,000 available in the TensorFlow training
     environment.
     Most models, though, only use a handful of these operators.
     This means it's important for binary size
     to only include the operators that you need for your model.
     And this is what op resolvers are for.
     They're way of specifying what operators you need at compile time,
     so the other code can be removed.
     Most models only use a handful of different operations
     so the majority of the space taken up by op implementations is wasted.
     In order to avoid linking in this unused code,
     we use this op resolver mechanism.
     This class is a way to let an application developer specify which ops
     they want to be included in the binary.
     For example, the micro speech model uses just four different kinds of layers.
     This means we can use the micro mutable op resolver class to pick just
     those four operator implementations.
     We create an instance, pass in how many operator types
     we'll want as a template argument, and then add in each layer that we'll need.
     Each layer in the graph corresponds to an operator type that's in the model.
     We're hoping to add more user friendly ways to create the op resolver
     code you need in the future.
     But for now, the simplest way to understand what ops are required
     is by using a visualizer program like Netron to view the graph
     and include the operator types that are visible.
     If you're not memory limited, you can also
     use the special all ops resolver, which pulls
     in all supported operator types at the cost of a larger library size
     footprint.
     Even though the code can be a bit fiddly, using op resolvers
     can dramatically reduce the binary size of your executable.
     So, I hope this has helped you understand them a little bit better.
     Ende des Transkriptes, zum Anfang springen
     Downloads und Transskripte
     Video
     Videodatei herunterladen
     Transskripte
     Download SubRip (.srt) file
     Download Text (.txt) file
     
     \section{Summary of TFLM -  TFLite Micro Developer Design Principles}
     
 There are four overarching design principles that TFMicro was built upon in order to address some of the challenges faced by developers when working with tinyML for embedded systems. This reading provides a synopsis of these core principles, as outlined in further detail in the TensorFlow Lite Micro paper.
 
\subsection{Principle 1: Minimize Feature Scope for Portability}

 This principle proposes that an embedded machine learning (ML) framework should assume, by default, that the model, input data, and output arrays are in memory, and do not need to be loaded into memory. In addition, accessing peripherals, such as an on-device camera, should not be the job of the ML framework. These functions still need to be fulfilled, but principally should not be fulfilled by the ML framework.
 
 While this may seem unimportant, some microcontrollers do not have memory management (e.g. malloc) and other capabilities. Thus, trying to accommodate all varieties of platforms would bloat the library in an attempt to provide sufficient portability. Fortunately, due to the self-contained nature of machine learning models, the model can be run on-device without the need to access peripherals and system functions.
 
\subsection{Principle 2: Enable Vendor Contributions to Span Ecosystem}

 Embedded devices come in all shapes and sizes, and require kernels to perform tinyML functions. The more optimized these kernels are for a particular device, the better performance will be achieved. However, because of the many differences between device platforms, there is no one-size-fits-all optimization solution. Consequently, the TFMicro team by itself is unable to support the wide variety of platforms that may want to run tinyML, and thus, vendors with strong motivation (i.e., those involved in microcontroller development) are encouraged to contribute to help bridge the gap. These vendors often have little experience with deep learning, and thus, TFMicro must provide sufficient resources to allow these teams to easily contribute. One way this is accomplished is by encouraging vendors to submit to a library repository and to provide tests and benchmarks for vendors to assess their hardware performance.
 
\subsection{Principle 3: Reuse TensorFlow Tools for Scalability}

 The third principle focuses on scalability. More than 1,400 operations (e.g. CONV2D) are supported by TensorFlow and other machine learning training frameworks. However, inference frameworks (i.e., those actually deploying the model) typically only support a fraction of these operations. For most use-cases, this will likely not cause issues since the most commonly used operations will likely be supported, but this inherent difference leads to a mismatch between the set of potential models produced by the training framework and the set of potential models that can be deployed by the inference framework.
 
 An exporter is used to convert a model from a training framework, such as TensorFlow, to a model for an inference framework, such as TFLite or TFLite for Microcontrollers. This model can then be deployed directly to a device and run using the library interpreter. Often, the training and inference frameworks are developed by different entities, which can present difficulties for developers when there are compatibility issues between the various stages of the developmental pipeline. This may render otherwise functional models unusable when trying to be deployed to a client device, especially when the incompatibilities are abstracted in high-level libraries such as Keras.
 
 Due to these concerns, the TFMicro developers decided to reuse as many TensorFlow tools available as possible to help minimize such complications and compatibility issues.
 
\subsection{Principle 4: Build System for Heterogeneous Support}

 The last principle focuses on promoting a flexible build environment. There are a large number of different types of embedded devices that may wish to use tinyML, and thus TFMicro should be designed without preference to any particular platform. This prevents vendor lock-in and also attracts a larger developer ecosystem due to improved portability. To combat this, TFMicro prioritizes code that can be built across a wide variety of integrated development environments (IDEs) and toolchains.
 
 These four principles help to facilitate a developer ecosystem that is oriented towards maximizing portability between various hardware platforms, architectures, frameworks, and toolchains.    
 
 
 
 \section{TinyML ``Keyword Spotting'' Workflow}
 
 In Course 2, you were introduced to the machine learning workflow. In this reading, we will expand on this, introducing the additional steps necessary to get our real-world deployment of the “Keyword Spotting” application running on our own microcontrollers. The reading gives you an overview of what is to come, while the videos that follow will help you understand the challenges of performing ML deployment effectively on-device.
 
 The first stages of the workflow are essentially the same as a traditional machine learning workflow: we collect and process our data, and then design and train a model for a specific machine learning task using the TensorFlow framework. The model will likely require an iterative design to meet performance goals. Following this, we move to TFLite to optimize our model (e.g., quantizing model weights into a suitable format) and then convert it into a suitable format, often a hexdump file (i.e., a binary file with its contents represented by hexadecimal values). Finally, we flash this hexdump file to the flash memory of the microcontroller, which is then able to perform inferences on the embedded device. The rest of this article will expand upon these stages with reference to our hello world application. The workflow is illustrated below.
 
 
 \includegraphics[width=\textwidth]{TinyML/3 Development/img3.png}  
 
 Machine learning workflow and pipeline. Begins with collection and preprocessing of data then the design and training of a model. These steps are done in TensorFlow. Then a model is optimized and converted and deployed using TensorFlow Lite. Finally inferences are made using TensorFlow Lite Micro.
 
 \subsection{Step 1: Data collection}
 
 As we explored in Course 2, collecting the dataset can often be the most challenging part of a tinyML application. Datasets must be large and specific in order to be applicable. For example, in the Keyword Spotting application we need a dataset aligned to individual words that include thousands of examples that are representative of real-world audio (e.g., including background noise). For our Keyword Spotting application, we will start off using the Speech Commands dataset from Pete Warden. Later on, we will show you how to collect your own data to train a new model.
 
 \subsection{Step 2 - Data Preprocessing}
 
 For efficient inference, we may need to extract features from the input signal to use for classification with our NN. First, we’ll need to take any analog input signals collected from our sensors and convert them into digital signals. Then we may need to apply additional transformations. For example, in the Keyword Spotting application, we will convert the audio signals into images through the use of spectrograms, all on-device!
 
  \subsection{Step 3 - Model Design}
  
 In order to deploy a model onto our microcontroller, we need it to be very small. We explored the tradeoffs of such models and just how small they need to be (hint: it’s tiny)! Since our Keyword Spotting application uses the “tinyConv” model, we should not need to worry about the size of the model, but parsimony should always be exercised when designing these models for the sake of efficiency - always have in mind, what is the smallest number of parameters that can be used to achieve a highly accurate model?
 
  \subsection{Step 4 - Training}
  
 When we build our own dataset, we will train our model using standard training techniques explored in Courses 1 and 2. Remember, training requires highly accurate computation (close to machine precision) since we are often working with small gradients during the training process, which largely precludes us from performing training on our embedded device. Because of this, we train our models using TensorFlow, as done previously.
 
 \subsection{Step 5 - Evaluation}
 We then test the final accuracy. In most industrial settings, initial testing is often poor and requires iterative design in order to match performance goals and the accuracy required by end-users. These performance goals are application-dependent and range from inference speed to model accuracy. For our Keyword Spotting application, we are only hoping to see a flashing light as our output, so there are no real performance goals beyond having a reasonable accuracy to detect a keyword of interest properly.
 
 \subsection{Step 6 - Conversion}
 
 In this stage, we take our functioning model and aim to compress and port it to our microcontroller device. First, we quantize our model weights and activations appropriately so that they can be represented by 8-bit or fixed-point arithmetic, and subsequently convert the model to a hexdump file. This stage is relatively simple to implement but is an important step to get right to ensure our model continues to function correctly when deployed on the microcontroller. The Keyword Spotting application needs no special treatment in this stage, all application models must go through this.
 
  \subsection{Step 7 - Write Deployment Code}
  
 Note that once we have our model ready, we also need all of the preprocessing code in C++ to be able to deploy our model to the Arduino with TFMicro. This involves the TFMicro runtime, as well as code to tell the Arduino how to respond to the information from the machine learning model. For our Keyword Spotting application, we have to tell it to turn on our LED and to report results over Serial!
 
  \subsection{Step 8 - Deploy the Full Application}
  
 The last thing to do is deploy the model. This stage involves loading the binary file into the application and using the Arduino IDE to flash the file to the microcontroller. Voila!
 
 \section{KWS Application Architecture}
 
     0:00 / 0:00Klicken Sie auf die Pfeiltaste nach oben, um das Geschwindigkeitsmenu aufzurufen. Regeln Sie dann mit den Hoch und Runter Pfeiltasten die Geschwindigkeit. Klicken Sie Enter, um die jeweilige Geschwindigkeit festzulegen.Geschwindigkeit1.0xKlicke auf diesen Knopf und diess Video stumm oder nicht stumm zu schalten, oder drücke die HOCH oder RUNTER Knöpfe um die Lautstärke zu erhöhen oder abzusenken.Maximal Lautstärke.Videotranskript
     Beginn des Transkriptes, zum Ende springen.
     VIJAY JANAPA REDDI: From this section on forward,
     one of the exciting things that we're going to be doing
     is actually learning how to deploy the applications onto our Tiny machine
     learning system.
     Now I'm going to start off with keyword spotting
     and then walk you through the other applications later on.
     Specifically in the context of keyword spotting,
     we're going to go pretty deep.
     I'm going to start off this video by talking about the overall application
     architecture because I'm going to use this
     as an example to help you understand other applications more systematically.
     As you know, keyword spotting is all about having
     a critical keyword like "Okay, Google" and getting that
     to wake the machine up.
     Now we have talked about this in the previous courses,
     and we've talked about all aspects of it outside of the deployment phase.
     We've talked about how to collect the data set properly,
     we learned how to train the models, we learn
     how to optimize them and convert them for deployment.
     But we never quite got around to actually deploying
     the model, which is what this course and this module is all about.
     So how do we go about doing this?
     Well, for that I need you to understand what the application architecture looks
     like for keyword spotting.
     For any application there are generally two main structures.
     There's the initialization phase and then there's the main loop.
     The main loop is really where the grunt of the work
     really sits because it's something that is constantly
     running, whereas initialization is something that you just do once.
     In our keyword spotting application, in the main loop
     there are many different phases.
     Don't be alarmed.
     I'm going to walk you through each and every single one of them in detail
     so that you understand what's going on.
     The way I've structured this video and the next couple of videos
     is such that I first introduce the core components that are there,
     and then in the next few videos I'll actually
     go deeper into each of those components and help
     you understand the constructs better.
     Before we get started, here is a demo of Brian
     doing an example of keyword spotting as you will be doing later
     on in the lab portion.
     BRIAN: Yes.
     No.
     And there we go.
     VIJAY JANAPA REDDI: As you can see, when Brian says the keyword,
     the machine lights up.
     And that's exactly what you want.
     Now let's try and understand how we actually do this.
     Starting from the very beginning where we have the microphone, where
     we actually have to capture the data and then provide it in some way
     so that we can then extract features out of that audio signal.
     Once we get the features, which are coming
     in the form of this continuous signal that will then
     convert into spectrograms, all of which you learned in Course 2,
     we will then be able to feed the spectrograms
     into the machine learning framework.
     That is specifically TensorFlow Lite for microcontrollers.
     We're going to set the interpreter up in order to execute the TinyConv
     model that you know all too well at this point, which
     uses a few different types of operators in order
     to make sure the model is nice and small while still
     being able to recognize keywords.
     Once the system has actually made an inference
     and it thinks it's recognized something or it has not recognized something,
     the command recognizer module kicks in.
     And the command recognizer module is basically
     to detect whether you said yes, no, up, down,
     or maybe there was really nothing that was being said.
     And once it makes a decision, it lets the command responder
     know where you actually actuate the physical device itself.
     Specifically in our context, we're going to be using the LEDs.
     So that effectively covers the full gamut of the different components
     that are involved.
     Now there's a lot of nuance into some of these components which
     I'm going to break down in the next couple of videos.
     Specifically, we're going to start off by understanding the initialization
     phase.
     What are the things that we have to do?
     Then we'll move into what we call is the preprocessing stage.
     The preprocessing stage is the stage that
     comes right before we do inference.
     Then we do inference.
     After we do inference on the network, then comes the postprocessing phase.
     And these are all important conceptual components
     that I want you to keep in mind because once you have that sort of a framework,
     then you can port that framework over to different types of applications
     because that structure is not going to change.
     Only the details within each of those structures change.
     With that said, let's get started by understanding how to initialize
     a keyword spotting application.
     Ende des Transkriptes, zum Anfang springen
     Downloads und Transskripte
     Video
     Videodatei herunterladen
     Transskripte
     Download SubRip (.srt) file
     Download Text (.txt) file
     
  \section{KWS Initialization}
  
      0:00 / 0:00Klicken Sie auf die Pfeiltaste nach oben, um das Geschwindigkeitsmenu aufzurufen. Regeln Sie dann mit den Hoch und Runter Pfeiltasten die Geschwindigkeit. Klicken Sie Enter, um die jeweilige Geschwindigkeit festzulegen.Geschwindigkeit1.0xKlicke auf diesen Knopf und diess Video stumm oder nicht stumm zu schalten, oder drücke die HOCH oder RUNTER Knöpfe um die Lautstärke zu erhöhen oder abzusenken.Maximal Lautstärke.Videotranskript
      Beginn des Transkriptes, zum Ende springen.
      VIJAY JANAPA REDDI: In this video, I'm going
      to walk you through keyword spotting's initialization step.
      Remember that the keyword spotting application
      has got many different components to it, going from the initialization
      into the main loop.
      Now, the initialization, which is the first step,
      has got many subcomponents to it.
      Don't be alarmed by the large number of steps we have to take.
      I'm going to walk you through each and every single one of these
      in greater detail.
      The most important thing to take away from this slide
      is that these are the steps you're going to be doing regardless
      of what application you're building.
      Let's get started.
      The very first step is to simply declare the variables.
      Specifically, you want to declare the variables that correspond to tflite.
      For example, that you need an error report,
      that there's a Model pointer, that there's
      a pointer to the MicroInterpreter, and so forth.
      These will tend to happen regardless of the application you're doing.
      I'm going to get more into the details in the screencast.
      For now, just kind of get the general idea.
      The second portion really focuses on the tensor arena.
      Now remember, you have to experimentally determine
      how big a tensor arena is going to be.
      Bigger is always OK from an execution standpoint, as long as the model fits.
      But remember, if you make the tensor arena too big,
      then you're effectively wasting space, which you might not be needing at all.
      So you want to try and figure out, what's the minimum amount of tensor
      arena size that you really want?
      And then, allocate it an array that can actually hold that space.
      After that comes loading the model, which
      you can do using the get model function, which takes a pointer called g\_model.
      G\_model is nothing but a pointer to an array that
      actually contains a hexadecimal dump of the actual neural network we're using.
      In our case, it's a tiny [? conv ?] network, which is effectively
      expressed here in a raw format.
      Now, out of these, you've actually got to determine, what are those operators?
      Now, you've got to know what the operators are
      that are used by a neural network.
      In the context of tiny [? conv ?] it happens
      to be Reshape, DepthwiseConv2D, FullyConnected, and Softmax.
      And you tell TF Micro that these are the ops you're using,
      as such by showing you the code on the right.
      Now, the key thing to remember here is that you could always
      use the all ops resolver.
      But you don't want to do that, yet again,
      because you want to keep the amount of binary space
      that's occupied by TF Micro, as an executable, small.
      And you can do that by only bringing in the ops that
      are used by the neural network, which in our case, DepthwiseConv2D,
      FullyConnected, Softmax, and Reshape.
      And that allows us to only use the code that's
      needed for actually executing the neural network.
      Now that we have all the bits and pieces,
      we can actually allocate or instantiate the interpreter,
      which we're doing here by passing in the model, the micro\_op\_resolver,
      the tensor\_arena, the size of the tensor arena, and the error\_reporter.
      Once we're done with that, we can actually allocate the memory
      for the tensor arena where the model is actually going to be executing out of.
      Next, comes actually defining the model inputs.
      Now, remember that you're going to have your input audio
      data coming in from someplace.
      Now, that might get processed.
      And then, it needs to be handed off to the interpreter, which
      is actually going to run that input through the network.
      So you have to link those two things together.
      And the first thing you do here is actually figure out
      the interpreter's input coming from.
      And then, you express that by making the connection to \PYTHON{model\_input}.
      More on this in the screencast, but the key takeaway
      here is that you need to be able to link the input that the model is
      going to be getting to the actual input that you're
      going to be providing it with.
      Now that we have effectively set up everything
      we need for TensorFlow Lite Micro itself,
      then we can setup other parts that are relevant to our application,
      such as the FeatureProvider and the RecognizeCommands.
      FeatureProvider is the one that, in keyword spotting,
      will generate the spectrograms.
      RecognizeCommands is the one that will actually
      do something based on the words that it might recognize, such as yes and no.
      All that said, I've just walked you through the steps
      from declaring the variables down to setup main loop.
      Now remember, again, the takeaway here has really
      got to be that these are the general steps you're going to be following
      regardless of the application.
      All that said, in the future slides, I'm just
      going to be using Initialize Interpreter to refer to many of the steps
      above because these are things that we will be doing repeatedly.
      So just kind of keep that in mind.
      With that said, in the next video, I'm going
      to pick up talking about the actual code in the screencast
      and help you understand how to navigate the large number of files
      we will be looking at.
      Ende des Transkriptes, zum Anfang springen
      Downloads und Transskripte
      Video
      Videodatei herunterladen
      Transskripte
      Download SubRip (.srt) file
      Download Text (.txt) file 
      
\section{KWS Initialization Screencast}

    0:00 / 0:00Klicken Sie auf die Pfeiltaste nach oben, um das Geschwindigkeitsmenu aufzurufen. Regeln Sie dann mit den Hoch und Runter Pfeiltasten die Geschwindigkeit. Klicken Sie Enter, um die jeweilige Geschwindigkeit festzulegen.Geschwindigkeit1.0xKlicke auf diesen Knopf und diess Video stumm oder nicht stumm zu schalten, oder drücke die HOCH oder RUNTER Knöpfe um die Lautstärke zu erhöhen oder abzusenken.Maximal Lautstärke.Videotranskript
    Beginn des Transkriptes, zum Ende springen.
    PROFESSOR: This is the first of many screencasts
    that I will be doing in order to help you understand some of the TensorFlow
    Lite Micro code in a bit more detail than we looked at in the videos
    previously.
    As you may recall, there are many different components
    involved in the Keyword Spotting application,
    starting with the initialization phase.
    In this particular screencast, I'm going to focus on the initialization stage.
    And I'm going to walk you through the different steps that are involved.
    I'm going to do this in a lot of detail this particular time around,
    but in subsequent applications, I will go a little bit
    faster just because it is a repetitive template that you'll find
    used in many different applications.
    So looking at our IDE, the first thing you will see here
    is that there are a couple of header files.
    Many of these header files are all about the TensorFlow Lite Micro.
    For example, line 26 through 31 are all about setting up the default TensorFlow
    Lite Micro ops.
    There are a couple of application-specific files here.
    Now, these correspond to activities that are
    part of the application architecture, which we'll get into later.
    Now, going back to my, slides one of the first steps
    is to actually declare the variables.
    So let's look at what the code says.
    Scrolling down a little bit, between line 34 and line 50, what you see
    are the actual variable declarations that are needed.
    The first part of the code is mostly centered around things
    that are needed for the TensorFlow Lite Micro interpreter itself.
    There is an ErrorReporter which is responsible for logging errors.
    You will find this in almost every single application.
    Then there is the model pointer, which actually holds the flat buffer
    model that we need to load that corresponds to application.
    For example, here we're interested in the tinycom
    network that does the keyword spotting.
    Then we have the interpreter itself, which
    we have to initialize with a certain amount of required state.
    Then we have an input pointer for the input
    that's going into the neural network.
    And the next two correspond to actually the application specifics,
    which are about getting the features that are needed, in other words,
    where is the spectrogram sitting, and the output that comes out
    of the network, which is stored in recognized commands, which tells us
    if a keyword was actually recognized or not--
    more on the previous time later.
    And the second half of the variable declaration scope
    says that we first have to define the tensorArena.
    Now, if you remember from what Pete said and what I have said before,
    the tensorArena is extremely critical.
    And you have to size it appropriately if you do not want your application
    to crash with an error.
    So here you have experimentally determined what the size needs to be.
    And then you go ahead and allocate the tensorArena buffer,
    which has got a certain size.
    Then we allocate the feature buffer, which
    contains the actual spectrogram features which we will be using in order
    to do the inference on.
    And it's of a certain size as well.
    And then we have the model input buffer.
    This is basically pointing to a buffer that
    will be initialized later in order to tell the interpreter where
    the inputs are coming from.
    The next step in our timeline is to actually load the model.
    So if I was to scroll down a little bit and go into the setup routine, which,
    as you should recall, is only executed once,
    you will find that here there's a function called GetModel.
    GetModel's job is effectively to take a pointer to a single-dimensional array
    which has actually got the flat buffer file
    schema, which defines the neural network model that we want to load and run.
    And the next couple of lines of basically doing some error checking
    to make sure that the schema version matches, because you always
    want to make sure that the network output file that you have
    in your schema is actually corresponding to the latest TensorFlow version
    that we will be using.
    So this is quite standard.
    And I expect you to do the same.
    After that comes resolving operators.
    Scrolling down a little bit further, you will
    see that between line 79 and line 99, we have a chunk of code which
    is all about the op resolver.
    The op resovler is taking in four different ops in here.
    Now, how do you know what the ops are?
    Well, you should have
    looked up what the ops are using the Netrun
    app that I told you about earlier.
    And if you look at the network architecture using Netrun,
    you would know that these are the different ops.
    Now, if you do not know what ops that network is actually using,
    let's say, hypothetically, you missed one of the ops,
    it's not the end of the world.
    TF Lite Micro will actually tell you through an error
    that you are missing one of the ops when it's actually trying to run it.
    So that's one of the ways you can easily figure out,
    what are all the options that you want to include in your ops resolver?
    As Pete and I both mentioned, you have the AllOps resolver,
    which is a cop-out way of doing op registration.
    In other words, if you do not want to explicitly register all these ops,
    you can easily say just use all ops, bringing all the ops.
    But of course, as a TinyML engineer, at this point
    you should know what the pros and cons are of bringing in all the ops.
    Bringing in all the DevOps is not necessarily a good thing,
    because it ends up wasting a lot of space.
    So keep that in mind.
    Then as we scroll down further, you will see
    that we are arriving at the next major step, which
    is initializing the interpreter.
    Initializing the interpreter requires us to pass a whole bunch of variables
    in here, starting with the model that we're actually bringing in,
    what are the ops, where is the tensorArena, how big is it,
    and what is the way by which we have to report errors.
    Once you have this bit of code here, you are effectively
    ready to start using the interpreter.
    Then going back to our timeline, we see that the next step
    is to allocate the arena, which is what is happening on line 107.
    We are actually allocating the tensors.
    This is using the tensorArena space that we had previously
    told the interpreter to use.
    Once again, we do error checking, as always.
    It's a good practice to have.
    Finally, after that, we tell the neural network
    where the model's inputs are really coming from.
    And that is happening between line 113 and line 123.
    So what we are asking the interpreter to do
    is to give us a reference to where its input buffers are actually stored.
    Now, you will see that it says the model input has got a dimension of 2.
    Now, you might be wondering, why does the input buffer have two dimensions?
    Well the first one is purely a wrapper.
    That's why you're seeing that here, when we
    look at the dimensions of the data of the first element,
    it's got only one entry in it.
    And that single entry, which you, then, if you poke at,
    will just have a single size one.
    And subsequently, the second entry that's in there
    is actually our spectrogram size.
    In other words, it is a k feature slice count times k feature slice.
    Now, if you remember how big our spectrogram is,
    it's going to be 40 columns wide.
    And it's going to be 49 rows in height.
    And that's what this is, slice count times slice size.
    So if you want to look this up, let's do a quick search.
    I'm going to do k feature slice size.
    I'm going to find it.
    It's used in a couple of places.
    But we'll get to the actual definition soon.
    There it is.
    It's declared as a constant.
    And you can see the k slice size is 40, which is effectively our column size.
    And k slice count is 49, which is the number
    of rows we have in our spectrograms.
    And so now if we go back up to our micro speech file,
    we're just making sure that the spectrogram dimensions which
    we are going to be feeding into the network
    are exactly the same as what the network is expecting.
    So this is actually what the network is expecting.
    And of course, we've got to make sure this is the size that we pass it in.
    And we want to make sure that the data type is correct, that it's not a fp32,
    it's actually an int8 value that we're passing in.
    And that's effectively what we're defining here.
    Finally, the next step in our process is to set up the main loop.
    In setting up the main loop, you can clearly
    see down here is that there are two stages here.
    The first one is called setting up the feature provider.
    The feature provider is the one that actually accesses the audio buffer
    and generates spectrogram.
    And that spectrogram is what we actually feed into the network later on.
    And after you do the inference, that's when
    we go into recognize commands, which will actually help us recognize.
    More on this later-- for now, all you need to know
    is that we're just simply setting these variables up.
    All right, with all of that said, you should at this point
    be quite comfortable having gone all the way from declaring variables,
    to loading the model, to resolving the ops, to initializing the interpreter,
    to allocating the arena, to defining the model inputs
    and setting up the main loop.
    This is a process that we will repeat over, and over, and over
    in all the different applications.
    So get comfortable with it.
    And you will naturally be feeling comfortable with it
    as you start writing a little bit more code.
    With that said, I'll pick you up in the next video.
    Ende des Transkriptes, zum Anfang springen
    Downloads und Transskripte
    Video
    Videodatei herunterladen
    Transskripte
    Download SubRip (.srt) file
    Download Text (.txt) file        
    
\section{KWS Pre-processing}

    0:00 / 0:00Klicken Sie auf die Pfeiltaste nach oben, um das Geschwindigkeitsmenu aufzurufen. Regeln Sie dann mit den Hoch und Runter Pfeiltasten die Geschwindigkeit. Klicken Sie Enter, um die jeweilige Geschwindigkeit festzulegen.Geschwindigkeit1.0xKlicke auf diesen Knopf und diess Video stumm oder nicht stumm zu schalten, oder drücke die HOCH oder RUNTER Knöpfe um die Lautstärke zu erhöhen oder abzusenken.Maximal Lautstärke.Videotranskript
    Beginn des Transkriptes, zum Ende springen.
    VIJAY JANAPA REDDI: Now, we're going to start
    getting into the pre-processing stage of keyword spotting.
    Pre-processing is a term that we typically
    use to refer to the set of tasks that we perform right before we do inference.
    Now, this is part of the main loop and not part of the initialization.
    So focusing on the main loop.
    Here, you see that the pre-processing contains of two steps.
    First is the audio provider.
    And second is the feature extractor.
    The job of the audio provider is really to deal
    with setting up the microphone to get the audio input from an analog
    into a digital representation.
    If you remember in course two, we talked real briefly
    about how to get an audio signal that is a continuous signal which
    we sample at about 16 kilohertz that we then
    convert into a digital representation.
    Now, all of this A-to-D is analog to digital conversion
    is done by libraries that we are piggybacking
    on that come equipped with Arduino.
    So the first step that we do is actually configure the Arduino's microphone.
    And there's a very simple function called that we have written
    that does exactly that.
    Now, don't worry about the details.
    I'll walk you through the screencast as to what
    this function is doing later on.
    But for now, just know that there is a function that
    actually sets up the microphone.
    That's the key takeaway.
    Typically, this function ends up being quite dependent upon the device.
    So, for instance, if you are using a SparkFun Edge 2 or a Himax board,
    then you're going to have to rewrite these functions to be quite
    specific for that particular platform.
    Keep that in mind as you try and extend outside the scope of the things
    that we are covering.
    So the audio provider's next job then is once the microphone has been configured
    to set up for a certain sampling [INAUDIBLE],,
    it's got to continuously read in the analog voltage signals that
    are coming in through the microphone and get a digital representation
    and dump the digital representation into an audio buffer.
    And that is a process that is repeated continuously
    as we need it because we need to continuously
    be sampling for audio signals to know when
    the user might have said a keyword.
    Now, the way we do that is by having this particular function here.
    This function effectively gets us the actual audio samples.
    To get our audio samples is populating a buffer
    that is going to store the audio signals as a 16-bit value inside each and every
    single of the entries.
    Now, once we have that audio buffer, now that
    is effectively the raw digital sampling of the audio signal.
    What we then have to do is extract features from that.
    Now, if you remember in course two, one of the key things that we said
    is that we take a continuous signal.
    And we convert that continuous audio waveform into a spectrogram-- something
    that a neural network can easily ingest and process
    as a convolutional neural network can for a given image.
    And out of that spectrogram is how we infer
    whether we actually have a certain key word that's being spoken--
    fairly simple.
    So how exactly do we do this?
    Well, in the implementation, we take a 30-millisecond clip.
    That is one snapshot.
    And out of that one snapshot, we generate as FFT.
    The FFT, recall, is a much more effective representation
    of the time series signal in the frequency domain.
    And we get 256 values out of that.
    And in order to make the FFT more representative and robust,
    we end up doing averaging which then leaves us with 43 unique values.
    These 43 unique values are effectively one slice of the spectrogram
    that we have.
    What we end up doing after we get one slice is then,
    we shift the red box that you are seeing on the screen a little bit.
    That's the second time step effectively.
    Those dotted lines on the waveform are the new set of recordings
    that we're going to be pulling in.
    So one of the key things to notice in the figure here right now
    is look at the overlap between the solid red box and the dotted red box.
    And what you will realize there is that we are
    sliding that window by 20 milliseconds.
    Now, remember, I said that the box itself
    is 30 milliseconds, which means if I slide it over by just 20 milliseconds,
    and that means that I'm going to have about 10 milliseconds worth of data
    still from the past.
    And that's actually a good thing because you
    want to have some sort of overlapping features.
    That's what we want to capture.
    And then we continue this process effectively
    until we fill up the spectrogram.
    And this behavior is actually handled by the feature provider.
    And the feature provider's job, which we defined earlier in the initialization
    stage, which then in the loop stage continues to repeat over and over,
    is to actually fill that buffer up.
    So what it's doing is it's getting audio samples and waiting for the buffer
    to actually replenish itself all the way to the FULL 1 second worth of data.
    Once we have 1 second worth of actual spectrogram data in there,
    then we can actually do an inference.
    So bottom line.
    The way things are tying up is that you've got this audio signal.
    That audio signal is actually a time to signal which we
    have converted into a frequency domain.
    And from the frequency domain, what we can do are additional optimizations.
    If you remember in course two, we talked about how we do MFCC.
    MFCC real simply is the fact that our ear naturally
    picks up low frequency signals much better than higher frequency signals.
    And therefore, if we focus our effort or energy
    more on the low frequency signals, then we're
    likely to pick up the words better.
    So what the bottom spectrograms are showing is like,
    what happens when you take a regular spectrogram
    and then you convert it with MFCC?
    It emphasizes the lower frequencies much more cleanly.
    And that allows our machine learning model
    to pick up the signals for the keywords much more clearly.
    Now, one thing I want you to keep in mind
    is this process might really vary depending
    on the capabilities off your microcontroller.
    For example, if you have a very powerful microcontroller,
    you might choose to run a much more complicated FFT and a much more
    sophisticated MFCC method.
    But if your microcontroller is rather simple and it's very power efficient,
    then you might not have the flexibility to do this kind of computationally
    in terms of test.
    So in which case, you might choose to implement something
    quite simple for your FFT and MFCC, which
    is something that we do in our TensorFlow like, micro example
    and micro speech.
    So all of this generating features going all the way
    from that audio signal to the spectrogram
    is actually done by this function called GenerateMicroFeatures.
    We'll look at the code that's involved within GenerateMicroFeatures later
    on in the screencast.
    But for now, conceptualize that this is the function that you should focus on
    for spectrogram generation aspects.
    This basically is taking the audio signals
    and generating spectrograms for us.
    Now, once you've got a spectrogram, then you
    feed it to the inference stage, which is where we will be going next.
    Now, to take a step back, this is effectively the keyword spottings
    preprocessing nugget.
    You've got the audio provider which knows how to take the analog signals
    and then convert them into a digital representation.
    And then you've got the feature extractor
    who takes those digital representations and converts them into a spectrogram.
    And this is the pre-processing that happens
    before we call it KWS's inference, which is effectively
    invoking the tiny [? COM ?] model, which I'll talk about next.
    Ende des Transkriptes, zum Anfang springen
    Downloads und Transskripte
    Video
    Videodatei herunterladen
    Transskripte
    Download SubRip (.srt) file
    Download Text (.txt) file    
    
 \section{KWS Pre-processing}
 
 
     0:00 / 0:00Klicken Sie auf die Pfeiltaste nach oben, um das Geschwindigkeitsmenu aufzurufen. Regeln Sie dann mit den Hoch und Runter Pfeiltasten die Geschwindigkeit. Klicken Sie Enter, um die jeweilige Geschwindigkeit festzulegen.Geschwindigkeit1.0xKlicke auf diesen Knopf und diess Video stumm oder nicht stumm zu schalten, oder drücke die HOCH oder RUNTER Knöpfe um die Lautstärke zu erhöhen oder abzusenken.Maximal Lautstärke.Videotranskript
     Beginn des Transkriptes, zum Ende springen.
     VIJAY JANAPA REDDI: Now, we're going to start
     getting into the pre-processing stage of keyword spotting.
     Pre-processing is a term that we typically
     use to refer to the set of tasks that we perform right before we do inference.
     Now, this is part of the main loop and not part of the initialization.
     So focusing on the main loop.
     Here, you see that the pre-processing contains of two steps.
     First is the audio provider.
     And second is the feature extractor.
     The job of the audio provider is really to deal
     with setting up the microphone to get the audio input from an analog
     into a digital representation.
     If you remember in course two, we talked real briefly
     about how to get an audio signal that is a continuous signal which
     we sample at about 16 kilohertz that we then
     convert into a digital representation.
     Now, all of this A-to-D is analog to digital conversion
     is done by libraries that we are piggybacking
     on that come equipped with Arduino.
     So the first step that we do is actually configure the Arduino's microphone.
     And there's a very simple function called that we have written
     that does exactly that.
     Now, don't worry about the details.
     I'll walk you through the screencast as to what
     this function is doing later on.
     But for now, just know that there is a function that
     actually sets up the microphone.
     That's the key takeaway.
     Typically, this function ends up being quite dependent upon the device.
     So, for instance, if you are using a SparkFun Edge 2 or a Himax board,
     then you're going to have to rewrite these functions to be quite
     specific for that particular platform.
     Keep that in mind as you try and extend outside the scope of the things
     that we are covering.
     So the audio provider's next job then is once the microphone has been configured
     to set up for a certain sampling [INAUDIBLE],,
     it's got to continuously read in the analog voltage signals that
     are coming in through the microphone and get a digital representation
     and dump the digital representation into an audio buffer.
     And that is a process that is repeated continuously
     as we need it because we need to continuously
     be sampling for audio signals to know when
     the user might have said a keyword.
     Now, the way we do that is by having this particular function here.
     This function effectively gets us the actual audio samples.
     To get our audio samples is populating a buffer
     that is going to store the audio signals as a 16-bit value inside each and every
     single of the entries.
     Now, once we have that audio buffer, now that
     is effectively the raw digital sampling of the audio signal.
     What we then have to do is extract features from that.
     Now, if you remember in course two, one of the key things that we said
     is that we take a continuous signal.
     And we convert that continuous audio waveform into a spectrogram-- something
     that a neural network can easily ingest and process
     as a convolutional neural network can for a given image.
     And out of that spectrogram is how we infer
     whether we actually have a certain key word that's being spoken--
     fairly simple.
     So how exactly do we do this?
     Well, in the implementation, we take a 30-millisecond clip.
     That is one snapshot.
     And out of that one snapshot, we generate as FFT.
     The FFT, recall, is a much more effective representation
     of the time series signal in the frequency domain.
     And we get 256 values out of that.
     And in order to make the FFT more representative and robust,
     we end up doing averaging which then leaves us with 43 unique values.
     These 43 unique values are effectively one slice of the spectrogram
     that we have.
     What we end up doing after we get one slice is then,
     we shift the red box that you are seeing on the screen a little bit.
     That's the second time step effectively.
     Those dotted lines on the waveform are the new set of recordings
     that we're going to be pulling in.
     So one of the key things to notice in the figure here right now
     is look at the overlap between the solid red box and the dotted red box.
     And what you will realize there is that we are
     sliding that window by 20 milliseconds.
     Now, remember, I said that the box itself
     is 30 milliseconds, which means if I slide it over by just 20 milliseconds,
     and that means that I'm going to have about 10 milliseconds worth of data
     still from the past.
     And that's actually a good thing because you
     want to have some sort of overlapping features.
     That's what we want to capture.
     And then we continue this process effectively
     until we fill up the spectrogram.
     And this behavior is actually handled by the feature provider.
     And the feature provider's job, which we defined earlier in the initialization
     stage, which then in the loop stage continues to repeat over and over,
     is to actually fill that buffer up.
     So what it's doing is it's getting audio samples and waiting for the buffer
     to actually replenish itself all the way to the FULL 1 second worth of data.
     Once we have 1 second worth of actual spectrogram data in there,
     then we can actually do an inference.
     So bottom line.
     The way things are tying up is that you've got this audio signal.
     That audio signal is actually a time to signal which we
     have converted into a frequency domain.
     And from the frequency domain, what we can do are additional optimizations.
     If you remember in course two, we talked about how we do MFCC.
     MFCC real simply is the fact that our ear naturally
     picks up low frequency signals much better than higher frequency signals.
     And therefore, if we focus our effort or energy
     more on the low frequency signals, then we're
     likely to pick up the words better.
     So what the bottom spectrograms are showing is like,
     what happens when you take a regular spectrogram
     and then you convert it with MFCC?
     It emphasizes the lower frequencies much more cleanly.
     And that allows our machine learning model
     to pick up the signals for the keywords much more clearly.
     Now, one thing I want you to keep in mind
     is this process might really vary depending
     on the capabilities off your microcontroller.
     For example, if you have a very powerful microcontroller,
     you might choose to run a much more complicated FFT and a much more
     sophisticated MFCC method.
     But if your microcontroller is rather simple and it's very power efficient,
     then you might not have the flexibility to do this kind of computationally
     in terms of test.
     So in which case, you might choose to implement something
     quite simple for your FFT and MFCC, which
     is something that we do in our TensorFlow like, micro example
     and micro speech.
     So all of this generating features going all the way
     from that audio signal to the spectrogram
     is actually done by this function called GenerateMicroFeatures.
     We'll look at the code that's involved within GenerateMicroFeatures later
     on in the screencast.
     But for now, conceptualize that this is the function that you should focus on
     for spectrogram generation aspects.
     This basically is taking the audio signals
     and generating spectrograms for us.
     Now, once you've got a spectrogram, then you
     feed it to the inference stage, which is where we will be going next.
     Now, to take a step back, this is effectively the keyword spottings
     preprocessing nugget.
     You've got the audio provider which knows how to take the analog signals
     and then convert them into a digital representation.
     And then you've got the feature extractor
     who takes those digital representations and converts them into a spectrogram.
     And this is the pre-processing that happens
     before we call it KWS's inference, which is effectively
     invoking the tiny [? COM ?] model, which I'll talk about next.
     Ende des Transkriptes, zum Anfang springen
     Downloads und Transskripte
     Video
     Videodatei herunterladen
     Transskripte
     Download SubRip (.srt) file
     Download Text (.txt) file   
     
\section{KWS Inference}

    0:00 / 0:00Klicken Sie auf die Pfeiltaste nach oben, um das Geschwindigkeitsmenu aufzurufen. Regeln Sie dann mit den Hoch und Runter Pfeiltasten die Geschwindigkeit. Klicken Sie Enter, um die jeweilige Geschwindigkeit festzulegen.Geschwindigkeit1.0xKlicke auf diesen Knopf und diess Video stumm oder nicht stumm zu schalten, oder drücke die HOCH oder RUNTER Knöpfe um die Lautstärke zu erhöhen oder abzusenken.Maximal Lautstärke.Videotranskript
    Beginn des Transkriptes, zum Ende springen.
    VIJAY REDDI: In this video, I'm going to talk about Keyword Spotting's
    inference phase.
    Remember, that there are two main stages, the initialization
    stage and the main loop.
    And specifically, when it comes to inference, it is in the main loop,
    and it comes after the pre-processing stages.
    Recall, that the pre-processing is really
    the audio provider and the feature extractor whose job is
    to take the raw audio signal and convert it into a format that
    is conducive for us to do inference on.
    You've done something like this, for example,
    in course one where, for instance, when you were dealing with the MNIST data
    set, you actually took the MNIST two-dimensional image
    and created a flattened one-dimensional tensor that you then fed into the MLP
    in order to do classification.
    Now, we're going to do something similar to that.
    We're going to take our spectrogram in a flattened format,
    and then feed that into the tiny comm network.
    Not necessarily the MLP here, but you get
    the idea about how we flatten a two-dimensional tensor
    and feed that into a neural network.
    Now, that's all done through this very simple piece of code.
    The top "for loop" is basically showing you the feature buffer data being
    copied into the model input buffer.
    The feature buffer is really where we're actually
    extracting the features and kind of populating that buffer
    with the key spectrogram data.
    And the model input buffer is what we had previously
    tied as input source for the model for TF Micro
    so, therefore, we have to make this data shuffling.
    Once the data has been shuffled, then we can actually invoke the interpreter.
    And that's all there really is.
    For the inference stage, it's quite simple.
    The key thing is to really focus on where
    is the input coming from and making sure that input
    from the pre-processing stage is ready for the interpreter to take it.
    Then, we call the neural network model, which is really simple.
    In our case, it's a simple tiny comm model that we're going to invoke.
    Now, the key takeaway that I want you to keep in mind
    is inference is really quite simple.
    And the most important phase here before you
    call inference is making sure that you're setting up the input buffers
    properly for TF Micro to consume.
    That's you need to remember.
    After this, we will be picking up with what do
    we actually do from the model's output.
    How do we actually process that?
    It's actually not as simple as it might seem.
    In order to detect yes or no, there are actually quite a lot of nuances.
    So I'll see you in the post-processing video.
    Ende des Transkriptes, zum Anfang springen
    Downloads und Transskripte
    Video
    Videodatei herunterladen
    Transskripte
    Download SubRip (.srt) file
    Download Text (.txt) file     
    
\section{KWS Inference Screencast}

    0:00 / 0:00Klicken Sie auf die Pfeiltaste nach oben, um das Geschwindigkeitsmenu aufzurufen. Regeln Sie dann mit den Hoch und Runter Pfeiltasten die Geschwindigkeit. Klicken Sie Enter, um die jeweilige Geschwindigkeit festzulegen.Geschwindigkeit1.0xKlicke auf diesen Knopf und diess Video stumm oder nicht stumm zu schalten, oder drücke die HOCH oder RUNTER Knöpfe um die Lautstärke zu erhöhen oder abzusenken.Maximal Lautstärke.Videotranskript
    Beginn des Transkriptes, zum Ende springen.
    VIJAY JANAPA REDDI: In this screencast I'm
    focused on how you do the inference for keyword spotting.
    Now you should at this point be able to recognize
    that inference is going to happen in the main loop
    because this is something that has to happen continuously.
    So as you can see here, if I were to scroll down just a little bit,
    after we do the preprocessing, then comes an inference stage.
    Now, if you remember, the preprocessing is
    done by PopulateFeatureData, which puts data into the feature buffer.
    In order to do the inference we actually have to get the data that's inside
    \PYTHON{feature\_buffer} and move that over to the \PYTHON{model\_input\_buffer}.
    And we do this by simply doing a raw data copy.
    Now one of the reasons that we actually write a manual loop
    and go through and iteratively copy it is
    so that we don't actually have to rely on library functions that
    actually do the copy.
    Why do you think we do not want to rely on library functions?
    Huh?
    Well, Pete said that we don't want to have library dependencies because when
    we have library dependencies, that increases
    the size of the actual executable, which in our case
    can be quite problematic because we're deploying things on a microcontroller.
    Therefore, often you will find us doing things
    such as manually copying the data over.
    After we have moved the input values into the model input buffer,
    this is where the spectrogram effectively is sitting right now.
    Because we have already set up the model input buffer previously
    and tied it to the interpreter, if you recall from the initialization phase,
    all we can do now is just call the interpreter through the Invoke function
    call, which effectively runs the model on the spectrogram
    data that's inside the model input buffer and then generates an output.
    I'll talk more about the output in the postprocessing stage.
    For now, when we call the interpreter, we simply
    want to make sure that there are no errors that
    are reported by the interpreter.
    At all times, I strongly encourage you to have these kinds of error checks.
    They will help you build a robust application that will help you
    debug the application really easily.
    With that said, I'll pick up in the postprocessing phase.
    Ende des Transkriptes, zum Anfang springen
    Downloads und Transskripte
    Video
    Videodatei herunterladen
    Transskripte
    Download SubRip (.srt) file
    
\section{KWS Post-processing}

    0:00 / 0:00Klicken Sie auf die Pfeiltaste nach oben, um das Geschwindigkeitsmenu aufzurufen. Regeln Sie dann mit den Hoch und Runter Pfeiltasten die Geschwindigkeit. Klicken Sie Enter, um die jeweilige Geschwindigkeit festzulegen.Geschwindigkeit1.0xKlicke auf diesen Knopf und diess Video stumm oder nicht stumm zu schalten, oder drücke die HOCH oder RUNTER Knöpfe um die Lautstärke zu erhöhen oder abzusenken.Maximal Lautstärke.Videotranskript
    Beginn des Transkriptes, zum Ende springen.
    VIJAY JANAPA REDDI: We are now going to talk
    about keyword spotting's post-processing stage, the last stage in the main loop.
    Again, as a quick reminder, we've got the initialization,
    and we've got the main loop.
    And the post-processing is something that comes about
    after you're done with the inference.
    And this is part of the main loop, which gets repeated over and over.
    Now, post-processing might seem like a little bit of a weird thing.
    You've already got your neural network that has made a prediction.
    What else do you really need to do?
    Well, often, when you're deploying things,
    there is a bit of post-processing that you have to do in order
    to get confidence in the results.
    This is something that is often not caught when you're just
    focused in the Colab where you're training a model,
    and you're not really seeing the full end-to-end picture.
    Let me give you a very concrete example.
    Let's say you're training your device to detect the keyword up.
    Here's a question for you.
    What happens if the user says upward?
    Is the machine supposed to wake up?
    Or is the machine supposed to ignore that?
    Well, obviously, your first reaction might
    be, well, the user said upward, not up.
    So yeah, the machine should not wake up.
    OK, well, that sounds fairly simple to you and me
    to be able to figure that out.
    But how does a machine know that?
    And the complexity really comes from the fact
    that it depends on where you read the sound.
    For example, if you look at this particular clip here,
    let's say the one second--
    if you've got a nice, clean one second of the full word
    that the user said and then you do your inference,
    well, is that going to be up or upward?
    Well, naturally, the network is not going
    to detect the word up here because it's nicely aligned for the full word.
    So it's going to say nothing.
    It's not a keyword at all.
    But then imagine if it so happened that, when
    you were creating your spectrograms in the preprocessing stage,
    you got your one-second word, the clip, right up
    until the mid-portion of the word, which means
    that you've only heard the word up, and you're taking a snapshot of that.
    And then you're feeding it into the network.
    What do you think is going to happen?
    Well, you're going to get a false positive there.
    The reason is because the machine only knows what it's just heard.
    It just saw the spectrogram from the word up in there.
    So it says yes, this is something that seems quite reasonable.
    This is totally possible.
    I'm not saying it happens all the time, but it's
    something that could totally happen and does in fact happen with our own model.
    So how can we improve this situation is the key question.
    Now, up is just one example.
    In reality there are many other examples like that.
    For example, think about no.
    Think about all the other words that are there, which happen
    to have no right at the beginning.
    The first one that comes to my mind is notable, "no-table."
    Well, if it picks up the word no and then you
    run it immediately through your inference,
    it's going to wake the machine up.
    So there are many other ones.
    In fact, I encourage you to kind of think of all the other words
    where such a weird situation might actually crop up.
    Now, how can we make things robust to this?
    Well, one important thing that we can do is that, after an inference is done,
    we don't immediately try and trigger the device
    to actually respond to something.
    Instead, we get the machine to listen to the spectrogram a couple of times.
    So we collect a couple of spectrograms in order to get confidence.
    In other words, we do preprocessing.
    We do inference.
    We collect some data.
    We don't immediately respond to it.
    We try and get some confidence.
    So we go back.
    We collect another bunch of data, do some preprocessing,
    get the new spectrogram, do the inference on that,
    see what the value ends up being, and repeat this process a few times.
    So the key takeaway is that, when we're doing keyword spotting,
    we actually do multiple inferences before we actually generate a result.
    Let me be a bit more specific here.
    So, for example, for every window, what we do
    is we store the results that come in from a new inference.
    And we calculate the new score for the keywords
    that the model is supposed to predict and keep a history of that.
    And we average them over multiple iterations.
    So, every time a new window comes in, we keep a running average
    of what the confidence is for a given keyword.
    If we do this enough times and we do this quickly enough,
    as in we make multiple inferences quickly enough,
    then we can actually have a very strong sort of a predictor.
    And that's the example I was showing here.
    So, for example, like let's say, at time step 1,
    you have some amount of confidence across these keywords like yes and no.
    Time step 2, after you do another inference,
    your confidence goes up that it's a yes.
    Time step 3, time step 4, time step 5, voila, in general, what you're seeing
    is that your confidence for that particular keyword is going up.
    So you're effectively averaging and building more hysteresis in there.
    And, in doing so, you can actually get a very good prediction.
    How do we do this in the code?
    Well, there is a function called ProcessLatestResults that actually
    implements this exact sort of behavior.
    So I'll go into the code later on when I do the screencast, but, for now,
    just understand that this is a core principle
    that you might want to leverage, as you build
    your own keyword spotting applications.
    You don't need to get into the details of this unless you really want to,
    but, by just calling this ProcessLatestResults,
    it's actually got the mechanisms to handle it.
    Now, once you've actually gotten some level of confidence,
    then you call the RespondToCommand.
    Now, the RespondToCommand's job is to do something
    based on a recognized command.
    In our case, it's all about triggering LEDs in different lights,
    depending on the keyword.
    For example, if it is a yes, it's going to be green.
    If it's a no, it's going to be red.
    If it's a word that it doesn't recognize, it's a blue.
    And that's how this example GIF that we're looking at here is working.
    The system is detecting a work.
    It's averaging things a couple of times by doing multiple inferences.
    And then it's actually seeing if it needs to respond.
    And that is how, my friend, keyword spotting really works.
    So post-processing is the key thing that I want
    you to walk away with in this video.
    It's not as simple as it seems.
    Now, not every application is going to have the same level of complexity.
    It really depends on what the use case is.
    For example, when we do the visual wake words or person
    detection model later on, you will realize that the post-processing
    is actually quite simple.
    In fact, you will find the keyword spotting is actually
    quite a complicated application by comparison to the other examples
    that we look at.
    And that's a good thing because I kind of
    want you to start with some of these core, important, fundamental building
    blocks rather than progressively increasing the complexity on you
    because then you would really not know where the ends would be.
    At least this way, you kind of know that keyword spotting
    is a rich enough application where you're
    getting exposed to some very important fundamental concepts
    about deploying an end-to-end TinyML application.
    In the next video, I'm going to try and summarize some of these core concepts
    for you so you know how to walk away in terms
    of understanding the code and the structure of these TinyML applications.
    And I'll see you there.
    Ende des Transkriptes, zum Anfang springen
    Downloads und Transskripte
    Video
    Videodatei herunterladen
    Transskripte
    Download SubRip (.srt) file
    Download Text (.txt) file    
    Download Text (.txt) file    
 
\section{KWS Post-processing Screencast}

    0:00 / 0:00Klicken Sie auf die Pfeiltaste nach oben, um das Geschwindigkeitsmenu aufzurufen. Regeln Sie dann mit den Hoch und Runter Pfeiltasten die Geschwindigkeit. Klicken Sie Enter, um die jeweilige Geschwindigkeit festzulegen.Geschwindigkeit1.0xKlicke auf diesen Knopf und diess Video stumm oder nicht stumm zu schalten, oder drücke die HOCH oder RUNTER Knöpfe um die Lautstärke zu erhöhen oder abzusenken.Maximal Lautstärke.Videotranskript
    Beginn des Transkriptes, zum Ende springen.
    SPEAKER: In this screen cast, I'm going to walk you
    through the post-processing phase of keyword spotting.
    Now if you remember the post-processing phase
    comes right after the inference stage and it's got two major components.
    The first stage is the command recognizer
    and the second stage is the command responder.
    As a quick recap, the command recognizer job
    is to make sure it's detecting the right key word.
    In other words, let's say if the word is up,
    then it's got to make sure it has detected the word up.
    Not only that, it's actually got a more important job,
    which is to make sure it robustly detects the key word up.
    So for example, the user says upward, then
    its job is to make sure that the system does not accidentally
    trigger a false positive, as i had previously explained.
    So let's get into the code and try and understand how this is actually done.
    Now, the post-processing stage comes right after the inference stage.
    Line 162 is exactly where we had done the inference.
    After the inferences happened, you get an output.
    The output is effectively the output that
    has come out from the neural network, the Tiny comm network that we run.
    What happens immediately after that, is that we pass the output to process
    latest results.
    As you can see here, the process latest results function
    is taking output as an argument, it's taking the current time,
    it's taking a reference to found command.
    If it determines that a command has indeed
    been found, even after it's done it's robustness checks,
    then it will set this variable to be true.
    And for the command, it's going to assign it a certain score
    and will determine if a new command was found or if it was an old command,
    it would simply be false.
    With that said, let's dig right into process latest results,
    and see what is actually going on here.
    So I'm going to do a quick search for it and that brings us
    into the recognized command CPP file.
    And you can see that the output tenser that we
    were looking at in this particular function is called latest results.
    The first thing that's happening here, is
    that we're doing a quick error check on latest results Tensor.
    We want to make sure that the output tensor that
    is coming in from the neural network, is exactly of the form
    that we expect it to be.
    In other words it got dimensions equal to 2
    and we want to make sure that it's only got
    a single element in the first entry.
    And in the second entry, it's got to have kcategory count.
    In other words, at this point you've got to make sure
    that the output that's coming out in data one,
    the number of softmax probabilities you have there,
    matches the exact number of categories we are trying to label.
    Let me be more specific and say what is kcategory count.
    Category count is talking about the yes, no, silence and unknown.
    That's four different categories.
    So the count is actually four.
    And each one of these here is actually going
    to be telling us if it's yes, no, silence or no.
    And so it's making sure that, that output Tensor has a right format
    and that it's off the right data type.
    After that, because we will actually be collecting
    a queue of different inference records in order
    to get robustness of confidence as I had previously explained,
    we are actually making sure here is doing yet
    another round of self checking.
    That any new entry that is coming in, is actually being put in the right order
    into the queue.
    So, if that's not the case, we say the results
    must be fed in increasing time order.
    This is a simple sanity check, to make sure that all the results are new
    and no old results are getting pushed in.
    Because of the order of the results is not being maintained
    and that can lead to confusion in how we're processing the data.
    Assuming that everything is clear, the very first thing we do
    is we store the new result that is just come in, by pushing it
    into the back of the queue.
    The next thing we want to do, is to make sure
    that we're not keeping an infinite number of prior guesses.
    What we're instead doing, is that we're only
    trying to keep a limited number of entries in there.
    And the way we do that, is by trying to restrict ourselves to about 1,000
    milliseconds or so, which is what this variable has been defined to do.
    So there's a certain time limit, and if the old values
    kind of fall outside of this time limit, then they're
    effectively edged out by popping them out of the queue.
    Pretty simple.
    So long story short, the goal of this piece of code
    is to make sure that we're only looking at some reasonable number of items,
    most recent items.
    After that, we want to check if we have sufficient results.
    If we don't have enough of a history in there,
    then there's no point of actually making confident result.
    This is part of the robustness that this particular function
    is offering, in order to make sure that the prediction that's going to make
    is robust against different sorts of input variables.
    The next thing we go into, assuming everything is good,
    is we want to calculate what's the average score for each of the entries
    that I previously talked about and compute a running average.
    This is effectively a way of building confidence
    by making sure you're done the inferences a few times and then
    over that time, even within a one second period,
    you've got some confidence, that yes, I am actually looking at the right word
    based on what the user has seen.
    And that's the average calculation that's happening here.
    After we have done the average calculation,
    we're simply trying to find which category actually
    has the highest score in the list of variables that we have.
    And that's being done by simply looking at the current top index based
    on the average scores.
    And we're seeing if the current average score,
    is greater than the previous average score that was recorded.
    Once we find that, we are determining that the current top index corresponds
    to a certain category label.
    For example, here if I was to look up what the category labels, where
    you can see that the category labels in the default scenario
    are silence, unknown, yes, no and so forth.
    If I go right back to the file that I was in, which is a recognized command
    CPP, back to where we were.
    Now, we are at the next stage.
    The next stage what we're trying to do, is
    to make sure that there is some level of confidence in the command that
    is actually being analyzed.
    In other words, if this time we got a no, but just recently, we got a yes,
    and that sort of a bouncing is happening back and forth,
    then we don't have confidence in the results.
    And that's effectively what this code is trying to do.
    It's trying to see if the labels are actually coming out consistently.
    And if it turns out that everything is good,
    then we will say that yes, there is a new command that we have detected.
    At that point, we can happily go back to our primary function.
    I'm going to scroll down real quickly to where we we're.
    We were right here.
    We were inside this function.
    At this point, we have got some process data,
    we have actually analyzed the history of the commands
    and so we should be quite confident, in terms of what the command output was.
    So then we go on to the second stage.
    If you remember the second stage is the actual response stage.
    In the response stage, we want to trigger that LED.
    So let's go ahead and see what's happening inside the respond
    to command.
    The response to command is taking a couple
    of different input variables, which we will
    look at inside the respond to command.
    So, here we are, inside the Arduino command responder file.
    We are going to be working with LED's, as a means of actuating the system
    in response to a certain command.
    You might choose to actuate a motor or something to that effect, that's
    up to your liking.
    And you can see here, that the first thing this code actually does,
    is actually check if the code has been previously initialized.
    This is only done the first time when you're inside
    the main loop doing the steady state.
    This is not actually getting executed.
    The core part of this function is really this region.
    If there is a new command, and if it turns out that the command is actually
    a yes, then we want to set the LED that is green in color,
    to a low value, which will turn the LED on.
    Similarly, if the key word that was heard was no, then
    we want to set the red LED on.
    We do that by simply writing low.
    We are digital, right?
    Last but not least, if the word is unknown,
    then we set the blue LED to low.
    That causes the system to actually turn the blue LED on.
    Finally, the last thing to do is to see if there is a need to turn the LED off.
    For example, let's say the first time you say yes, now that turns the LED on.
    Well that's great.
    And then it stays on for a while.
    So what we're doing here, is that if the previous command happened
    to be like about three seconds ago and we turned the LED on, then
    here we want to simply turn it off.
    And that's what's going on here.
    And that's all there really is.
    At this point, you should be quite familiar with the overall structure.
    Preprocessing in summary, involves making sure
    that we're confident about the result that is being generated
    and that's being handled by process latest results.
    And after process latest results gives us a confidence response,
    then we determine what to do with that respond.
    And here we are turning on LEDs.
    I hope at this point you're feeling quite
    comfortable with the overall structure of the code.
    As always, take a step back, think about the code
    in terms of what's the initialization step, what's the preprocessing
    step, what's the inference step, and finally
    what's the post-processing step.
    And that's all there is.
    Ende des Transkriptes, zum Anfang springen
    Downloads und Transskripte
    Video
    Videodatei herunterladen
    Transskripte
    Download SubRip (.srt) file
    Download Text (.txt) file    
    
 \section{KWS Summary}
 
     0:00 / 0:00Klicken Sie auf die Pfeiltaste nach oben, um das Geschwindigkeitsmenu aufzurufen. Regeln Sie dann mit den Hoch und Runter Pfeiltasten die Geschwindigkeit. Klicken Sie Enter, um die jeweilige Geschwindigkeit festzulegen.Geschwindigkeit1.0xKlicke auf diesen Knopf und diess Video stumm oder nicht stumm zu schalten, oder drücke die HOCH oder RUNTER Knöpfe um die Lautstärke zu erhöhen oder abzusenken.Maximal Lautstärke.Videotranskript
     Beginn des Transkriptes, zum Ende springen.
     VIJAY JANAPA REDDI: You're almost to the end of understanding
     the keyword spotting application.
     After this, you're going to be working with Brian
     and actually deploying the keyword spotting model
     onto your embedded microcontroller.
     But, before I let you go over to Brian, I
     want to help you understand how to digest the vast amount of complexity
     that you might be feeling at this point.
     I know that we stepped through many different functions.
     Then we discussed about all the different topics
     that are involved around the code constructs.
     And we looked at a lot of code in the form of raw files.
     So, if you're feeling overwhelmed, don't worry.
     There's a method to the madness, and I'm going to teach you
     how to think about this systematically.
     First and foremost, there's a whole notion
     around setting up the environment.
     So there's a set of hows that are involved there.
     Then you go into setting up the ops and the interpreter.
     And that requires you to look at a couple of functions,
     usually in the setup portion of the code.
     Then we get into actually ingesting data from the microphone, which
     is the preprocessing part of the application.
     That can be a little bit complicated.
     And remember there, you have a device-specific portion,
     which is where you are dealing with the specifics
     of a particular microcontroller, for instance,
     and the way it interfaces with the I/O subsystem,
     specifically from the embedded sensor part.
     And there are non-device-specific aspects
     where you are just simply processing raw data,
     converting that input buffer data of audio signals into the spectrogram.
     Keep that in mind because those two things are a little different.
     And, once you have that, you invoke the interpreter with the right data.
     So you're typically copying that raw signal
     that you've perfected into features and then giving that over
     to the interpreter, which then performs the inference.
     And then it generates outputs.
     And that output is something that we act on
     by looking at what the actual output is from the network,
     be it a recognized word or not, and then decide
     if we're going to actuate some sort of LED or another subsystem.
     So most of the code is kind of largely sprinkled along these axes.
     So this could very well be a framework for the way
     you need to think about digesting this vast volume of code.
     That's lesson number one.
     Now, another thing around this is that there are many different components
     that we step through.
     Now, how do you kind of make sense of what to change and what not to change?
     Let's say you want to make some sort of tweak to the application.
     The way you've got to think about this is not thinking of the whole code
     because that's overwhelming.
     What you should really do is recognize the fact
     that the end-to-end application is really broken up
     into these four big building blocks.
     There's the setup.
     There's the preprocessing.
     There's the inference.
     And then there's the post-processing.
     And you tend to repeat the cycle between preprocessing and post-processing
     over and over and over.
     So pay attention to the different parts that I've
     listed out as examples that fall underneath each one
     of these components.
     For example, in preprocessing, you have the KISS FFT,
     or that's where you would do the MFCC.
     So let's say if you want to enhance the spectrogram that you're actually
     generating.
     Well, then you just need to focus on the preprocessing,
     and you don't need to really worry about the other pieces.
     But keep in mind though, if you change your preprocessing step here before you
     call the network, you've got to make sure your network has
     been trained accordingly in the earlier stages, like in course 2.
     Keep that in mind.
     You cannot simply change the preprocessing step and expect
     the network to perform better because the network has learned the features
     that preprocessing is actually doing.
     So it's extremely important to keep in mind
     that you would have to go back and retrain
     the network with a new set of preprocessing steps that you're doing,
     something that is actually quite subtle, but very important.
     Similarly, if you want to improve the network's performance
     through optimizations, well, then focus on what the network is.
     And, if you want to improve the confidence of the system,
     then focus on the post-processing side of things.
     You don't have to look at the whole code as one big, monolithic block
     sprinkled across different files.
     There is structure in the chaos that you might be witnessing.
     So, if you want to modify something, identify
     what you want to modify, and only focus on that part is the key lesson here.
     One of the reasons that I have specifically broken things down
     into the setup, preprocessing, inference, and post-processing
     is because it's a rather important message from the end-to-end application
     and design perspective that I want to convey over to you.
     Now, remember, right at the beginning of this course,
     I said that we're going to do things a little different here.
     We're not just interested in the model capabilities,
     and we're not just interested in coming up with the smallest neural network.
     We're actually interested in orchestrating the full flow
     and being able to design that end-to-end application
     so that it runs effectively.
     And, to that note, know that setup, preprocessing, and post-processing
     are just as important if not, in some sense,
     more important than the actual neural network model
     that we would actually do inference on.
     Why do I say this?
     Think about it from the execution latency perspective.
     You have to do setup once.
     And then, after that, you have to do the preprocessing, inference,
     and post-processing all in real time while the device is actually running,
     trying to detect your keyword.
     What that means is that they all add latencies.
     Now, this is just a pictorial representation
     of what the latencies might be.
     I'm not necessarily saying that preprocessing is always
     the most biggest chunk.
     It might be a smaller chunk.
     Post-processing might be bigger.
     Either way, the point that I'm trying to convey here
     is that they have different amounts of overhead.
     And those overheads can be sometimes more larger than the inference itself.
     So, when you're trying to optimize your device for performance,
     you need to keep this in mind.
     Similarly, from memory as well, you need to keep this in mind.
     Preprocessing, post-processing, setup, they all
     consume that limited amount of memory that we have on our microcontroller.
     So we need to keep that in mind, as we try and optimize the system.
     This is also the reason that earlier I talked to you about the embedded system
     software stack because that Mbed OS and all of those things
     actually eat up memory.
     So you need to be conscious about where you're
     going to tune the memory resource requirements if you have to do so.
     Bottom line is that, both in terms of execution latency,
     as well as memory usage, there might be different overheads that
     are coming from all the non-ML aspects.
     In other words, if I take inference out, there
     is some AI tax that is left behind, as I like to call it.
     And that AI tax is a pure overhead.
     And so, when we're designing applications
     to be efficient from an end-to-end perspective,
     as we need to do in keyword spotting to actually get it to work,
     well, then we need to also pay attention to how quickly we're
     able to do this processing.
     This is why, for instance, our version of the FFT in the preprocessing stage
     is actually very simple.
     It's actually not the best FFT that you can run.
     It's actually called KISS, which stands for Keep It Simple, Stupid.
     That's literally what it stands for.
     But you understand the notion of keeping things simple and why it's important.
     So keep AI tax in mind because that'll help you be a better TinyML
     engineer because you will understand what it really means
     to design an end-to-end application.
     Now, with all that said, I'm going to leave you
     in the good hands of Brian who's going to help you take that model
     and deploy it onto your embedded system.
     And also, later on, we will be changing that model in response
     to a new data set.
     So a lot of exciting stuff is coming up.
     So stay tuned for that, and have fun deploying the model.
     Ende des Transkriptes, zum Anfang springen
     Downloads und Transskripte
     Video
     Videodatei herunterladen
     Transskripte
     Download SubRip (.srt) file
     Download Text (.txt) file   
     
\section{Deploying the Pretrained KWS Model}

    In this reading, we are going to first generate a binary file representing the pretrained KWS model and then deploy that model to our Arduino using the Arduino IDE.
    
\subsection{Converting the TFLite Model File into a Binary Array}

    The first step in deploying the pre-trained KWS model is to take the quantized TensorFlow Lite model you created in Course 2 and convert it into a binary array that we can load onto our microcontroller. And don’t worry if you don’t still have the .tflite file from Course 2; we have provided a direct link to the staff copy in the following Colab.
    
    Please run this Colab and then, either leave the tab open with the printout of the final binary file and/or download the final .cc model file and open it up in a separate text editor. In either case, we will need to copy the data in that file into our Arduino sketch in the next section, so make sure you keep it in an easy to access place!
    
    https://colab.research.google.com/github/tinyMLx/colabs/blob/master/4-5-16-KWS-PretrainedModel.ipynb
    
\subsection{Deploying the Pretrained KWS Model}

    0:00 / 0:00Klicken Sie auf die Pfeiltaste nach oben, um das Geschwindigkeitsmenu aufzurufen. Regeln Sie dann mit den Hoch und Runter Pfeiltasten die Geschwindigkeit. Klicken Sie Enter, um die jeweilige Geschwindigkeit festzulegen.Geschwindigkeit1.0xKlicke auf diesen Knopf und diess Video stumm oder nicht stumm zu schalten, oder drücke die HOCH oder RUNTER Knöpfe um die Lautstärke zu erhöhen oder abzusenken.Maximal Lautstärke.Videotranskript
    Beginn des Transkriptes, zum Ende springen.
    BRIAN PLANCHER: Hello again, everybody.
    We're now going to take the binary model that you just generated in Colab
    and deploy it to your Arduino using the IDE.
    We're going to do this by opening up the micro speech example, which
    is the keyword spotting example.
    We'll find that in the Harvard TinyMLX library.
    So we'll go back to file examples, scroll down to Harvard TinyMLX
    and open up micro speech.
    We're then going to go over here to the right and there's this little arrow
    and we're going to click on that and that's
    going to open up all of the files, a lot of which Vijay just mentioned.
    Now the file we're interested in right now
    is called micro features model dot CPP.
    And that's where the actual model data lives.
    So if we click on that file we're going to see
    a file that has a couple of pound to finds at the top,
    and then it has this huge binary array.
    And that's what we want to copy into from what we just generated from Colab.
    So if you go back to wherever you have saved that model file,
    whether that's in Colab itself or in some other tab
    or as some downloaded file, if you just highlight all of this binary data--
    now make sure you don't get all of the syntax up at the top.
    We just want the binary bytes.
    So if we go down to the bottom and copy all of that,
    again ignoring the semicolons, parentheses,
    all of that kind of stuff--
    we can then go back to the Arduino IDE, and we could highlight, again,
    just the binary data.
    And the reason why that's important is because the syntax is actually
    different in the dot CC model file and in the Arduino file.
    So let's highlight all of this.
    And then we can paste that new model data into this file.
    Now, the last thing we have to do is we have to go back
    and we have to get the model length.
    Copy that.
    And paste that as well into the file.
    Now you might notice that in this case, the model length is the same.
    And that's because there was already a different pre trained yes/no
    model in here.
    But we wanted to make sure you saw how to copy a model file over.
    And so we're just going to overwrite that pre trained copy with our own pre
    trained copy.
    So if you can just save that, what you'll notice
    is that a pop up comes up that says that this is a read only file.
    And that's because all of the examples are read only by default.
    So you can retain them as examples.
    And so what we're going to do is click OK to find a new location to save it.
    One thing that might be helpful for this course, for your own organization,
    is to generate maybe a folder called, say, TinyML.
    Save it in there.
    And I actually am going to call this micro speech underscore pre
    trained, just to keep things organized.
    And that's it.
    With that all in place, you can get ready for your standard deployment.
    So go to Tools, select the board and select the port.
    And then you can just click Upload to start compiling and uploading
    your code.
    And, again, don't be worried if you see some warnings that show up there
    or if it takes a while to compile, we're working with leading edge code
    and sometimes you get a couple of warnings
    and sometimes the first time it compiles there's a lot of libraries to compile,
    but pretty soon it should load onto your board.
    And then we can explore the result of deploying your first TinyML model.
    True, hardcore TinyML application to your Arduino.
    All right now that your model is deployed, let's test it out.
    We can do that by opening up the serial monitor,
    and we should probably immediately start seeing a bunch of Heard unknowns,
    as I am not saying one of the key words, yes or no.
    And I said yes or no too fast there for it to necessarily be able to pick up.
    And we're still seeing lots of unknowns.
    But, anyway, you can also see this happening on the board.
    And so what you'll see is that this blue LED down here
    is signifying what it found.
    And so blue is for unknown.
    Green is for yes and red is for no.
    And what you can also do is, if I get quiet for a second--
    on the top left you saw that blinking yellow LED.
    And that is sort of when it says it's primed
    and really ready to take a new keyword.
    So why I give it a yes and a no and we'll
    see if it can turn the LED green and red and show up in the serial monitor.
    Yes.
    No.
    And there we go.
    You have a pretty decent keyword spotting
    model deployed to your own Arduino right in front of you.
    Hopefully you can have some fun playing around with it.
    So I'll see you in the next video.
    Ende des Transkriptes, zum Anfang springen
    Downloads und Transskripte
    Video
    Videodatei herunterladen
    Transskripte
    Download SubRip (.srt) file
    Download Text (.txt) file
    
    
\subsection{Updating the Arduino Code}

\begin{enumerate}
  \item Open the \FILE{micro\_speech.ino} sketch, which you can find via the File drop-down menu. Navigate, as follows: File → Examples → Harvard\_TinyMLx → micro\_speech.
    
  \item Navigate to the \FILE{micro\_features\_model.cpp} file and update the model. You can find that file by selecting it from tabs across the top of the Arduino IDE. If that file is not visible, you can navigate to it (or other additional files) by clicking on the downward facing triangle at the end of the tabs, which will open up a dropdown showing all of the files.
    
  \includegraphics[width=\textwidth]{TinyML/3 Development/4-5.1-re.png}
    
    A Screenshot of the Arduino IDE showing that the main code starts on line 35 the type of the model and ends with a \{. When you are pasting in your model all of that definition needs to remain, You are only replacing the hexidecimal values which follow. Note all hexidecimal values are values in base 16 represented by 0-9a-f preceded by 0x so for example, 0x20, 0x00, 0x5f etc.
     
      \begin{itemize}
        \item Copy the binary model file contents from the \FILE{KWS\_yes\_no.cc} file into the \FILE{micro\_features\_model.cpp} file. Make sure to only copy the binary data inside the \{ \} as the variable type is different in the downloaded or printed \FILE{KWS\_yes\_no.cc} file (it is of type unsigned char in the \FILE{.cc} file, but it needs to be of type const unsigned char with the \PYTHON{DATA\_ALIGN\_ATTRIBUTE} in the \FILE{.cpp} file). If you lost your \FILE{KWS\_yes\_no.cc} file, don't worry, you can use the staff's copy!

        \includegraphics[width=\textwidth]{TinyML/3 Development/4-5.2.png}

        A Screenshot of the Arduino IDE showing the end of the hexidecimal model values (as described in the previous image), ending with a \};. After that the next line says const int \PYTHON{g\_model\_len = 18712;}. Again you'll want to make sure to just update the numerical value.
    
        \item Next, scroll all the way down to the bottom of the file and replace the model length. Again, note that the .cpp file needs the variable to be of type const int while the .cc file will show unsigned int. Our suggestion again is to simply copy the numerical value.
        
        \includegraphics[width=\textwidth]{TinyML/3 Development/4-5.3.png}
        
        A Screenshot of the Arduino IDE showing where to copy and paste the model length.
      \end{itemize}        
    
    \item Next save your changes. You will most likely see a popup as shown below asking you to save a copy of the example as all examples are treated as “read-only” by default. treated as “read-only” by default.
    
     \includegraphics[width=\textwidth]{TinyML/3 Development/4-5.4.png}

      Screenshot of popup asking to save the sketch to a new location as it is ``read-only''.
    
      We suggest that you make a folder called e.g., TinyML inside of your Arduino folder. You can find your main Arduino folder either inside of your Documents folder or in your Home folder, and save it in that folder with a descriptive name like \PYTHON{micro\_speech\_edx\_pretrained}. That said, you can save it wherever you like with whatever name you want!

      \includegraphics[width=\textwidth]{TinyML/3 Development/4-5.5.png}
    
      Screenshot showing a TinyML folder inside of the main Arduino folder.
\end{enumerate}


    
\subsection{Deploying the Pretrained Model}

\begin{enumerate}
    \item Use a USB cable to connect the Arduino Nano 33 BLE Sense to your machine. You should see the green LED power indicator come on when the board first receives power.
    
    \item As always, use the Tools drop-down menu to select appropriate Port and Board.
    
      \begin{itemize}
        \item Select the Arduino Nano 33 BLE as the board by going to Tools → Board: <Current Board Name> → Arduino Mbed OS Boards (nRF52840) → Arduino Nano 33 BLE.

        \item Then, select the USB Port associated with your board. This will appear differently on Windows, macOS, Linux, but will likely indicate ‘Arduino Nano 33 BLE” in parenthesis. You can select this by going to Tools → Port: <Current Port (Board on Port)> → <TBD Based on OS> (Arduino Nano 33 BLE). Where <TBD Based on OS> is most likely to come from the list below where <\#> indicates some integer number.

          \begin{itemize}
            \item Windows → COM<\#>
            \item macOS → /dev/cu.usbmodem<\#>
            \item Linux → ttyUSB<\#> or ttyACM<\#>
        \end{itemize}
      \end{itemize}

    \item     Use the rightward arrow next to the ‘upload’ / flash the code. You’ll know the upload is complete when you see red text in the console at the bottom of the IDE that shows 100\% upload of the code and a statement that says something like “Done in <\#.\#> seconds.”
    
    If you receive an error you will see an orange error bar appear and a red error message in the console (as shown below). Don’t worry -- there are many common reasons this may have occurred. To help you debug, please check out our FAQ appendix with answers to the most common errors!
    
    \item You should now have a working Keyword Spotting model for “Yes” and “No” deployed to your Arduino! To test it out, open the Serial Monitor. You should start to see the result of the model “Yes, No, Unknown, Silence” begin to display in the Serial Monitor.

      \includegraphics[width=\textwidth]{TinyML/3 Development/4-5.6.png}
    
    Screenshot of the Serial Monitor showing the recognition of some utterances of yes and no as well as some unknown words.
    This should match the multicolor LED near the Bluetooth module on the board changing colors as follows (we’ve shown a picture of it turning green below):
    
    \bigskip
    
    Yes     = Green LED
    
    No = Red LED
    
    Unknown = Blue LED
    
    Silence = Nothing
    
      \includegraphics[width=\textwidth]{TinyML/3 Development/4-5.7.png}

    Screenshot of the Serial Monitor showing the recognition of some utterances of yes and no as well as some unknown words. It follows the patter Heard X (Y) @Zms where X is what it heard (in our case yes, no, unkown) and Y is the score output by the model and Z is a timestamp of when this occured from when the Arduino was turned on.
\end{enumerate}

\section{Deploying a KWS Model with Your Favorite Keyword(s)}

    In this reading, we are going to first generate a binary file representing the KWS model with our favorite keywords from Course 2 and then deploy that model to our Arduino using the Arduino IDE (which will include some code changes).
    
\subsection{Converting the TFLite Model File into a Binary Array}

    Again we are going to first convert the .tflite file into a .cc file in Colab for use with the Arduino IDE. This time let’s use a model with different keywords (and a different number of keywords), so we can explore more of the changes you will need to make to deploy custom models. We’ve provided you with one option for a model in the Colab below. That said, we also invite you to use your favorite KWS model you created in Course 2 if you happen to still have the .tflite file! If you’d like to go back and train a new .tflite file, you can find the training Colab at this link.
    
    As with the pretrained model, we will be using the resulting .cc file, so make sure to download it or leave the tab open with the printout!
    
    \url{https://colab.research.google.com/github/tinyMLx/colabs/blob/master/4-5-18-KWS-FavoriteKeywords.ipynb}
    
    
\subsection{Deploying a KWS Model with Your Favorite Keyword}

    0:00 / 0:00Klicken Sie auf die Pfeiltaste nach oben, um das Geschwindigkeitsmenu aufzurufen. Regeln Sie dann mit den Hoch und Runter Pfeiltasten die Geschwindigkeit. Klicken Sie Enter, um die jeweilige Geschwindigkeit festzulegen.Geschwindigkeit1.0xKlicke auf diesen Knopf und diess Video stumm oder nicht stumm zu schalten, oder drücke die HOCH oder RUNTER Knöpfe um die Lautstärke zu erhöhen oder abzusenken.Maximal Lautstärke.Videotranskript
    Beginn des Transkriptes, zum Ende springen.
    SPEAKER: All right.
    Let's now deploy a keyword spotting model with your favorite keywords.
    We'll start again with another basic micro speech example, just
    the default one that we opened up from the examples menu.
    And again, we will go and we will take our model file
    and we will copy just the binary contents of that file.
    And we will paste it into the file called micro features model.cpp.
    And we will just replace this giant binary model here
    with the binary model of the model with our favorite keywords.
    Again, be careful not to overwrite any of the syntax.
    And then this time, be extra careful to make sure
    that you copy over the model length, because it's
    very likely that the model length has changed since we no longer have
    just a yes/no model.
    We have something else.
    When we hit save-- again, it's going to ask us to find a new location
    to make a copy, since it's read only.
    And I'm going to call this one micro speech favorite.
    And I'll save it in that same folder that I
    saved the pre-trained model as well.
    So now we have this project and we've replaced the model.
    But we're not done yet because we've changed
    the number and the actual keywords themselves.
    So first we're going to go to micro features micro model settings.h.
    And in that file on line 40, right here.
    We're going to update the variable kCategory count
    and we're going to update that to be two plus the number of keywords we have.
    So that's silence unknown and then in our case, up, down, go.
    So that'll be five.
    We'll make sure to save that file.
    Then we will open micro features micro model settings.cpp.
    And here we'll see the actual list of the keywords
    where we'll need to change yes and no to up, down, and go.
    And we'll make sure to save that file too.
    OK.
    So now the last thing we need to do now that we have changed the model
    and changed the labels for that model, we're
    going to need to change the command responder.
    So we're going to open Arduino command responder.cpp.
    This is the file that actually does something when it finds a keyword.
    So in our case, we've been turning on LEDs.
    And you'll see right here in these if statements starting on line 54,
    it's looking to see if there's a variable
    found command, which is a string.
    So it's an array of characters, which are
    the letters of each of the key words.
    So found command 0 is the first letter.
    So it's looking for y for yes, n for no, and u for unknown.
    And that's how we're turning things green, red, and blue.
    So let's update this to our keywords.
    And the first problem you might immediately realize
    is, that up also starts with a u.
    But that's OK.
    What we can do is we can add a second part to this if statement.
    And we can say if the first letter was a u and if the second letter is a p,
    then we have a up.
    And so it's green for up.
    Down is simple.
    It's just an n.
    Red for down.
    And the go would also be easy, but I want to keep this unknown one.
    So I'm going to add actually, another if statement in here.
    And I'm going to say, if we find the d, up, down, go.
    If we find the g, then we're going to turn things blue.
    Blue for go.
    And then, what I'm going to do to unknown
    is I'm going to do again the double clause here
    to make sure that we are finding the correct word unknown and not just up.
    And what I'm going to do is I'm actually going to turn on all three LEDs,
    red, blue, and green.
    And together it will look kind of white.
    You'll be able to see the red and the green a little bit
    and the blue around the outside, but it will look mostly white.
    And so we'll call this white for unknown.
    And that's it.
    We save this file.
    So we've updated the Arduino command responder
    to respond correctly to our command.
    We've updated the micro model settings to the exact words
    that we actually are using.
    And one note here again, the order of these words matters.
    So if you trained it as up , down, go.
    It must be up, down, go.
    If you change the order, it will break things.
    So be careful when you train a model.
    Make sure you remember what order you put the key words in.
    We changed the command responder.
    This was the micro features micro model settings.cpp.
    We also changed the micro model settings.h.
    We changed the kCategory count and we changed the actual model
    itself, the micro features model.cpp.
    And once you've changed all of those, we are
    ready to first check and make sure that the board and the port
    are set correctly.
    And once they are, we're ready to compile and upload that to our device.
    And that's it.
    That's all you need to do and you'll be able to test out
    your model with whatever of your favorite keywords
    that you have deployed to it.
    So I hope you have fun with that.
    Ende des Transkriptes, zum Anfang springen
    Downloads und Transskripte
    Video
    Videodatei herunterladen
    Transskripte
    Download SubRip (.srt) file
    Download Text (.txt) file
    
\subsection{Updating the Arduino Code}

\begin{enumerate}
  \item Open the \FILE{micro\_speech.ino} sketch, which you can find via the File drop-down menu. Navigate, as follows: File → Examples → \FILE{Harvard\_TinyMLx} → \FILE{micro\_speech}.
    
  \item As before, navigate to the \FILE{micro\_features\_model.cpp} file and update the model. You can find that file by selecting it from tabs across the top of the Arduino IDE. If that file is not visible you can navigate to it (or other additional files) by clicking on the downward facing triangle at the end of the tabs which will open up a dropdown showing all of the files.
    
  \includegraphics[width=\textwidth]{TinyML/3 Development/4-18.1-re.png}
    
  Screenshot of the IDE showing the downward arrow and the list of files that it opens up as described in the reading.
  
  \begin{itemize}
  \item  Copy the binary model file contents from the \FILE{KWS\_favorite.cc} file into the \FILE{micro\_features\_model.cpp} file. Make sure to only copy the binary data inside the \{ \} as the variable type is different in the downloaded or printed \FILE{model.cc} file (it is of type unsigned char in the .cc file, but it needs to be of type const unsigned char with the \PYTHON{DATA\_ALIGN\_ATTRIBUTE} in the \FILE{.cpp} file). If you lost your \FILE{KWS\_favorite.cc} file, don't worry, you can use the staff's copy!
   
  \includegraphics[width=\textwidth]{TinyML/3 Development/4-18.1-re.png}
    
    A screenshot of the Arduino IDE show the arrow on the top right corner of the screen that has been clicked to show all of the additional files in this Arudino sketch. The file \FILE{micro\_features\_model.cpp} is highlighted as it needs to be selected to allow us to update the model.
    
    \item Next, scroll all the way down to the bottom of the file and replace the model length. Again, note that the .cpp file needs the variable to be of type const int while the .cc file will show unsigned int. Our suggestion again is to simply copy the numerical value.
    
    \includegraphics[width=\textwidth]{TinyML/3 Development/4-18.3.png}
    
    A Screenshot of the Arduino IDE showing where to copy and paste the model length.
 
     \end{itemize}

  \item   Next, save your changes. You will most likely see a popup as shown below asking you to save a copy of the example as all examples are treated as “read-only” by default.
  
    \includegraphics[width=\textwidth]{TinyML/3 Development/4-18.4.png}

  Screenshot of popup asking to save the sketch to a new location as it is “read-only”.
  
  Again, we suggest that you make a folder called e.g., TinyML inside of your Arduino folder. You can find your main Arduino folder either inside of your Documents folder or in your Home folder, and save it in that folder with a descriptive name like \FILE{micro\_speech\_favorite}. That said, you can save it wherever you like with whatever name you want!

    \includegraphics[width=\textwidth]{TinyML/3 Development/4-18.5.png}

    Screenshot showing a TinyML folder inside of the main Arduino folder.
    
  \item    The two other things we need to change are the count and list of keywords in the model settings files and the output response file. As we do that in the rest of this document, we are going to assume that you used our model with the Keywords “up,down,go.” If you choose to use different keywords, make sure that you update the following steps accordingly.
   
  \item  Navigate to the \FILE{micro\_features\_micro\_model\_settings.h} file and scroll down to line 40 and change the value of the \PYTHON{kCategoryCount} variable to be two more than the number of keywords you selected (one extra for silence and one extra for unknown). For example, if you had 3 keywords (as we do in our example model), we need to update the variable to be $3 + 2 = 5$ as shown:
  
  \PYTHON{constexpr int}
  
  \PYTHON{kCategoryCount = 5;}
    
  \item  Then, navigate to the \FILE{micro\_features\_micro\_model\_settings.cpp} file. You’ll find that it includes the .h file and otherwise has a single array. Update the values of that array to match your keywords. Make sure to leave ``silence'' and ``unknown'' in the list as well, and make sure to list your keywords in the order that you listed them on the training script. So, in the case of our example model, we would update the array to be:
  
  \begin{lstlisting}
    const char* kCategoryLabels[kCategoryCount] = {
        
        "silence",
        
        "unknown",
        
        "up",
        
        "down",
        
        "go",
        
    };
  \end{lstlisting}  
    
  \item  Finally, navigate to the \FILE{arduino\_command\_responder.cpp} file. Scroll down to line 54. There, you will see a series of if statements that control how the Arduino responds to each command. They are each structured to look at the value of the \PYTHON{found\_command} variable. This variable will be the string of the keyword (or ``silence'' or ``unknown'') you entered into the \PYTHON{kCategoryLabels} array in the previous step. So, by comparing the [0] value in that string, the if statements are simply looking at the first letter! In each if statement, you’ll also see a digitalWrite to either the LEDG, LEDR, or LEDB variable which is used to turn on the appropriate color LED. So, if we wanted to adapt this to our three keywords and, for example, only turn on the single colors for our keywords and turn all of the colors on for unknown, which will come out look white, we could update those three if statements to the following:
  
  \begin{lstlisting}
    // Red for up -- note here that you do not need to index
    
    //               into the first letter only, just a unique
    
    //               letter combination in the keyword! That
    
    //               said make sure you do not index beyond the
    
    //               end of ANY keyword or you will get an error!,
    
    if(found_command[1] == 'p') {
        
        last_command_time = current_time;
        
        digitalWrite(LEDR, LOW);
        
    }
  \end{lstlisting}  
    
  \begin{lstlisting}
    // Green for down
    
    if(found_command[0] == 'd') {
        
        last_command_time = current_time;
        
        digitalWrite(LEDG, LOW);
        
    }
  \end{lstlisting}  
    
  \begin{lstlisting}
    // Blue for go
    
    if(found_command[0] == 'g') {
        
        last_command_time = current_time;
        
        digitalWrite(LEDB, LOW);
        
    }
  \end{lstlisting}  
    
  \begin{lstlisting}
    // All three for unknown (white)
    
    if(found_command[1] == 'n') {
        
        last_command_time = current_time;
        
        digitalWrite(LEDR, LOW);
        
        digitalWrite(LEDG, LOW);
        
        digitalWrite(LEDB, LOW);
        
    }
  \end{lstlisting}  
  
\end{enumerate}

    
\subsection{Deploying the New Model}

\begin{enumerate}
  \item Use a USB cable to connect the Arduino Nano 33 BLE Sense to your machine. You should see the green LED power indicator come on when the board first receives power.
  
  \item  As always, use the Tools drop-down menu to select the appropriate Port and Board.
    
     \begin{itemize}
        
        \item Select the Arduino Nano 33 BLE as the board by going to Tools → Board: <Current Board Name> → Arduino Mbed OS Boards (nRF52840) → Arduino Nano 33 BLE.
   
        \item    Then, select the USB Port associated with your board. This will appear differently on Windows, macOS, Linux, but will likely indicate ‘Arduino Nano 33 BLE” in parenthesis. You can select this by going to Tools → Port: <Current Port (Board on Port)> → <TBD Based on OS> (Arduino Nano 33 BLE). Where <TBD Based on OS> is most likely to come from the list below where <\#> indicates some integer number.
    
      \begin{itemize}
        \item Windows → COM<\#>
        \item macOS → /dev/cu.usbmodem<\#>
        \item Linux → ttyUSB<\#> or ttyACM<\#>
       \end{itemize}
    \end{itemize}

    \item Use the rightward arrow next to the ‘upload’ / flash the code. You’ll know the upload is complete when you see red text in the console at the bottom of the IDE that shows 100\% upload of the code and a statement that says something like “Done in <\#.\#> seconds.”
    
    If you receive an error, you will see an orange error bar appear and a red error message in the console (as shown below). Don’t worry -- there are many common reasons this may have occurred. To help you debug, please check out our FAQ appendix with answers to the most common errors!
    
    \item You should now have a working Keyword Spotting model for either your favorite keywords from Course 2 or our model for “up,down,go” deployed to your Arduino! To test it out, open the Serial Monitor. You should start to see the result of the model begin to display in the Serial Monitor and the LED’s light up according to how you instructed them in the previous section!
   
\end{enumerate}
        
        
% 1.6: Custom Dataset Engineering for Keyword Spotting

\section{Recap: Dataset Engineering}

    In this module, you will learn about the nuances of data engineering, and develop your own dataset which you will use to train and deploy a custom keyword spotting model!
    
    Within these courses so far, we have used numerous datasets to train our algorithms. The existence of a dataset spared us from having to tackle the most time-consuming components of the machine learning pipeline: the collection and wrangling of data. In this section of the course, we are going to create our own dataset to perform keyword spotting, and expose you to some of the challenges that come along with the process, such as dataset bias, the need for edge cases, and the difficulty associated with procuring large quantities of data. This procedure is aptly named dataset engineering.
    
    Data engineering is an important aspect of supervised learning and entails the specification of the dataset, as well as the procedures for collecting and processing data. Naturally, this involves defining what features will be present, the structure of the data (e.g., images, time series, row tuples), and other characteristics.
    
    One of the key features of data engineering is identifying potential sources of data, which may come from external sources (free or purchased), crowdsourcing, existing users of a service, sensors, or other means. When developing our dataset, it is imperative that we consider our use case at all times, since this is our ultimate goal.
    
    A natural corollary of this is that we must determine the environment our device will be in once deployed. For example, if an algorithm processing sound data expects to be used on the street, it is necessary to make recordings in the presence of adversarial sounds such as construction work, loud vehicles, as well as other background noise. We may have to decide what time of day images are taken, and also make sure that each of the classes in our dataset has roughly equal representation to prevent biasing to the majority class. 
    
\section{Introduce Custom Dataset for KWS}

    0:00 / 0:00Klicken Sie auf die Pfeiltaste nach oben, um das Geschwindigkeitsmenu aufzurufen. Regeln Sie dann mit den Hoch und Runter Pfeiltasten die Geschwindigkeit. Klicken Sie Enter, um die jeweilige Geschwindigkeit festzulegen.Geschwindigkeit1.0xKlicke auf diesen Knopf und diess Video stumm oder nicht stumm zu schalten, oder drücke die HOCH oder RUNTER Knöpfe um die Lautstärke zu erhöhen oder abzusenken.Maximal Lautstärke.Videotranskript
    Beginn des Transkriptes, zum Ende springen.
    VIJAY JANAPA REDDI: Welcome back.
    At this point, I hope you're getting comfortable deploying the TinyML
    models onto your embedded device.
    Now one of the key things that I'm sure you must have experienced
    as you were deploying the model is that sometimes the model works well,
    and sometimes the model does not work well.
    Now what's the reason?
    Well, there are many different reasons, but one
    of the key areas where you might be systematically having this problem
    is with the data set itself.
    Now we have come a long way from understanding
    how data is collected to actually deploying the model
    and making inferences, but what I want to do now
    is to kind of circle you back all the way to the beginning
    and think about what it means to collect data for a new machine learning task.
    Now we're going to do this specifically in the context of keyword
    spotting because since you have already deployed it,
    you have a good experience now about how the models are actually working.
    And I want to take you back to the source
    and help you understand how the data collection process actually influences
    the deployment and inference stage.
    Let's look at this big picture that I showed you
    in the previous course, which is all about the data set.
    In this big picture, you have an original data set
    that is chunked down into the training set and the test set.
    The training set is really split up further into the actual training
    set, which is actually used to fit the model,
    and the validation set, which is actually used to test and optimize
    the machine learning algorithm's performance
    or the machine learning algorithm's accuracy
    by tuning its hyper parameters.
    Once you have a final model that you're really happy with
    and you want to claim a certain level of accuracy,
    that's when you use the test set.
    Great.
    So over of course one and course two, we spent a lot of time
    talking about the training set, the validation set, and the test set,
    and using them repeatedly.
    But the key question is this, who actually
    generates the original data set?
    Where does that really come from?
    Now we kind of touched on a little bit about this in the previous course,
    but I want to do a recap on that portion because I
    want to help you understand the systemic issues that are involved with actually
    collecting a data set before you go on to actually building your own data set.
    Now recall that data sets require a significant amount of effort,
    both in terms of the labeling that's involved,
    as well as the manual effort that's actually involved in collecting
    all of that data together.
    For example, with Common Voice, we actually
    have to crowdsource data in order to be able to pull together data, speech
    data specifically, from across many different people spread
    across the planet.
    MS COCO, from Microsoft, has over 2.5 million
    labeled images specifically to deal with detection related use cases.
    ImageNet, which is used widely for image classification tasks,
    was created by Fei-Fei Li at Princeton with her students,
    and that has over millions of images at this point.
    Similarly, Waymo and KITTI, which are specifically
    tuned for actually deploying autonomous vehicles on the road,
    have had to have a massive amount of data set collected, all of which
    costs a lot of money, time, and effort.
    So the question is this.
    All these folks have gone through the experience
    of actually collecting the data set and know the nuances behind it.
    So the question here for us is, how do we build our own data set?
    What are the data engineering principles that are involved?
    In the previous course, I touched on the core data engineering principles,
    which fall largely into 4 buckets.
    Requirements, gathering, refinement, and sustainment.
    Requirements, briefly, is about identifying
    what you want to have in the data set.
    Gathering is about going out and actually figuring
    out and getting that data.
    Refinement is making sure that the data you've collected is, in fact, the data
    that you actually want.
    Sustainment is about improving the data to meet the various demands
    of the model over time.
    Let's talk about this very specifically in the context of speech
    commands, which Pete put together a couple of years ago.
    And let's talk about what requirement, gathering, refinement, and sustainment
    mean in the context of this particular TinyML data set.
    For requirements, Pete thought that it would
    be useful to have common use cases, like yes, no, and so forth.
    That's something that a data set should have.
    He also thought, well, robotics is on the rise,
    so it would be interesting to have words like left, right, go, stop,
    and so forth useful in the data set.
    He also thought numbers might be useful.
    0 through 9 would prove quite useful for general purpose applications.
    So those were the requirements that Pete went with.
    Then he went about the art of gathering all this data.
    Using crowdsourcing, he pulled together over 2,000
    some volunteers who consented to have their voices recorded and distributed.
    And they come from many different sources,
    so it's a rich data set, which has got many different accents in it.
    Luckily, that led to having over 1,000 samples for each of the keywords
    collected easily through a web based interface.
    Then came the refinement stage because not all of that data
    is actually useful.
    There are often words in there that were a little quiet.
    There were, in fact, wrong words.
    Someone was having a little bit of fun saying yes instead of no,
    saying no instead of yes, and so forth.
    So there had to be a process of actually cleaning up all of that data.
    Well, we can't sit around and actually manually clean the data.
    We need automated infrastructures and tools.
    For example, to remove low volume recordings because the network
    will not be able to pick them up.
    Or extract just a key portion of the word and not the entire sound
    that you might be collecting.
    For example, if there's a five second recording,
    well, within those 5 seconds, the keyword
    might actually just be a one second part in some snippet.
    So we have to cut out all the quiet noise.
    Then, of course, all of this needs to be validated in some sense.
    Luckily, with the help of crowdsourcing, we
    able to end up with over 100,000 utterances, which is great.
    All that said, you will have to do similar things when
    you start collecting your own data set.
    We'll talk about that more in the readings
    that come up that will supplement all the knowledge you need
    to have to collect a proper data set.
    Of course, even after you collect the data set,
    there's always the need to keep it evolving.
    You've got to meet new demographics, for instance.
    You've got to worry about biases.
    Who are the people who are actually going to use it?
    Do you have a nice balance between male and female voices?
    For example, in a previous course, I mentioned that in the Common Voice data
    set, for instance, they found that most of the recordings
    were male based recordings, and very few female voices were actually there
    in the data set, which can lead to a model that
    only tends to be more biased towards males, which is not a good thing.
    There's also diversity in age groups, the kind of accents,
    the dialects, the languages, and so forth.
    There are many different needs.
    So you have to understand where you're deploying the model,
    and then think about how you are going to sustain the data set.
    In other words, how are you going to keep evolving that data set?
    These are all great concepts to think and know about,
    but you've got to put this into practice.
    And that's what the goal is in your next assignment.
    The challenge that I'm setting out for you
    is, I want you to go and collect your own custom data
    set for building a keyword spotting model.
    You already know how to deploy a keyword spotting model on your embedded system.
    Next question is, can you collect the data
    in order to go through the training process, the evaluation
    process, and the deployment process to then understand
    how what you do right at the beginning ends up impacting
    how well the model performs at the end?
    This is a really unique chance, and I want
    you to pay attention and work through this
    because this is one of those rare opportunities where
    you have this experience.
    So, you've got to start off by understanding
    what the requirements are.
    And you've got to then go out, and you've got to gather the data.
    The question is, how much data do you need?
    I don't know what kind of use case you are thinking in your head,
    so you're the one who's going to be able to answer that question.
    Now what sort of false positive and false negative rates are acceptable?
    I can't answer that question for you.
    You're going to have to answer that question because I don't know which
    scenario you're imagining this particular custom data set based model
    is going to be working.
    Is it in a life critical sort of situation,
    when someone's reaching out for help and you're
    deploying this model in a medical emergency, for instance?
    Well, then you cannot have a really poor response rate.
    You're going to have to be careful.
    Or is it just a device that occasionally needs a sample of people around,
    or something like that, in which case maybe
    a little bit of false positive and false negative
    is something that's actually acceptable?
    I don't have the answers to these questions.
    You're going to have to tell me what you thought is your use case scenario,
    and how you went about identifying what you think are the right metrics.
    Now after you have gone through the process of actually collecting
    the data, you're going to definitely find
    that you've had all kinds of issues.
    You thought you collected the key word, but your microphone was not plugged in.
    You thought you collected the key word, but the microphone's volume
    was too low.
    You thought you collected the keyword, but, by the way,
    you read the wrong keyword instead of what you were planning on doing.
    There are many issues you're going to run into.
    And in addition to all of this, one of the key things
    you're going to have to do is you're going
    to have to supplement the keywords that you've
    collected with the right kind of background noise.
    Remember, background noise is very critical.
    We don't live in an isolated bubble in a noise proof environment.
    There's a lot of noise around us continuously that we are just
    used to picking up and ignoring.
    But we're going to have to teach our neural network
    to be able to do the exact same thing.
    So, you're going to have to augment the data set
    with a little bit of background noise.
    The readings tell you how to go about doing this.
    There's a shortcut that we are providing where
    you can reuse some of the speech commands background noise if needed.
    In terms of sustainment, when you evaluate the data set with the trained
    model, the question becomes, are others able to use the data
    set that you have put together?
    For example, if you are trying to get the model to recognize your word, well,
    does it work for everybody else?
    I don't know.
    Is it supposed to work for everybody else?
    You know the answer to that.
    Then you've got to think about where you're actually
    deploying your embedded system.
    Are you're deploying it in a quiet environment, where
    you have recorded it, so the deployment environment is representative?
    Or are you actually deploying this in a noisy environment, like in a kitchen,
    for instance, where the family tends to often gather?
    Or perhaps in the car, while you're actually playing your background music?
    Or is it next to the TV or the radio, where
    there's a lot of background noise?
    Or in a very crowded room, where it's very hard to pick things up?
    Well, you're going to have to think through all of these things.
    And I promise you that it's a very rich experience that you're going to have.
    It might be a little frustrating if the model doesn't work.
    Don't get caught up on the end result. I want
    you to appreciate the process because when you appreciate the process,
    and you go out, and you start building systems in the real world,
    then you will know the mechanisms and the issues you're going to come into.
    And that is the real learning.
    It's not whether the model simply works or not.
    So keep that in mind.
    And with that said, I'm going to turn you over
    back to Brian, who's going to help you walk through the process of collecting
    the custom keyword data set.
    Have fun with it.
    That's the most important thing.
    Ende des Transkriptes, zum Anfang springen
    Downloads und Transskripte
    Video
    Videodatei herunterladen
    Transskripte
    Download SubRip (.srt) file
    Download Text (.txt) file    
    
    Another important concern when performing data engineering is data provenance. We must be cognizant of the origin of our data, most notably when we are using information procured from an external source, such as images from the internet. These images may be subject to copyright laws and other restrictions related to privacy. To preclude the possibility of legal issues, the origin of data should also be assessed, consulting with domain experts where necessary.
    
    Certain types of data may require labeling, which may be done using crowdsourcing platforms such as Amazon Mechanical Turk in some cases. However, in other cases, domain experts may be required to label data. For example, radiologists may be necessary to assess x-rays or mammograms for abnormalities that would be vague or unnoticeable to the layman.
    
    A further concern is the shifting of our objectives. If a plan is not sufficiently detailed or implemented correctly, dataset creep may occur, wherein more features are added in an attempt to improve the utility of the dataset. Alternatively, certain feature variables may be omitted or inputted incorrectly, such as with incorrect labeling. 
    
    Clearly, there are many ways in which the engineering of a dataset can go awry, and it is our job as a data engineer to ensure that the dataset has been properly collected with consideration of the deployment environment, as well as considerations of labeling, dataset bias, copyright, and other aforementioned concerns.
    
    In the remainder of this section, you will learn to design, collect and develop your own dataset for the creation of a custom keyword spotting model. Good luck!

\section{Things to Consider for your Data Collection Plan}

    As you make a plan to collect data for a custom keyword spotting application, there are several important data engineering questions to consider in designing your data collection and testing scheme:
    
\begin{itemize}
  \item Who are the anticipated end-users?
    \begin{itemize}
      \item What will the age range be?
      \item What languages will the user be speaking? Will the users have accents?
    \end{itemize}    
  \item What are the goals for using the application? How will this impact the requirements/needs of the ML model's performance?
    \begin{itemize}
      \item For example, will the user be trying to turn specific lights on/off vs. turn the thermostat up/down vs. arm/disarm a security system vs play your favorite playlist from a smart speaker, etc.?
    \end{itemize}    
  \item  What false positive or false negative rates can the application handle / what would the consequences be of such an event?
  \item In what environments will the users employ the application?
    \begin{itemize}
      \item How much background noise do we expect?
      \item How far will the users be from the device, and from noise sources in the environment?
      \item Will these values vary over time?
    \end{itemize}    
  \item    In what situations will the users use the application? Will the user’s speech sentiments vary?
    \begin{itemize}
      \item Will the user be stressed vs. calm vs. panicked?
      \item Will the user use different volumes of voice (whispered/normal/loud/shouted)?
      \item Will different users have different sentiments?
    \end{itemize}    
  \item   How do all of these previously mentioned factors affect how much and what kinds of target, non-target, and imposter/adversarial data you will need to train with?
    \begin{itemize}
      \item How likely is it that these wake words will be triggered unintentionally during a conversation?
      \item Who/how many people will be talking in the application environment?
    \end{itemize}    
  \item   Given all of these previously mentioned factors, what will you collect?
    \begin{itemize}
      \item What custom wake words will you select?
      \item Are the background noise samples from the Speech Commands dataset sufficient for model robustness? Would it be helpful to collect additional background noise samples from specific environments?
      \item Do you need to collect other words to ensure the model learns the difference between them and your wake words?
      \item How much data will you need to collect? We all know more is better (usually), but you also live in the real world with time constraints so how much do you think will be enough?
      \item In what environments will you collect these samples?
    \end{itemize}    
  \item   How will you test your model?
    \begin{itemize}
      \item What words will you try? How many times will you try these words? Will any of them be adversarial? Remember even current production systems struggle with close words at times!
      \item Who will test the model? How diverse will the testers be?
      \item Will the testers vary their speech sentiments?
      \item Where will the testers be in relation to the device?
      \item In what environments will the testing happen?
    \end{itemize}    
  \item Finally, based on your initial testing, how will you improve your results:
    \begin{itemize}
      \item What didn't work the first time?
      \item Will you need more or different data?
      \item How can you iteratively improve your results?
    \end{itemize}    
\end{itemize}    
         
\section{Your Data Collection Plan}

\begin{itemize}
    \item What is the goal of your KWS application going to be?
    
    \item How are you going to go about collecting your data (e.g., who, what, where, when)? 
    
    \item How many recordings are you going to do initially?
    
    \item Do you have a testing and revision plan?
    
    \item Do you have any suggestions for other students about their plans?
\end{itemize}            

\section{Building a Custom Dataset}

\subsection{Collecting a Custom Dataset for KWS}

NOTE: Our URL has changed from the video! See below for details!

There is another change from the video below -- please only record 10 words per recording session as Chrome will only allow you to download 10 files at a time. We are working to debug the zip library which will allow us to get back to ~30 files per recording session.

If you would like to run the open-speech-recording application locally (and therefore not have a limit on simultaneous downloads) please see our appendix document. Do note that you must be comfortable running python and python-pip from the command line to take this approach.

If you record using your own tools make sure to name your files keyword\_\#\#\#.ogg where keyword is whatever word you are recording and \#\#\# is any integer number. This is needed for the training scripts to process your files correctly.

\subsection{Collecting a Custom Dataset for KWS}

0:00 / 0:00Klicken Sie auf die Pfeiltaste nach oben, um das Geschwindigkeitsmenu aufzurufen. Regeln Sie dann mit den Hoch und Runter Pfeiltasten die Geschwindigkeit. Klicken Sie Enter, um die jeweilige Geschwindigkeit festzulegen.Geschwindigkeit1.0xKlicke auf diesen Knopf und diess Video stumm oder nicht stumm zu schalten, oder drücke die HOCH oder RUNTER Knöpfe um die Lautstärke zu erhöhen oder abzusenken.Maximal Lautstärke.Videotranskript
Beginn des Transkriptes, zum Ende springen.
BRIAN PLANCHER: Hi, everyone.
In this reading, we're going to learn how
to use the Open Speech Recording app to record our own custom keyword data set.
As shown in the reading, we can get to the Open Speech Recording app
at the tinymlx.org website.
Make sure you use HTTPS, with an S, otherwise, this app won't work.
When we first get to the web page, we'll have
to give microphone access to our browser.
So click Allow, and now we can actually use our microphone.
And for those of you that are concerned about security and privacy,
do know that the open speech recording app does not send anything to a server.
All of the data remains local, in your browser,
the entire time you're using it.
So let's talk about how to actually use the app.
All you need to do is input the words you
want to record, where it says, these are example words,
as a comma separated list.
So I could, say, do Brian, Vijay, Pete.
And then, in the second entry, we can put how many of each word we want.
So if we left this as 10, it would be 10 Brians, 10 Vijays,
and 10 Petes we would record.
So let's just, for sample purposes, say three.
Now, this actually brings up a good point,
which is, how many words should you be recording?
And the answer is that, in general, more words are better.
More samples are better.
Pete's data set has thousands of samples in it.
But don't panic.
I actually just trained a Brian, Vijay, Pete model
that worked OK, as in it recognized my voice, saying
those words in this room, most of the time, using only about--
I think it was 30 to 50 samples for each of those three.
So if you have around, say, 40 samples of each word, for two to four keywords,
there's a good chance that you can train a model that will work for you.
If you want this to generalize beyond that--
maybe you want your family members to be able to use it,
you'll probably want to get your family members to also record some words.
Now, one other note of caution, since this is all running in your browser,
you'll only want to record about 30 words total per each recording,
or you might have issues with your browser's memory.
And the downloaded file you'll get may only
still include around 30 recordings, even if you record 40 or 50 recordings.
So we suggest just simply refreshing the page--
doing your recording, getting the download,
and then just simply refresh the page.
It'll start from scratch.
Put your keywords back in.
Put the amount back in, that you want to use, and record again.
And you should just do that a couple of times until you collect your full data
set.
That's the safest method, at least.
So to reuse the app, once you've typed all that stuff in, just click Record.
You'll see a countdown timer going.
And then, the words will start to appear right here,
in the middle of the screen, with a number counting you
through how many words you have done and to go.
And the recordings will start to appear down on the bottom.
Let's stop this for a minute because I've been just recording
what I've been saying.
But let me just show you.
So if I play the first one--
[AUDIO PLAYBACK]
BRIAN PLANCHER (ON RECORDING): And then the words will start--
[END PLAYBACK]
BRIAN PLANCHER: So that's just recording everything
that I am saying because I wasn't doing it thoughtfully.
So let's delete those.
So we can go to this red delete button here.
We can go through and we can delete all these recordings
because there was nothing there.
And if we now hit record again, it should start back at the beginning,
asking again for the first word because it knows that there's still nine to go.
Now, just to time the recordings, you'll see this record button goes red
and then it goes gray again.
And every time it goes red, that's when it's recording.
So if I wanted to actually get one to work well, I would just start this up,
and then I would wait for that moment.
Pete.
Vijay.
Pete.
And so again, I'm waiting, watching it turn gray.
And then, when it turns back red, I'm saying the word that it shows there.
And then, you'll get this pop-up that'll ask you if you're ready to download.
But let's say, I'm not ready yet.
So if you click OK, it will download it.
But if I click Cancel, it will let you keep playing with these here.
And so I could say--
[AUDIO PLAYBACK]
BRIAN PLANCHER (ON RECORDING): Ask me again.
[END PLAYBACK]
BRIAN PLANCHER: Right, so these first couple were wrong.
So I don't want those.
And maybe down at the bottom, I think, I probably,
hopefully got one of these right.
[AUDIO PLAYBACK]
BRIAN PLANCHER (ON RECORDING): Vijay.
[END PLAYBACK]
BRIAN PLANCHER: Yeah, so that one worked right.
And let's say I want to download the set that is here now.
I can just click the blue Download button.
And you'll get a nice zip file with all your words from that recording in it.
And we suggest when you're done doing all of your recordings,
to unzip all of those zip files, and put all of the .ogg files that are in there
in one folder.
And that'll make it a lot easier for you to upload them en masse in the Colab,
for use with training and deployment in the next reading.
So good luck, and have fun recording your own custom data set.
Ende des Transkriptes, zum Anfang springen
Downloads und Transskripte
Video
Videodatei herunterladen
Transskripte
Download SubRip (.srt) file
Download Text (.txt) file


We are going to use a modified version of Pete’s open-speech-recording project that will keep all of your recording data local to your browser (and thus it will never leave your computer)! You can open up the application in your browser by navigating to: \url{https://tinyml.seas.harvard.edu/open_speech_recording}
 NOTE: Our URL has changed from the video!


Please make sure to open this application on a device that you can easily download and upload files from (e.g., a laptop as compared to a phone).

If you see a warning that says: Your device does not support the HTML5 API needed to record audio (this is a known problem on iOS).

\begin{enumerate}
  \item If you are not using Chrome, we suggest you switch to Chrome. We have tested all of our applications on Chrome and cannot guarantee compatibility with other browsers.

  \item Make sure you are navigating to https://. For some reason, if you do not use the secure protocol, then the HTML 5 API will not work.

  \item Make sure you approve access to your microphone! The first time you load the webpage (or if you load it again in an incognito window), you should get a popup on the top left of your screen as shown below. Make sure to click the “allow” button!

  \includegraphics[width=\textwidth]{TinyML/3 Development/4-6.1.png}
  
  A screenshot of the Chrome browser that has navigated to \url{https://tinymlx.org/open_speech_recording/} which has generated a warning popup which says “tinymlx.org wants to use your microphone” and below it has a button to “block” or to “allow.” This is where you have to click “allow” for recording to work. Please defer to the text for changes that have been made to the url and suggestions of how to use the webpage as it has evolved over time.

\end{enumerate}


\subsection{Using the Open Speech Recording App}

\begin{enumerate}
  \item In the middle of the screen, you will see the instructions for the App. Make sure you follow the steps in order. First, in step 1, type in as a comma separated list, the words you would like to record. And in step 2, type in as a single integer, the number of times you would like to record those words.

  Note that we have found that ~40 recordings each of 2-4 words often will be enough data to successfully train a moderately decent KWS application (i.e., it will work for your voice in whatever exact room you are in but may not generalize beyond that and may only work well for some of the keywords).

  Also, note that since all recordings will be saved in your browser memory and you can record multiple times by refreshing the page to reset the application, we do not suggest recording too many words in any one round of recording. We have found that Chrome will only let you download 10 recordings per session. So do not try to record more than 10 words before refreshing the page. Note that this is a change from the video as at the time of recording we were using a zip library that allowed ~30 words to be downloaded at a time. Unfortunately, that library, JSZip, now does not seem to be working. If you happen to be a web development guru and have ideas for a stable way to zip and download more files please do let us know!

  As a final note, remember that you can record both keywords and other distractor words (e.g., if your keyword was “slow” you could also record “snow”) to help. refine the model.

  \includegraphics[width=\textwidth]{TinyML/3 Development/4-6.2-re.png}
  
  Another screenshot of the Open Speech Recording web page, this time showing the inputs to the UI. Namely step 1 says “Input the words you would like to record as a comma seperated list:” and in the input box below it has the example “These,are,example,words”. Step 2 says “Input the number of samples of each word you would like to record:” and the input below it has the example value 10. Please defer to the text for changes that have been made to the url and suggestions of how to use the webpage as it has evolved over time.

  \item Then, click the blue record button. You will see a 3,2,1 timer countdown occur and then recording will start.
  
   You will see a counter of X/Y appear that will show you your progress in recording all of the words. You will also see the current word you are supposed to record appear in the white box. The words will come shuffled in a random order, so sometimes they will alternate and sometimes, you will get the same word multiple times in a row. You will also see the waveform of your current speech appear in the grey box below that.

   It may take you a couple of tries to get the exact timing working. As a tip, the “Record” button will alternate between grey (as shown below) and red. Record the word shown in the white box each time the word record turns red.

   After each word is recorded, it will appear as a new row in the bottom of the screen.

  You can pause the recording session by clicking “Stop” at any time. The “Stop” button will then turn grey and the “Record” button will turn blue again. Click “Record” to restart the recording session.

  \includegraphics[width=\textwidth]{TinyML/3 Development/4-6.3-re.png}
  
  Another screenshot of the Open Speech Recording web page, this time showing “record”, “stop”, and “download” buttons that appear in the middle of the page. Above those buttons is the wave form viewer that displays the sounds being recorded and above that is a counter showing (in this case) 1/40 and the word “words” which implies that this session is recording 40 words and the first word to be recorded is the word “words”. Below the buttons the recorded samples appear. Each has an audio play window that allows you to click play to hear the recording and a volume slider as well as a delete button to delete bad recordings as well as the label of what that recording's word is supposed to be. Please defer to the text for changes that have been made to the url and suggestions of how to use the webpage as it has evolved over time.

  \item Once you complete recording the final word in this session, a pop-up will appear asking if you’d like to download your recordings now. If you are satisfied, click “OK” otherwise, click “Cancel” and you can inspect the recordings one by one and download them later.

  \includegraphics[width=\textwidth]{TinyML/3 Development/4-6.4.png}

  Another screenshot of the Open Speech Recording popup that appears once recording is done. It says “ tinymlx.org says Are you ready to download your data? If not, press cancel now, and then press Download once you are ready.” It then has a cancel and an OK button. Please defer to the text for changes that have been made to the url and suggestions of how to use the webpage as it has evolved over time.

  Note that if you use the application multiple times, you may get a pop-up asking for your permission to download multiple files. Click “Allow” and it should then automatically download all of your recordings again in the future without asking.

  \includegraphics[width=\textwidth]{TinyML/3 Development/4-6.5.png}

  Another screenshot of the Chrome browser that has navigated to https://tinymlx.org/open\_speech\_recording/ which has generated a new warning popup which says “tinymlx.org wants to download multiple files” and below it has a button to “block” or to “allow.” This is where you have to click “allow” for downlaods to work. Please defer to the text for changes that have been made to the url and suggestions of how to use the webpage as it has evolved over time.

  \item To inspect a recording, simply press the play button on the left of its row. You can delete any recordings you don’t like by clicking the red “delete” button on the bottom right of its row. You can also adjust the volume of the recording as you listen to it by using the speaker icon on the right of each row.

  \includegraphics[width=\textwidth]{TinyML/3 Development/4-6.6.png}

  Another screenshot of the Open Speech Recording web page, this time showing a zoom in on the previously mentioned audio play window that allows you to click play to hear the recording and a volume slider as well as a delete button to delete bad recordings as well as the label of what that recording's word is supposed to be. Please defer to the text for changes that have been made to the url and suggestions of how to use the webpage as it has evolved over time.

  \item Once you have removed any recordings, you do not like simply click the “Download” button in the middle of the screen that will have turned blue. Note the warning mentioned above about Chrome needing approval to download multiple items.

  \item And that’s it! You’ve successfully collected a KWS dataset. Feel free to record multiple times if you’d like to increase the size of your dataset and remember where your files are being downloaded. If you happen to be recording at a time where your download is a single zip file (instead of many .ogg audio files), you’ll want to unzip all of your recordings to a folder from which you can upload them into the next Colab for training! On most operating systems, right-clicking a zip file will reveal an option that says either “extract” or “open archive manager” or something along those lines. Clicking that will either immediately unzip the files or open up a window which will have an option to unzip the files.
\end{enumerate}


\section{Train and Deploy Your Custom Dataset KWS Model}

In this reading, we are going to walk through the steps for training a custom KWS model based on the dataset you just collected and then deploy that model onto your Arduino (including the necessary code changes in the Arduino IDE).
    
\subsection{Train and Deploy Your Custom Dataset KWS Model}

    0:00 / 0:00Klicken Sie auf die Pfeiltaste nach oben, um das Geschwindigkeitsmenu aufzurufen. Regeln Sie dann mit den Hoch und Runter Pfeiltasten die Geschwindigkeit. Klicken Sie Enter, um die jeweilige Geschwindigkeit festzulegen.Geschwindigkeit1.0xKlicke auf diesen Knopf und diess Video stumm oder nicht stumm zu schalten, oder drücke die HOCH oder RUNTER Knöpfe um die Lautstärke zu erhöhen oder abzusenken.Maximal Lautstärke.Videotranskript
    Beginn des Transkriptes, zum Ende springen.
    SPEAKER 1: All right, now that you've finished collecting your data set,
    it's time to train a model and then deploy it onto your Arduino.
    The Colab below will help you to take all of your existing recordings
    that you just did, and make sure that they're formatted correctly,
    so that the training script that we used in course two will work with them.
    You'll see, essentially, a copy of the assignment
    from keyword spotting when you retrained a model in course two.
    Paste it at the bottom of that Colab.
    We've kept the data set formatting and the training in the same Colab,
    so you don't have to upload the data set more than once.
    This will be especially important if you have a slow internet connection.
    Once you finish training your model, you'll see this cell at the bottom that
    runs that XXD step and will generate out your .cc model file.
    You can then take that model file and follow
    the steps we used in the last section, when we deployed
    your favorite keyword, to update the various files in the Arduino IDE
    to deploy that model.
    And that's it!
    Well, the training will take about two hours.
    But once the training is done, within a couple of minutes,
    you should have your own custom keyword spotting
    model deployed onto your Arduino.
    I hope you have a lot of fun with that.
    Ende des Transkriptes, zum Anfang springen
    Downloads und Transskripte
    Video
    Videodatei herunterladen
    Transskripte
    Download SubRip (.srt) file
    Download Text (.txt) file

\subsection{Training with your Custom Dataset}

 The first thing you’ll need to do is to upload your custom dataset into Colab and train a KWS model. Then, in that Colab, we’ll need to convert that model first into a quantized .tflite file and then into a .cc file for use with the Arduino IDE.
    
 As with the pretrained model and favorite keyword model, we will be using the resulting .cc file, so make sure to download it or leave the tab open with the printout!
    
    \url{https://colab.research.google.com/github/tinyMLx/colabs/blob/master/4-6-8-CustomDatasetKWSModel.ipynb}
    
\subsection{Updating the Arduino Code}

\begin{enumerate}
  \item Again, open the micro\_speech.ino sketch, which you can find via the File drop-down menu. Navigate, as follows: File → Examples → Arduino\_TensorFlowLite →   micro\_speech.
  
  \item As before, navigate to the micro\_features\_model.cpp file and update the model.
    \begin{itemize}
      \item Copy the binary model file contents from the KWS\_custom.cc file into the micro\_features\_model.cpp file. As always, make sure to only copy the binary data inside the { } as the variable type is different in the downloaded or printed KWS\_custom.cc file.
      \item Next, scroll all the way down to the bottom of the file and replace the model length. Again, note that the .cpp file needs the variable to be of type const int while the .cc file will show unsigned int. Our suggestion again is to simply copy the numerical value.
    \end{itemize}
  \item Next, save your changes. Again, we suggest that you make a folder called e.g., TinyML inside of your Arduino folder. You can find your main Arduino folder either inside of your Documents folder or in your Home folder, and save it in that folder with a descriptive name like micro\_speech\_edx\_custom\_keywords. That said, you can save it wherever you like with whatever name you want!
 
  \item  The two other things we need to change are the count and list of keywords in the model settings files and the output response file. As we do that in the rest of   this document, we are going to assume that you used our model with the keywords “brian,” “vijay,” “pete.” If you choose to use different keywords, make sure   that you update the following steps accordingly.

  \item  Navigate to the micro\_features\_micro\_model\_settings.h file and scroll down to line 40 and change the value of the kCategoryCount variable to be two more than the number of keywords you selected (one extra for silence and one extra for unknown). For example, if you had 3 keywords (as we do in our example model), we need to update the variable to be 3 + 2 = 5 as shown below:
 
   \PYTHON{constexpr int kCategoryCount = 5;}
    
  \item  Then, navigate to the micro\_features\_micro\_model\_settings.cpp file. You’ll find that it includes the .h file and otherwise has a single array. Update the values of   that array to match your keywords. Make sure to leave “silence” and “unknown” in the list as well, and make sure to list your key words in the order   that you listed them on the training script. So, in the case of our example model, we would update the array to be:
  
  \begin{lstlisting}
   const char* kCategoryLabels[kCategoryCount] = {
        
        "silence",
        
        "unknown",
        
        "brian",
        
        "vijay",
        
        "pete",
        
    };
  \end{lstlisting} 
    
  \item  Finally, navigate to the arduino\_command\_responder.cpp file. Scroll down to line 54. There, you will see a series of if statements which control how the Arduino responds to each command. They are each structured to look at the value of the found\_command variable. This variable will be the string of the keyword (or “silence” or “unknown”) you entered into the kCategoryLabels array in the previous step. So, by comparing the [0] value in that string, the if statements are simply looking at the first letter! In each if statement you’ll also see a digitalWrite to either the LEDG, LEDR, or LEDB variable, which is used to turn on the appropriate color LED. So, if we wanted to adapt this to our three keywords and, for example, only turn on the single colors for our keywords and turn all of the colors on for unknown, which will come out look white, we could update those three if statements to the following:

  \begin{lstlisting}
    // Red for Vijay
    
    if(found_command[0] == 'v') {
        
        last_command_time = current_time;
        
        digitalWrite(LEDR, LOW);
        
    }
  \end{lstlisting} 
    
  \begin{lstlisting}
    // Green for Pete
    
    if(found_command[0] == 'p') {
        
        last_command_time = current_time;
        
        digitalWrite(LEDG, LOW);
        
    }
  \end{lstlisting} 
    
  \begin{lstlisting}
    // Blue for Brian -- note here that you do not need to index
    
    //                   into the first letter only, just a unique
    
    //                   letter combination in the keyword! That
    
    //                   said make sure you do not index beyond the
    
    //                   end of ANY keyword or you will get an error!
    
    if(found_command[0] == 'b' && found_command[1] == 'r') {
        
        last_command_time = current_time;
        
        digitalWrite(LEDB, LOW);
        
    }
  \end{lstlisting} 
    
  \begin{lstlisting}
    // All three for unknown (white)
    
    if(found_command[0] == 'u') {
        
        last_command_time = current_time;
        
        digitalWrite(LEDR, LOW);
        
        digitalWrite(LEDG, LOW);
        
        digitalWrite(LEDB, LOW);
        
    }
  \end{lstlisting} 

\end{enumerate}    
    
    Remember, if you used different keywords make sure to update the if statements accordingly. And that’s it! You’re Arduino code should be all set to handle your new model!
    
\subsection{Deploying the New Model}

\begin{enumerate}
  \item   Use a USB cable to connect the Arduino Nano 33 BLE Sense to your machine. You should see the green LED power indicator come on when the board first receives power.
  
  \item   As always, use the Tools drop-down menu to select appropriate Port and Board.
    \begin{itemize}
      \item Select the Arduino Nano 33 BLE as the board by going to Tools → Board: <Current Board Name> → Arduino Mbed OS Boards (nRF52840) → Arduino Nano 33 BLE.

      \item Then, select the USB Port associated with your board. This will appear differently on Windows, macOS, Linux but will likely indicate ‘Arduino Nano 33 BLE” in parenthesis. You can select this by going to Tools → Port: <Current Port (Board on Port)> → <TBD Based on OS> (Arduino Nano 33 BLE). Where <TBD Based on OS> is most likely to come from the list below where <\#> indicates some integer number.
        \begin{itemize}
          \item Windows → COM<\#>
          \item macOS → /dev/cu.usbmodem<\#>
          \item Linux → ttyUSB<\#> or ttyACM<\#>
        \end{itemize}
      \end{itemize}

  \item   Use the rightward arrow next to the ‘upload’ / flash the code. You’ll know the upload is complete when you see red text in the console at the bottom of the IDE that shows 100\% upload of the code and a statement that says something like “Done in <\#.\#> seconds.”

  \item If you receive an error, you will see an orange error bar appear and a red error message in the console (as shown below). Don’t worry -- there are many common reasons this may have occurred. To help you debug, please check out our FAQ appendix with answers to the most common errors!

  \item You should now have a working Keyword Spotting model trained on your own custom dataset deployed to your Arduino! To test it out, open the Serial Monitor. You should start to see the result of the model begin to display in the Serial Monitor and the LED’s light up according to how you instructed them in the previous section!
\end{enumerate}    
 
\section{How Well Does your Custom Dataset Model Work?}

How did your Custom Dataset Model work? Does it work for you in different locations? Does it work for any of your family and friends? How much data did you collect? Did you collect data from anyone other than yourself? If you made multiple models (or re-trained with additional data) did that improve performance? How well did your data collection plan actually work out?    


% 1.7: Visual Wake Words

\section{Recap: What are Visual Wake Words?}
    In this module, you will learn to build and implement an algorithm for detecting visual wake words on your local microcontroller.
    
    In the previous section focused on keyword spotting, we empowered our microcontroller with a lightweight auditory modality, allowing the system to monitor the microphone input in real-time to be able to recognize and respond to several keywords of interest to users. Visual wake words (VWW) are the natural extension of this technique to visual data, providing a lightweight visual modality to our microcontroller system.
    
    Image information is inherently more complex than auditory information to analyze using machine learning due to the higher dimensionality and complex spatial correlations. Consequently, image data requires greater computational and memory resources, making it challenging to deploy on resource-constrained microcontrollers. An additional complication, similar to keyword spotting, is the difficulty in procuring sufficient data to train our algorithm. In some scenarios, we may have to generate the dataset ourselves.
    
    In the previous classes, we have looked at the VWW task in some detail, from the implementation in Colab to discussions of image privacy and copyright. We analyzed a specific set of architectures, known as the MobileNet architectures, which leverage depthwise separable convolutions to minimize required memory and computational resources. We also looked at how transfer learning can be used to build upon pre-trained models that were used for similar tasks in the same domain to reduce our need to procure large datasets and train large deep learning models.
    
    The main VWW task we focused on was person detection, which focuses on determining the presence or absence of an individual in an image. In this task, we used the VWW dataset to train our person detection model and then performed transfer learning using a relatively small amount of images to adapt our model to perform mask detection. However, we have yet to demonstrate how to port and deploy such a model to our end-device, which will be the main focus of this section.
    
    By the end of this section, you should have a fully functioning VWW detection system. This shares some similarities to keyword spotting, but the deployment procedure will require some pre-and post-processing (e.g., image downscaling, encoding/decoding, and interfacing with the onboard camera) unique to the VWW detection task.
    
    Not only this, but we will also demonstrate the deployment of keyword spotting and VWW simultaneously using model “multitenancy.” More often than not, the typical focus is on running individual models in isolation, independent of other contextual activities. But increasingly there is a growing need to combine information from multiple sensors together to perform an intelligent action. To that end, we will learn how to perform audio-visual recognition, using both the microphone and camera. Our goal is to simultaneously operate two models at the same time in order to mimic sensor fusion (i.e., the ability to tie together inputs from different sensors in close temporal proximity).
    
    We will learn how to do this and more in TensorFlow Lite Micro.
    
\section{Person Detection Application Architecture}

    0:00 / 0:00Klicken Sie auf die Pfeiltaste nach oben, um das Geschwindigkeitsmenu aufzurufen. Regeln Sie dann mit den Hoch und Runter Pfeiltasten die Geschwindigkeit. Klicken Sie Enter, um die jeweilige Geschwindigkeit festzulegen.Geschwindigkeit1.0xKlicke auf diesen Knopf und diess Video stumm oder nicht stumm zu schalten, oder drücke die HOCH oder RUNTER Knöpfe um die Lautstärke zu erhöhen oder abzusenken.Maximal Lautstärke.Videotranskript
    Beginn des Transkriptes, zum Ende springen.
    VIJAY JANAPA REDDI: In this module, I'm very
    excited to be introducing you to a whole new application that
    involves the camera module.
    Specifically, we will be talking about person detection
    or sometimes referred to as the visual wakewards application.
    As I've done with the keyword spotting application,
    I will be talking about the overall application architecture.
    Let's first start off by talking about what the application is all about.
    Person detection is really simple.
    We want to know if there is a person in the camera's viewpoint or not.
    For example you are seeing Dylan Newton here and on the right hand side,
    you're not seeing him.
    So if there's a person, then we want to be notified that yes there is a person.
    If there is no person, I simply want to be notified that there is no person.
    This simple task has many diverse applications in the real world.
    For example, if you are deploying this sort of person detection in office
    buildings, you can use this to automatically turn
    lights on lights or off depending on the occupancy.
    That could be quite useful for conserving energy,
    especially in office buildings that are often large
    and consume a significant amount of electricity.
    Let's get started in understanding what's
    actually involved in putting in such an application together.
    There are two components as usual; there's initialization phase
    and there's a steady state main loop.
    The steady state main loop, has got multiple components
    involved the pre processing, the inference and the post-processing.
    Hopefully, by this point, you're quite comfortable with what those terms mean.
    And I'll step through them in detail soon.
    Let's get started with the initialization face.
    The initialization phase has got to subcomponents,
    you've got to initialize the camera and then you've
    got to set up things for your model.
    By and large, this breaks down into these smaller steps that
    are covered in a lot of detail, in the key word spotting application example,
    declaring variables all the way down to setting up
    things with the main loop, which is going to execute repeatedly.
    So I'm going to let things be at this point.
    Later when I do the screen cast, I'll talk about these in a bit more detail.
    Now, when you do the initialization, you've
    got initialize the camera, which is as simple as making this one
    single function call, camara.begin.
    As you can see here, this function call takes a couple of arguments.
    The first argument, QCIF, is what's called a Quarter Common Intermediate
    Format.
    That's a standard that means something in terms of the resolution.
    Second argument is GRAYSCALE.
    In order to do person detection, you do not need a color image.
    So we simply need a grayscale image that consumes less number of pixels
    and takes up less space.
    Third argument, is a number of frames per second.
    How often do we want the data coming in from the camera.
    And the final argument, is just a specific camera module itself.
    Which in our case is the OV7675.
    Then we move on to allocating the model in the initialization phase, which
    involves declaring the variable, setting up your interpreter,
    tying up the interpreter's input to where the buffer that
    is associated with the camera module.
    So the camera dumps the data into a buffer,
    that buffer has got to then get passed into the hands of like micro,
    because that's where the neural network is
    going to be seeding its inputs strong.
    And then you set things up for the main Loop.
    Pretty straightforward.
    Let's get into the pre processing part.
    The pre processing stage involves two components, which is,
    reading the image, cropping and converting.
    Reading the image involves simply invoking the function call that
    gets a frame of data from the OV7675.
    The data that you get back is in the format that I mentioned earlier, QCIF.
    QCIF has a precise resolution that's associated
    with it, which is 176 in number of columns
    and 144 in terms of the number of rows.
    This data with this function call, will end up on [INAUDIBLE] system.
    The next part, before we can actually go into the inference stage,
    is really cropping and converting the model from QCIF,
    into a format that's rendered for our actual neural network.
    In this case, it is about scaling that original QCIF image, into a 96
    by 96 square pixel.
    This can easily be done through this little function
    that we have already provided for you, which I will talk about more in detail,
    let me do the screen test.
    But it's very simple, is the key takeaway here.
    You're cropping the image down in size and then
    you're converting it into a format that's ready for the neural network
    to consume.
    Given that, once we have the input data that's ready,
    we simply invoke the model.
    In our case, the model that we're invoking here,
    is a mobile nets model, that actually runs
    on top of the image to detect whether a person is there or not.
    And in order to invoke the model, this is all we do.
    At this point you've got to be quite comfortable with actually
    looking at this APIs knowing exactly what's happening.
    The next step, is figuring out what needs to happen
    during the post processing phase.
    Remember the post processing phase is all about understanding
    what to do with the data that has come out of the neural network.
    So in our application, it's all about responding to
    whether a person was detected or not, which
    we are going to do simply through the LEDs that
    are present on our embedded system.
    So for example, if there is a person in the view,
    then we simply want to light that green LED,
    if there is no person in the viewpoint, then we simply
    want to light the red LED.
    And that's all there is.
    So hopefully one of the key things that you're seeing now,
    is that compared to keyword spotting.
    The visual wakewards to the person detection example,
    is remarkably simple.
    So not every application has the same level of complexity
    in terms of the initialization, in terms of the preprocessing,
    in terms of inference and in terms of the post-processing.
    So just keep that in mind.
    One of the thing you should probably think about,
    is the fact that the stages tend to stay the same.
    So there's a certain pattern that repeats
    in terms of how you want to process.
    And if you want to look for something that you want to change,
    keep that in mind.
    Just as I said to do the same in the context
    of the keyword spotting application.
    With that said, let's get a little more creative
    in terms of how we can actually deploy an interesting application using
    the camera module, which incorporates person detection,
    but in addition it'll actually incorporate
    getting another sensor, which makes it much more exciting, more on that soon.
    Ende des Transkriptes, zum Anfang springen
    Downloads und Transskripte
    Video
    Videodatei herunterladen
    Transskripte
    Download SubRip (.srt) file
    Download Text (.txt) file    
    
\section{Person Detection Screencast}

    0:00 / 0:00Klicken Sie auf die Pfeiltaste nach oben, um das Geschwindigkeitsmenu aufzurufen. Regeln Sie dann mit den Hoch und Runter Pfeiltasten die Geschwindigkeit. Klicken Sie Enter, um die jeweilige Geschwindigkeit festzulegen.Geschwindigkeit1.0xKlicke auf diesen Knopf und diess Video stumm oder nicht stumm zu schalten, oder drücke die HOCH oder RUNTER Knöpfe um die Lautstärke zu erhöhen oder abzusenken.Maximal Lautstärke.Videotranskript
    Beginn des Transkriptes, zum Ende springen.
    VIJAY JANAPA REDDI: In this screencast, I'm
    going to walk you through the person detection example.
    Specifically, I'm going to walk you through the entire application
    architecture, starting with the initialization phase
    and then stepping through the various components of the main loop.
    Now, one of the things that I'm going to do differently
    in this particular screencast is that I'm not going to break down
    the screencast into multiple screen casts
    to explain each one of the stages.
    That is the initialization, preprocessing, inference,
    and postprocessing, and different levels.
    The reason is because it's a fairly simple example.
    Moreover, I want to help you see the connections in a single screencast
    so you can feel confident that you understand
    the general template that any TinyML application is going to follow.
    So let's start off with the initialization step,
    which you may recall has got a whole bunch of small steps involved.
    So let's get started with that.
    The first step is to declare your variables.
    As you can see here from line 30, we're declaring the variables.
    The first chunk of variable declarations should
    look quite familiar to what we did in keyword spotting.
    We declare a handle to the error reporter,
    a handle to the model that we'll be loading,
    a handle to the microinterpreter, and a handle to the input
    tensor that will be used as the source for doing the inference.
    And the second portion is really simply declaring the tensor arena
    and allocating some memory for it.
    Once we're done with this, we get into the next step of loading the model.
    Now that is actually done inside the setup routine, which is right here.
    So as you can see here, we're going in and getting the model.
    G\_person\_detect\_model\_data is the file that actually
    contains the actual neural network.
    That is actually held in here, which if I scroll to the top, you can see
    is g\_person\_detect\_model\_data, which is an array which has
    got the neural network architecture.
    Going back here, we do a quick error check
    to make sure that the schema of the model
    actually matches the schema of what we're
    expecting in this particular version of TensorFlow Lite Micro.
    Following that, we have to resolve the operators
    and that's what we're doing here.
    We've got five operators that we will be using.
    Each one of those includes the AveragePool2D, the Conv2D,
    the DepthWiseConv2D, the m and the Softmax.
    Alternatively, we could have used AllOpsResolver, which we are not
    because we want to conserve the memory.
    After initializing the OpsResolver, our next step
    is to actually initialize the interpreter using all the variables we
    have set up, such as what is the model file, what's the ops that we'll
    be using, what's the tensor arena, what's its size,
    and how do we report errors.
    Finally, we will then allocate the raw memory that's
    required for the tensor arena, and then we
    will set up the input buffer for the neural network.
    And that's all there is for the initialization phase.
    This should look very similar to what you did with keyword spotting.
    The key takeaway here that this, there is a certain template that repeats over
    and over regardless of what the application is.
    Keep that in mind.
    It will help you make forward progress much faster.
    Now we're going to get into the actual main loop of the application.
    Remember, the main loop has preprocessing, inference,
    and postprocessing.
    So that is broken down here into three simple parts.
    The first part is the preprocessing, where
    the job is simply to get an image.
    That means you've got to read an image, you've got a crop and convert it.
    Both those are tucked away nicely inside GetImage.
    After you've finished the preprocessing, we
    do the inference, which is simply through this function call.
    After that, we do the postprocessing, where
    we process the data that's coming out of the neural network.
    Let's go through each of these stages in detail.
    Let's go and see what's happening inside GetImage.
    The job of GetImage is quite simple.
    We want to get an image from the camera module.
    If this is the very first time we're calling the camera module,
    then we're going to initialize it.
    As you recall from the slides I showed you,
    the camera module uses a QCIF format, requesting pixels in grayscale
    at five frames per second for this camera module, which
    we have interfaced with the TinyML [? shield. ?]
    This is only done the very first time.
    Subsequently, we read the data, and after that, we
    crop and scale the image down into a resolution that
    makes sense for our particular neural network, which is 96 by 96 pixels.
    After we are done with the preprocessing,
    we then get into the inference stage, and the inference stage is very simple.
    We simply call the interpreter, which gets its input data from the input
    buffer.
    After the neural network is run, then we get into the postprocessing stage
    where we get the output from the neural network, which is a tensor,
    and then we look at the Softmax probabilities
    for its outputs, which are in the form of a person's score or no person score.
    We pass that off to the responder.
    The responder's job is to light an LED depending on
    whether it detects a person or not.
    So let's go ahead and look at what's actually
    happening inside RespondToDetection.
    This function, which is inside the Arduino detection responder,
    performs a couple of small bookkeeping tasks
    here, which I will let you read at this point
    because you should be familiar with the code.
    The key nugget of this particular function is really this part.
    If it thinks that the probability--
    "it" being the neural network-- if the neural network guessed
    that the probability that the person actually is in the picture
    is greater than the probability that there's no person in the picture,
    then we will ask the green LED to turn on by writing digitalWrite
    as a LOW on the LEDG.
    That turns on the green LED.
    Or, vice versa.
    If there is no person, then it will turn on the red LED.
    And that is pretty much all there is to the postprocessing stage.
    Again, once this is done, we go back to the start of the loop.
    We go through the preprocessing stage, we run the inference,
    then we run the preprocessing.
    What I'm trying to get you to realize here
    is the same general structure of initialization followed
    by main loop coupled with preprocessing, inference,
    and postprocessing inside the main loop is the overall pipeline
    that we'll be rinsing and repeating regardless
    of what application you're doing.
    And that is the key takeaway that I want you to have
    from this particular screencast.
    Ende des Transkriptes, zum Anfang springen
    Downloads und Transskripte
    Video
    Videodatei herunterladen
    Transskripte
    Download SubRip (.srt) file
    Download Text (.txt) file    
    
\section{Person Detection with Keyword Spotting: MultiModal}

    0:00 / 0:00Klicken Sie auf die Pfeiltaste nach oben, um das Geschwindigkeitsmenu aufzurufen. Regeln Sie dann mit den Hoch und Runter Pfeiltasten die Geschwindigkeit. Klicken Sie Enter, um die jeweilige Geschwindigkeit festzulegen.Geschwindigkeit1.0xKlicke auf diesen Knopf und diess Video stumm oder nicht stumm zu schalten, oder drücke die HOCH oder RUNTER Knöpfe um die Lautstärke zu erhöhen oder abzusenken.Maximal Lautstärke.Videotranskript
    Beginn des Transkriptes, zum Ende springen.
    VIJAY JANAPA REDDI: As we get ready to deploy the person detection example,
    I want us to kind of reflect back on the way we have been
    thinking about the usage of sensors.
    Specifically, when I talked about keyword spotting,
    I only mentioned the microphone.
    When I talked about the person detection example,
    I only mentioned the camera module.
    But is there something interesting that we
    can do if we combined multiple sensors together to build richer applications?
    For example, can our personal detection example
    be richer if we use the microphone data as well?
    The answer is absolutely yes.
    For instance, if we combined the visual input with the audio input,
    we will be able to predict moods, which is quite powerful.
    Now, how would we go about this?
    This is where the concept of multimodal comes in.
    Multimodal means combining input that's coming in from multiple sensors
    into a single classification, and there are loads of benefits around this.
    Specifically, they fall into two categories, which
    is called complementary or redundancy.
    Redundancy means there are multiple sources of the same information,
    such as lip movement and parts of speech,
    that can be used in conjunction with one another
    to increase the reliability and the prediction accuracy of the model.
    Complementary means that we can detect things through data fusion
    that we would have otherwise not been able to detect
    if we were only using a single sensor.
    So this sounds like a really awesome thing to be able to do.
    We all want multimodal because if we have multimodal, then
    we might be able to build really interesting applications.
    How do we go about doing that?
    Well, let's go back to the ML workflow.
    It all starts with the data set.
    If you want to be able to do multimodal, that
    means we need to have multiple streams of input data.
    So here, we need to have a visual stream for the facial expressions
    and then we need to have an audio stream.
    Now, these are data sets, basically, but these data sets need to be aligned.
    For example, in this first figure, the person has a neutral expression
    and there's a set of words that are associated
    with that neutral expression.
    This person here is happy and there's something
    that this person is saying at that particular point that
    is aligned with the happy feeling.
    Here, the person looks sad, so they're probably saying something
    that's associated with a sad word.
    And here, they seem excited, which means that there's probably something
    that they're saying that is actually associated with the excitement.
    Now, these things have to be perfectly aligned
    for us to be able to train the model.
    Now, let's assume that we actually have such a wonderful data set where there's
    perfect alignment between the two different input sensors,
    because multimodal tends to work best for tasks
    where each sensor provides context for a single classification.
    For example, predicting a user's mood based
    on their voice and facial expression.
    So let's assume that we have this data set,
    then how would we actually go about training the model?
    Well, we have to effectively choose an embedding in order to train the model.
    Now, the embedding has to be a joint representation that
    comes in both from the camera stream as well as the audio stream.
    So for images that are coming in from the camera,
    we would have an image encoder that spits out an encoding that
    corresponds to the image.
    Similarly, we would have an encoding that corresponds to the audio input.
    Now, these encoders can be trained separately
    and their output is effectively concatenated and then
    used to train the multimodal classifier.
    Now, let's assume that we know how to do this and everything goes well.
    When it comes to the evaluation conversion process,
    it's exactly the same as what we have done before.
    Now comes the deployment phase.
    Wow, we are very excited about deploying a multimodal network
    on our embedded device.
    Wouldn't that be exciting for person detection, where we were
    able to tell the mood of the person?
    Let's talk about some of the complexities
    that we're going to run into because there
    are some gotchas here where there's a fundamental decoupling
    between the deployment stage and the rest of the pipeline,
    and that's extremely important to understand because this does happen
    in production scenarios where the people who are training the model
    do not realize how the model is going to be deployed
    and that leads into a fundamental mismatch that breaks the application
    use case.
    The complication with deployment really comes
    about as a result of having to deal with synchronizing the input
    buffers from multiple sensors.
    The camera capture is happening at five [? fips ?] at a steady state,
    but the actual speed really depends on the inference latency
    on the embedded microcontroller.
    On the Adreno, one iteration of the whole person detection loop
    takes 740 milliseconds for us.
    And the keyword spotting audio track filled up the one second buffer
    and the amount of data there depends on how long
    it has been since the last inference that was performed
    on the keyword spotting pipeline.
    In either case, if we take too long between the inferences,
    we're going to miss proper alignment for the keyword.
    Which means even though we have this network that
    has been trained with aligned data and actually
    is capable of multimodal execution, the problem you're going to have
    is that the inputs, when they're coming in on the embedded system,
    if they are off and they are not aligned,
    then the network is going to have missed predictions
    and therefore it would simply not work.
    Now, that puts us in a bit of a pickle because this sounds like something
    that we would ideally like to have, but is there
    something that we can do in order to unlock
    this sort of an interesting application use case
    without having to deal with the issues that we are just learning about,
    which is the issue of having to deal with the synchronization
    between the sensors on the deployment aspect of it?
    On that note, I'm going to stop here and I'll
    come back and talk about what I call is multitenancy, which
    gets us a step closer to being able to realize
    the utilization of multiple sensors to unlock
    these kinds of creative applications where
    we can benefit from having rich input coming
    in from multiple sensors in order to build an interesting application.
    So I'll see you there in the next video.
    Ende des Transkriptes, zum Anfang springen
    Downloads und Transskripte
    Video
    Videodatei herunterladen
    Transskripte
    Download SubRip (.srt) file
    Download Text (.txt) file    
    
\section{Person Detection with Keyword Spotting: Multi-Tenancy}

    0:00 / 0:00Klicken Sie auf die Pfeiltaste nach oben, um das Geschwindigkeitsmenu aufzurufen. Regeln Sie dann mit den Hoch und Runter Pfeiltasten die Geschwindigkeit. Klicken Sie Enter, um die jeweilige Geschwindigkeit festzulegen.Geschwindigkeit1.0xKlicke auf diesen Knopf und diess Video stumm oder nicht stumm zu schalten, oder drücke die HOCH oder RUNTER Knöpfe um die Lautstärke zu erhöhen oder abzusenken.Maximal Lautstärke.Videotranskript
    Beginn des Transkriptes, zum Ende springen.
    VIJAY JANAPA REDDI: In this video, I'm going
    to introduce the concept of multi-tenancy, which
    is different than multimodal execution.
    Multi-tenancy is a proxy way of getting close
    to some aspects of multimodal execution.
    Remember, in multimodal execution, what we wanted to do
    was that we wanted to take the inputs that are coming from multiple sensors
    and fuse them together into a combined encoding
    and feed that into a single neural network
    in order to perform a classification.
    However, when it comes to deployment, due
    to rate mismatches between the input sensors,
    be it the microphone or the camera or whatever your favorite sensor is,
    there's a high probability that the data is not going to be aligned.
    Now, if the input data coming into your trained neural network is not aligned,
    then the network might not be able to recognize the things
    that it was originally trained on.
    Now, remember, during the training phase in a multimodal network,
    the data is nicely aligned.
    However, when you're deploying things onto an embedded microcontroller
    where there are many different kinds of sensors
    and there are many different kinds of platforms on which you can deploy,
    and that heterogeneity can potentially lead into problems where
    the rate at which the input data is coming in from one sensor
    might not match the rate at which the other sensor is coming in at.
    And that can lead to failure problems.
    I'm not saying that this is not something that we can never realize.
    I'm just saying that this often tends to happen.
    However, there are ways to get around this.
    Now, let me take a very specific application
    to introduce the concept of multi-tenancy.
    Let's say we want to build a contact-free elevator, where the goal
    is to enforce mask wearing at offices.
    This will fundamentally require us to build two types of neural networks
    that work together.
    Because we're saying contact-free, well, then there
    has to be some way for the user to communicate
    which floor they want to go to.
    Maybe the keywords here are level one, level two, level three, level four,
    and so forth.
    And now we also want the system to enforce mask wearing.
    In other words, we wanted to open the door only
    when it recognizes that the person is wearing a mask.
    Well, in that particular case, after it's received the keyword,
    it would have to check to see if the person is actually wearing a mask
    and then allow the doors to open.
    Now, how would we do this?
    Well, in many ways, you've already been exposed to all the fundamental building
    blocks to enable you to do this.
    We want to run two separate models on the same system, effectively
    at the same time.
    You want the camera module to be running to look
    for masks, which we did in course two.
    And you want the keyword spotting pipeline
    to be running in order to detect the keyword, which you've already done
    and, in fact, gone to the extent of deploying on your microcontroller.
    Now, as these two networks run in parallel, in multi-tenancy,
    you effectively combine the outputs that are coming out of the networks
    into a decision-making process.
    Now let's look at what this means from a deployment pipeline
    sort of perspective.
    If you look at the machine learning workflow for multi-tenancy,
    we don't have to go all the way back to collecting the data.
    Remember, in multimodal, we actually have to go all the way back
    to the data set part.
    And we have to make sure we have data sets
    where things are actually perfectly aligned so that we can train a model
    and then be able to deploy it.
    In the multi-tenancy perspective, you don't have that problem.
    Because you're running the networks independently,
    there is no need for the data to be aligned.
    As such, the left-hand part of the entire pipeline can stay intact
    and you can reuse whatever you have done.
    Now, when it comes to the deployment stage, what we would do
    is that we would cascade the neural network architectures together
    in order to make a combined decision.
    For example, here, you would take the audio that's
    coming in from the microphone, make a decision
    on whether the input is yes or no, and then you would feed that,
    if it is a yes, into the camera pipeline,
    and say that, well, what is the camera saying?
    If the keyword is detected and the camera
    says that the person is wearing a mask, voila, then you can open the doors.
    So hopefully what you're learning here is
    that you can run more than one model.
    Very often when people talk about [INAUDIBLE] applications,
    they tend to deal with a single sensor.
    And what I'm trying to do here is to introduce you to more than that.
    I'm trying to get you to think beyond just a single sensor at a time.
    Hopefully you get that concept.
    And in the next video, what I'm going to do
    is I'm going to show you some of the nuances in the deployment stage
    about deploying the multi-tenancy effectively.
    And I'll see you there.
    Ende des Transkriptes, zum Anfang springen
    Downloads und Transskripte
    Video
    Videodatei herunterladen
    Transskripte
    Download SubRip (.srt) file
    Download Text (.txt) file    
    
 \section{Multi Tenancy in TensorFlow Lite Micro}
 
     0:00 / 0:00Klicken Sie auf die Pfeiltaste nach oben, um das Geschwindigkeitsmenu aufzurufen. Regeln Sie dann mit den Hoch und Runter Pfeiltasten die Geschwindigkeit. Klicken Sie Enter, um die jeweilige Geschwindigkeit festzulegen.Geschwindigkeit1.0xKlicke auf diesen Knopf und diess Video stumm oder nicht stumm zu schalten, oder drücke die HOCH oder RUNTER Knöpfe um die Lautstärke zu erhöhen oder abzusenken.Maximal Lautstärke.Videotranskript
     Beginn des Transkriptes, zum Ende springen.
     SPEAKER 1: In this video, I want to talk about how
     you achieve Multi Tenancy in TensorFlow Lite Micro
     so that you can effectively deploy two neural networks on the device,
     at the same time, to be able to achieve something close to a [? multimodal ?]
     experience.
     Now how would we go about that, and what are
     some of the issues that will come up, are the things
     that I want to talk to you about.
     Remember that we want to deploy two models.
     We want to deploy a vision model and we want to deploy an audio model.
     Now, whenever you're deploying things, a lot of the problems
     will really come down to the amount of memory you actually have.
     We have SRAM and we have Flash.
     The vision model takes up some amount of SRAM and Flash, as shown here.
     Then you have the audio model, which in addition to the vision model,
     will take up a bit more space.
     Now, whether this is a problem or not really
     depends on the device's limit for both the SRAM and the Flash.
     In this particular situation, the SRAM proves
     to be problematic because you only have 256 kilobytes,
     but your combined usage of memory for the audio and vision processing
     is going to exceed that, which means we will have a problem during execution
     time.
     On Flash, this is not a problem.
     Our one megabyte of available storage does not
     seem to be a concern because the vision and audio models
     are not exceeding that device limit.
     So let's try and understand what the problem really is with the SRAM.
     For any given model, be it the vision model or the audio model,
     there are two parts.
     There is the persistent part, and there's the temporary part.
     Think of the temporary as all the temporary storage
     that is required for the input feature, maps, and activations, and so forth.
     The persistent part of the part that will
     correspond to the actual weights of the network, for instance.
     Now, these two combined together-- that is,
     the vision model's temporary plus the vision model's persistent part,
     in addition to the audio model's temporary and the audio model's
     persistent part-- is what is causing us to exceed the SRAM that we actually
     have available on our device.
     Now is there something we can do here?
     Well, there is in fact a trick that you can
     play when you're implementing your code in TensorFlow Lite Micro, which
     is, you can combine the memory usage for the audio model's temporary state
     and the vision model's temporary state.
     By combining them, you reduce the total amount of memory that is being used.
     Because the models are not effectively going to be running concurrently,
     what we can do is that we can actually get the memory
     to be reused across the models.
     Now, I said that the models are not running concurrently.
     Keep in mind here, that we are implementing a cascade architecture.
     A cascade architecture means you have one model that runs first.
     And then based on its input, you might trigger a second model.
     Now because of that, the two models are effectively not competing
     for resources at the same time so you would actually
     be able to do a little bit of optimization here.
     Now, I just want to give you a quick view of what the code op
     changes are going to look like.
     In a standard allocator, for example, what we
     would have is the above definition shown on the right.
     We would have a tensor arena that is set up for the visual [INAUDIBLE] model,
     for the person detector.
     And then, when we look at the audio for the keyword spotting model,
     we would have an associated amount of keyword spotting tensor
     arena dedicated for it.
     And then we would create an interpreter that
     corresponds to the keyword spotting, and then we would allocate its own tensors.
     And we're doing the exact same thing for the vision model up on top.
     Now because of this, we're exceeding the amount of memory allocation
     that's available on the SRAM.
     Now with a shared allocator, however, what we are doing is
     we're actually creating a combined tensor arena,
     where the total amount of memory is actually
     going to be much less in total usage.
     The reason is because again, when we create interpreters
     to run for the audio model and the vision model,
     they're actually running back to back in the cascade style.
     And therefore, they're not going to be competing for the temporary storage,
     and that is something that we can actually reuse.
     And to that end, TensorFlow Lite Micro has actually
     got APIs that allow us to do this very effectively, as shown on the right.
     And we'll talk about this in a bit more detail later on.
     But specifically, the key takeaway has got
     to be that TensorFlow Lite Micro allows us to implement a shared allocator that
     will reuse temporary allocations between models, which helps us conserve space.
     And in doing so, we have the ability now to support Multi Tenancy in execution.
     And this means that we can enable the rich kind of applications
     that we started talking about in the person detection use case.
     With that said, later on, Brian is going to show you
     how we're going to deploy a Multi Tenancy model where the goal is
     to actually be able to use the keyword spotting model that you know,
     and to couple it with the person detection model, which
     is getting introduced to you for the first time now.
     And together, we're not only learning to use the camera model,
     but we're actually unlocking a whole new level of TinyML application space.
     Have fun, and I will see you later.
     Ende des Transkriptes, zum Anfang springen
     Downloads und Transskripte
     Video
     Videodatei herunterladen
     Transskripte
     Download SubRip (.srt) file
     Download Text (.txt) file   
    
\section{Deploying the Pretrained Person Detection Model}

    In this reading, you will deploy the pretrained person detection model onto your Arduino and walk through a little bit of the code to better understand how it works.
    
\subsection{Deploying the Pretrained Person Detection Model}

    0:00 / 0:00Klicken Sie auf die Pfeiltaste nach oben, um das Geschwindigkeitsmenu aufzurufen. Regeln Sie dann mit den Hoch und Runter Pfeiltasten die Geschwindigkeit. Klicken Sie Enter, um die jeweilige Geschwindigkeit festzulegen.Geschwindigkeit1.0xKlicke auf diesen Knopf und diess Video stumm oder nicht stumm zu schalten, oder drücke die HOCH oder RUNTER Knöpfe um die Lautstärke zu erhöhen oder abzusenken.Maximal Lautstärke.Videotranskript
    Beginn des Transkriptes, zum Ende springen.
    BRIAN PLANCHER: Hi, everyone.
    In this reading we're going to be deploying the stock person detection
    example.
    You can find that example by going to the Harvard Tiny MLX library
    and opening up the person detection example.
    We won't be making any changes to the person detection example,
    so we can just simply deploy it.
    Check the board and check the port as always,
    and then you can just click Upload and deploy it onto your device.
    Now that you've gotten the model deployed,
    let's open up the serial monitor.
    In the serial monitor, you're going to see a series
    of person score and no person score.
    And what this is indicating with each integer value,
    is how strongly it believes that there is a person or not
    a person in the image.
    And so negative numbers are bad positive, numbers are good.
    And so what you're seeing, is right now it believes
    there is not a person in the image.
    And if we look at the Arduino, we can see that the light is red
    and that confirms that.
    Now you see these little blue flashes, and the blue light
    flashes every time the camera takes a picture and it does a single inference.
    Now, if I prop this up here on the box and I look down at it from a little bit
    farther away, you can see that I can turn the light green.
    And that's because I'm now far enough and in the right field of view
    that the camera can see my face and it now thinks there's a person.
    And if we look in the serial monitor, you
    can see that the person score is now positive,
    while the no person score is negative.
    And that's it.
    That's how the person detection example works.
    Ende des Transkriptes, zum Anfang springen
    Downloads und Transskripte
    Video
    Videodatei herunterladen
    Transskripte
    Download SubRip (.srt) file
    Download Text (.txt) file
    
\begin{enumerate}
  \item Use a USB cable to connect the Arduino Nano 33 BLE Sense to your machine. You should see the green LED power indicator come on when the board first receives power.
  
  \item  Open the person\_detection.ino sketch, which you can find via the File drop-down menu. Navigate, as follows: File → Examples → Harvard\_TinyMLx→ person detection. Our library version of the example is very similar to the one you’ll find in the TensorFlow Lite for Microcontrolers library, but includes a couple of important tweaks to work better with the TinyML Shield and both the OV7675 and OV7670 so make sure you open the example out of our library.

  \item  Use the Tools drop-down menu to select the appropriate Port and Board. This is important as it is telling the IDE which board files to use and on which serial connection it should send the code. In some cases, this may happen automatically, but if not, you’ll want to select: 

    \begin{itemize}
      \item Select the Arduino Nano 33 BLE as the board by going to Tools → Board: <Current Board Name> → Arduino Mbed OS Boards (nRF52840) → Arduino Nano 33 BLE. Note that on different operating systems the exact name of the board may vary but/and it should include the word Nano at a minimum. If you do not see that as an option, then please go back to Setting up the Software and make sure you have installed the necessary board files.
      
      \item Then, select the USB Port associated with your board. This will appear differently on Windows, macOS, Linux, but will likely indicate ‘Arduino Nano 33 BLE” in parenthesis. You can select this by going to Tools → Port: <Current Port (Board on Port)> → <TBD Based on OS> (Arduino Nano 33 BLE). Where <TBD Based on OS> is most likely to come from the list below where <\#> indicates some integer number:
        \begin{itemize}
          \item Windows → COM<\#>
          \item macOS → /dev/cu.usbmodem<\#>
          \item Linux → ttyUSB<\#> or ttyACM<\#>
        \end{itemize}
      \end{itemize}

  \item   Use the rightward arrow to upload / flash the code. Do not be alarmed if you see a series of orange warnings appear in the console. This is expected as we are working with bleeding edge code. You’ll know the upload is complete when you red text in the console at the bottom of the IDE that shows 100% upload of the code and a statement that says something like “Done in <#.#> seconds.”
    
    If you have the OV7670 camera, make sure to first make the changes indicated at the end of this document!
    
    If you receive an error, you will see an orange error bar appear and a red error message in the console. Don’t worry -- there are many common reasons this may have occurred.
    
    To help you debug other issues, please check out our FAQ appendix with answers to the most common errors!
    
  \item    Open up the serial monitor. You will start to see it print a series of messages about the person score and the no person score. This is indicating the model’s confidence in whether there is either a person or no person in the image!

  \includegraphics[width=\textwidth]{TinyML/3 Development/4-7.png}

    Screenshot of the Serial Monitor showing a series of Person vs. No Person scores. We see that the values can be either positive or negative integers. This is indicating the confidence that the model thinks and image is either a person or not a person. We see it start with positive values for person and negative values for not a person and then switch. This is because Brian started in view of the camera and then moved.

    At this point, you’ll want to look at the board itself. The LED should be flashing blue when an image is being captured and then either red or green. Green indicates that the model detects that there is a person in the image and red indicates that it does not detect a person. Try holding the Arduino at arm's length and point it toward your face. The LED should flash green. As you may see, the field of view is very narrow, so if it’s not working, try propping up the Arduino and standing further back.
    
  \includegraphics[width=\textwidth]{TinyML/3 Development/4-7.2}% .gif}
    
    A GIF of the board LED turning from red to green as Brian offscreen angles it at his face by propping it up on the box.
\end{enumerate}    


\subsection{Understanding the Code in the Person Detection Example}

    Now that you have gotten the person detection application deployed to your microcontroller, let's explore the code a little bit. The main person\_detection.ino file consists of the standard Arduino setup() and loop() functions. These are both quite similar to the hello\_world and micro\_speech example. To really start to understand the code, let's take a look at a couple of the other files.
    
    Open the arduino\_image\_provider.cpp file. Here, we will find the GetImage() function, which captures an image and loads it into the model’s input.
    
    First, let’s take a look at the camera initialization. Here, we set the resolution of the image we receive from the camera to QCIF (176x144). In order to save space and complexity, the input to our model is a grayscale image, therefore, we set the camera to transmit a grayscale image. We also set our camera model to the OV7675.
    
    Note: As shown in the changes to the OV7670 section below, this is where you’ll need to change the OV7675 to OV7670.
    
  \begin{lstlisting}
    // Initialize camera if necessary
    
    if (!g_is_camera_initialized) {
        
        if (!Camera.begin(QCIF, GRAYSCALE, 5, OV7675)) {
            
            TF_LITE_REPORT_ERROR(error_reporter, "Failed to initialize camera!");
            
            return kTfLiteError;
            
        }
        
        g_is_camera_initialized = true;
        
    }
  \end{lstlisting}    
    
    Next, we can capture an image with Camera.readFrame(). Since the input of our model is a 96x96 image, we have to crop the 176x144 image we receive from the camera before passing it to the model. To do this, we copy over the desired pixel values from the center of our original image into the input vector of our model.
    
  \begin{lstlisting}
    // Crop 96x96 image. This lowers FOV, ideally we would downsample
    
    // but this is simpler and more efficient on device.
    
    for (int y = min_y; y < min_y + 96; y++) {
        
        for (int x = min_x; x < min_x + 96; x++) {
            
            // convert TF input image to signed 8-bit
            
            image_data[index++] = static_cast<int8_t>(data[(y * 176) + x] - 128);
            
        }
        
    }
  \end{lstlisting}    
    
    Now that we are familiar with capturing and loading the image, let’s take a look at how the person detection application handles the output of the model. To do that, open the arduino\_detection\_responder.cpp file.
    
    The primary function of the detection responder is to control the LED output. In a real application, this function would drive some other core function that is triggered by the presence of a person.
    
    The blue LED will flash every inference, demonstrating how often an inference is run. Then we will compare the person no\_person scores that are the outputs of the model. In this case, we flash the green LED when the person\_score is higher than the no\_person\_score (both of which are reported over serial). In some applications, we might be more concerned with the false positives or false negatives, in which case, we can adjust the criteria for detecting a person, potentially by using a threshold as we have discussed in Course 2.
    
  \begin{lstlisting}
    //Flash the blue LED after every inference.
    
    digitalWrite(LEDB, LOW);
    
    delay(100);
    
    digitalWrite(LEDB, HIGH);
  \end{lstlisting}    
    
  \begin{lstlisting}
    // Switch on the green LED when a person is detected,
    
    // the red when no person is detected
    
    if (person_score > no_person_score) {
        
        digitalWrite(LEDG, LOW);
        
        digitalWrite(LEDR, HIGH);
        
    }
    
    else {
        
        digitalWrite(LEDG, HIGH);
        
        digitalWrite(LEDR, LOW);
        
    }
  \end{lstlisting}    
    
    TF\_LITE\_REPORT\_ERROR(error\_reporter, "Person score: \%d No person score: \%d", person\_score, no\_person\_score);
    
    And at a high level, that is how the person\_detection example works! We simply take and crop the input images, pass them to the neural network, and then report the higher scoring result!
    
    Changes for the OV7670
    To use the OV7670 instead of the OV7675, simply make the change you made in the camera test! Change the fourth argument of the call to Camera.begin() as highlighted below, from OV7675 to OV7670. That’s it! The library will handle the rest! You can find the camera initialization in:
    
  \begin{lstlisting}
    // Initialize the OV7675 camera
    void if(!Camera.begin(QCIF, GRAYSCALE, 5, OV7675)) {
        Serial.println("Failed to initialize camera");
        While (1);
    }    
  \end{lstlisting}    

%1.7: Visual Wake Words -  Deploying the Multi Tenant Example

\section{Deploying a Multi-Tenant Application}
In this reading, you will deploy the pretrained multi-tenant KWS/person detection model onto your Arduino and walk through a little bit of the code to better understand how it works.

\section{Deploying the Multi Tenant Example}
0:00 / 0:00Klicken Sie auf die Pfeiltaste nach oben, um das Geschwindigkeitsmenu aufzurufen. Regeln Sie dann mit den Hoch und Runter Pfeiltasten die Geschwindigkeit. Klicken Sie Enter, um die jeweilige Geschwindigkeit festzulegen.Geschwindigkeit1.0xKlicke auf diesen Knopf und diess Video stumm oder nicht stumm zu schalten, oder drücke die HOCH oder RUNTER Knöpfe um die Lautstärke zu erhöhen oder abzusenken.Maximal Lautstärke.Videotranskript
Beginn des Transkriptes, zum Ende springen.
BRIAN PLANCHER: All right, let's now deploy the multi-tenant example.
In this example, we're combining the person detection model which you just
deployed, with the pre-trained yes no keyword spotting model,
using the approach described by [? fuge ?] in this section.
To get the multi-tenant example, we're going
to go back to the examples for the Harvard TinyMLx library
and open up multi-tenant Again, just like with the person detection example,
we're going to deploy its stock as is.
So let's just make sure the board and the port are right,
and then we can click Upload, to start compiling and uploading the sketch.
So now let's test out the multi-tenant model.
The way the multi-tenant model works, is that it first
runs keyword spotting as fast as it can, looking for the yes keyword.
Once it hears that YES keyword, it'll turn on the person detection model
and then look for a person.
So let's open up that serial monitor and then let's
prop up this board here, just like I did in the person detection example.
So as i hope that it will see me.
All right.
And then let's trigger the keyword Spotting.
And what you should see is that--
well first of all, you're seeing heard unknown.
And once you hear the key word that we're
looking for, you'll see it appear there and then immediately after that
in the serial monitor, you'll see the results from personal detection.
So let's trigger that.
Yes.
Did you see that?
So on the yes, the light turned blue, and then because it saw me,
person score was bigger than no person score, it turned green.
Let's try that again.
Yes.
Now we can actually also get it to go red if there's not a person.
So we can just leave it up this way and go.
And we're looking again at that LED on the bottom there, yes.
Oops!
Yes.
There and it went blue and then red.
And you can see that even though the YES NO model actually works quite well,
it's not always perfect, it did miss that one YES from me.
But Yeah, that's it.
You now have a working multi-tenant model
and we hope you enjoy using these sorts of ideas
to design your own sensor fusion projects in the future.
Ende des Transkriptes, zum Anfang springen

\begin{enumerate}
  \item Use a USB cable to connect the Arduino Nano 33 BLE Sense to your machine. You should see the green LED power indicator come on when the board first receives power.

  \item Open the multi\_tenant.ino sketch, which you can find via the File drop-down menu. Navigate, as follows: File → Examples → Harvard\_TinyMLx→multi\_tenant.

  \item Use the Tools drop-down menu to select the appropriate Port and Board. This is important as it is telling the IDE which board files to use and on which serial connection it should send the code. In some cases, this may happen automatically, but if not, you’ll want to select: 

    \begin{itemize}
      \item Select the Arduino Nano 33 BLE as the board by going to Tools → Board: <Current Board Name> → Arduino Mbed OS Boards (nRF52840) → Arduino Nano 33 BLE. Note that on different operating systems the exact name of the board may vary but/and it should include the word Nano at a minimum. If you do not see that as an option, then please go back to Setting up the Software and make sure you have installed the necessary board files.

      \item Then, select the USB Port associated with your board. This will appear differently on Windows, macOS, Linux, but will likely indicate ``Arduino Nano 33 BLE'' in parenthesis. You can select this by going to Tools → Port: <Current Port (Board on Port)> → <TBD Based on OS> (Arduino Nano 33 BLE). Where <TBD Based on OS> is most likely to come from the list below where <\#> indicates some integer number.

        \begin{itemize}
          \item Windows → COM<\#>
          \item macOS → /dev/cu.usbmodem<\#>
          \item Linux → ttyUSB<\#> or ttyACM<\#>
        \end{itemize}
    \end{itemize}

  \item Use the rightward arrow to upload / flash the code. Do not be alarmed if you see a series of orange warnings appear in the console. This is expected as we are working with bleeding edge code. You'll know the upload is complete when you see red text in the console at the bottom of the IDE that shows 100\% upload of the code and a statement that says something like ``Done in <\#.\#> seconds.''

If you have the OV7670 camera, make sure to first make the changes indicated at the end of this document!

If you receive an error, you will see an orange error bar appear and a red error message in the console. Don’t worry -- there are many common reasons this may have occurred.

To help you debug other issues, please check out our FAQ appendix with answers to the most common errors!

  \item Open up the serial monitor. You will start to see the output from the ``Yes, No'' keyword spotting model. We can then trigger it with a ``Yes''. Once triggered, the person detection model will run and we will then see the output from that printed to the serial monitor as well.

    \includegraphics[width=\textwidth]{TinyML/3 Development/4-7.3.png}
  
  A screenshot of the serial monitor showing a lot of unknowns and then a yes triggering a person detection result -- one of which found a person and the other a no person. This is done by first seeing outputs from Keyword Spotting as described in the earlier sections (e.g., “Heard yes (score: X)”). This triggers the person detection model which outputs a value as described in the previous section (e.g., ``Person score: X No person Score: Y''). If a ``yes'' is heard and the person score > no person score then the result is triggered.

At this point, you’ll want to look at the board itself. The yellow LED in the top left should be blinking rapidly. Try saying ‘yes’ and see if the board responds. If the micro\_speech model detects ‘yes,’ it will turn the middle LED blue, and then, it will try to detect if a person is in the camera frame. If the person\_detection model detects a person, it will flash the LED green, otherwise, it will flash the LED red. The field of view is quite narrow, so try holding the device at arm’s length and then say ‘yes’ with an emphasis on the ‘sss’ sound.

  \includegraphics[width=\textwidth]{TinyML/3 Development/4.7-4} % .gif}

  Gif of the mutli tenant model triggering pointing at Brian and turning blue then green (right) and pointing up into the air and turning blue then red (left)

\includegraphics[width=\textwidth]{TinyML/3 Development/4.7-5} %.gif}

  Gif of the mutli tenant model triggering pointing at Brian and turning blue then green (right) and pointing up into the air and turning blue then red (left)
\end{enumerate}

\subsection{Understanding the Code in the Multi-Tenant Example}

Now that you have gotten the multi\_tenent application deployed to your microcontroller, let's explore the code a little bit. The main multi\_tenant.ino file is similar to the previous person detection example. The difference here is that we use the shared allocator from TFLite Micro to allocate space for two models using the same tensor arena.

Unlike previous examples where we directly passed the tensor\_arena and kTensorArenaSize to the interpreter, we need to directly create an allocator and pass it to our two separate interpreters.

In the setup() function, we create our allocator that will be shared between both of our interpreters as follows:

\begin{lstlisting}
tflite::MicroAllocator* allocator =

tflite::MicroAllocator::Create(tensor_arena, kTensorArenaSize, error_reporter);
\end{lstlisting}


Then we create the interpreter for the person detection model, passing the allocator to it as a parameter, and then calling AllocateTensors().

\begin{lstlisting}
static tflite::MicroInterpreter vww_static_interpreter(

vww_model, micro_op_resolver, allocator, error_reporter);

vww_interpreter = &vww_static_interpreter;
\end{lstlisting}


\begin{lstlisting}
// allocate the VWW model from tensor_arena

TfLiteStatus allocate_status = vww_interpreter->AllocateTensors();
\end{lstlisting}

Now that both interpreters are created using the same allocator, we can use them as we would normally.

In the loop() function, we first run the micro\_speech model as we did in the KWS example, except this time if the model detects ‘yes’ we capture an image and run the person\_detection model. The LED control is performed in arduino\_detection\_responder.cpp and follows the patterns used in both the micro\_speech and person\_detection examples.

\begin{lstlisting}
bool heardYes = RespondToKWS(error_reporter, found_command,
\end{lstlisting}

\begin{lstlisting}
is_new_command, score);
if(heardYes){
    //Our keyword spotting model heard 'yes' so we detect if a person is visible
    // Get image from provider.
    if (kTfLiteOk != GetImage(error_reporter, kNumCols, kNumRows, kNumChannels,
    vww_input->data.int8)) {
        TF_LITE_REPORT_ERROR(error_reporter, "Image capture failed.");
    }
    
    // Run the model on this input and make sure it succeeds
    if (kTfLiteOk != vww_interpreter->Invoke()) {
        TF_LITE_REPORT_ERROR(error_reporter, "Invoke failed.");
    }
    
    // Process the inference results.
    int8_t person_score = vww_output->data.uint8[kPersonIndex];
    int8_t no_person_score = vww_output->data.uint8[kNotAPersonIndex];
    RespondToDetection(error_reporter, person_score, no_person_score);
    
}
\end{lstlisting}


\subsection{Changes for the OV7670}

To use the OV7670 instead of the OV7675, simply make the change you made in the camera test! Change the fourth argument of the call to Camera.begin() as highlighted below, from OV7675 to OV7670. That’s it! The library will handle the rest! You can find the camera initialization in:

\begin{lstlisting}
// Initialize the OV7675 camera
void if(!Camera.begin(QCIF, GRAYSCALE, 5, OV7675)) {
    Serial.println("Failed to initialize camera");
    While (1);
}
\end{lstlisting}

\section{Recap: Time Series for Anomaly Detection}

\subsection{Recap: Time Series}

In this reading, we will do a quick recap of the time series from Course 2 as we will be leveraging that knowledge for a slightly different application than what we discussed previously (i.e., anomaly detection): gesture recognition with a custom-built magic wand.

\subsection{Supervised vs. Unsupervised Learning}

So far, we have tackled many interesting machine learning tasks that cover a variety of sensory modalities (e.g., vision, audition). All of these tasks required a labeled dataset - a known output value for a given set of inputs - and were thus comfortably situated within the supervised learning paradigm, which we first discussed in Course 1.

However, frequently, we encounter datasets in which our output variable is unknown. Whether because we do not have access to the output variable, or we are merely interested in the internal structure of our dataset, we cannot resort to our standard machine learning framework, which relies on training and test sets. Instead, we must adopt an alternative, equally useful, a paradigm known as unsupervised learning to approach these tasks, which we discussed in Course 2. 

There are still similarities between the supervised and unsupervised learning paradigms. Firstly, both assume the existence of a dataset that is used to train our machine learning algorithm. The performance of both generally improves as more data becomes available. Nevertheless, unsupervised learning is distinct in that we do not require the existence of a predictor variable. Unsupervised learning looks at the internal structure of a dataset and attempts to either group this data systematically, known as clustering, or look for abnormalities within the dataset, known as outlier analysis, or anomaly detection.

\subsection{Anomaly Detection}

In course 2, we were mostly interested in anomaly detection, predominantly due to its potential for predictive maintenance (i.e., noticing abnormalities in machines that can be resolved before catastrophic events). Anomaly detection can be performed spatially or temporally, either looking at associations between different physical locations, or the same location at different points in time.

Most sensory devices are fixed in discrete locations, such as those monitoring the temperature and vibrational characteristics of a piece of machinery. Inherently, these devices output a temporal signal that can be monitored and used to look for anomalies over time. These anomalies manifest in myriad ways: a sudden temperature change, a sudden increase in vibration or sound, or a sudden drop in pressure. Given sufficient explanatory variables, anomaly detection algorithms allow us to detect outliers and perform corrective actions in a timely manner to mitigate the possibility of catastrophic events, such as an error, a machine breakdown, or even an explosion.

Many algorithms exist for the detection of anomalies in time series data. One algorithm that was discussed in the previous course was k-means, which looks at the distance of a data point in feature space from a set of discrete clusters that represent the majority of data points. We saw that one drawback of k-means was its poor performance on high-dimensional data due to the curse of dimensionality. However, we were able to improve performance by reducing the dimensionality of the data with t-SNE. Anomaly detection can also be achieved using neural network structures known as autoencoders, which attempt to collapse knowledge of the dataset in a small set of neurons located at the mid-layer of the autoencoder (this formulation is very similar to principal component analysis). Autoencoders tend to have superior performance to traditional methods due to their ability to capture complex non-linear associations.

\subsection{Magic Wand}

In the remainder of this section, we will leverage this prior knowledge for a slightly different purpose: building in gesture recognition using a custom-built magic wand.

\subsection{TinyML Sensor Ecosystem}

    0:00 / 0:00Klicken Sie auf die Pfeiltaste nach oben, um das Geschwindigkeitsmenu aufzurufen. Regeln Sie dann mit den Hoch und Runter Pfeiltasten die Geschwindigkeit. Klicken Sie Enter, um die jeweilige Geschwindigkeit festzulegen.Geschwindigkeit1.0xKlicke auf diesen Knopf und diess Video stumm oder nicht stumm zu schalten, oder drücke die HOCH oder RUNTER Knöpfe um die Lautstärke zu erhöhen oder abzusenken.Maximal Lautstärke.Videotranskript
    Beginn des Transkriptes, zum Ende springen.
    VIJAY JANAPA REDDI: Hi, in this module, I'm
    very excited to be embarking on a whole new application journey with you which
    is around the magic wand.
    Now let me take this opportunity to talk about the kind of sensors
    that we have been using and what the prospects are
    if we expand the landscape of sensors that we look at specifically for TinyML
    applications.
    Now generally when you talk about machine
    learning as it's applied to sensors, it often
    ends up being focused on two domains.
    It happens to be speech, and it happens to be vision.
    However, when you're talking about embedded sensors,
    there's a large number of sensors out there
    from being able to measure temperature to pressure to humidity.
    To being able to tell the strength of the magnetic field, and so forth.
    There are many different kinds of sensors out there.
    And so if we can tap into some of these rich array of sensors,
    then we might be able to uncover a whole new range of applications that are
    specifically unique to TinyML itself.
    And not just the usual machine learning applications
    from just speech and vision.
    Now we will specifically be talking about the IMU.
    The inertial measurement unit is a really rich sensor.
    It generates a lot of data in a very short amount of time.
    And it gives you very rich information as well
    if you know how to deal with the data that's coming through it.
    Specifically, a 9 degree of freedom IMU typically has three sensors in it.
    It's got the accelerometer, the gyroscope, and the magnetometer.
    You can do really interesting things by collecting
    the data that's coming from the sensors, and fusing them together.
    For example, the three axis gyroscope has the ability
    to give us rotation data around a known point.
    While the accelerometer has the ability to tell us about
    displacements in the x-axis, y-axis, and z-axis.
    Now a post processing step down on the microcontroller
    might actually allow us to infer some interesting properties
    about the experience of the sensor as a whole.
    And do so very robustly because each of these sensors
    has some pitfalls that I'll talk about later.
    Now why am I going on and on about the IMU?
    Well, I'm trying to get you excited about the possibilities of using
    other kinds of sensors.
    Now a real practical reason as to why you would consider other sensors out
    of things like microphones and just cameras
    is really because of power consumption.
    Now remember TinyML is all about running on ultra low power consumption devices
    as much as using just a small coin cell battery.
    Well, let's look at a sensor here.
    If you look at the image sensor, it consumes around 98 milliwatts.
    But if you look at a standard IMU sensor,
    it only consumes about 9.5 milliwatts.
    That's an order of magnitude difference in the power consumption.
    So if you think about the TinyML system or the device,
    then the sensor that you use can actually make a very big difference.
    Now, of course, this depends on the application
    that you're choosing to build.
    But the general point I'm getting across here
    is that these very small little sensors have the ability
    to consume very low power.
    And still they can enable very rich application experiences.
    Moreover, you also want to think about the size of the device
    because you want these sensors to be really small so you can put them
    everywhere.
    And they have got to be lightweight, especially for robotic applications.
    And you want them to be really cheap so you can use a lot of them
    and do very interesting TinyML applications using them.
    Because of their low cost and because of the low power consumption,
    there are countless applications.
    For example, you can do embedded sensing on bodies for activities like sports
    or activity recognition or detecting medical [INAUDIBLE] before they
    actually become problematic.
    You can enable a rich number of robotic applications.
    Specifically IMUs are very useful for things like stabilization.
    And, of course, IMUs are already used today in things like the Apple Watch
    for instance, which actually detects a person suddenly falling down.
    And then can actually alert an emergency health worker.
    Now there are many different applications.
    Now you can think of many other ones well beyond the ones
    that I'm actually showing you here.
    All that sounds exciting.
    But what are some challenges?
    Well, there are some gotchas.
    And they're largely about interpretability, sensor drift,
    and deployment sensitivity.
    Interpretability is the notion that when you look at the sensor data,
    can you really make sense of it?
    For example, we can see images.
    For example, we can hear the audio that we're actually playing.
    But when you look at some of the sensor data that's
    coming in from things like the IMU, we don't really
    have an easy way to interpret that kind of data.
    And that kind of makes it a little challenging
    to be able to easily build applications.
    That's the first thing.
    The second thing is that the sensors are not quite ideal.
    In other words, sensors tend to behave differently over time.
    Now an ideal sensor would be one that has
    the exact same behavior on day 1,000 as it does for instance, on day zero.
    Well turns out, that many of these small little sensors
    tend to have problems like drift for example.
    Gyroscopes are specifically subject to bias instabilities.
    Now, that means that over time their initial zero reading ends up drifting.
    And that can be problematic if you have made some sort of assumptions about how
    to engineer the application based on day zero
    and then deployed it into the field.
    But 1,000 days afterward, the sensor has a certain drift.
    That means you've got to suddenly compensate for that.
    But how would you do it if you already deployed into the field?
    Now, accelerometers commercially, they're sensitive to vibration.
    That any kind of non gravity accelerations can also affect them.
    So you have sensors that have this capability,
    but they also come with the pitfalls.
    And so you want to understand some of this in order
    to be able to use them effectively in applications.
    Last one is really around deployment sensitivity.
    You often have to calibrate or fine tune some of these sensors.
    There tend to be these minor differences that
    can actually have a major effect on the application outcome.
    Some of which we'll actually see as we build a magic wand application.
    These differences often end up manifesting
    in terms of biases in the data that you have to compensate
    for in the pre-processing stage.
    And we'll look at that in our application.
    All that said, the thing that I'm trying to get you excited about
    is the fact that there are countless applications that can use one or more
    of these sensor for purposes that are well beyond the scope of what
    we are covering here.
    We will specifically focus on the magic wand application
    that uses the accelerometer and the gyroscope inside the IMU.
    But keep in mind that that's just the tip of the iceberg,
    and I'm looking forward to seeing some of the projects you will all build
    after you're done with this course.
    That said, let's get started with the magic wand application.
    Ende des Transkriptes, zum Anfang springen
    Downloads und Transskripte
    Video
    Videodatei herunterladen
    Transskripte
    Download SubRip (.srt) file
    Download Text (.txt) file

\section{Anatomy of an IMU}

\includegraphics[width=\textwidth]{TinyML/3 Development/4-8.8.png} 

Icons for Accelerometer, Magnetometer, and Gyroscope from the previous video. The acceleraometer shows a square with arrows in three directions indicating the XYZ planes. The magnetometer shows a standard U-shaped magnet. The gyroscope shows a 3D gyroscope, indicating rotations in 3 dimmensions.
    
An inertial measurement unit (IMU) is a system composed of sensors that relay information about a device’s movement, such as accelerometers, gyroscopes, and magnetometers. In concert, each of the sensors provides a detailed account of a device’s instantaneous motion, orientation, and acceleration, which are often used for navigational and stabilizing purposes in vehicles and aircraft as part of electronic feedback control systems. We will be using it for our Magic Wand application. The Arduino Nano 33 BLE sense has a 9-axis IMU.
    
    More recently, applications of the IMU to wearable devices have been stimulated by the quantified self-movement. The data generated from IMU’s can be used for machine learning purposes in classifying different types of actions associated with a device. For example, by attaching an IMU to a specific body part, such as an arm, it may be possible to discern different gestures purely from time-series information, such as waving, shaking hands, or typing on a keyboard. This makes the IMU particularly useful for embedded machine learning applications that are gesture-activated, such as lighting up your phone screen when picking it up, when checking the time on a smartwatch, or even fall detection.
    
\subsection{Structure of an IMU}

    An IMU consists of at least two of the following sensors: an accelerometer, a gyroscope, and a magnetometer. In principle, these devices can be designed to provide single- or multi-axis measurements, but most offer multi-axis measurements to provide an accurate set of information associated with the position and orientation of a device in three-dimensional space. The importance of the IMU for sensing applications is highlighted by our Arduino Nano BLE Sense, which comes packaged with a built-in 9-axis IMU.
 
\includegraphics[width=\textwidth]{TinyML/3 Development/4-8.9.png} 
    
    Image of the Arduino with arrows for acceleration and rotation drawn on it. We see three acceleration and three rotation arrows showing that we are measuring XYZ acceleration and roll, pitch, yaw rotations.
    
    The three sensors that constitute an IMU afford for the following capabilities:
    
    \textbf{Accelerometer.} Measures changes in velocity (acceleration) and position (velocity), as well as absolute orientation. The accelerometer is the device in tablets and smartphones which ensures the image on-screen remains upright regardless of orientation. By itself, the accelerometer provides information about the linear and rotational X-, Y-, and Z- directions. The accelerometer allows the 3-axis of motion to be captured. An onboard accelerometer is modeled as a micromechanical damped mass-spring system, wherein the compression or extension of the mass-spring system can be mathematically related to an object’s acceleration.
    
    \textbf{Magnetometer.} Establishes cardinal direction (directional heading). A simple example of a magnetometer is a compass, which is used to measure the direction of the Earth’s magnetic field. Smaller versions called microelectromechanical magnetic field sensors can be incorporated into integrated circuits, allowing them to be combined with other sensors as part of an IMU. The magnetometer measures orientation to magnetic north on the X, Y, and Z axes. Together with the accelerometer, the 6 degrees-of-freedom of a system can be accounted for, which fully describes the kinematics of a system. Oftentimes, magnetometers work via the Hall effect, which involves the creation of a potential difference across a conductor as a result of a perpendicular applied magnetic field.
    
    \textbf{Gyroscope.} Measures changes in orientation (rotation) and rotational velocity. Microelectromechanical gyroscopes, often called gyro meters, are present in many consumer electronics such as gaming controllers. A gyroscope provides information about the rotational X- (roll), Y- (pitch), and Z- (yaw) directions. The 6 degree-of-freedom system created by combining a magnetometer and accelerometer suffers from several shortcomings, such as the sensitivity of magnetometers to time-varying magnetic disturbances, and the corruption of accelerometer readings due to the presence of linear acceleration distorting the Earth gravity vector.
    
    To alleviate these shortcomings, a 3-axis gyroscope can be added, creating a 9 degree-of-freedom, or “gyro-stabilized” system. The gyroscope provides the system with an independent measurement of instantaneous rotation speed, complementing the original 6 degree-of-freedom system. The gyroscope works by the conservation of angular momentum; a spinning wheel or disc rotates at high speed, while the spin axis is freely allowed to rotate and assume any orientation. The inertia of the spinning wheel allows it to remain unperturbed in the same direction during tilting or rotation, allowing rotational information to be discerned.
    
\subsection{Next Steps}
    In the following sections, we will focus on obtaining time-series information from the onboard IMU, and subsequently, use this information to train a machine learning model capable of detecting specific actions of a magic wand.
    
\section{Magic Wand Application}

    0:00 / 0:00Klicken Sie auf die Pfeiltaste nach oben, um das Geschwindigkeitsmenu aufzurufen. Regeln Sie dann mit den Hoch und Runter Pfeiltasten die Geschwindigkeit. Klicken Sie Enter, um die jeweilige Geschwindigkeit festzulegen.Geschwindigkeit1.0xKlicke auf diesen Knopf und diess Video stumm oder nicht stumm zu schalten, oder drücke die HOCH oder RUNTER Knöpfe um die Lautstärke zu erhöhen oder abzusenken.Maximal Lautstärke.Videotranskript
    Beginn des Transkriptes, zum Ende springen.
    VIJAY JANAPA REDDI: In this video, I'm excited to introduce
    the magic wand application.
    First, I'm going to talk about what the goal of the application is.
    And then, I'm going to talk about how we can actually use the IMU data along
    with the TinyML system we have in order to realize the application.
    Now, here is a very simple magic wand.
    Now this wand has got a very simple LED attached to the end of it.
    Now, what I want this wand to be able to do
    is to be able to recognize a certain gesture that I'm making.
    That's the ultimate goal of the magic wand application.
    Let me first start off by asking you to recognize the number that I'm actually
    writing in front of you.
    If you guessed that the number is zero, you're right.
    Now let's say, I'm writing this number.
    What is that number that you see?
    If you guessed that the number is eight, excellent.
    Now, the question is this, how is it that you know, just
    from my simple gestures, that the number is either a zero or an eight?
    That's exactly what we're going to try and figure out
    how to realize in our TinyML system using the initial measurement
    unit that I talked about earlier.
    Now, before we get into something quite complicated like that,
    let's take a step back, and let's just say
    that this is a very simple magic wand.
    The goal of this magic wand is quite simple.
    Whenever I swing it in a particular direction, I want the LED to turn off.
    Whenever I swing it again in that exact direction, I want the LED to turn on.
    Now, stop for a moment and think about how you
    might be able to do this using an IMU.
    Now, an IMU has got an accelerometer inside it.
    In other words, the IMU can actually tell
    when there's sudden changes in a certain direction of the IMU.
    Now, I could very well use this to be able to say if the wand was
    swung in a certain direction.
    Now, if it turns out that I happen to swing the wand in a certain direction
    with a certain level of acceleration, and if that acceleration exceeds
    a certain threshold, then I would be able to, say, turn the LED on.
    And if I wanted to turn it off, I would simply
    look for the exact same level of acceleration,
    and then realize that the LED was on, and then simply turn it off.
    This works great for very simple gestures, like simply waving it.
    What if I wanted to make something more sophisticated and complicated,
    like the example that I'm showing you here?
    What if I wanted to be able to draw a gesture, like the number zero,
    or write the alphabet Z?
    How would we go about doing that?
    Let's see what would happen if we fall back
    to our non-ML based methods or our very simple accelerometer-based methods.
    Well, we could try and threshold things on a sequence.
    We could try and build a sequence of those certain actions we're making
    and then be able to see if those sequences match
    to some sort of predetermined categories that we might have.
    For example, if I want to draw the number zero,
    how would I go about doing that?
    Well, I might move the wand to the right.
    Therefore, the acceleration in the x-axis has increased quite a bit.
    Therefore, I record that.
    And then, naturally, I would move the wand down,
    which means that the y-axis acceleration increases.
    So I would record that.
    And then, I would move it again in the x-axis, so I record that.
    And finally, I record the Y.
    Then, I can take the sequence and then try and pair it or match it
    to some sort of known history and say, does this sequence
    line up with the previous sequence that I know that corresponds to the number
    zero?
    Well yes, conceptually, that would work.
    But there are lots of nuances in actually getting this to work robustly.
    This is challenging because this ends up being very sensitive to the speed
    and shape of the gesture.
    Even small differences here could lead to major issues.
    In other words, even though you're writing the number zero,
    well, small variations in the way you've done it could lead to mismatches,
    and that could be problematic.
    Can we do better?
    The answer is actually yes.
    You can use angular position, which gives you much higher fidelity in terms
    of the recording you're making.
    An improvement of this would be to use the angular values
    that we get from the gyroscope to be able to map the position of the IMU
    onto the spherical coordinate system.
    This also helps mitigate the problem with differences in gesture speed
    and direction since we're only estimating
    the shape of the actual gesture that we're making.
    Well, that's a step forward.
    However, there is a gotcha.
    And the gotcha is it's kind of hard to visualize.
    And therefore, it's actually hard to classify
    what that number or alphabet might be.
    Well, what can we do?
    In order to visualize it, then, what if we could turn that angular rotation
    that we have in the gyroscope into a two-dimensional plot of the gesture?
    This allows us to more easily visualize the data.
    Now we can take those points that we have collected
    and then turn them into a 2D image through a process called rasterization.
    Images with defined size are more standard
    and are easier to compare to each other.
    And so what makes it easy is that we can effectively
    be able to save the references that we create into our disk.
    And in the future, when we do a certain sequence,
    we can do the rasterization again and then we can compare them.
    Specifically, we can see what the root mean square error is
    between the two of them.
    And if it turns out the root mean square error in both of them
    is actually quite small, then we can classify that they're actually a match.
    This seems great.
    Now, let's try and understand what the full pipeline would look like.
    The full pipeline would start with vectorization of the IMU data
    that's coming in from the input side, which would then rasterize and flatten
    into a two-dimensional image, which we would then save as a reference.
    And then future inputs, we would rasterize and then
    compare to that reference that we actually have in order
    to classify the gesture as an output.
    Well, that seems great.
    It's certainly a step forward.
    But let's see what ends up happening.
    The problem with this method is that it directly
    compares to a single saved reference, meaning our classification is
    very sensitive to any changes in shape.
    And therefore, we're sort of back at the square one problem,
    which is that subtle changes don't really
    help us address the fundamental issues that we're facing.
    Now, you know one place where machine learning is very useful.
    It's very useful in dealing with small variations as noise.
    And it learns to ignore those small variations
    and just tries to find the general structure that's
    actually encoded in the image.
    So what if one approach was to take just the raw sensor
    data that's coming in from the IMU, and we
    feed that straight into the neural network
    and ask it to predict the gesture?
    Well, that would be an end-to-end learning method.
    Now, this would require a lot of data.
    And it would require a fairly large model
    in order to get good accuracy, especially in complicated gestures.
    So this is not really a preferred method if you
    want to do things in a TinyML system.
    However, if we go back to the original proposal that we had,
    which was certainly an improvement, where
    we were able to take the image input, and then we
    were able to flatten and rasterize it-- and then
    the problem that we had previously was simply that our comparison
    method was kind of flawed.
    We needed a more robust comparison method.
    Now, what if we could instead use a neural network in place of the root
    mean square error check?
    We already know that neural networks are awesome at being
    robust to small little variation.
    So therefore, we could use a neural network to actually do the comparison.
    Now, this is awesome because this means that suddenly we're
    able to take the existing pipeline, of being able to take the input,
    doing the flattening, use the machine learning method
    for what it's really good at, in doing the classifications.
    In this way, we'll be able to realize a rather robust magic wand
    application that is capable of recognizing gestures--
    specifically, gestures that correspond to numbers or alphabets.
    And so in the next video, I'm going to talk a little bit more
    about how we actually go about implementing that
    from a code perspective.
    Hopefully, the key takeaway you've had here are twofold.
    One is IMU sort of data is really rich, and so you
    can do very cool applications, like recognizing certain gestures.
    And two, you don't necessarily want to use
    machine learning end to end because it's not necessarily the best approach.
    What you want to do is you want to understand when
    you want to use traditional methods.
    And you want to couple those traditional methods
    with the right place where that machine-learning method is actually
    effective, which, in our case, was to robustly
    be able to classify images which are coming from the rasterization
    process, which is totally non-machine learning based.
    So keep that in mind.
    And I'll see you in the next video.
    Ende des Transkriptes, zum Anfang springen
    Downloads und Transskripte
    Video
    Videodatei herunterladen
    Transskripte
    Download SubRip (.srt) file
    Download Text (.txt) file    

\section{Magic Wand Application Architecture}

    0:00 / 0:00Klicken Sie auf die Pfeiltaste nach oben, um das Geschwindigkeitsmenu aufzurufen. Regeln Sie dann mit den Hoch und Runter Pfeiltasten die Geschwindigkeit. Klicken Sie Enter, um die jeweilige Geschwindigkeit festzulegen.Geschwindigkeit1.0xKlicke auf diesen Knopf und diess Video stumm oder nicht stumm zu schalten, oder drücke die HOCH oder RUNTER Knöpfe um die Lautstärke zu erhöhen oder abzusenken.Maximal Lautstärke.Videotranskript
    Beginn des Transkriptes, zum Ende springen.
    VIJAY JANAPA REDDI: In this video, I'm going to introduce the Magic Wand
    Application Architecture.
    Now, as you may recall, as with any other application,
    there are two major components, which is the initialization
    phase followed by the main loop.
    And the main loop has many different components that I'll talk about soon.
    Let's talk about the initialization phase.
    The initialization phase job is primarily to set up the IMU
    and then set up all the resources needed to run the TensorFlow like macro model.
    The first step of the initialization phase,
    is the IMU initialization, which is done using this set up IMU routine.
    When you go into the set up IMU routine, you
    will find device specific calls that tap into the IMU functions that
    are provided by the library.
    So please make sure that if according to whatever device you're using,
    you're setting these IMU's up.
    Now, if you're taking this course, you should really just
    be focused on using the Arduino, as none of 33 bearly sense
    IMU, that is actually on device.
    So there's no physical integration.
    You should readily be able to use this.
    The second component of it, is really setting up all the resources
    needed to run the model.
    This might be the pointer to the model, which
    you would get using Git model, the initialization of the interpreter using
    the Tensor arena, the model, the observer and so forth.
    I trust you know how to do all of these things,
    because we've talked about them in a lot of detail in keyword spotting as well,
    as well in the visual wakewards and person detection one.
    If these don't ring a bell, please go back to those videos
    and get yourself a refresher.
    After this we're inside the preprocessing loop or the main loop,
    where we start off with the IMU provider, whose job is
    to get data in from the gyroscope and the accelerometer
    and then be able to process that data.
    There are function calls that are readily available,
    that will actually allow us to read the data
    and from the gyroscope and the accelerometer.
    So as long as there's data available from the IMU,
    we're going to go ahead and process that data.
    If you remember, I said that the gyroscope ends up
    having a little bit of drift, so what we have to do
    is we have to compensate for that drift.
    And that's what this function estimate gyroscope drift is really doing.
    When we determine that the IMU is not moving using the accelerometer,
    we can then estimate the drift of the gyroscope
    and then account for it subsequently.
    Next, we want to integrate the gyroscopes incremental angular
    changes that are coming in over time, because that's going to really give us
    a part of the gesture in the spherical coordinate system
    and that's what this function update orientation is doing.
    It's trying to continuously capture that part that's coming in.
    Next we effectively want to project that inside this physical system,
    into a two dimensional plane.
    And if you remember why we're doing this.
    Well that's because it's much more easier for us
    to be able to understand a 2D gesture than it is a complex 3D gesture.
    And therefore, update stroke is going to do that flat mapping.
    Then comes the process of actually processing the accelerometer data.
    We want to estimate the direction of the gravity,
    so that we can actually control for the role of the sensor
    in the gyroscope readings.
    Now, remember everyone is going to be holding the wand in a certain angular
    momentum.
    So which means that you have to be able to neutralize or normalize
    for that effect.
    For example, you might be holding the wand like this,
    I might be holding the wand like this.
    However, the gesture that we're performing is the exact same thing.
    In either way, we've got the exact same number being written,
    so we get to compensate for that.
    And the way we do that is by effectively trying to figure out what's
    the role from the gyroscopes reading.
    And then we update the velocity, so we know when the sensor is still
    and we can update the sensor drift.
    Now the step that comes after doing that,
    which is effectively capturing the data, is to then rasterize that stroke.
    You already know why we do this.
    Because it's easy for us to feed something
    that's rasterizing into the image, into a convolutional neural network.
    And there is a function that helps us do that, which is rasterized stroke.
    The next part, after we have done all that preprocessing
    is then to hand that rasterized image directly
    to our convolutional neural net.
    In this case, we're actually going to be passing in an RGB image, which
    is a red, green, and blue image.
    So there are three challenge to the image that we're passing into the net.
    And that input is then going to be run using
    a convolutional neural network, which is going
    to make a prediction on the gesture.
    In order to invoke the model, as always we
    have to set up the input buffers, which is exactly what I'm showing you here.
    And following that, we simply get the output from the neural network,
    to see what it has determined as actual gesture.
    And in terms of processing the output, we
    simply actuating LED or we also print the output to the screen
    or there's something else you might be able to do, depending on your interest
    after we have recognized the gesture.
    And that's all there is to the overall application architecture.
    We've gone from understanding how to initialize the code,
    to actually post-processing them.
    So at this point, I would strongly encourage you to take a breather
    and go and look at the actual example code
    and try and understand some of these bits and pieces from the comments
    that we have provided.
    With that said, I hope you enjoy implementing
    the magic wand application.
    And I'm looking forward to seeing all small variants
    that you're likely to do on top of it, for example
    using the magic wand gesture recognition,
    to be able to produce useful applications out of.
    That said, have fun and I'll see you later.
    Ende des Transkriptes, zum Anfang springen
    Downloads und Transskripte
    Video
    Videodatei herunterladen
    Transskripte
    Download SubRip (.srt) file
    Download Text (.txt) file
    
\section{Understanding the Magic Wand Application}

    This reading is about understanding the components of the magic wand application and how they generalize to other IMU-based projects. By gaining familiarity with the data flow, from the sensor to classification, you can apply similar concepts to other applications.
    
    \includegraphics[width=\textwidth]{TinyML/3 Development/1642303958726932648973.png}
    
    A dataflow visualization of the magic wand application. The accelerometer data is fed into the estimate gyroscope drift and estimate gravity direction methods. The gyroscope data is passed into estimate gyroscope drift and update orientation. Update orientation is used to produce the rotational orientation data which is, in turn, passed to the update stroke method. Update stroke produces the x/y coordinates of the gesture in a flat 2D plane. The x/y coordinates are sent over BlueTooth to the data capture. The x/y coordinates are rasterized into a 2d image and then displayed and passed to the model which outputs a classification.
    
    The Magic Wand application accomplishes a fairly complicated task (gesture recognition) by carefully crafting a 2D image from 3D IMU data. The dataflow displayed in the diagram above consists of 14 steps:
    
  \begin{enumerate}
    \item The accelerometer data is read and passed to the EstimateGravityDirection() where it is used to determine the orientation of the Arduino with respect to the ground.
    \item The Accelerometer data is passed to UpdateVelocity() where it is used to calculate the velocity of the Arduino.
    \item The direction of gravity is passed to UpdateVelocity() and is used to cancel out the acceleration due to gravity from the accelerometer data.
    \item The velocity is then passed to EstimateGyroscopeDrift() where it is used to determine if the Arduino is stationary or moving.
    \item The gyroscope data is passed to EstimateGyroscopeDrift() where it is used to calculate the sensor drift of the gyroscope if the Arduino is not moving (velocity is 0).
    \item The gyroscope data is passed to UpdateOrientation() where it is integrated to determine the angular orientation of the Arduino.
    \item The gyroscope drift is also passed to UpdateOrientation() and subtracted from the gyroscope reading to cancel out the drift.
    \item The 3D angular orientation is passed to UpdateStroke() transformed into 2d positional coordinates. UpdateStroke() also handles whether the current gesture has ended or if a new gesture has been started by analyzing both the length of the gesture and testing whether the orientation data is still changing.
    \item The direction of gravity is also passed to UpdateStroke() to determine the roll orientation of the Arduino.
    \item The 2d positional coordinates are sent over BlueTooth to the data collection application in the browser.
    \item The 2d positional coordinates are passed to RasterizeStroke() which takes the 2D coordinates and draws lines between them on a 2D image. The color of the lines shifts from red to green to blue to indicate the direction of motion during the gesture.
    \item The 2D image of the gesture is converted to ASCII art and printed on the serial monitor.
    \item The 2D image of the gesture is passed to the model.
    \item The model predicts the label of the gesture and the label is printed to the serial monitor.
  \end{enumerate}

    For a deeper understanding of the Magic wand application, try following the flow diagram while reading the code. The comments should provide the context needed to understand the function of each section. Try to piece the application together, one piece at a time, as there is a lot going on in this application. That said, each piece by itself is fairly self-contained.    
    
\section{Deploying the Magic Wand Project}

    In this document, we are going to deploy the pre-trained magic wand project, which can recognize gestures of the digits 0-9 by analyzing the time series data produced by the IMU!
    
\subsection{Deploying the Pretained Magic Wand Project}

    0:00 / 0:00Klicken Sie auf die Pfeiltaste nach oben, um das Geschwindigkeitsmenu aufzurufen. Regeln Sie dann mit den Hoch und Runter Pfeiltasten die Geschwindigkeit. Klicken Sie Enter, um die jeweilige Geschwindigkeit festzulegen.Geschwindigkeit1.0xKlicke auf diesen Knopf und diess Video stumm oder nicht stumm zu schalten, oder drücke die HOCH oder RUNTER Knöpfe um die Lautstärke zu erhöhen oder abzusenken.Maximal Lautstärke.Videotranskript
    Beginn des Transkriptes, zum Ende springen.
    BRIAN PLANCHER: In this video, I'm going to provide you with some tips
    for building your own magic wand and deploying the default Magic Wand
    project.
    Now before I talk about building the magic wand,
    I do want to note that you don't actually have to build the magic wand.
    You can use the Arduino inside of the shield,
    or on a breadboard attached to a camera, or however you like it just as is.
    And that will totally work for all the gestures
    you need to do for the magic wand.
    But we think it's a little fun to build a magic wand.
    And we, the core staff, has found that the default Magic Wand
    example tends to work a little bit better
    if you have your Arduino on the end of the magic wand.
    So in that spirit, I have three whiteboard markers here.
    And my magic wand was just going to be to connect the three
    whiteboard markers together, and just use that as a little bit of a wand.
    And so just a little quick tip.
    As you're taking the Arduino in and out of the shield or breadboard or what
    have you, the best way to do it is, again, just like when we put it in,
    to grab it by the outside and sort of just
    gently rock it back and forth to slowly ease it out of there.
    Just being careful that you don't bend any of the pins.
    And then in terms of attaching your Arduino to a magic wand,
    just one thing really to note is that we want
    to have the BLE module up here facing out along the end,
    and the port for the USB coming down along
    so you can run the cable closer to you.
    And then you can take something like just some simple Scotch tape
    to attach it.
    And that's sort of really it.
    And two good places to that I found nice to
    tape it on as, one, around that BLE module.
    Because that won't get in the way or touch any of the sensors in the middle
    here, like the microphone and whatnot.
    And then another great place to put a second piece of tape
    is around the USB cable.
    Because it doesn't really matter if we get a bunch of tape on the USB cable.
    And you know, just two simple pieces of tape
    should probably be enough to hold it pretty steady on there.
    And that would be all you need to have a nice little magic wand.
    Now, in terms of deploying the magic wand code--
    Well, first of all, we now need to reconnect our Arduino here
    to the computer.
    But all we're going to do is we're going to open up the Magic Wand
    example, which just like all the other examples,
    we can find in the hardware TinyMLx library set of examples.
    So we're going to open up Magic Wand.
    And then this pre-built example is ready to go.
    So we're just going to check the board and the port.
    And then we're just going to upload it to our Arduino.
    And once it uploads, we can talk about how to explore gestures
    with the default example.
    And then in the next reading, we'll collect some of our own gesture data.
    And then we can train and deploy our own custom magic wand.
    All right.
    Now we can explore the Magic Wand example using the Serial Monitor.
    In the Serial Monitor, what we're going to see
    is both an ASCII arc plot of the gesture as you move the wand around,
    as well as a prediction of what gesture you drew,
    as well as a confident score ranging from negative 128
    to positive 127 of how likely that gesture was to be the gesture that you
    thought it was.
    So for example, if I just take this wand and I bring it straight up,
    then that's pretty close to a perfect 1.
    And as you can see, you get this drawing of a mostly straight line.
    I guess I erred a little bit in the beginning and the end.
    But it did think that's a 1, with a score of 112--
    quite likely.
    If I kind of just scribble around something quite random,
    we'll get this blob here.
    And it says, well, maybe that's a 3 because it kind of looks
    like it has sort of two humps to it.
    But my score is a zero.
    I'm not I'm not very confident that that's a 3.
    Again, if I just lower it back down with a straight 1, then we got a 1.
    Although, I guess I erred a bit more.
    Now it's only 88 on its confidence.
    But play around with it.
    See which digits you can draw, which ones you think are easy,
    which ones you think are hard.
    And hopefully, you can start to explore the magic wand.
    And then I'll see you in the next reading, where
    we will collect a custom data set that we can use to then train our own custom
    magic wand project.
    Ende des Transkriptes, zum Anfang springen
    Downloads und Transskripte
    Video
    Videodatei herunterladen
    Transskripte
    Download SubRip (.srt) file
    Download Text (.txt) file

\subsection{Building the Wand}

    Note that you do not have to build the full magic wand and can simply hold the Arduino (with or without the shield) in your hand. However, you will likely find that when you deploy the example with Pete’s pre-trained digit model, since he trained it using a ‘wand,’ it will be more accurate if you use a ‘wand’ too.
    
    If you’d like to build a ‘wand,’ the 'wand' itself can be as simple as a stick. It doesn't need to do anything other than keep the board at its end as you hold the other end and wave it about. For fun, the course staff members have found cheap wands from online retailers and used them, but a simple piece of wood, a ruler, or even a thick piece of cardboard should work just as well.
    
    You should place the board at the end of the wand, with the USB socket facing downwards, towards where you hold it, so that the cable can run down the handle. The sketch is designed to compensate for any rotation of the board around the wand's shaft, so as long as it's parallel to the wand's length, the board's twist won't matter. Use some sticky tape or some other easy-to-remove method to attach the board to the wand, and hold the cable in place along the shaft. The end result should look something like the image below.
    
    \includegraphics[width=\textwidth]{TinyML/3 Development/4-8.1.png} 
    
    The board attached to a “wand.” The board is taped around the USB cable at the bottom of the board and around the BLE module at the top of the board to a wooden stick.
    
    If an ASCII-art diagram is more helpful, here's what you should aim for:
    
    -----
    |      | <- Arduino board
    |      |
    |  ()  |   <- Reset button
    |      |
    -TT-      <- USB port
    ||
    ||<- Wand
    ....
    ||
    ||
    ()
    
\subsection{Deploying the Magic Wand Project}

\begin{enumerate}
    \item Use a USB cable to connect the Arduino Nano 33 BLE Sense to your machine. You should see the green LED power indicator come on when the board first receives power.
    
    \item Open the \FILE{magic\_wand.ino} sketch, which you can find via the File drop-down menu. Navigate, as follows: File → Examples → Harvard\_TinyMLx → magic\_wand.
    
    \item As always, use the Tools drop-down menu to select appropriate Port and Board.
      \begin{itemize}
        \item Select the Arduino Nano 33 BLE as the board by going to Tools → Board: <Current Board Name> → Arduino Mbed OS Boards (nRF52840) → Arduino Nano 33 BLE. Note that on different operating systems the exact name of the board may vary but/and it should include the word Nano at a minimum. If you do not see that as an option, then please go back to Setting up the Software and make sure you have installed the necessary board files.

        \item Then select the USB Port associated with your board. This will appear differently on Windows, macOS, Linux but will likely indicate ‘Arduino Nano 33 BLE” in parenthesis. You can select this by going to Tools → Port: <Current Port (Board on Port)> → <TBD Based on OS> (Arduino Nano 33 BLE). Where <TBD Based on OS> is most likely to come from the list below where <\#> indicates some integer number:
          \begin{itemize}
            \item Windows → COM<\#>
            \item macOS → /dev/cu.usbmodem<\#>
            \item Linux → ttyUSB<\#> or ttyACM<\#>
           \end{itemize}
      \end{itemize}

  \item Use the rightward arrow to upload / flash the code. Do not be alarmed if you see a series of orange warnings appear in the console. This is expected as we are working with bleeding edge code. You’ll know the upload is complete when you see red text in the console at the bottom of the IDE that shows 100\% upload of the code and a statement that says something like “Done in <\#.\#> seconds.”
    
  If you receive an error, you will see an orange error bar appear and a red error message in the console. Don’t worry -- there are many common reasons this may have occurred.
    
    To help you debug other issues, please check out our FAQ appendix with answers to the most common errors!
    
  \item Now open the serial monitor. The default magic wand will recognize the digits 0-9 based on data recorded by Pete, so your accuracy may vary depending on how close your “gesture handwriting” is to his! The serial monitor will output first ASCII art showing the gesture you just performed and below it will be the best match label as well as a confidence score between -128 and 127. The confidence score indicates how strongly the model believes you performed the gesture. Do note that every time you move the board and then stop, a new gesture will be processed, so don’t be surprised to get some odd results as you move the board to prepare for a gesture.
  
  \includegraphics[width=\textwidth]{TinyML/3 Development/pasted_image_0__1_.png} 


\end{enumerate}
  
  
\section{Collecting Data for Your Custom Magic Wand Project}

 In this document, we are going to collect custom gestures, which we can then later use to train a custom magic wand project.
    
    NOTE: Our URL has changed from the video! See below for details!
    
  \subsection{Collecting Data for Your Custom Magic Wand Project}
  
    0:00 / 0:00Klicken Sie auf die Pfeiltaste nach oben, um das Geschwindigkeitsmenu aufzurufen. Regeln Sie dann mit den Hoch und Runter Pfeiltasten die Geschwindigkeit. Klicken Sie Enter, um die jeweilige Geschwindigkeit festzulegen.Geschwindigkeit1.0xKlicke auf diesen Knopf und diess Video stumm oder nicht stumm zu schalten, oder drücke die HOCH oder RUNTER Knöpfe um die Lautstärke zu erhöhen oder abzusenken. Maximal Lautstärke.Videotranskript
    Beginn des Transkriptes, zum Ende springen.
    BRIAN PLANCHER: All right, now that we've
    played with the default magic wand project, it's time
    to record some of our own gestures so we can train our own custom model.
    To record some gestures, you're going to want to go to HTTPS,
    and again the S is important or this won't work, tinymlx.org/magic\_wand.
    And you'll get to a page that looks something like this.
    And here the first thing you're going to want to do
    is make sure that your magic wand deployed example code is running
    on an Arduino, and it is powered.
    We're then going to connect to it over Bluetooth
    by clicking the Bluetooth button.
    You're going to see an ask to pair with a BLE
    sense and some sort of number, which is a serial number or type ID of your BLE.
    So just click pair.
    And you should see connecting to device, and then connected.
    And then now, the Bluetooth is green.
    And you're going to see that things are starting to appear on the right here.
    And essentially, we're doing the same thing that was happening before
    on the serial monitor, except for now it's just taking the raw drawing
    and saving them, and not actually running it through the model,
    but allowing us to generate a data set.
    So Every time we move, you're going to see a drawing
    on sort of the-- on this left box here.
    And then once we stop, it's going to become
    some sort of drawing on the right.
    If you do not want any of these gestures,
    you can just click this little trashcan up
    in the top right of each gesture to delete them.
    And then you can also--
    let me make a gesture.
    You can also label-- you'll also need to label each gesture,
    here where that question mark is, you'll need
    to give each gesture that you want a name.
    So let's get rid of this one.
    For example, let's say I wanted to record some Z's.
    I could draw out a Z, and that would be a nice gesture on the right there.
    Now, this brings up the question of, actually,
    how many gestures do you need to record?
    And the answer is that, again, just sort of like with keyword spotting, more
    is generally better and more robust.
    But we found that if you record about 20 or so gestures of, say, two
    to three or four different kinds of gestures,
    and then you train a model with that data,
    it will actually work quite well for you holding the wand the same way, et
    cetera.
    But you can actually get something that'll work quite decently.
    So that that's a good place to start.
    So let's say I wanted to record some more.
    Now, one thing you're going to have to keep in mind
    is that every time you move it, it's going to record a gesture.
    So you may end up having to delete a lot as you try to go,
    because every time you want to go--
    reset before I draw another Z, I'm going to get something.
    Now, one thing you could decide is, maybe one gesture
    you want to add to your data set is sort of like a diagonal line
    or something like that.
    And then, suddenly, you now wouldn't have to delete these.
    But remember again, before you're done, you're
    going to have to go through and clear out the ones you don't want.
    So maybe I don't want these diagonal lines,
    I'd have to make sure to go through and delete all of them.
    And you can see up here the count is telling you
    how many gestures you've recorded, and it updates after you delete.
    And I do need to go through and label these all as the letter Z here.
    Because this is labeling your data set for training later.
    And once you're done collecting all your data,
    again, you probably want more than four.
    And the good news here is we've found that you can record 40, 60, 80,
    100 different gestures into this system, and it should just
    download as one file.
    But also, we've set it up so that if you have multiple files in the training
    script, you can just upload multiple different gesture recording
    sessions, files at the same time, and just it'll
    automatically go through all of them and pull in all of the relevant gestures.
    But again, when you're ready to download,
    just click the Download Data button.
    And it should download a WandData.JSON file.
    And that's the file you're going to use in the training
    script in the next reading to train a model
    and then deploy it onto your Arduino.
    Ende des Transkriptes, zum Anfang springen
    Downloads und Transskripte
    Video
    Videodatei herunterladen
    Transskripte
    Download SubRip (.srt) file
    Download Text (.txt) file
    
\begin{enumerate}
  \item Make sure the default Magic Wand project is deployed as described in the previous reading!
  
  \item Open up your browser and navigate to: \url{https://tinyml.seas.harvard.edu/magic_wand}. NOTE: Our URL has changed from the video!
    
    \begin{itemize}
      \item If you see a warning that says: Error: This browser doesn't support Web Bluetooth. Try using Chrome. If you are not using Chrome, we suggest you switch to Chrome. If you are already using Chrome, then navigate to the following and make sure to “enable” the “experimental web platform features.” We have found this step to be necessary on older versions of Chrome and users running the Linux operating system. \url{chrome://flags/#enable-experimental-web-platform-features}.
    
      \includegraphics[width=\textwidth]{TinyML/3 Development/4-8.2.png} 
    
      A screenshot of the webpage you land on after navigating to “\url{chrome://flags/#enable-experimental-web-platform-features}” which has “Experimental Web Platform Features” highlighted in yellow and some explanatory text below. Importantly, on the right is a blue drop down which says “Enabled” or “Disabled” which needs to be set to “Enabled”.

      If this problem persists, make sure you are navigating to https://. For some reason, if you do not use the secure protocol, then the web Bluetooth will not work.
    
      \item You should then arrive at a webpage that looks something like the following:
    
      \includegraphics[width=\textwidth]{TinyML/3 Development/4-8.3.png} 
    
       screenshot of the magic wand data collection webpage. It shows the text instructions that state: “To get started recording magic wand gestures: 1) Upload the Magic Wand Capture sketch to an Arduino Nano 33 BLE Sense board, 2) Connect to the board using the Bluetooth button below. 3) Wave the wand to make gestures. They'll be recorded and displayed on the right. 4) Review the gestures, add labels by clicking on the '?', and remove mistakes. 5) Download the gestures as a JSON data file, ready for model training.” Below that are two buttons that say “Download Data” and “Bluetooth”. Both should be blue in color. Next to the Bluetooth button it should say “Click the bluetooth button to connect to your board.” Please defer to the text for changes that have been made to the url and suggestions of how to use the webpage as it has evolved over time.
     \end{itemize}
  \item You’ll then need to connect your device over Bluetooth. To do that, simply click the blue Bluetooth button and a pop-up will appear asking to pair. Select your BLE Sense, which should be called something like “BLESense-XXXX” where XXXX will vary, and click the pair button. Do note that the course staff has found that sometimes you have to repeat this step twice. Once you are connected, the Bluetooth button will turn green.
    
      \includegraphics[width=\textwidth]{TinyML/3 Development/4-8a.png} 
    
    Another screenshot of the chrome browser navigating to the magic wand data collection application. It has a popup that says “tinymlx.org wants to pair” and shows below it the option “BLESEnse-6E00-Paired” at the bottom of the popup it has a cancel and pair button. Note that the “-Paired” indicates that something is paired although we have found that you may need to pair a device twice to get it to work. Please defer to the text for changes that have been made to the url and suggestions of how to use the webpage as it has evolved over time.
    
  \item Once your Arduino is paired, it’s time to record some gestures. You’ll notice that every time you move the Arduino around and then stop, a new gesture is recorded. This is because the gestures are automatically split up by times when the wand is kept still. These pauses act like spaces between words, and so when you've finished, a gesture you should stop moving the wand, so that it ends cleanly. Note that the direction of the gesture matters (e.g., a clockwise circle is different from a counterclockwise circle)!
    
    These gestures you are drawing will start to show up in the list on the right side of the screen. You can look at the shapes shown there to understand whether the gestures came out cleanly. A good rule of thumb is that if you can't tell what the gesture is by looking at it, then a model will have a hard time recognizing it too.
    
    If you want to delete a recording, simply click the trashcan icon at the top right of each gesture recording (You may need to delete a lot of spurious recordings that were made as you moved the Arduino into position between each gesture).
    
    Also, make sure to label all of your gestures for training. You can label each recording by clicking on the question mark at the top left of each gesture and typing in your label. For example, the screenshot below shows the label “Z” added to the gesture of the letter “Z” recorded by the course staff. Make sure you label all of your gestures and note that case matters (i.e., "z" and "Z" are different)! Some students have noted that it can be easier to create one gesture per recording session in order to make sure all of the gestures are labeled identically but/and the staff has also gotten it to work with multiple gestures per recording session.
    
    The staff has found that collecting ~20 examples each of 2-3 different gestures often will be enough data to successfully train a moderately decent magic wand application (i.e., it will work often for you may not generalize to other users). To help you keep track of how many gestures you recorded, there is a number in the top right of the screen (e.g., the number 1 as shown in the image below). The staff has also found that gestures like a circle (O) or the (Z) for Zoro tend to work quite well!
    
    Finally, you can upload multiple JSON files to the training script so don’t feel pressured to do all of your gesture recordings in one shot!
    
      \includegraphics[width=\textwidth]{TinyML/3 Development/4-8.4.png} 
    
    Screenshot of the interface showing the letter Z gesture labeled as the letter Z and the number 1 denoting 1 recording so far in the upper right of the UI.
    
    \item When you are done collecting all of your data, simply click the blue “Download Data” button and a JSON file with all of the gestures will be automatically downloaded! We’ll use that file in the Colab in the next section to train a custom model! Be careful, when you leave or refresh the web page, your recorded gestures will be lost, so make sure you use the "Download Data" link to save them!
\end{enumerate}
    
  Note: If you encounter an error where the number of gestures does not match the number of labels, check to ensure the data collection tool correctly labeled your gestures. Open the JSON file and check for “label”:“” and replace all of the mislabeled gestures.        
  
\section{Training and Deploying Your Custom Magic Wand Project}

In this document, we are going to train and then deploy a custom magic wand model based on the custom gestures we just collected.
    
\subsection{Training and Deploying Your Custom Magic Wand Project}

  0:00 / 0:00Klicken Sie auf die Pfeiltaste nach oben, um das Geschwindigkeitsmenu aufzurufen. Regeln Sie dann mit den Hoch und Runter Pfeiltasten die Geschwindigkeit. Klicken Sie Enter, um die jeweilige Geschwindigkeit festzulegen.Geschwindigkeit1.0xKlicke auf diesen Knopf und diess Video stumm oder nicht stumm zu schalten, oder drücke die HOCH oder RUNTER Knöpfe um die Lautstärke zu erhöhen oder abzusenken.Maximal Lautstärke.Videotranskript
    Beginn des Transkriptes, zum Ende springen.
    BRIAN PLANCHER: In this reading, we're going
    to train and deploy a custom magic wand project
    based on the data you just collected.
    In the Colab below, you'll be able to upload
    all of the JSON files of gestures you just recorded
    and use them to train and model that we can then deploy onto the Arduino.
    In the remainder of this video, I'm going
    to show you the changes you need to make to the Arduino code
    in order to successfully deploy that model.
    So let's assume you've finished that Colab,
    and you've got your fully trained model dot cc file.
    What do we need to do next?
    Well, we need to do two things.
    The first thing we need to do is in this main magic wand tab
    of the magic wand sketch.
    We need to scroll down to about line 55 right around here.
    And we need to update these two variables with, first,
    the number of gestures and the names of the gestures.
    And the order of the gestures matters, and so you're
    going to want to pay attention to that one step in the Colab
    where it says in bold make, a note of this
    ordering as you will need to enter the labels in order in the Arduino code.
    And so you see the output here.
    I have zero, one, and two and that's the order
    that I need to enter in those labels.
    You might also note here that I've used a special character, a backslash.
    And because we just treat these like strings,
    it actually just did a nice little escape there,
    you'll see it's actually a backslash in the task, but it's doing two of them.
    And that was what I used to reset because I actually recorded a z,
    and then I did a reset as backslash, and then I did some circles.
    Anyway, those are the three labels that I need, and that's the order.
    And so what I need to do is I need to take these and put them
    into the Arduino code.
    And so what I'll do is I'll update the count to three,
    and I'll update the labels to slash o and z.
    And then I will delete the rest of these.
    When we hit Save, we're going to get that same warning about read-only,
    so let's go and let's save down a copy.
    And let's save it in that TinyML folder, and let's call this \FILE{magic\_wand\_custom}.
    All right.
    So now we have our custom labels and the custom count of the number of labels
    that we need.
    So the only thing we now have to change is we need to go
    to the \FILE{magic\_wand\_model\_data.cpp}.
    That's where all the model data is.
    So we could go back to Chrome, and we could get that model file,
    and we can just, again, copy the model file.
    And remember from keyword spotting, that we only want to get the binary data,
    we don't want to get any of the surrounding syntax.
    So I don't want that semicolon there, so we can copy that.
    And we can go back to the Arduino IDE, and we can paste that in there.
    And then we can also go and get the model length, which
    is most likely different in this case.
    Save that, and that's it.
    That's all you need to do.
    We just need to update those labels, the count, and the actual model file,
    and then we're ready to deploy it.
    So we can just make sure the board and the port are set,
    and then we can click Upload to compile and upload that to our magic wand.
    And then you can start playing around with your own magic wand
    and see how the model works.
    And if you want, you can go back and collect more data,
    try training it for longer, change the model a little bit,
    do whatever you like.
    So hopefully, you have fun playing with your magic wand.
    And that's it for me.
    This is the last lab, so I hope you enjoyed
    all of the lab parts of this course.
    Ende des Transkriptes, zum Anfang springen
    Downloads und Transskripte
    Video
    Videodatei herunterladen
    Transskripte
    Download SubRip (.srt) file
    Download Text (.txt) file
    
\subsection{Training the Magic Wand Model}

    The first thing you’ll need to do is to upload your gesture dataset into Colab and train a new magic wand model. Then in that Colab, we’ll need to convert that model first into a quantized .tflite file and then into a .cc file for use with the Arduino IDE.
    
    As with the KWS examples, we will be using the resulting .cc file, so make sure to download it or leave the tab open with the printout!
    
    https://colab.research.google.com/github/tinyMLx/colabs/blob/master/4-8-11-CustomMagicWand.ipynb
    
\subsection{Deploying the Trained Model}

\begin{enumerate}
  \item Use a USB cable to connect the Arduino Nano 33 BLE Sense to your machine. You should see the green LED power indicator come on when the board first receives power.
  
  \item Open the \FILE{magic\_wand.ino }sketch, which you can find via the File drop-down menu. Navigate, as follows: File → Examples → Harvard\_TinyMLx → magic\_wand.
  
  \item You’ll then need to make two changes to the \FILE{magic\_wand.ino} file to alert it of your number of gestures and gesture labels. These changes occur on lines 55-59, which currently read as:
  
  \begin{lstlisting}
    // ------------------------------------------------------------------------------- //
    // UPDATE THESE VARIABLES TO MATCH THE NUMBER AND LIST OF GESTURES IN YOUR DATASET //
    // ------------------------------------------------------------------------------- //
    constexpr int label_count = 10;
    const char * labels[label_count] = {"0","1","2","3","4","5","6","7","8","9"};
  \end{lstlisting}  

  \begin{itemize}
      \item Update the \PYTHON{label\_count} to reflect the number of gestures in your dataset.
      \item Update the list of labels to reflect the gestures in your dataset. Note: the order matters! Make sure it matches the alphanumeric order as printed out in the training script!
  \end{itemize}    
      
 \item You'll also need to update the model data as we've done with the previous examples. The model data can be found in the magic\_wand\_model\_data.cpp file. As we've done in the past, make sure to only update the binary values and leave the cpp syntax constant.
 
 \item When you save, just like with the KWS examples, you will be asked to save a copy. Again, we suggest that you make a folder called e.g., TinyML inside of your Arduino folder. You can find your main Arduino folder either inside of your Documents folder or in your Home folder, and save it in that folder with a descriptive name like magic\_wand\_custom. That said, you can save it wherever you like with whatever name you want!

  \item  As always, use the Tools drop-down menu to select appropriate Port and Board.

    \begin{itemize}    
      \item Select the Arduino Nano 33 BLE as the board by going to Tools → Board: <Current Board Name> → Arduino Mbed OS Boards (nRF52840) → Arduino Nano 33 BLE. Note that on different operating systems the exact name of the board may vary but/and it should include the word Nano at a minimum. If you do not see that as an option, then please go back to Setting up the Software and make sure you have installed the necessary board files.

      \item Then. select the USB Port associated with your board. This will appear differently on Windows, macOS, Linux but will likely indicate ‘Arduino Nano 33 BLE” in parenthesis. You can select this by going to Tools → Port: <Current Port (Board on Port)> → <TBD Based on OS> (Arduino Nano 33 BLE). Where <TBD Based on OS> is most likely to come from the list below where <\#> indicates some integer number:
        \begin{itemize}    
          \item Windows → COM<\#>
          \item macOS → /dev/cu.usbmodem<\#>
          \item Linux → ttyUSB<\#> or ttyACM<\#>
        \end{itemize}    
     \end{itemize}    
  \item   Use the rightward arrow to upload / flash the code. Do not be alarmed if you see a series of orange warnings appear in the console. This is expected as we are working with bleeding edge code. You’ll know the upload is complete when you see red text in the console at the bottom of the IDE that shows 100\% upload of the code and a statement that says something like “Done in <\#.\#> seconds.”

    If you receive an error, you will see an orange error bar appear and a red error message in the console. Don’t worry -- there are many common reasons this may have occurred.
    
    To help you debug other issues, please check out our FAQ appendix with answers to the most common errors!
    
  \item     Now, open the serial monitor and test out your custom model. As a reminder, the serial monitor will output first ASCII art showing the gesture you just performed and below it will be the best match label as well as a confidence score between -128 and 127. The confidence score indicates how strongly the model believes you performed the gesture. Do note that every time you move the board and then stop, a new gesture will be processed, so don’t be surprised to get some odd results as you move the board to prepare for a gesture.

      \includegraphics[width=\textwidth]{TinyML/3 Development/4-8.5.png} 
    
    A screenshot of the serial monitor showing ASCII art of a gesture and the prediction and confidence interval below it as described in the reading. There are a series of ”.“ that fill the screen and wherever the wand went there is a “\#”. This visually shows the pattern the wand made. In this case a wobbly 0. On the bottom we see the result “Found 0 (126)”.  
\end{enumerate}
    
    
\section{How Well Does Your Custom Magic Wand Model Work?}

How did your Custom Dataset Magic Wand Model work? Does it work for all gestures equally well? Does it work for any of your family and friends? How much data did you collect? Did you collect data from anyone other than yourself? If you made multiple models (or re-trained with additional data), did that improve performance? How does your experience compare to others in the forum?
    
    
% 1.9: Responsible AI Deployment

\section{Privacy}

0:00 / 0:00Klicken Sie auf die Pfeiltaste nach oben, um das Geschwindigkeitsmenu aufzurufen. Regeln Sie dann mit den Hoch und Runter Pfeiltasten die Geschwindigkeit. Klicken Sie Enter, um die jeweilige Geschwindigkeit festzulegen.Geschwindigkeit1.0xKlicke auf diesen Knopf und diess Video stumm oder nicht stumm zu schalten, oder drücke die HOCH oder RUNTER Knöpfe um die Lautstärke zu erhöhen oder abzusenken.Maximal Lautstärke.Videotranskript
Beginn des Transkriptes, zum Ende springen.
SUSAN KENNEDY: Revisiting our human-centered design framework
for building responsible AI.
We're now going to focus on the ethical challenges that
arise during deployment.
The specific questions we'll think about are why is privacy valuable
and what should we do to protect user's privacy?
Now, when it comes to privacy, there's actually
a huge benefit to working with TinyML.
On-device machine learning allows you to process sensitive data locally.
And by harnessing the power of edge computing, where
data is being processed at the source, is going to decrease latency,
but more importantly, it minimizes the transmission
of data up to the cloud, which would otherwise create room
for potential privacy violations.
So with edge computing, there's just going
to be less opportunities for malicious actors
to intercept data that's being transmitted.
And since the data is being primarily stored on the device,
it's more secure than if it were warehoused
in a singular location like the cloud.
These features of TinyML are definitely a step in the right direction.
But there's still a long way to go when it
comes to adequately protecting privacy.
Why?
Well, TinyML is going to be contributing to the exploding
number of IoT devices, which means that we're just going
to be collecting more and more data.
And this matters because even seemingly innocuous data,
like the time a coffee machine is turned on,
can reveal deeply intimate aspects of a person's life
when integrated and aggregated with other information.
Now, there's a lot of different steps you can take to protect privacy.
But before jumping into solutions, you'll
first want to understand what the problem is that you're trying to solve.
So to that end, we're going to turn our attention to what moral reasons there
are for data protection and what counts as a violation of privacy.
The first reason to value privacy is to prevent information-based harms.
For example, if a hacker were to gain inappropriate access
to someone's passwords, characteristics, whereabouts, or other information,
they could potentially use this information to cause harm.
And while TinyML reduces opportunities for hackers to gain access,
it's still a possibility.
And in cases like this, privacy is closely
connected to security concerns, which we'll discuss in an upcoming video.
But why else might we be concerned about privacy, aside from malicious actors?
Well, because informational injustice and discrimination
can result from an inappropriate transfer of personal data
across the boundaries of a context.
For example, we think if a medical professional asks
about your health status, that's fine.
But if a potential employer had access to this information,
you might be discriminated against or disadvantaged in some way.
So you want to give careful consideration
to things like the context your device will be deployed in,
what type of information you'll be collecting,
and who has access to this information.
A third reason to value privacy is to preserve
moral autonomy and human dignity.
In essence, autonomy refers to the idea of being self-governing and living life
according to one's own reasons, considerations, and desires.
When privacy is preserved, people are able to control
how they present themselves to others.
But when there's a lack of privacy, individuals
might be exposed to outside forces that influence their choices
and bring them to make decisions that they wouldn't have otherwise made.
For example, if a TinyML device captures information about a user's driving
habits, they might feel pressured to change their habits depending
on who has access to this information.
So in order to preserve autonomy and human dignity,
it's important to obtain informed consent from users.
This way, they have a say in what data is being collected
and how it's being used.
So what should we do to protect privacy?
The first thing is minimize the data that's being transmitted and stored.
Don't collect what you don't need and set up a timeline for deleting data
periodically.
After all, you won't have to worry about others gaining inappropriate access
to data that you don't have.
Another reason it's important to minimize data handling wherever
possible is that government pressure can sometimes
force you to hand over information.
And while there might be some cases where this is for the greater good,
there's a possibility that governments will
abuse this power, which could lead to serious harms to your users.
The second thing to do is protect what information you do have.
Even if you remove personal identifiers like name and address,
the safest thing to do is treat all data as if it's personal data
and protect it with encryption techniques.
But since privacy's about more than just preventing hackers from gaining access,
you want to map out clearly how data and personal information flows.
Be specific about what you're collecting, who will have access,
and whether or not the information is crossing boundaries from one context
into another.
This will help you spot potential concerns relating
to informational injustice.
And lastly, you want to obtain informed consent, which
requires being transparent with users about how their data is being collected
and used so that they can make an informed
decision about their involvement.
And you can revisit course two for more details about what's
required for informed consent.
If you've done a good job of minimizing, you
should be able to effectively communicate to your users
why your product's benefits are worth the privacy implications.
Since each device will raise a unique set of privacy concerns,
there isn't a standard privacy checklist that will apply across all cases.
So you'll want to record and document potential concerns
for your particular product and then spend
some time thinking about what privacy preserving measures you should take.
Ende des Transkriptes, zum Anfang springen
Downloads und Transskripte
Video
Videodatei herunterladen
Transskripte
Download SubRip (.srt) file
Download Text (.txt) file

    
\section{Security}

    0:00 / 0:00Klicken Sie auf die Pfeiltaste nach oben, um das Geschwindigkeitsmenu aufzurufen. Regeln Sie dann mit den Hoch und Runter Pfeiltasten die Geschwindigkeit. Klicken Sie Enter, um die jeweilige Geschwindigkeit festzulegen.Geschwindigkeit1.0xKlicke auf diesen Knopf und diess Video stumm oder nicht stumm zu schalten, oder drücke die HOCH oder RUNTER Knöpfe um die Lautstärke zu erhöhen oder abzusenken.Maximal Lautstärke.Videotranskript
    Beginn des Transkriptes, zum Ende springen.
    SUSAN KENNEDY: Returning to our human-centered design
    framework for responsible AI, we're now going
    to focus on the ethical issues surrounding security.
    The specific questions we'll think about are who values security and why,
    and how should we address security concerns?
    Diving in, we'll first discuss different types of security threats
    alongside some real-world examples that made headlines in recent years.
    The first type of security threat are data leaks,
    where there's an unauthorized transmission or exposure of data.
    In most cases, this happens unintentionally,
    although data leaks can include cases where
    a whistleblower from inside the company purposefully
    leaks sensitive information.
    For example, an accidental data leak occurred in 2018
    when the popular fitness app Strava leaked sensitive information
    about the location of military bases, owing to heat maps that
    show patterns and fitness behavior.
    Another closely related type of security threat
    are data breaches, where personal information is intentionally
    exposed or accessed by hackers.
    For example, in 2018, hackers abused Alexa and Google Home smart assistance
    to eavesdrop on users without their knowledge
    and even worse, trick some users into handing over sensitive information.
    A third kind of security threat are adversarial attacks.
    One kind of attack is DDoS, or Distributed Denial
    of Service, where cyber criminals aim to overload a system that
    leads to a disruption of service.
    For example, in 2016, cyber criminals shut down the heating of two buildings
    in Finland.
    In this case, the attack left the remote system
    in an endless cycle of rebooting in an attempt to reconnect.
    This left residents with no heating, a serious concern,
    given that temperatures in Finland could be well below freezing
    at that time of year.
    Another kind of attack relies on exploiting vulnerabilities.
    For example, in 2019, researchers found six vulnerabilities
    in the smart deadbolt device from Hickory Hardware, three of which
    could be exploited to remotely unlock doors and break into homes.
    Of course, exploiting vulnerabilities is a broad category.
    So to see what this might mean for a TinyML application
    like keyword spotting, for instance, let's consider
    two types of attacks with very different outcomes.
    In the first kind of attack, noise is introduced
    to fool the machine, so to speak, with the result being
    that the model doesn't recognize the wake word
    and fails to activate the device.
    But now let's look at the DolphinAttack case,
    where researchers at Zhejiang University published a paper showing
    that speech commands that are inaudible to humans
    can still succeed in triggering a wake word on an ML device.
    This case is especially concerning because the attack can go undetected,
    and it leaves the user susceptible to various harms.
    Given the wide range of potential security threats,
    it should be easy to see why security is such a hard problem.
    The truth is, building 100% secure devices is near impossible.
    And this is because breaking into a system
    is often easier than protecting it.
    The burden lies on developers to protect against all present and future threats,
    whereas attackers only need to find one point of weakness.
    And predicting security threats is even more difficult
    when you consider how attackers are constantly
    coming up with new strategies.
    So in response, you might think that we should just
    postpone the deployment of vulnerable ML devices
    until they're made adequately resilient.
    But given the potential of these devices to improve our lives,
    the strategy might do more harm than good.
    So already, we're confronted with a matter of ethics.
    Do the benefits of deploying the device outweigh the potential harms
    from security risks?
    When evaluating the potential types of security threats you might face,
    you also want to identify who is implicated
    and why to figure out the best way to handle security issues.
    Because as we'll see, there can be competing reasons
    that pull us in different directions.
    So who might value security?
    Well, the most obvious is individual users.
    A close second is the service provider.
    If your device happens to rely on transmitting data to the cloud,
    then perhaps a third party cloud service provider.
    And finally, governments and state-level agents.
    And why is security considered ethically valuable?
    Well, one reason is to preserve privacy for all the reasons
    that we discussed in the previous video.
    Another reason is to prevent harms to users.
    And finally, to protect citizens and communities.
    Now, it's important to map out the different stakeholders and reasons
    to value security because these might come into conflict.
    So let's suppose we're trying to address security issues in data storage.
    One option is to delete all the data after a certain time frame.
    But another option might be to encrypt it and let some trusted
    party in the company have access.
    If we care about the privacy for users, deleting the data
    might be the best option.
    But if we consider the government's interest
    in protecting citizens and communities, being
    able to share access to encrypted data might be the best option.
    So deciding what security measures to implement
    might require some critical thinking about how to weigh these competing
    considerations.
    So what can we do to improve security on TinyML devices?
    The first thing we can try and do is limit opportunities for attackers
    by minimizing your hardware design.
    If you don't need Wi-Fi or particular sensors,
    build a device that doesn't have these.
    And if you recall from a discussion forum in course 1,
    users were surprised when they discovered microphones
    on their Nest devices.
    While the microphones might have allowed the development of future capabilities,
    it also opened up unnecessary risks.
    So try not to do this.
    In some cases, having multiple sensors could work to your advantage.
    For example, earlier in this course, you learned about sensor fusion
    in the case of AV speech recognition.
    This could make your device more resilient because attackers
    would need to find a way to fool both the audio and visual inputs.
    A third way to improve security is encryption techniques.
    So even if hackers gain access, encrypted data
    will preserve privacy of users by hiding certain details of the information.
    But you still need to make some tough decisions
    about who's going to have access to the encryption key.
    And finally, the last way to improve security
    is to map the stakeholders who have an interest in security
    and identify conflicts and values.
    This will help you navigate different options
    and figure out what security measures are most important to implement
    for your particular product.
    Ende des Transkriptes, zum Anfang springen
    Downloads und Transskripte
    Video
    Videodatei herunterladen
    Transskripte
    Download SubRip (.srt) file
    Download Text (.txt) file
    
\section{Attacking a KWS Model}

Now that we’ve talked about the need for security considerations with TinyML applications it's time for you to get some hands-on experience attacking a KWS model! Open up the Colab below to maliciously trigger the “Yes, No” KWS model with some seemingly innocent static noise!

\url{https://colab.research.google.com/github/tinyMLx/colabs/blob/master/4-9-6-Attacking-KWSModel.ipynb}

In the Colab, just before this, you gained some hands-on experience attacking a KWS model where you can trigger a “yes” (false positive) with static noise.

Now that you have identified a security vulnerability for the KWS model.

\begin{itemize}
  \item What kinds of applications can you imagine where exploiting this vulnerability would be especially concerning? (Do you think the risks associated with potential false positives are serious enough to outweigh the potential benefits of the device?)
  \item What kind of applications can you imagine where this security vulnerability might be okay?
  \item Do you agree or disagree with what others have said so far in the forum?
\end{itemize}

\section{Why do ML Models Fail after Deployment?}

  Imagine you train and validate a machine learning (ML) model with high predictive accuracy and you have determined that it’s ready for deployment. Once your model is deployed in the real world, you might think that this is cause for celebration as your hard work is now complete! However, the journey isn’t over quite yet. When it comes to machine learning, you will need to continuously monitor your model’s performance in the real world to ensure that it is still performing as well as it did when you validated its accuracy in production.
    
    The reason ML models must be monitored after deployment is owing to the possibility of model drift, sometimes referred to as ‘model decay’, where the model degrades over time leading to a drop-off in predictive power. Why does model drift occur? Put simply, every model is created with the assumption that the way the world will be in the future is similar to how we interpreted the world at the time the model was developed. Of course, this assumption is not always warranted, as the real world is a dynamic environment that is constantly evolving.
    
    When a model is deployed in an environment that changes in unforeseen ways, the model’s performance is likely to change as well. In order to maintain the predictive accuracy of a model, successful deployment requires continuous monitoring over time to guard against model drift. Otherwise, the deployment of a model may lead to failure, with unhappy customers complaining that the model doesn’t work as promised.
    
    David Talby recounts his own experience with model drift in the article, “Why Machine Learning Models Crash and Burn in Production,” in hopes that others may learn from his mistakes. After he deployed an ML model in hospitals to predict 30-day readmissions, he discovered a drop-off in performance after just 2-3 months. Moreover, the particular causes of model drift varied between hospitals, “or even buildings within the same hospital.” For example, in some cases, the problem was ‘missing values’ owing to changes in the electronic health record system. In other cases, the population of patients presenting to the ER had shifted after the hospital began accepting a new type of insurance.
    
    Talby’s experience is informative as it shows how a variety of factors can contribute to model drift. At this point, we will shift our attention to learning how to distinguish between the two most common causes of model drift: data drift and concept drift.
    
\subsection{Data Drift}

    When data drift occurs, the properties of the data inputs are different from the data that was used to train the model. Similar to how Talby saw a change in the population of patients presenting to the ER, it’s easy to imagine how the distribution of data may change with respect to TinyML applications.
    
    Let’s suppose you are creating a keyword spotting application to activate a device. Since the intended context for deployment consists almost entirely of native English speakers, you acquire data from this population to train your model. Fast forward to a few months after deployment--the model’s accuracy has held steady at around 85%. Now, let’s suppose the context has suddenly changed, as the environment the device was deployed in has become a popular tourist destination. Since your model is encountering a large number of non-native English speakers with thick accents, the model’s performance has dropped to 60%. This example shows how an unforeseen change in the real world led to a change in the distribution of data inputs, subsequently causing data drift.
    
    In addition to changes in the data inputs due to externalities, upstream data changes can also contribute to data drift. This is especially relevant to TinyML which exists at the intersection of machine learning and embedded IoT devices. On-device machine learning requires data inputs from sensors, and these sensors can raise a number of problems. For example, if a sensor is broken this might result in missing values that are required for making a prediction. Alternatively, the sensor itself may still be functioning properly but nevertheless result in faulty data inputs, e.g. if a camera lens is smudged and produces blurry images. Moreover, the sensor is used on the deployed device may be different than the sensor used to acquire the training data, and this can result in measurement changes, e.g. if the training data came from a high-resolution camera but the model is deployed on a device with a low-resolution camera.
    
\subsection{Concept Drift}

    When concept drift occurs, the statistical properties of the target variable (which is what the model has been designed to predict) have changed. In other words, the actual relationship between the data inputs and outputs is different from what the model has learned. It is worth highlighting that concept drift can occur even when the distribution of incoming data remains the same and still accurately resembles the data used to train the model.
    
    In Course 2, we discussed how defining the target variable is more like an art form than a perfect science, because it is a subjective process wherein one attempts to define some concept in quantitative measures e.g. defining ‘healthy’ in terms of measurable features like heart rate, number of steps, etc. Similar to how we might worry about how bias can influence the definition of the target variable, we also need to consider how the definition of the target variable might change or evolve over time in unforeseen ways.
    
    For example, imagine we create a TinyML application to predict whether a plant is diseased by identifying features like yellow spotting and wilting leaves (see the left panel of the figure below). Now let’s suppose that advances in research have led to a better understanding of this plant disease by identifying a new feature for accurately diagnosing it. As a result, some plants that would normally be classified as ‘healthy’ are now classified as ‘diseased’. Since the concept of ‘disease’ has changed, our model’s decision boundary is no longer accurate (see the right panel of the figure below). As a result, our model will incorrectly label data that falls between the decision boundary our model has learned and where the decision boundary is in reality.
    
    \includegraphics[width=\textwidth]{TinyML/3 Development/4.9img1.png} 
    
    A pictorial representation of the following scenario. The original decision boundary between the populations shown in the left panel (a diagonal line splitting a square in half) becomes inaccurate as the true boundary shifts to the left. This leaves a space bordered by the true boundary on the left and the decision boundary on the right that is misclassified. This is shown by the original decision boundary in blue and red being shown on the left and on the right we see the new decision boundary and the original red and blue sections that are still correct and a yellow section in the middle that was misclassified.
    
\subsection{How should model drift be addressed?}

    Now that you have gained familiarity with model drift and its common causes, the next step is learning how to detect it and address it. Proceed to the next video to learn more about model drift with a special focus on how to address it in the context of TinyML.

\section{Monitoring after Deployment}

    0:00 / 0:00Klicken Sie auf die Pfeiltaste nach oben, um das Geschwindigkeitsmenu aufzurufen. Regeln Sie dann mit den Hoch und Runter Pfeiltasten die Geschwindigkeit. Klicken Sie Enter, um die jeweilige Geschwindigkeit festzulegen.Geschwindigkeit1.0xKlicke auf diesen Knopf und diess Video stumm oder nicht stumm zu schalten, oder drücke die HOCH oder RUNTER Knöpfe um die Lautstärke zu erhöhen oder abzusenken.Maximal Lautstärke.Videotranskript
    Beginn des Transkriptes, zum Ende springen.
    PROFESSOR: The final topic we'll cover for responsible AI is
    the importance of monitoring your model after deployment.
    The specific questions we'll address are how can we detect model drift
    and what should we do about it, what considerations should we
    keep in mind for TinyML applications, and where should ethical considerations
    enter the picture.
    Starting with how to detect model drift, one of the most reliable indicators
    of drift is a drop off in your model's performance
    assuming ground truth labels are available and sufficiently real time.
    If the model was 85% accurate, but drops to 60% accuracy
    after just a few months, we have a problem.
    Seems simple enough, right?
    Well, one thing that's unclear is just how much of a drop in accuracy
    is considered a problem, and this is what
    we need to know if we want to set a threshold for accuracy where
    we can get an alert for model drift when performance falls outside this range.
    The acceptable threshold will depend on the particular context and device.
    So you want to do is actually repeat the process you performed in Course 1
    where you carefully consider the consequences for the user
    when your model fails.
    For example, a wider threshold might be acceptable for a weak word where
    failure means an inability to activate a voice assistant,
    but unacceptable for a device on a smart car that's
    responsible for activating the brakes.
    It's worth noting that simply tracking your model's accuracy
    may not be enough, and so you'll want a more detailed analysis of performance
    that takes into account measures like precision and recall.
    Because, depending on the application, a shift
    in the rate of false positives or false negatives
    might be just as problematic as a drop in overall accuracy.
    For example, imagine if a smart smoke detector still
    has the same accuracy rate when deployed in the real world,
    but suddenly there's a shift in the distribution of inaccurate predictions.
    So instead of false positives, the model produces
    false negatives or missed detections.
    This would be a serious issue.
    So it's worth revisiting the ethical considerations you worked through
    for precision and recall in Course 1.
    This way you have a clear picture of what's
    important for maintaining good performance for your model.
    Similarly, you'll want to check whether there's a shift in performance
    for different populations of users.
    For example, if your model's performance remains the same for male users
    but has degraded for female users, the model drift
    might raise problems of fairness, and you can revisit these in Course 2.
    Sometimes, you may not have access to real time ground truth
    labels that let you monitor the predictive accuracy of your model.
    So there's two other avenues for detecting drift, albeit less reliable.
    The first is that you can keep an eye out
    for a shift in the distribution of predictions.
    So suppose you have a model to classify images as either class A and class B,
    and during testing the predictions were a pretty even split.
    But now that your model's in the real world,
    it's classifying almost all images as class B. This might
    be a potential indicator of drift, but it's also
    possible to see a shift in predictions with an accurately
    modeled decision boundary as it might actually just
    be the case that all images belong in class B.
    Another option for detecting drift requires
    monitoring the changes between your baseline data
    set, which is what you used for training,
    and the target data set, which is the incoming data.
    But again, this is only a potential indicator of drift
    because changes in the target data set may not result
    in drop off in the model's performance.
    For example, a feature might have shifted that doesn't actually
    cause a meaningful change in the model output,
    and this is known as spurious drift.
    So after you detect or suspect model drift,
    you'll want to identify what caused it.
    If we diagnose concept drift, the affected old data
    needs to be relabeled as well as the model retrained.
    If we diagnose data drift enough of the new data
    needs to be labeled to introduce new classes and the model retrained.
    Now, when we consider model drift for TinyML devices,
    there are a couple of things worth noting.
    The first is that retraining might not always be a feasible option.
    For example, in order to conserve battery
    you might have set up a low power communication protocol,
    and this means that we might not be able to send a new model to the device
    if it's large in size.
    Or you may have designed a TinyML application to operate without Wi-Fi,
    for instance, but this means that we won't have the capability of loading
    the new model--
    at least not without retrieving all the devices
    and manually attaching them to a computer to retrain.
    So you want to revisit the process for deciding what am I building
    and who is it for, as we did in Course 1, and try
    and see if the benefits of your device outweigh the potential harms that
    might result from model drift.
    The second thing to keep in mind is that retraining opens up security risks.
    For example, a data poisoning attack could
    occur that might make your model performance worse than what
    you were initially trying to fix.
    So if the harm's associated with the security breach or privacy
    violation are more concerning than a drop in your model's performance,
    you may want to forgo the option of retraining
    altogether and instead just try and run supervised tests of your model
    in the real world and incorporate whatever problems you encounter
    into the normal development process.
    Now the process of monitoring your model after deployment
    basically presents a culmination of all of the ethical concerns
    that you've been dealing with throughout design and development.
    Thinking about the consequences for the user when your model fails
    with special attention paid to precision and recall and metrics for fairness,
    these can help you decide what's an acceptable performance threshold
    and when should model drift trigger retraining.
    In addition, weighing the benefits of your model against privacy and security
    concerns can help you decide whether retraining should even be an option.
    And I know it might sound like a lot, but if you
    persist in taking ethics seriously from start to finish with your project,
    you'll be more likely to see success at the deployment stage and beyond.
    Ende des Transkriptes, zum Anfang springen
    Downloads und Transskripte
    Video
    Videodatei herunterladen
    Transskripte
    Download SubRip (.srt) file
    Download Text (.txt) file
    
\section{Congratulations! You Made it to the Finish Line!}

    0:00 / 0:00Klicken Sie auf die Pfeiltaste nach oben, um das Geschwindigkeitsmenu aufzurufen. Regeln Sie dann mit den Hoch und Runter Pfeiltasten die Geschwindigkeit. Klicken Sie Enter, um die jeweilige Geschwindigkeit festzulegen.Geschwindigkeit1.0xKlicke auf diesen Knopf und diess Video stumm oder nicht stumm zu schalten, oder drücke die HOCH oder RUNTER Knöpfe um die Lautstärke zu erhöhen oder abzusenken.Maximal Lautstärke.Videotranskript
    Beginn des Transkriptes, zum Ende springen.
    VIJAY JANAPA REDDI: Congratulations on making
    it all the way to the end of Course 3.
    It's very exciting that you've come so far.
    You have learned a massive amount of material
    from collecting data and making inferences through Course 1, 2, and 3.
    In Course 3 specifically, we focused very heavily
    on teaching you how to deploy models and make inferences efficiently.
    If we take a step back, I would say that the three key takeaways that you should
    have from this particular course are learning
    about the range of TinyML applications, specifically
    learning how to deploy them from an end-to-end sort of standpoint
    on a real [? ambient ?] microcontroller.
    And doing so while learning how to program in a new machine learning
    framework, which in our case, was TensorFlow Lite for Microcontrollers.
    Now we did this by having you get hands-on experience
    with keyword spotting, visual wake words, and the magic wand example.
    Now stop for a moment and I want you to think
    about what are the interesting things that you uniquely learned
    in each of these application examples.
    We thought very carefully about the things
    that you're going to pick up in each and every single one of them.
    Believe it or not, in every single application,
    there was something more than what you had
    learned in the previous applications.
    It was not simply that we repeated the same pattern.
    We were actually increasing your level of knowledge
    with every single application.
    So stop and think about that for a moment
    because it's going to make you feel good about yourself,
    that you've actually really picked up quite a bit of material.
    Let me focus on some of the takeaways that you should keep in mind
    from a hands-on sort of standpoint.
    Specifically, we talked about how to structure
    an application for deployment.
    What do I mean by that?
    We introduced you to the standard template
    that you'll be using in any TinyML system,
    be it with TensorFlow Lite Micro or otherwise.
    So it's a general takeaway.
    Of course, in the context of TensorFlow Lite Micro,
    we taught you how to use its API so you can run inference models under it.
    Of course, in that same avenue, we also taught you about multi-tenancy
    and how TensorFlow Lite Micro can support multiple models
    so you can realize rich applications.
    For example, combining keyword spotting with visual wake words
    is something that you actually did, which
    actually opens up a new creative avenue for new applications.
    So think about what kinds of sensors you would like to fuse
    and then go ahead and try and deploy multiple models onto the same device.
    Now while you know how to deploy those multiple models,
    we also told you about a system architecture-- a cascade architecture
    that you would actually use, which is used in fact in production
    where inputs from one sensor based on outputs
    of that sensor in the output of the neural network
    would naturally trigger a follow on sensor.
    So you actually did that in practice in this course.
    We also told you about multi-modal sort of execution,
    where you might have sensor data coming in from two different streams
    and you combine them into a single classification, which
    has got pros and cons.
    And hopefully you remember some of these because there might be examples
    that you might go on to build which could very well fall into this realm.
    But hopefully the concept that we told you kind of
    sticks with you, though we did not do the actual example itself.
    Next, we talked a lot about how to actually
    carefully deploy a TinyML application.
    Now that you know the structure, how do you actually deploy it in practice?
    It's much more than just training a neural network
    and getting accuracy and knowing the math behind it.
    There is more to it.
    And that's what we focused on.
    We went all the way from understanding the characteristics
    of the sensors, from a microphone standpoint,
    from an image camera standpoint, and then an IMU standpoint,
    and then feeding them into corresponding networks,
    and then learning how to assess them effectively using appropriate metrics.
    Let me give you a quick reminder.
    We did this exact full flow in the context of keyword spotting,
    where we took the microphone data, which was an analog signal which we converted
    into a frequency domain that we then represented
    in the context of a spectrogram, where it's like frequencies effectively
    changing over time.
    And then we fed that into the neural network
    in order to do keyword spotting.
    So we went through many different stages of pre-processing here.
    But then, of course, we realized that if you just blindly deploy that,
    you might have false positives because keywords like "up" and "upward"
    might have overlaps.
    So if the keyword is "up," but the user says "upward,"
    you might end up having a false positive.
    So you need to be very careful in terms of how you do post-processing.
    So you need to look at the whole signal that's coming in
    and then gain confidence in what the user said
    and then actually go on to making a prediction.
    This introduced us to the concept of actually making multiple predictions
    through our neural network before actually
    triggering the green LED to indicate that we
    have in fact detected the keyword.
    Now this is the kind of thing that you would not ever
    learn if you were just doing things in the Colab.
    This requires you to do physical deployment.
    And that's what you have done.
    And so you are actually advancing in terms
    of knowing how to deploy models onto the microcontrollers.
    Also, we touched a little bit on what are some interesting and emerging
    application use cases.
    And to this end, I specifically focused us
    on the notion of using other types of sensors
    instead of just focusing on a microphone or in sort of the camera,
    because the visual and speech modalities have been heavily exploited
    in the machine learning space.
    For TinyML, it's all about sensors.
    And just do me a favor and go ahead and look up the number of sensors
    that you would have if you're just looking
    for mechanical sensors or electromagnetic sensors, and so forth.
    There are lots of different sensors out there.
    And if you can use machine learning on them,
    then you would be able to unlock all kinds of cool new applications.
    So to this end, we specifically happened to touch
    on the IMU sensor, which allowed us to do
    integration of the [? accelerometer ?] data along with the gyroscope data.
    And together, we were able to create the magic wand application.
    Now, having talked about all this exciting stuff,
    one of the key things we have done is also
    focus on some of the responsible AI aspects
    as we start putting sensors along with smarts all over the place.
    We talk about three key things in this context.
    Susan told us about the issues around model drift.
    We had to worry about what it means to have a concept drift, in other words,
    you train a model and you deploy it in the field at a certain time.
    But over time, your concept of what you're trying to do changes.
    And according to that, you would actually have to adapt your model.
    So you cannot just deploy it once and be done.
    The second aspect of that is the actual data drift.
    You might have trained the model on one class of data
    that you actually have access to, but then as the field sort of evolves,
    a new user base evolves, things end up changing.
    So that means your model then needs to be adjusted accordingly.
    That's just one example.
    Then we worried about security breaches.
    How do we actually design the hardware properly?
    How do we actually get the right policies in place
    so there are security mechanisms are put in place right from the get go?
    And of course, there's a notion of how do we actually protect privacy.
    Fundamentally understanding what privacy means
    and who the actual stakeholders are in privacy is actually a critical aspect.
    So as you go on to build your TinyML applications,
    I strongly urge you to actually think about the responsible AI aspects
    as an integral component of engineering your solution.
    Keep that in mind, because if you forget that,
    you might build the coolest application.
    But if that has serious consequences down the road,
    well, all that effort is going to be lost.
    So taking a step back again, the three key takeaways
    that you've gotten at this point are that you've understood,
    wow, I can actually do cool things with TinyML
    because I can look at interesting sensors.
    And not only that, I can take those sensors
    and I know how to integrate this into the context
    of the full end-to-end pipeline going from pre-processing
    through inference to post-processing, and what the role of pre and post is.
    Because that's often not talked about by people.
    But whether your ML model can really work or not is really
    a function of the pre-processing.
    And even if your model did work well, whether it can actually
    deliver useful information in the real world
    is a function of your post-processing.
    So please keep that in mind and think of that as the end-to-end pipeline.
    Finally, you really did hands-on development
    using to TensorFlow Lite Micro in order to realize those concepts.
    With that said, I hope you enjoyed Course 1, 2, 3 very much.
    And in Course 3 specifically, I hope you got all the nuanced details
    around deploying TinyML.
    In the next video, I'm going to show you how you can move ahead
    with all of the information you have gathered from Course 1, 2, 3,
    so you can determine where you're going next.
    Ende des Transkriptes, zum Anfang springen
    Downloads und Transskripte
    Video
    Videodatei herunterladen
    Transskripte
    Download SubRip (.srt) file
    Download Text (.txt) file    
    
\section{Advanced Topics in TinyML}

Throughout this course, you have developed multiple tinyML projects on your local microcontroller. You have designed a keyword spotting algorithm that was trained using your own customized dataset. You have designed a visual wake words algorithm and then, through transfer learning, adapted this to create a mask detection algorithm. You have even created your own gesturing magic wand!
    
Although you have now come to the end of this class, your journey as a tinyML engineer is only just beginning. You are now equipped with a set of tools and techniques that will be invaluable in the coming years as the field of tinyML reaches ubiquity and is incorporated beyond academia into commercial markets. The exponential increase of IoT devices, increasing prevalence of sensor networks, and transition towards industry 4.0 will continue to drive further development of tinyML at an accelerating pace. Consequently, keeping up-to-date with state-of-the-art methods, as well as advanced and novel techniques will be a necessity for your success as a tinyML engineer.
    
In this reading, we will touch on some of the major topics that we consider “advanced” as well as potential future developments and applications which may become relevant in the near future. Such predictions are, at best, informed guesses about how the field may evolve and may deviate depending on the conflicting needs of industry, end-users, and academia. However, we believe that these topics may act as a helpful stepping stone to you in your continuing development as a tinyML engineer.
    
\subsection{Profiling and Benchmarking}

\textbf{Profiling} is used in software engineering to optimize software performance. Similarly, in tinyML, profiling can be used to optimize a system for its deployment. This might include finding bottlenecks, computational inefficiencies, along with other properties which may negatively impact algorithmic performance. Profiling is typically used to determine where speedup can be achieved to aid in minimizing the computational costs of running an algorithm.
    
To accompany this, \textbf{benchmarking} is used to provide like-for-like comparisons between an algorithmic implementation running on different hardware, different implementations running on the same hardware, and so on. Benchmarking typically involves the use of comparative metrics, which for tinyML might include inference time, operations per second, and memory utilization.
    
    Benchmarking and profiling will become increasingly important for commercial applications where large numbers of devices are running the same code since inefficiencies lead to higher costs and other undesirable effects. TinyMLPerf is a recent attempt to help provide a set of benchmarks for performing comparative measurements of typical tinyML tasks.
    
\subsection{ML Operations (MLOps)}

MLOps is very closely related to DevOps. DevOps focuses on seamless integration of automated testing and benchmarking capabilities to a software engineering environment, allowing updated code and algorithms to be automatically checked for errors and assessed for performance. MLOps would apply this to our tinyML algorithms, which would significantly increase the productivity of developers by providing real-time feedback which can be tailored to the deployment environment visible to the end-user. Several companies focusing on MLOps already exist, and there will likely be more as market penetration of tinyML applications increases.
    
    Edge Impulse is one such example, which makes it easy to gather your dataset, train models, optimize and deploy them. Now that you know what is underneath the hood, you can go ahead and use such tools with good progress.
    
\subsection{Cross-platform Interoperability}

    Algorithm interoperability across various hardware platforms remains an important barrier to tinyML. Microcontroller devices built upon different instruction sets may have (1) different sensor capabilities, (2) be bare-metal as opposed to having a lightweight operating system, (3) different memory capabilities, with or without a memory controller, and (4) different computational capabilities (e.g., clock speed, word lengths). These differences lead to a panoply of challenges for creating algorithms that can run on various hardware platforms.
    
    Several attempts at overcoming this issue have already been attempted. OctoML provides a compiler stack that acts as a middleman that provides interoperability between different deep learning platforms and microcontroller devices. This challenge will become vital to overcoming once the tinyML ecosystem begins to propagate since multiple vendors will have to have products that can work together seamlessly, similar to how different computers and devices are able to interact on the internet today.
    
\subsection{Neural Architecture Search}

    As we have seen throughout the course, it is difficult to find the optimal network configuration for a given application and dataset. Neural Architecture Search (NAS) is a common method used in machine learning to search and test multiple architectures in an attempt to find an optimal configuration automatically. However, we are more resource-constrained with tinyML devices, which must be incorporated into our NAS. Therefore, for tinyML devices, we must recast NAS to produce an optimal configuration given the hardware constraints imposed by our system, such that we can eke out the maximal performance.
    
    Several attempts have been made so far to do this, such as with MCUNet and MicroNets. The commercial drivers for this are strong since devices with higher performance capabilities are preferable in any given application.
    
    No doubt, there will be more items that could be added to this list in the near future, but we hope this will provide you with a valuable starting point for finding additional topics that may be useful for you in future applications.    
    
    
\section{What Do I Do Now?}

 0:00 / 0:00Klicken Sie auf die Pfeiltaste nach oben, um das Geschwindigkeitsmenu aufzurufen. Regeln Sie dann mit den Hoch und Runter Pfeiltasten die Geschwindigkeit. Klicken Sie Enter, um die jeweilige Geschwindigkeit festzulegen.Geschwindigkeit1.0xKlicke auf diesen Knopf und diess Video stumm oder nicht stumm zu schalten, oder drücke die HOCH oder RUNTER Knöpfe um die Lautstärke zu erhöhen oder abzusenken.Maximal Lautstärke.Videotranskript
    Beginn des Transkriptes, zum Ende springen.
    VIJAY JAPANA REDDI: Given that you've now completed course one, two,
    and three, you might be asking yourself this question of what do I do now?
    Well, first and foremost, I want you to stop and reflect on the fact
    that you've come a very long way.
    Recognize that in going from fundamentals of TinyML,
    to learning about the applications of TinyML, to actually deploying TinyML,
    you went through this complete machine learning workflow,
    going from collecting data to making inferences.
    What I want to do now is to step you through each one of those cells
    in order to help you reflect back on the critical items,
    and then I'll use that as a seed to give you new ideas about things
    that you might be able to do.
    Recall that when it comes to like dealing with data,
    it's extremely important to know how you're going to collect the right data.
    In this vein, we actually taught you how to do your own custom data
    collection in course three.
    Hopefully, that exercise taught you about all the things
    that you have to think about in terms of carefully identifying
    the requirements for your data set.
    Now, it's one thing to hear about the requirements, which we explained
    to you through speech commands, for example,
    but it's a whole different thing when you actually have
    to go about in collecting your data.
    So hopefully, that exercise empowers you into understanding the value of data
    sets and how you might be able to go ahead in building new data sets.
    Second, we talked a lot about preprocessing data.
    Now, we talked about preprocessing data in the form of data
    augmentation in course one, but especially
    in course two and course three, we emphasize the role
    of preprocessing, where we take signals in one domain
    and might convert them into a totally different representation.
    In other words, we took an audio signal and converted that
    into a spectrogram image, and then used that to be
    able to actually detect what the audio word was which is a remarkable process
    that we did.
    And that's only possible because we preprocess the data very intelligently,
    and that's quite crucial for any ML application.
    But especially in the context of TinyML, I
    want you to think about what it means to preprocess sensor data that's
    coming in, because if you don't think about this carefully,
    you might accidentally make the mistake of getting
    a rich amount of data that would not really
    work well through a neural network.
    For example, we talked about end-to-end learning
    in the context of taking IMU sensor data and feeding it directly
    into the neural network instead of going through the preprocessing stage, where
    we use traditional methods to rasterize the image into a 2D plane.
    That's an example.
    Next, we talked about designing a model.
    Well, mostly, I hope you understood that, when
    you're thinking of deployment scenarios, you really
    have to think very carefully about where you're deploying the model.
    So because we have very limited amount of microcontroller memory,
    we thought very carefully about making sure
    that we had a small neural network that would easily fit
    on both the SRAM as well as the flash.
    So we need to be conscious about actually designing the models,
    until as you go onto build new TinyML models, keep this in your mind.
    And there's also other stuff you might be
    able to pick up, like neural architecture search, which
    helps you automatically try and find the right, small, little neural network
    that would perform best from an accuracy standpoint.
    So there's much more out there, but hopefully, you
    realize that the key nugget here is that you've
    got to have a small model that achieves reasonable accuracy.
    Then, of course, you learned about how to train models, especially
    in course one, and the key takeaway here is really
    understanding how we get the original data set which again goes back
    to the notion of collecting data.
    And then you learned how that data can actually
    be split into training set and the test set,
    and the training set can further be broken into really the training set
    that's actually used for evaluating the model,
    doing the hyperparameter tuning space, and then
    how the test set is ultimately used to validate
    the accuracy of the final model.
    And you also got exposed to doing training using Google Colab
    rather than using your own resources.
    Then, of course, we talk a lot about evaluating and optimizing models.
    Specifically in course one, for instance,
    we talked about underfitting and overfitting.
    You don't want the network to underfit.
    You don't want the network to overfit.
    You got to try and find that right amount of balance
    in order to get the best experience from your neural network.
    And in that process, as we went into course two,
    we looked at how the accuracy of the model
    changes when you start throwing in new data
    and how it works across different kinds of accents and so forth.
    And we also learned about optimizing the models.
    Remember, you typically end up with a float 32 model,
    but float 32 models are not good for microcontrollers
    which sometimes might not, in fact, have floating point support.
    And you also learned that floating point is much more expensive
    in terms of computation, as compared to integer operations.
    And so one of the things you want to do is
    you want to quantize the models, both for size as well as latency.
    And in that process, you learned about PTQ and QAT,
    where PTQ stands for Post-Training Quantization,
    and QAT stands for Quantization Array Training.
    In other words, after you quantize a model,
    you retrain the model in order to get back the accuracy lost that
    happens as a result of quantization.
    Now, that's a very powerful concept.
    By and large, if you look at TinyML systems,
    you're almost always going to have to have a quantized model.
    Right now, we're talking about in 8 quantized model,
    but their are architectures out there that are starting
    to support in 4 models, for instance.
    So keep this in mind, because you're definitely
    going to want to reuse quantization with whatever application you do.
    Because all microcontrollers, at the end of the day,
    are going to run much more efficiently with [INAUDIBLE]..
    And then you learned the process of actually converting
    the model into a flat buffer format, so it's actually
    widely deployable across a wide range of devices.
    Specifically at this point, you learn the differences
    between the different types of frameworks you
    have for machine learning, going all the way
    from TensorFlow, to TensorFlow Lite, to then TensorFlow Lite Micro.
    Hopefully, if one was to ask you, hey, what
    is the difference between the three of them,
    you're able to explain this from the perspective of whether you
    need to train or not, how many different operators you need to support,
    what kind of operating system support is needed to run these systems,
    and so forth.
    And so there's a rich amount of knowledge you picked up on this.
    And so when you go on to build your own application, reflect on this,
    and see what is a right tool that you need to have.
    While we talk a lot about TensorFlow Lite Micro in this particular course,
    I did try and give you hints about how other systems also work.
    So keep an open mind, and go and explore other frameworks
    is the real key takeaway here.
    And of course, when it came down to actually deploying things,
    we did a lot around TensorFlow Lite Micro
    which is a production-grade inference framework for TinyML.
    So go ahead and use the APIs that you learn,
    because you can actually build production-grade systems using
    the exact code that we have actually shown you.
    Now, you would also have picked up the fact
    that, if you wanted to build systems like TensorFlow Lite Micro,
    you would have an appreciation for understanding the complexities.
    Why do we need to build a system like TensorFlow Lite Micro?
    Well, that's because there's lots of problems
    when it comes to the embedded ecosystem which is really fragmented.
    So reflect back on some of those readings we gave you
    which explain the hardware and software differences,
    the heterogeneity, the resource constraints, the missing operating
    system support, all of the things that we typically take
    for granted in much larger systems.
    The lack of that support makes it extremely crucial
    for us to learn how to build these lean and mean inference engines.
    Finally, you learned the art of making inferences through the experiences
    that you've had with keyword spotting, Visual Wake Words, and the Magic Wand.
    Through these application deployments, I hope
    you realize some of the pitfalls of these application.
    They don't always work well, but that's not the point of this course.
    It's not about getting things working well.
    It's about getting things working, because the lessons here really
    are about when you get the things working,
    then you start experiencing when something works
    and when something does not work.
    That is what's going to make you a strong TinyML engineer, because you
    have an appreciation for where the issues are really going to stem from.
    All that said, now that you've got this vast amount of knowledge going
    all the way from collecting data to making inferences,
    the question in your head might be, hey, what
    are some new applications I can unlock?
    Well, that's a great question to have.
    I just want you to keep a certain mentality in mind.
    I want you to focus on the Golden Circle approach,
    when you're trying to build your next TinyML application.
    It does not matter how big or small your application is, but always focus first
    on the question of why.
    Why do you want to build this application?
    What is it that is actually interesting?
    What's the purpose?
    Why is this particular TinyML application
    going to unlock the future in a way that you see it?
    Then, focus on how.
    What is your value proposition?
    What is the unique approach that you bring to the table that separates you
    from someone else who might be interested in building a same TinyML
    solution?
    And finally, then think about what, how you actually went about it.
    So don't focus on how just because you know all the different things
    we have thought you, don't start there.
    Instead, I'm asking you to flip your mindset,
    and go into identifying the root cause of wanting to build an application,
    and then work your way outwards.
    Because if you do that, you will end up building a successful TinyML
    application, and that's exactly what I want to see.
    I personally can guarantee you that the success from this course for me
    comes not so much from how many students are taking it but much more
    from what are the things that you're building, now that you have learned all
    these things that we have taught you.
    With that said, remember that the future of ML is tiny and bright.
    I'm looking forward to all the amazing things
    that are going to come out of all the exciting projects
    that you're going to be building.
    Have fun.
    Ende des Transkriptes, zum Anfang springen
    Downloads und Transskripte
    Video
    Videodatei herunterladen
    Transskripte
    Download SubRip (.srt) file
    Download Text (.txt) file    
    
    
\section{What Do I Do Now?}

    0:00 / 0:00Klicken Sie auf die Pfeiltaste nach oben, um das Geschwindigkeitsmenu aufzurufen. Regeln Sie dann mit den Hoch und Runter Pfeiltasten die Geschwindigkeit. Klicken Sie Enter, um die jeweilige Geschwindigkeit festzulegen.Geschwindigkeit1.0xKlicke auf diesen Knopf und diess Video stumm oder nicht stumm zu schalten, oder drücke die HOCH oder RUNTER Knöpfe um die Lautstärke zu erhöhen oder abzusenken.Maximal Lautstärke.Videotranskript
    Beginn des Transkriptes, zum Ende springen.
    VIJAY JAPANA REDDI: Given that you've now completed course one, two,
    and three, you might be asking yourself this question of what do I do now?
    Well, first and foremost, I want you to stop and reflect on the fact
    that you've come a very long way.
    Recognize that in going from fundamentals of TinyML,
    to learning about the applications of TinyML, to actually deploying TinyML,
    you went through this complete machine learning workflow,
    going from collecting data to making inferences.
    What I want to do now is to step you through each one of those cells
    in order to help you reflect back on the critical items,
    and then I'll use that as a seed to give you new ideas about things
    that you might be able to do.
    Recall that when it comes to like dealing with data,
    it's extremely important to know how you're going to collect the right data.
    In this vein, we actually taught you how to do your own custom data
    collection in course three.
    Hopefully, that exercise taught you about all the things
    that you have to think about in terms of carefully identifying
    the requirements for your data set.
    Now, it's one thing to hear about the requirements, which we explained
    to you through speech commands, for example,
    but it's a whole different thing when you actually have
    to go about in collecting your data.
    So hopefully, that exercise empowers you into understanding the value of data
    sets and how you might be able to go ahead in building new data sets.
    Second, we talked a lot about preprocessing data.
    Now, we talked about preprocessing data in the form of data
    augmentation in course one, but especially
    in course two and course three, we emphasize the role
    of preprocessing, where we take signals in one domain
    and might convert them into a totally different representation.
    In other words, we took an audio signal and converted that
    into a spectrogram image, and then used that to be
    able to actually detect what the audio word was which is a remarkable process
    that we did.
    And that's only possible because we preprocess the data very intelligently,
    and that's quite crucial for any ML application.
    But especially in the context of TinyML, I
    want you to think about what it means to preprocess sensor data that's
    coming in, because if you don't think about this carefully,
    you might accidentally make the mistake of getting
    a rich amount of data that would not really
    work well through a neural network.
    For example, we talked about end-to-end learning
    in the context of taking IMU sensor data and feeding it directly
    into the neural network instead of going through the preprocessing stage, where
    we use traditional methods to rasterize the image into a 2D plane.
    That's an example.
    Next, we talked about designing a model.
    Well, mostly, I hope you understood that, when
    you're thinking of deployment scenarios, you really
    have to think very carefully about where you're deploying the model.
    So because we have very limited amount of microcontroller memory,
    we thought very carefully about making sure
    that we had a small neural network that would easily fit
    on both the SRAM as well as the flash.
    So we need to be conscious about actually designing the models,
    until as you go onto build new TinyML models, keep this in your mind.
    And there's also other stuff you might be
    able to pick up, like neural architecture search, which
    helps you automatically try and find the right, small, little neural network
    that would perform best from an accuracy standpoint.
    So there's much more out there, but hopefully, you
    realize that the key nugget here is that you've
    got to have a small model that achieves reasonable accuracy.
    Then, of course, you learned about how to train models, especially
    in course one, and the key takeaway here is really
    understanding how we get the original data set which again goes back
    to the notion of collecting data.
    And then you learned how that data can actually
    be split into training set and the test set,
    and the training set can further be broken into really the training set
    that's actually used for evaluating the model,
    doing the hyperparameter tuning space, and then
    how the test set is ultimately used to validate
    the accuracy of the final model.
    And you also got exposed to doing training using Google Colab
    rather than using your own resources.
    Then, of course, we talk a lot about evaluating and optimizing models.
    Specifically in course one, for instance,
    we talked about underfitting and overfitting.
    You don't want the network to underfit.
    You don't want the network to overfit.
    You got to try and find that right amount of balance
    in order to get the best experience from your neural network.
    And in that process, as we went into course two,
    we looked at how the accuracy of the model
    changes when you start throwing in new data
    and how it works across different kinds of accents and so forth.
    And we also learned about optimizing the models.
    Remember, you typically end up with a float 32 model,
    but float 32 models are not good for microcontrollers
    which sometimes might not, in fact, have floating point support.
    And you also learned that floating point is much more expensive
    in terms of computation, as compared to integer operations.
    And so one of the things you want to do is
    you want to quantize the models, both for size as well as latency.
    And in that process, you learned about PTQ and QAT,
    where PTQ stands for Post-Training Quantization,
    and QAT stands for Quantization Array Training.
    In other words, after you quantize a model,
    you retrain the model in order to get back the accuracy lost that
    happens as a result of quantization.
    Now, that's a very powerful concept.
    By and large, if you look at TinyML systems,
    you're almost always going to have to have a quantized model.
    Right now, we're talking about in 8 quantized model,
    but their are architectures out there that are starting
    to support in 4 models, for instance.
    So keep this in mind, because you're definitely
    going to want to reuse quantization with whatever application you do.
    Because all microcontrollers, at the end of the day,
    are going to run much more efficiently with [INAUDIBLE]..
    And then you learned the process of actually converting
    the model into a flat buffer format, so it's actually
    widely deployable across a wide range of devices.
    Specifically at this point, you learn the differences
    between the different types of frameworks you
    have for machine learning, going all the way
    from TensorFlow, to TensorFlow Lite, to then TensorFlow Lite Micro.
    Hopefully, if one was to ask you, hey, what
    is the difference between the three of them,
    you're able to explain this from the perspective of whether you
    need to train or not, how many different operators you need to support,
    what kind of operating system support is needed to run these systems,
    and so forth.
    And so there's a rich amount of knowledge you picked up on this.
    And so when you go on to build your own application, reflect on this,
    and see what is a right tool that you need to have.
    While we talk a lot about TensorFlow Lite Micro in this particular course,
    I did try and give you hints about how other systems also work.
    So keep an open mind, and go and explore other frameworks
    is the real key takeaway here.
    And of course, when it came down to actually deploying things,
    we did a lot around TensorFlow Lite Micro
    which is a production-grade inference framework for TinyML.
    So go ahead and use the APIs that you learn,
    because you can actually build production-grade systems using
    the exact code that we have actually shown you.
    Now, you would also have picked up the fact
    that, if you wanted to build systems like TensorFlow Lite Micro,
    you would have an appreciation for understanding the complexities.
    Why do we need to build a system like TensorFlow Lite Micro?
    Well, that's because there's lots of problems
    when it comes to the embedded ecosystem which is really fragmented.
    So reflect back on some of those readings we gave you
    which explain the hardware and software differences,
    the heterogeneity, the resource constraints, the missing operating
    system support, all of the things that we typically take
    for granted in much larger systems.
    The lack of that support makes it extremely crucial
    for us to learn how to build these lean and mean inference engines.
    Finally, you learned the art of making inferences through the experiences
    that you've had with keyword spotting, Visual Wake Words, and the Magic Wand.
    Through these application deployments, I hope
    you realize some of the pitfalls of these application.
    They don't always work well, but that's not the point of this course.
    It's not about getting things working well.
    It's about getting things working, because the lessons here really
    are about when you get the things working,
    then you start experiencing when something works
    and when something does not work.
    That is what's going to make you a strong TinyML engineer, because you
    have an appreciation for where the issues are really going to stem from.
    All that said, now that you've got this vast amount of knowledge going
    all the way from collecting data to making inferences,
    the question in your head might be, hey, what
    are some new applications I can unlock?
    Well, that's a great question to have.
    I just want you to keep a certain mentality in mind.
    I want you to focus on the Golden Circle approach,
    when you're trying to build your next TinyML application.
    It does not matter how big or small your application is, but always focus first
    on the question of why.
    Why do you want to build this application?
    What is it that is actually interesting?
    What's the purpose?
    Why is this particular TinyML application
    going to unlock the future in a way that you see it?
    Then, focus on how.
    What is your value proposition?
    What is the unique approach that you bring to the table that separates you
    from someone else who might be interested in building a same TinyML
    solution?
    And finally, then think about what, how you actually went about it.
    So don't focus on how just because you know all the different things
    we have thought you, don't start there.
    Instead, I'm asking you to flip your mindset,
    and go into identifying the root cause of wanting to build an application,
    and then work your way outwards.
    Because if you do that, you will end up building a successful TinyML
    application, and that's exactly what I want to see.
    I personally can guarantee you that the success from this course for me
    comes not so much from how many students are taking it but much more
    from what are the things that you're building, now that you have learned all
    these things that we have taught you.
    With that said, remember that the future of ML is tiny and bright.
    I'm looking forward to all the amazing things
    that are going to come out of all the exciting projects
    that you're going to be building.
    Have fun.
    Ende des Transkriptes, zum Anfang springen
    Downloads und Transskripte
    Video
    Videodatei herunterladen
    Transskripte
    Download SubRip (.srt) file
    Download Text (.txt) file    
    
 \section{TinyMLx Project Extension (Optional)}
 
     At this point, you’ve amassed an incredible amount of knowledge on Tiny Machine Learning. You have seen how models are developed using the machine learning workflow, from data collection to training and deployment. You have seen how models are produced using TensorFlow and can be ported into light-weight frameworks, like TensorFlow Lite and Micro. You have also seen how these models can interact with various forms of input data, such as time-series data from sensors and microphones, to images from a camera.
     
     In Course 3, we have looked closely at how to deploy TinyML models within embedded systems, using the Arduino Nano 33 BLE Sense microcontroller development board as a representative host. Further still, you have already leveraged TensorFlow Micro to deploy a range of TinyML applications on heavily resource-constrained hardware,  creating a roadmap for end-to-end deployment.
     
     We want everyone who completes the TinyML edX course to be confident that they can build custom TinyML applications using the knowledge they have gained. As such, we invite you all to join the public TinyMLx community on Discourse and to design your own TinyML projects beyond the scope of this class. 
     
     To help motivate you and foster assimilation to the wider TinyMLx community, we’ve set up an open TinyMLx project competition with a number of prizes for exemplary work. This project represents an opportunity to demonstrate what you have learned and we hope will give you an incentive and a forum to push beyond the course examples to introduce new use cases and applications of TinyML. We can’t wait to see what you’ll design and build collaboratively as a part of the growing TinyMLx community.
     
     For more information about the optional project extension including FAQs, prizes, and submission details, please follow this link to Vijay’s Open Call for Projects on Discourse.   
     
\section{Appendix}

You may have noticed throughout the course that we have provided appendix links with additional supplementary material on various topics (e.g., powering your TinyML kit with a battery, details on serial protocols). All of those documents can be found in our appendix GitHub repository: \url{https://github.com/tinyMLx/appendix}
